% Options for packages loaded elsewhere
\PassOptionsToPackage{unicode}{hyperref}
\PassOptionsToPackage{hyphens}{url}
\PassOptionsToPackage{space}{xeCJK}
\documentclass[
  10pt,
  a4paper,
]{article}
\usepackage{xcolor}
\usepackage[margin=18mm]{geometry}
\usepackage{amsmath,amssymb}
\setcounter{secnumdepth}{-\maxdimen} % remove section numbering
\usepackage{iftex}
\ifPDFTeX
  \usepackage[T1]{fontenc}
  \usepackage[utf8]{inputenc}
  \usepackage{textcomp} % provide euro and other symbols
\else % if luatex or xetex
  \usepackage{unicode-math} % this also loads fontspec
  \defaultfontfeatures{Scale=MatchLowercase}
  \defaultfontfeatures[\rmfamily]{Ligatures=TeX,Scale=1}
\fi
\usepackage{lmodern}
\ifPDFTeX\else
  % xetex/luatex font selection
  \setmainfont[AutoFakeBold=1.5,AutoFakeSlant=0.15]{Arial Unicode MS}
  \setmonofont[]{Menlo}
  \ifXeTeX
    \usepackage{xeCJK}
    \setCJKmainfont[]{Songti SC}
  \fi
  \ifLuaTeX
    \usepackage[]{luatexja-fontspec}
    \setmainjfont[]{Songti SC}
  \fi
\fi
% Use upquote if available, for straight quotes in verbatim environments
\IfFileExists{upquote.sty}{\usepackage{upquote}}{}
\IfFileExists{microtype.sty}{% use microtype if available
  \usepackage[]{microtype}
  \UseMicrotypeSet[protrusion]{basicmath} % disable protrusion for tt fonts
}{}
\usepackage{setspace}
\makeatletter
\@ifundefined{KOMAClassName}{% if non-KOMA class
  \IfFileExists{parskip.sty}{%
    \usepackage{parskip}
  }{% else
    \setlength{\parindent}{0pt}
    \setlength{\parskip}{6pt plus 2pt minus 1pt}}
}{% if KOMA class
  \KOMAoptions{parskip=half}}
\makeatother
\usepackage{longtable,booktabs,array}
\newcounter{none} % for unnumbered tables
\usepackage{calc} % for calculating minipage widths
% Correct order of tables after \paragraph or \subparagraph
\usepackage{etoolbox}
\makeatletter
\patchcmd\longtable{\par}{\if@noskipsec\mbox{}\fi\par}{}{}
\makeatother
% Allow footnotes in longtable head/foot
\IfFileExists{footnotehyper.sty}{\usepackage{footnotehyper}}{\usepackage{footnote}}
\makesavenoteenv{longtable}
\setlength{\emergencystretch}{3em} % prevent overfull lines
\providecommand{\tightlist}{%
  \setlength{\itemsep}{0pt}\setlength{\parskip}{0pt}}
\usepackage{xurl}
\usepackage{etoolbox}
\emergencystretch=3em
\setlength{\parindent}{0pt}
\setlength{\parskip}{5pt}
\AtBeginEnvironment{longtable}{\fontsize{8.7}{12.2}\selectfont\setlength{\tabcolsep}{3pt}}
\tracinglostchars=3
\AtBeginDocument{\hypersetup{unicode=true,pdftitle={回到1803：拿破仑帝国的可能性边界}}}
\usepackage{bookmark}
\IfFileExists{xurl.sty}{\usepackage{xurl}}{} % add URL line breaks if available
\urlstyle{same}
\hypersetup{
  hidelinks,
  pdfcreator={LaTeX via pandoc}}

\author{}
\date{}

\newcommand{\literalquote}{{\addfontfeatures{Mapping=}\char34}}
\setCJKmonofont{Songti SC}
\pdfstringdefDisableCommands{\def\literalquote{"}}
\begin{document}

\setstretch{1.15}
\section{回到1803：拿破仑帝国的可能性边界}\label{ux56deux52301803ux62ffux7834ux4ed1ux5e1dux56fdux7684ux53efux80fdux6027ux8fb9ux754c}

\subsection{一部反事实的欧洲史，1803--1848}\label{ux4e00ux90e8ux53cdux4e8bux5b9eux7684ux6b27ux6d32ux53f218031848}

\textbf{终审报告 · 运行 r\_5b1357a9c6（v4，53卡）· Last Order 撰} （红队
10043 审查后由 Last Order 亲自修订；不再派发子节点）

\begin{center}\rule{0.5\linewidth}{0.5pt}\end{center}

\subsubsection{凡例}\label{ux51e1ux4f8b}

\begin{enumerate}
\def\labelenumi{\arabic{enumi}.}
\tightlist
\item
  \textbf{三段式标签。}
  全书每一处推演都区分：【锚】＝史料或文书（附作者、书名、页码或原文）；【机】＝史实类比或理论机制（附适用条件）；【推】＝在情景下的具体演变，标高／中／低置信。锚点与机制约束推演，推演是被标注的推断。禁止以\literalquote{}无来源\literalquote{}拒绝推演，也禁止无锚点的推演冒充史实。第一部各章按单位卡原结构行文，【锚】与【推】逐段标出；【机】在援引一个可命名的类比或理论时单独标出，其余场合机制即锚点与推演之间的那句话（如\literalquote{}反应函数总律\literalquote{}\literalquote{}四闸门\literalquote{}\literalquote{}里德机制\literalquote{}），不另加标签。第二至五部为总装，【锚】仍逐处标出，机制以定理与开关的形式命名（T1--T3、开关一／二、α／β终局）。
\item
  \textbf{情景代号。}
  S0史实；S1不列颠中立化（含入侵线）；S2大陆整固（不废西王、不打俄、对英武装共存------以提尔西特后处境为前提）；S3封锁绞杀；S4东方转向；S5秩序固化1815--48（含一次继承）；S6行省化最大版。政策菜单B1入侵英国与否；B2对奥普宽严；B3提尔西特三版；B4伊比利亚三版（a史实／b保波旁并联姻／c兼并至埃布罗）；B5奥地利肢解／婚姻；B6
  1810兼并封锁vs许可证自由化、俄婚vs奥婚；B7对俄三版；B8秩序建构（仅采当时提过的方案）。
\item
  \textbf{\literalquote{}行省化\literalquote{}三层。}
  A＝并合省（法国省制直辖）；B＝法典化藩属（保有本国军政机器、行法国法典）；C＝强制同盟（条约约束）。本书区分\literalquote{}字面版\literalquote{}（A层覆盖大陆）、\literalquote{}边界版\literalquote{}（A层限于内圈）、\literalquote{}控制版\literalquote{}（大陆四层控制＋俄受约束＋英被击败或中立化）。
\item
  \textbf{概率。}
  全书概率均为研究者裁量的\textbf{条件判断带}，不是历史频率、统计置信区间或市场赔率；不得相乘。凡引用模型输出，注明\literalquote{}模型非观测\literalquote{}。
\item
  \textbf{引用诚信。}
  只引读到的版本；转引标\literalquote{}转引\literalquote{}；圣赫勒拿与回忆录标\literalquote{}回溯性\literalquote{}；档案系列名只作谱系不假装亲阅；索引库查无≠文献不存在。引文保留原文语言。
\item
  \textbf{单位与货币。}
  法郎（m＝百万）、英镑（£）、比索（pesos）、里拉、盾（fl.）、塔勒、卢布按原文；人数\literalquote{}在册\literalquote{}≠\literalquote{}在场\literalquote{}（在场≈在册×0.85）；战列舰分H（舰体）／I（现役）／A12（十二个月内可修复武装配员）。
\end{enumerate}

\begin{center}\rule{0.5\linewidth}{0.5pt}\end{center}

\subsubsection{目录}\label{ux76eeux5f55}

\textbf{导论}　问题、方法、上一轮的三病、本轮的修正、红队的十项异议与处置

\textbf{第一部　起点：1803年的世界与每一个行动者的能力包络} -
第一章　法国（一）决策中枢：拿破仑本人、大臣与秩序蓝图 -
第二章　法国（二）陆军：兵员、马匹、军械、后勤与占领军上限 -
第三章　法国（三）海军与海事经济 -
第四章　法国（四）财政、银行与统治阶级 -
第五章　法国（五）社会与政权存续：征兵容忍、教会、保王派、继承 -
第六章　不列颠（一）政治：内阁、议会、王室、和平党与登陆后的决策 -
第七章　不列颠（二）财政与信用 - 第八章　不列颠（三）皇家海军与海军部 -
第九章　不列颠（四）陆军与本土防御 -
第十章　不列颠（五）工业、商业与劳工 - 第十一章　不列颠（六）海外帝国 -
第十二章　爱尔兰 - 第十三章　俄国（一）宫廷与外交精英 -
第十四章　俄国（二）贵族、农奴经济与财政 -
第十五章　俄国（三）军队与多线战争 - 第十六章　波兰-立陶宛 -
第十七章　奥地利 - 第十八章　普鲁士 - 第十九章　南德诸邦 -
第二十章　北中德意志 - 第二十一章　荷兰与瑞士 -
第二十二章　丹麦-挪威与瑞典 - 第二十三章　意大利王国 -
第二十四章　那不勒斯与西西里 - 第二十五章　教廷与欧洲诸教会 -
第二十六章　西班牙（一）宫廷与国家 - 第二十七章　西班牙（二）社会与抵抗
- 第二十八章　葡萄牙与巴西 - 第二十九章　奥斯曼帝国 -
第三十章　埃及、穆罕默德·阿里与红海 - 第三十一章　波斯、阿富汗与旁遮普 -
第三十二章　东印度公司 - 第三十三章　印度诸邦与法属印度洋 -
第三十四章　美国 - 第三十五章　海地与加勒比 - 第三十六章　新西班牙与古巴
- 第三十七章　南美 - 第三十八章　国际金融家 -
第三十九章　大陆体系作为执行机器 - 第四十章　科技与工业能力 -
第四十一章　思潮、民族主义、宗教复兴与社会运动 -
第四十二章　生态、人口与资源储量 - 第四十三章　并合省内圈的整合实绩 -
第四十四章　并合省外圈的断裂点 -
第四十五章　反抗的决定因素：三十五个地区---时段的比较 -
第四十六章　整合绩效面板 - 第四十七章　市场的证词：公债、汇率与保费

\textbf{第二部　如何赢：路径的对弈} -
第四十八章　对弈规则、奇迹的计价与三条互斥定理 -
第四十九章　被否决的三条路：入侵、绞杀、东方 -
第五十章　主线的重构：S2本义（史实至提尔西特＋两个开关）与路径B（1803年即切换）
- 第五十一章　两个开关的解剖 -
第五十二章　英国如何\literalquote{}不败而接受\literalquote{}：条件、窗口、安特卫普 -
第五十三章　概率：从1803年出发，这条路有多现实

\textbf{第三部　行省化：最重要的问题} - 第五十四章　三层回答与五本账 -
第五十五章　为什么外圈必断 - 第五十六章　乙层：技术可闭合、政治未证 -
第五十七章　官员、宪兵、语言与教会

\textbf{第四部　法胜欧洲的运作，1815--1848} -
第五十八章　骨架年表（修正后） - 第五十九章　经济秩序 -
第六十章　社会、思想、教会与继承 - 第六十一章　各国在秩序中的位置 -
第六十二章　1848

\textbf{第五部　分问} - 第六十三章　英法海军差距 -
第六十四章　西班牙的终局 - 第六十五章　罗马王与意大利 -
第六十六章　普鲁士与奥地利 - 第六十七章　美洲 - 第六十八章　印度

\textbf{结论}　主线、备选结论群、置信、改判证据、未决问题
\textbf{附录}　一、材料总目（已读／转引／未获）；二、口径冲突表；三、红队意见处置表；四、数字总表

\begin{center}\rule{0.5\linewidth}{0.5pt}\end{center}

\section{导论}\label{ux5bfcux8bba}

\subsection{一、问题}\label{ux4e00ux95eeux9898}

回到1803年。亚眠和约签署十四个月后，英法在马耳他问题上再度决裂；波拿巴是终身第一执政，帝国尚未称帝，大军团在布洛涅集结，皇家海军在海峡外等待。从这一刻出发，有没有一条路，能让法国\textbf{赢得战争}、\textbf{控制整个欧洲}、并建立一个\textbf{长久的}法国霸权秩序？最重要的一问是：\textbf{法国把欧洲行省化，是否真有合理、现实、可行的路径？}
随之而来的是六个分问：一个法国胜利的欧洲在十九世纪上半叶如何运作；英法海军的差距会缩到什么程度；西班牙会成为共和国、波旁立宪君主国还是波拿巴立宪君主国；拿破仑是否有意愿和能力让罗马王统一中南意大利；普鲁士还会比奥地利强势吗；法国能否帮助西班牙盟友保住美洲，拉美格局如何变化；拿破仑对收回法属印度乃至夺取英属印度是否有意愿，如何做。

这不是一个可以用\literalquote{}拿破仑做了什么\literalquote{}来回答的问题，而是一个\literalquote{}欧洲的每一个行动者最多能做到什么、在什么条件下会那样做\literalquote{}的问题。上一轮研究（r\_b55f2c1cf5，十八卡，四百七十八项材料）之所以失败，正是因为它把\literalquote{}细节有来源\literalquote{}误读成\literalquote{}只说有来源的事\literalquote{}，承认别国不会在法国改变政策时按史实行动，却没有去研究别国的国情，并且拒绝给出区间与终局。本轮把研究单位下沉到\textbf{历史参与主体}：四十四个国家机器、军队、阶级、教会、公司、金融网络、思想场与生态系统，各按十二个维度穷尽国情------决策结构、法律与行政渗透、财政货币、经济与资源、撷取与动员上限、军事、社会结构、抵抗传统、思潮与舆论、科技、生态人口、交流与流动------先给\textbf{能力包络}（历史峰值、崩溃点、未用潜力），再给对法国政策菜单的\textbf{反应函数}，再推演到1848年。九张跨单位卡把这些单位放到同一张桌上：反抗的决定因素、整合绩效、市场隐含预期、行省化总判、海军差距、大陆经济、胜利路径对弈、世界年表、概率聚合。

\subsection{二、方法与纪律}\label{ux4e8cux65b9ux6cd5ux4e0eux7eaaux5f8b}

本书的每一条推演都受\literalquote{}三段式\literalquote{}约束：先给史料锚点，再给机制及其适用条件，最后才是在情景下的推演，并标注置信。锚点不是装饰。它们决定了机制能否成立，机制决定了推演的边界。反过来，我们也拒绝上一轮的怯懦：\textbf{\literalquote{}史实如此\literalquote{}不是结论的理由}。每一个单位都被迫回答\literalquote{}它最多能做到什么\literalquote{}。

\literalquote{}胜利\literalquote{}被操作化为三层：军事胜利；政治控制（四层：A并合省／B法典化藩属／C强制同盟／D受约束中立）；长久（跨过拿破仑的寿限并完成一次继承）。\literalquote{}整个欧洲\literalquote{}被定义为：大陆四层控制成立，俄国受约束，英国被击败或中立化------其中\literalquote{}英国不败而接受\literalquote{}是允许的形态。\literalquote{}行省化\literalquote{}是三个不同的东西：A层覆盖大陆的字面版，A层限于内圈的边界版，以及不问省制只问控制的控制版。本书对三者的判决相反，读者不可混读。

概率是本书最容易被误读的部分。全书所有百分数都是研究者在明示假设下的\textbf{条件判断带}：它们不是频率，不是统计置信区间，也不是市场赔率。四十四张卡各自给出的区间不能相乘；两张综合卡（P6对弈、P8聚合）用了不同的口径，本书并列保留，不投票裁决。

\subsection{三、上一轮的三病与本轮的四条契约}\label{ux4e09ux4e0aux4e00ux8f6eux7684ux4e09ux75c5ux4e0eux672cux8f6eux7684ux56dbux6761ux5951ux7ea6}

三病：以史料拼凑替代推演；不研究他国国情；不给区间与终局。四条契约：三段式；先能力包络后反应函数；每卡给主线＋备选结论群＋置信＋至少两条改判证据；引用诚信（只引读到的版本、转引标注、圣赫勒拿材料标回溯性、档案系列名为谱系不假装亲阅、索引库查无不等于文献不存在）。

\subsection{四、红队的十项异议与本书的处置}\label{ux56dbux7ea2ux961fux7684ux5341ux9879ux5f02ux8baeux4e0eux672cux4e66ux7684ux5904ux7f6e}

节点综合完成后，红队（10043）抽核了八处引文（全部在下载件中命中）、复算了F2/P2/P8/B2的数字（相符）、复跑了p8\_tree.py、p2\_joint\_model.py、verify\_p4.py（与已声明CSV一致），并提出十项实质异议。用户裁定不再派发子节点，由本书作者亲自修正。处置如下（详见附录三）：

{\def\LTcaptype{none} % do not increment counter
\begin{longtable}[]{@{}
  >{\raggedright\arraybackslash}p{(\linewidth - 2\tabcolsep) * \real{0.5000}}
  >{\raggedright\arraybackslash}p{(\linewidth - 2\tabcolsep) * \real{0.5000}}@{}}
\toprule\noalign{}
\begin{minipage}[b]{\linewidth}\raggedright
异议
\end{minipage} & \begin{minipage}[b]{\linewidth}\raggedright
处置
\end{minipage} \\
\midrule\noalign{}
\endhead
\bottomrule\noalign{}
\endlastfoot
\textbf{K1\literalquote{}大陆四层控制\literalquote{}在路径B下无条约支撑}：P6把奥地利（申布伦限军150,000）、普鲁士（1808巴黎公约限军42,000）列为C层，把1808年三十九邦莱茵邦联与华沙公国列为B层------四项都是路径B明文移除的1805/1806/1807/1809战争的和约产物
& \textbf{接受，改判。}
第五十章重构主线：区分\literalquote{}S2本义\literalquote{}（史实打到1807年提尔西特，再切换两个开关）与\literalquote{}路径B\literalquote{}（1803年即切换）。控制版在S2本义下有条约支撑；在路径B下只以\literalquote{}primacy\literalquote{}（西欧霸权固化）形态存在，\literalquote{}control\literalquote{}未证。第五十章逐国重填K1，重算版图表。 \\
\textbf{最强竞争路径缺席}：史实至提尔西特＋两个开关从未被P6计分 &
\textbf{接受，本书作者亲自补跑P6引擎}（不改原卡文件，新脚本与输出存
\texttt{final/r\_5b1357a9c6-calc/}）。结果见第五十章。 \\
\textbf{P6与P7对1803--12英法战争状态矛盾}；M-B1借自战时窗口 &
\textbf{接受。}
S2本义世界中英法战争自1803持续（史实），B1b的1811--14窗口直接适用；路径B的\literalquote{}无战世界\literalquote{}需按1802--03检验另行定价。第五十二章。 \\
\textbf{七处\literalquote{}保留史实果实而移除原因\literalquote{}} & \textbf{接受。}
S2本义把1805--07的果实（莱茵邦联、巴符王号、普雷斯堡、提尔西特、华沙公国、残普、1810奥婚的动机）还给它们的原因；路径B的相关条目在第五十八章降为低置信并注明无锚点。 \\
\textbf{安特卫普}：1814巴黎和约第XV条未被任何卡使用；K3与③可能互斥 &
\textbf{接受，新增第五十二章第四节与第六十三章第四节。}
本书作者读了两个文本，给出两条终局带。 \\
\textbf{标题偷换}：\literalquote{}逻辑可构造\literalquote{}≠\literalquote{}现实可行\literalquote{} & \textbf{接受。}
结论表把控制版拆为条件／无条件两行，写入P6自认最可能的终局（英国不签约的武装共存）。第五十三章。 \\
F4的T=750锚于1810年6月疆域，路径B疆域更小 & \textbf{接受，作者重算。}
第五十四章第五节：T≈710--725（路径B／S2本义）；\literalquote{}陆300＋海195\literalquote{}在无盟约转移下不闭合，缺口约60--110m。 \\
IT1概率带以含威尼托的史实王国为前提；art.3分冠触发条件不满足 &
\textbf{接受。} 第六十五章按两条路径分别重估。 \\
完整普鲁士＋北德主持权未估；P7自相矛盾 & \textbf{接受。}
第六十六章按两条路径分别重估；北德主持权条目降为无锚点。 \\
无胜利帝国的合法性经济未入耐久模型 & \textbf{部分接受。}
第六十章第四节作为独立侵蚀项讨论，给方向与幅度区间，不改概率主表。 \\
\end{longtable}
}

局部修正（海军数字、白银两口径、Branda精度、MAN-01标签、独立束数、\literalquote{}普遍准入\literalquote{}非当时提案、K3手填）均已在相应章节改正并标明。

\subsection{五、读法}\label{ux4e94ux8bfbux6cd5}

本书第一部按行动者逐一铺开能力包络与反应函数；这是全书的证据基座，也是上一轮所缺的东西。第二部把它们放上同一张桌子对弈。第三部回答最重要的问题。第四部写1815--1848的秩序如何运作。第五部答六个分问。结论给主线、备选结论群、置信与改判证据。急于知道答案的读者可以从第五十章、第五十四章与结论开始读；但请记住，本书的一切判断都只在第一部的证据上成立。

\begin{center}\rule{0.5\linewidth}{0.5pt}\end{center}

\section{第一部　起点：1803年的世界与每一个行动者的能力包络}\label{ux7b2cux4e00ux90e8-ux8d77ux70b91803ux5e74ux7684ux4e16ux754cux4e0eux6bcfux4e00ux4e2aux884cux52a8ux8005ux7684ux80fdux529bux5305ux7edc}

\subsection{第一章　法国（一）决策中枢：拿破仑本人、大臣与秩序蓝图}\label{ux7b2cux4e00ux7ae0-ux6cd5ux56fdux4e00ux51b3ux7b56ux4e2dux67a2ux62ffux7834ux4ed1ux672cux4ebaux5927ux81e3ux4e0eux79e9ux5e8fux84ddux56fe}

\subsubsection{一、\literalquote{}拿破仑仍是拿破仑\literalquote{}是有结构的约束，不是均匀的顽固}\label{ux4e00ux62ffux7834ux4ed1ux4ecdux662fux62ffux7834ux4ed1ux662fux6709ux7ed3ux6784ux7684ux7ea6ux675fux4e0dux662fux5747ux5300ux7684ux987dux56fa}

任何反事实都要先回答一个问题：这个人会不会做那件事。本轮把1801--1814年间拿破仑的全部\literalquote{}克制\literalquote{}实例摆到一张表上，发现它们不是随机的。四个真实发生的克制------1801年与教廷的专约、1802年亚眠和约、1810年的许可证制度、1810年的奥地利联姻------共享三个特征：\textbf{工具性}（换取王朝、财政或包围收益）、\textbf{保留单边裁量权}（组织条款、停战化用法、亲签可废、不改备战）、\textbf{不承认对等约束}。凡克制意味着交出支配权或承认外部约束之处------1803年马耳他、1805年斯特拉斯堡备忘录、1808年西班牙、1810年荷兰与汉萨、1811--12年对俄、1813--14年三连拒------无一反例。

由此得出全书反复使用的两条先验：\textbf{非工具性克制0.05--0.25（中心≈0.15），工具性克制0.45--0.70（1809--10对奥为实证上界0.50）}；以及一条可操作的规则：\literalquote{}克制可以被购买，条件是让步被包装为王朝／财政／包围收益，且不要求承认外部约束。\literalquote{}表级规则更严：凡要求拿破仑\textbf{主动持续}非工具性克制的节点，单节点转移概率不应高于0.2；经工具性包装的同型节点可用0.45--0.70。

这张表最有证据等级的一格是B8（秩序建构）：1802--03与1810--11两个和平窗口的显示偏好一致------都选择继续扩张（皮埃蒙特、瑞士；荷兰、汉萨与预定的对俄战争）。主动自缚型的秩序建构先验只有0.05--0.15；被动型（灾难触发、如1813年2月的摄政法）一旦触发条件出现则约0.6。

【锚】1805年10月17日，塔列朗在斯特拉斯堡向皇帝呈交备忘录，主张不粉碎奥地利、以东方补偿换取奥法同盟、法国回归自然疆界、放弃意大利王冠；12月5日重申，未被采纳。\literalquote{}Votre
Majesté peut maintenant briser la monarchie autrichienne ou la relever.
Une fois brisée, il ne serait pas au pouvoir de Votre Majesté elle-même
d\textquotesingle en rassembler les débris
épars.\literalquote{}（Dard转录，本轮亲读）奥斯特里茨之后，拿破仑对他说：\literalquote{}我赢了一场大胜，你是一位大部长。\literalquote{}------克制被挖苦掉了。B2（对奥普宽严）的克制先验因此只有0.15（0.05--0.25）。

【锚】1810年春存在幅度最大的真实缓和窗口：莱茵境内撤军至两个师、加速对普财政了结、中止荷兰兼并、重启对英试探------\literalquote{}la
paix ne peut venir qu\textquotesingle en faisant d\textquotesingle abord
la guerre d\textquotesingle une façon moins
acerbe\literalquote{}。每一项在十二个月内被反向覆盖：七月并荷兰，十二月并汉萨与奥尔登堡。缓和是婚约红利的一次性支付，不是政策函数的改变。

【锚】1811年6月5日，圣克卢，科兰古向皇帝当面复述亚历山大的话：西班牙人屡败而未屈服；俄国有空间与气候；\literalquote{}If
the fighting went against me, I should retire to Kamchatka rather than
cede provinces and sign, in my capital, treaties that were really only
truces\literalquote{}；\literalquote{}奇迹只在皇帝亲临处发生，而他不能年复一年离开巴黎\literalquote{}。谈话持续五小时。拿破仑听后仍以列举兵力自醉：\literalquote{}one
good battle would knock the bottom out of my friend
Alexander\textquotesingle s fine
resolutions\literalquote{}。这不是情报缺失，是信念压制：真实信息完整送达，被\literalquote{}兵力列举\literalquote{}的心理机制压下。

【锚】1813年6月26日，德累斯顿，拿破仑对梅特涅（梅特涅1829年追述，回溯性，未与Fain／Maret版本对照）：\literalquote{}Vos
Souverains, nés sur le trône, peuvent se laisser battre vingt fois et
rentrer toujours dans leurs capitales; moi, je ne le puis pas, parce que
je suis un soldat
parvenu.\literalquote{}克制的约束不只是性格，也是其政权合法性经济的结构条件------一个靠自己爬上来的军人皇帝，其统治\literalquote{}在他不再强大、因而不再被畏惧之日即告终结\literalquote{}。这句话在第六十章会再次出现：它同时是\literalquote{}0.7/18个月冲动\literalquote{}的原因，也是一个无胜利帝国的合法性问题。

\subsubsection{二、中枢的容量：一个人能治理多大的欧洲}\label{ux4e8cux4e2dux67a2ux7684ux5bb9ux91cfux4e00ux4e2aux4ebaux80fdux6cbbux7406ux591aux5927ux7684ux6b27ux6d32}

秩序蓝图的文本群全部落锚：1806年fédératif训令原文、1808年致路易的法币统一信、1810年罗马王法令第7／10条、尚帕尼1810年12月8日报告（考出系拿破仑亲自改定的自画像）、1813年摄政法。蓝图的特征是：分层帝国、司法与货币统一主义、零横向协调、零自缚。

容量是可以数的。《通信集》通量：1806年2,685封→1811年折年峰值≈4,190封→1812年2,551封。环路时延：西欧8--12天，伊比利亚3--4周，俄境4--5周。1812年10月，莫斯科：巴黎来的信使14--15天一趟，\literalquote{}The
despatch-case for the Emperor and his headquarters arrived regularly
every day from Paris in fifteen --- often fourteen ---
days\literalquote{}；单节点治理半径两千八百公里，决策环路约一个月。1811年4月至12月的九个月里，皇帝写了3,144封信，其中陆军部长克拉克一人收到1,149封（36\%），军务占六成；而西班牙国王约瑟夫只收到\textbf{一封}。同年宗教、经济、备战三线危机下出现结构性放弃：半岛授权克拉克与贝尔蒂埃代管，凡尔赛工程停摆。1812年1月27日他自陈\literalquote{}这些准备用了一年\literalquote{}。

【推·中高】单人中枢在1810--12年已在极限运行。S6行省化所需的治理密度，估为现有容量的2.5--4倍。\textbf{没有分权改造的S6是自相矛盾的。}

分权有没有可能？行政性分权当时存在且可运转：副王欧仁、总督埃莉萨、达武、勒布伦、博尔盖塞、贝特朗。政治性分权零当时提案（西哀士1799年的方案被本人扼杀于起点），且与其成文意志正面冲突：1810年7月9日给勒布伦的训令，\literalquote{}Mon
intention est de gouverner moi-même le pays. Mon lieutenant général sera
là pour tout voir, prendre des renseignements,
m\textquotesingle instruire de tout, recevoir directement mes ordres
pour les faire
exécuter\literalquote{}；1811年11月27日致欧仁，\literalquote{}意大利王国的政策只由我一人决定\literalquote{}。唯一的自缚是1813年的摄政法，由灾难触发。因此S6在世期的主动分权路径无先例，只能按\literalquote{}奇迹\literalquote{}计价；摄政期（α路径）制度化可行性中等。

信息系统的失真不在采集层------警务日报、黑室、二十九个使馆、军事情报是欧洲一流------而在整合层：单节点整合无红队，激励失真（1808年西班牙：大使逢迎、密探盲目、缪拉的野心污染，且中枢反向注入假信息\literalquote{}俄国对瑞典宣战\literalquote{}以稳住缪拉），以及信念压制（前引五小时召对）。对弈中，法国的\literalquote{}情报失误\literalquote{}应建模为整合失真而非采集缺失。

\subsubsection{三、寿限}\label{ux4e09ux5bffux9650}

死因重建为散发型胃癌，至少T3N1M0（IIIA期），瘤体自贲门至幽门逾十厘米；Lugli等（2007）判断幽门螺杆菌路径比家族遗传更可能------\literalquote{}Napoleon\textquotesingle s
main risk factor might have been Helicobacter pylori infection rather
than a familial
predisposition\literalquote{}。父亲1785年死于幽门胃癌（尸检\literalquote{}冬梨\literalquote{}状硬化肿瘤）。家族长寿（母85岁、四兄弟65--76岁）托高非胃癌基线。综合三分支：中位死亡1826--1835年，活到罗马王1829年3月20日成年的概率≈0.50--0.65。1810--12年已有体能衰退的独立目击（de
Kock：\literalquote{}我期待见到一位神，只看到一个胖子\literalquote{}，转引Lentz
T11导言注5）；1815年后运营强度按0.6--0.8折减。

\subsubsection{四、给对弈的接口}\label{ux56dbux7ed9ux5bf9ux5f08ux7684ux63a5ux53e3}

\begin{itemize}
\tightlist
\item
  胜利和约签订后十八个月内，中枢启动新单边行动的概率≈0.7（对象按情景：西属美洲、奥斯曼、印度方向姿态）。\textbf{这是长久霸权最大的内生威胁}------只作拿破仑侧先验，不可移作他国参数。
\item
  1812年是S2的生死关：S2若活过1808，仍要在1811--12把\literalquote{}接受1810年俄国关税＋配额\literalquote{}包装成拿破仑的胜利（例如俄国公开重申封锁＋象征性联合舰队），才能压住开战；设计得当，避战概率可抬到0.4--0.5。
\item
  任何情景中首次重大军事挫败都触发\literalquote{}摄政法时刻\literalquote{}：自缚只在存亡压力下供给。
\item
  未决：罗德雷\literalquote{}权力靠光荣\literalquote{}引语的原卷未获；梅特涅德累斯顿版本未与Fain／Maret对照；《通信集》1807、1809--10年通量缺；沙普塔尔所记俄方许可证配额提议为回忆录孤证。
\end{itemize}

\begin{center}\rule{0.5\linewidth}{0.5pt}\end{center}

\subsection{第二章　法国（二）陆军：兵员、马匹、军械、后勤与占领军上限}\label{ux7b2cux4e8cux7ae0-ux6cd5ux56fdux4e8cux9646ux519bux5175ux5458ux9a6cux5339ux519bux68b0ux540eux52e4ux4e0eux5360ux9886ux519bux4e0aux9650}

\subsubsection{一、能养多少兵：包络}\label{ux4e00ux80fdux517bux591aux5c11ux5175ux5305ux7edc}

法国陆军的能力包络不是一个数字，是一组相互制约的账。\textbf{在册}（sous
les
drapeaux）40--55万，加上有自己军队的盟邦15--25万；在场可用率0.80--0.90，取0.85。可分配到外国驻防的上限25--35万，成熟宽案40--45万。年抽人上限（陆海合计）9.8--12.7万------陆军9--11.5万，海军0.8--1.2万，同一个人不双算。年轻男性年级约30万（Houdaille）；年征20--30万即吞掉多数年级；和平年可持续的年征≤8--10万。

【锚】征召的崩溃点是实测过的。皮加尔（Alain Pigeard，《La conscription
sous le Premier Empire》，Fondation Napoléon网页，原刊RSN
420，1998）：\literalquote{}Du 1er septembre 1812 au 20 novembre 1813, 1 527 000
hommes avaient été appelés en quinze
mois.\literalquote{}1813年五道元老院令合计约一百一十四万；帝国全期表列征召总额2,432,335。这十五个月的分项可复算之和为1,317k，皮加尔聚合的1,527k含未明细的210k。

【锚】1813年10月9日的一批：敕令征127,433人（1808--1814年级），\literalquote{}yet only
72.265 men reached their units by the end of
December\literalquote{}------到营率56.71\%（Grab 2013印p.118转Leggiere
2007印p.69）。这只是一个分批截止日的证据，不是全国全年到营率；全国1803--1814年逐年到营、死亡、退出的连续序列本轮未取得。

\subsubsection{二、姿态与余额：筛查地图}\label{ux4e8cux59ffux6001ux4e0eux4f59ux989dux7b5bux67e5ux5730ux56fe}

以在册450k＋盟军200k为中心，在场率0.85，给出各姿态的中心余额（千人在场）：

{\def\LTcaptype{none} % do not increment counter
\begin{longtable}[]{@{}lll@{}}
\toprule\noalign{}
姿态 & 需求 & 余额 \\
\midrule\noalign{}
\endhead
\bottomrule\noalign{}
\endlastfoot
S2（大陆整固，对英武装共存） & 540 & \textbf{+55} \\
S5（和平秩序） & 490 & \textbf{+147.5} \\
S6-lite（并外圈） & 566.25 & +28.75 \\
S6-max（并至南德、意大利、西班牙北部） & 922.25 & \textbf{−122.25} \\
S6-full & 1,102.25 & \textbf{−302.25} \\
\end{longtable}
}

西班牙的增量（22.5万）占S6-max缺口的63\%。这些数值是\literalquote{}原范围筛查地图\literalquote{}，不是最终地图：最终的行省化范围要读P2的逐地区联合账（第五十四章），它做了地区去重、盟军退出、税余与陆海财政的联算。

财政联算的约束更紧。按F4的陆军线300m、700法郎／兵年、独立盟军20万、在场率85\%，可用在场兵53.43万：S2中心缺5,714，S5余44,286，S6-lite缺31,964。\textbf{人力可行不等于财政闭合。}
首阶段的主线因此是：法国40--45万在册＋可靠自养盟军20--25万，并延缓造舰。

\subsubsection{三、新省供兵：第十年才转正，而且取决于役期}\label{ux4e09ux65b0ux7701ux4f9bux5175ux7b2cux5341ux5e74ux624dux8f6cux6b63ux800cux4e14ux53d6ux51b3ux4e8eux5f79ux671f}

\literalquote{}多吞＝多兵\literalquote{}是上一轮从未检验的假定。本轮的整合模型给出：占领成本≤4‰（人口）的新省可以净供兵，15--25‰的高安全负担区是黑洞；断裂窗口------1810年启动的并合在1812--14年断，1815年启动的在1817--20年断。新省的净供兵在\textbf{第十年}才转正------而这只在\literalquote{}hazard6\literalquote{}（常退役率）口径下成立。二轮复核给出对照：同样指定征额2‰、受训转化90\%、可用85\%、可靠60\%、安全需求4‰，原口径成熟净供兵+0.790‰；若改为\literalquote{}总役期五年含一年训练、四批受训兵、训练后年留存97.5\%\literalquote{}，则净−0.463‰，需指定征额2.262‰或可靠率67.9\%才转零。这是制度敏感性，不是新识别的历史临界率；不得对全版图机械扣减。

【锚】Rowe《From Reich to
State》表3（印p.179，本轮看图核验）：莱茵四省1806--10，\literalquote{}Mobilised
24,186; Draft dodgers 8,138 (34\%); Deserters 725
(3\%)\literalquote{}。征前逃避33.65\%与入伍后逃亡3.00\%并存------分母与行为阶段不同。莱茵的征前逃避永未趋同阿尔萨斯（17.12\%），收敛的只是入伍后服从与照额交付（第四十三章）。

\subsubsection{四、马、军械、后勤}\label{ux56dbux9a6cux519bux68b0ux540eux52e4}

马不是同质存货。一次军马损失13--17.5万（1812年俄征量级，非年折旧）：次年仍能采购，但新骑手、装具、运输、马龄的恢复不同步，至少跨过下一个战役季。1813年全国马群350万，其中83.8万未满四岁（E1）。分兵种军马在役数与军官年龄分布仍缺。

【锚】军械（Brun，《Des armes et de la
poudre》）：1812年4月枪械库存811,000；至1813年7月新造159,000、库出880,000（非全战损）；常态月产8,000--10,000，1813年肩射武器月产难过18,000；火药约3,000吨／年（作者认为1813年仍够）。1794年法国硝石峰值817万公斤是刮墙含硝土的应急，不是人工硝床的稳态；1828年约50万公斤（E1）。

后勤没有普适半径。范·克里维尔德《Supplying
War》第二章的著名算例------\literalquote{}即使以原六十万的三分之一抵达俄国首都、用六十天（实际八十二天）\ldots\ldots\literalquote{}二十万人、六十天、日耗三百吨、距基地六百英里、运输纵队日行二十英里------是作者的\literalquote{}Even
if\literalquote{}反事实算例，不是莫斯科实发军粮；原段未指明吨制，本书只把1.5公斤／人日作公吨近似的规划输入（若为长吨则≈1.524公斤）。粮秣说明案：五十万人＋十二万五千马＝2,000公吨／日（E1按十万马为1,750吨，差额由2.5万马×10公斤解释）------均为模型说明案，非全帝国军队，也不意味着全部从法国本土长运。

\subsubsection{五、占领密度的类比锚}\label{ux4e94ux5360ux9886ux5bc6ux5ea6ux7684ux7c7bux6bd4ux951a}

奥地利驻伦巴第-威尼西亚1848年约七万地区军团（Hilleprandt，转引），当地供三万至三万五千意大利兵且1848年过半仍忠于奥（Sondhaus）；法属阿尔及利亚1846／47年约十万驻军，1848--70年多数年份六万五千（Frémeaux）。3‰治安＋1‰边防只是低压纯治安的测试参数。西班牙两个切片冻结为占领态的实测：1810年1月15日境内324,996／在伍262,051（Oman
III附录VII）；1811年7月15日在册354,461／在伍286,414--291,414（Oman
IV附录XVIII）。占领西班牙吞掉29--36万------这是SP1\literalquote{}军事价值二值函数\literalquote{}的另一半。

\subsubsection{六、外籍兵员的三个口径不得混合}\label{ux516dux5916ux7c4dux5175ux5458ux7684ux4e09ux4e2aux53e3ux5f84ux4e0dux5f97ux6df7ux5408}

Grab
p.26：1812年俄征六十万中过半非法籍------这是战役国际军队的比例，不是全国陆军的比例。Elting第三章：1807年普波战区约324,000，可能三分之一为莱茵邦联与波军，含医院两万，不含失踪擅离约两万五千。Houdaille（印p.49）：111,000莱茵＋105,000弗拉芒＋69,000瓦隆＋104,000意大利＝389,000，属累计服役且按1804年前并合地区来源校正，与1,915,000名1815年疆界法国人相加后还须扣共和十二年以前的服役------不是1812年非法籍总数，也不能加到某一天的战区兵额上。

\subsubsection{七、二轮变动与遗留缺口}\label{ux4e03ux4e8cux8f6eux53d8ux52a8ux4e0eux9057ux7559ux7f3aux53e3}

F2在重开时做了合同内的有界复核：主线不变（法国核心军＋可靠自养盟军＋有限并合），接口数值不变；新增Rowe表3六行、1813年一批到营证据、五行役期口径桥、两行粮秣范围桥；四生成器与验证器复跑通过。通过只证明结构、算术与引文字面存在，不是历史或因果认证。遗留：全国逐年到营／退出／净存量；五厂年产；军官年龄与分兵种军马；地方可靠率与省级净财政实收。

\begin{center}\rule{0.5\linewidth}{0.5pt}\end{center}

\subsection{第三章　法国（三）海军与海事经济：船厂、木材、水手、军官团}\label{ux7b2cux4e09ux7ae0-ux6cd5ux56fdux4e09ux6d77ux519bux4e0eux6d77ux4e8bux7ecfux6d4eux8239ux5382ux6728ux6750ux6c34ux624bux519bux5b98ux56e2}

\subsubsection{一、舰体不是舰队}\label{ux4e00ux8230ux4f53ux4e0dux662fux8230ux961f}

【锚】托多罗夫（Todorov）按帝国控制口径给出的战列舰体：1810年夏50艘（旧法国港29、安特卫普12、荷兰9）；1812年72艘（旧法国港41、安特卫普20、荷兰10、热那亚-威尼斯-科孚1）；1814年4月74艘（旧法国港42、安特卫普21、荷兰9、热那亚／威尼斯／科孚2）。\literalquote{}Été
1810 ... 50; 1812 ... 72; Avril 1814 ... 74\literalquote{}（PDF
p.15）。这些是舰体，不是配员可战的舰队；\literalquote{}104 vaisseaux à flot ou en
chantier\literalquote{}含在台。

【锚】1810年7月13日的命令提出从小艇到战列舰的技能转移，规划\literalquote{}un corps de
66,000 matelots ... 50 vaisseaux de ligne et 30
frégates\literalquote{}及400艇------可作非战列舰人力预留的量级，不能当实到人数。1825年法国海事登记94,611人，其中73,053为军士、水手、见习水手及少年------是登记职业池，不是海军在役或全熟练池。

【锚】1811年2月27日海军部报告：此前不进口外国铁，铜已有储备，\literalquote{}les
produits de la France et de l\textquotesingle Italie suffisent aux
besoins du
service\literalquote{}（大麻）。这是当期部门报告，不是永久无约束的现代审计。

\subsubsection{二、木材的约束是地理，不是储量}\label{ux4e8cux6728ux6750ux7684ux7ea6ux675fux662fux5730ux7406ux4e0dux662fux50a8ux91cf}

1805年的橡木普查显示安特卫普补给流域集中了约62\%的合规格大橡木：\literalquote{}Le seul
bassin d\textquotesingle Anvers concentre presque 62 \% de
l\textquotesingle ensemble des arbres de chêne d\textquotesingle au
moins 5 pieds de
tour\literalquote{}。法国森林绝对量并不短缺------E1核得1810--12年七组指定船厂逐年下水7／9／8艘，推翻\literalquote{}年六舰物理不可能\literalquote{}；733.8万株大橡树也不等于熟材库存。主约束是\textbf{通道与船厂的地理错位}：安特卫普的浅水航道、分港训练、人员征募、财政拨付。木材以\literalquote{}年造所需到厂合格材流量\literalquote{}表示（24格需求表），不虚构立方米运量；丧失安特卫普通道的敏感性另列。1811年3月5日的命令只限安特卫普：\literalquote{}18船台、舰留三年\literalquote{}；10月2日因可再分配荷兰与布洛涅舰员而拟\literalquote{}construire
chaque année huit vaisseaux au lieu de
six\literalquote{}------不能与全国满负荷\literalquote{}二十艘\literalquote{}混比。指定七组船厂1807--13年至少下水50个不同战列舰体（逐舰表52行，改名与俘获转籍已去重；多经Winfield--Roberts／Roche转引，原书未阅）；年均下水≥7.14。托多罗夫所称年二十艘未从最小逐舰名单复现，本书不以二十／年为长期生产率。

蒸汽不是封锁下无代价的捷径：1829年的Sphinx号，\literalquote{}sa machine à vapeur fut
construite par une entreprise de Liverpool, W. Fawcett\literalquote{}。

\subsubsection{三、1830年的三档}\label{ux4e091830ux5e74ux7684ux4e09ux6863}

以持续海战（W）与1815年后海上停战（P）两种态，1830年十二个月内可修复武装配员的战列舰（A12）：主线84（W）／111（P），包络53--124／84--140；常备在役I主73／61。可拨总员10.4万／12.9万，合格海员6.5万／9.4万，年新增可拨人员0.8--1.2万（部分可来自年长商渔民）。敏感性：新下水减30\%→60／83舰；总员与技能流入减30\%→57／69；放弃安特卫普（起始少12舰、此后年少3舰）→50／71。财政减30\%不能直接等同流量减30\%。1811--15年下水的年轻舰群退役率1.5--3.5\%（Todorov
p.15年龄近似）。

财政联算：海军170--220m／年＝主线档；1830年高档（A12
121--140、总员13.5--15.2万）需220--260m／年，\textbf{不在F4任何平衡案内}；钱与舰不可线性折算。

\subsubsection{四、结论}\label{ux56dbux7ed3ux8bba}

法国有成为强大第二海权的路径，没有靠兼并沿海＋多造船自动取代皇家海军的路径。\textbf{最大行省化≠最大有效海权}：资源控制的收益必须扣掉陆军驻防、商航萎缩与征募抵抗。最优是保持安特卫普-荷兰合作、大西洋与地中海基地、停止吞噬海军资源的大陆扩张，以及以护商航行做长期训练。法国能否在外交上取得海上停战，由B1／P6裁定，本章不把它当已赢得。未决：1814年帝国afloat／chantier／armé口径与荷意包含关系；Glover
1967与Winfield--Roberts原书；材料表给需求非实测运量；1848年蒸汽舰队未数值拟合。

\begin{center}\rule{0.5\linewidth}{0.5pt}\end{center}

\subsection{第四章　法国（四）财政、银行与统治阶级：贡赋帝国能否变成税收国家}\label{ux7b2cux56dbux7ae0-ux6cd5ux56fdux56dbux8d22ux653fux94f6ux884cux4e0eux7edfux6cbbux9636ux7ea7ux8d21ux8d4bux5e1dux56fdux80fdux5426ux53d8ux6210ux7a0eux6536ux56fdux5bb6}

\subsubsection{一、法国是税收国家加战争特别账户，不是贡赋空壳}\label{ux4e00ux6cd5ux56fdux662fux7a0eux6536ux56fdux5bb6ux52a0ux6218ux4e89ux7279ux522bux8d26ux6237ux4e0dux662fux8d21ux8d4bux7a7aux58f3}

净经常收入以750m为中心（F4自给宽下行650、高收入案800）。S5过渡期的陆海联合预算按恒等式：750−300（非军与既有息养）−25（新增补偿）＝425（陆海合计）；海军取195则陆军余230。\textbf{230是陆军的余量，不是\literalquote{}陆300已闭合\literalquote{}}：F2按700法郎／兵年计，50万法军需350m，缺口120m／年；固定350m陆军时中心案只剩75m海军。要海军170--220m，需可靠的额外经常收入或净转移95--145m，或等额减支。\literalquote{}陆300＋海195闭合\literalquote{}是\textbf{高收入＋可兑现盟约}的另一组输入（495m），不能抹掉中心案的120m缺口。

【锚】布兰达（Pierre
Branda，附表，本轮据缓存页逐格开采；缓存页尾的\literalquote{}Discuss the
implications\ldots\literalquote{}问答段是平台生成摘要，不是布兰达原文）：共和十一年（1803）至1814年的非常融资来源全表4,284m内部自洽------国库小计1,272＋224＋137＋62＋69＋95＝1,859；战争小计809＋154＋453＋383＝1,799；总计1,859＋1,799＋503＋123＝4,284；所列百分比29.69／43.39／41.99／11.74／2.87\%全部复算相符。\literalquote{}Total
paid by the French Treasury 1,859 43.39\% ... Total financing through
the war 1,799 41.99\% ... Grand total 4,284 100.00\%\literalquote{}。

除以约11.5年：战争\textbf{附加}成本年均372.5m＝外部156.4＋本土国库161.7＋递延赤字43.7＋银行10.7（复算相符）。外部156.4含避免成本33.3（非现金，不得进入依赖率分子）；现金类123.1＝（普通贡赋809＋Domaine对国库贷款154＋条约外部收入453）÷11.5。\textbf{贡赋现金依赖度＝123.1／750＝16.4\%}；剔除西班牙就地消耗的350m后，1,066÷11.5＝92.7，92.7／750＝\textbf{12.4\%}。和平释放的本土资源＝161.7＋43.7＋10.7＝216.1m／年，是所失外部现金的1.75--2.3倍。

【推·高（狭义会计命题）】\literalquote{}贡赋帝国一旦停止征服就财政崩溃\literalquote{}是一个会计错误：它把融资腿误认为净资源。贡赋所融资的是战争的附加成本；和平取消该成本时，四条腿一同消失。\textbf{前提}：和平建制（常备陆军＋海军）必须装进经常收入------即仍完全受F2／F3联算约束。禁止读作\literalquote{}任何规模的和平军备都能自养\literalquote{}。政治命题（军队、承包商、渴望晋升的军官团是否接受和平）置信维持中。

【锚】\literalquote{}Spain 1808-1813 350 350
0\literalquote{}：西班牙1808--13年的普通贡赋350m全由军队就地取用，付予法国国库为0。红队指出该行只有三个数，第四格（Domaine）空白或丢失，\literalquote{}两列皆0（原表明确格）\literalquote{}过度精确；实质结论不变------\textbf{占领地榨取额不等于中央可支配资源}。这是本书对S6\literalquote{}并入即得税\literalquote{}最强的单一史实反证。

【锚】1814--18年战胜国对法国的反向榨取：\literalquote{}Total paid by France
1,412\literalquote{}＝条约前普通贡赋144＋一次性赔款700＋1815-11至1818-10占领军费328＋对帝国军队所夺之补偿240。在法国不战败的S5／S6中，这是不需任何假设的净避免额。

【锚】布兰达正文：\literalquote{}The cost had in fact been estimated by Mollien at 700
francs a year, i.e., 1.91 francs a day.\literalquote{}注15出处Simon de la
Rupelle，Revue
Politique，1892，p.662。但作者自述只计战役活动日、不计驻营日，故1.91法郎是边际战役成本，不是全年维持费。附表A盟军兵日节省253,748,180法郎＋附表B盟国支付驻屯费129.1m＝382.85≈383；正文同段先写\literalquote{}over
380 million\literalquote{}后写\literalquote{}352 million francs in
savings\literalquote{}，352无法复现------本书采383。附表B的1806--09各国逐年驻屯值因OCR丢格不可恢复（合计行仅九数、累加98.43≠129.1），仅共和十一、十二、十三年三列经交叉校验可用。正文与附表的276／278、454／453、80／62、232／224、59／69差异一律保留不调和。

\subsubsection{二、S6的分地区账：举证责任在扩张者}\label{ux4e8cs6ux7684ux5206ux5730ux533aux8d26ux4e3eux8bc1ux8d23ux4efbux5728ux6269ux5f20ux8005}

【锚】Marion（Histoire financière de la France，IV，1925，p.321，引AF IV
1072，本轮页图核）：1812年并合地的\textbf{预期}毛税342.260044m／净226.389345m------比利时83、荷兰66.5、汉萨38.791、莱茵37.5、皮埃蒙特33、托斯卡纳22.5、利古里亚16、罗马16.5。这是预算指令的预期，不是实收。Gaudin
1811年4月30日：\literalquote{}至多105\literalquote{}。1810年罗马预算指令每百万人≈14m主权费（目标非实收）；荷兰≈60m门槛。S6阈值表中仅两个独立可核格（荷兰、罗马）皆净耗。

由此给下游的规划参数：S6每新增五百万人口三档+55／+10／−55m（每百万人+11／+2／−11），每百万人继承100m债务×6\%＝−6m／年（规划假设，非观测借款利率），S3税基损伤−60至−110m／年；S3与S6在海岸带相互抵消。1810年的\literalquote{}非常关税\literalquote{}61.046m由荷兰存货50\%税、提税、许可证与罚没构成；严格S3下总关税≈97m→15--35m。Droits
réunis
1807年85m→1812年144m；1818年起毛／净断点。赠产（Domaine）年收益15.165m（1809-08）→18m＋（1810-01）→24.73m（1814-07重建），仅经常收入的2--3.5\%、战争附加成本的4--6.6\%------不是不可补的财政支柱。

\textbf{财政终局}：一个\literalquote{}大而有限的直接核心＋法典化藩属＋条约贡献国\literalquote{}，每批200--500万人，3--5年整合。

\subsubsection{三、金融机器的自伤（预告第三十八章）}\label{ux4e09ux91d1ux878dux673aux5668ux7684ux81eaux4f24ux9884ux544aux7b2cux4e09ux5341ux516bux7ae0}

法国在阿姆斯特丹与汉堡一笔未借，不是因为市场拒绝，而是因为它选择了没收：热那亚1805年（1.44亿里拉→1,850万法郎、1.5\%）、荷兰1810年tiercering（78→26m年金）、汉堡1813年（4,800万法郎罚金并搬空银行7,506,956
Mark
Banco）。决定性反例：1816年拉菲特断言法国在欧洲任何地方借不到一个小埃居，次年一个战败被占领的法国以发行价55募到3.15亿法郎。金融技术可用；在拿破仑执政下被采纳的概率约0.15--0.30。

\subsubsection{四、未闭合}\label{ux56dbux672aux95edux5408}

布兰达《Le prix de la
gloire》（2007）逐年表、Gaudin回忆录附录、各年《Compte
général》未取------它们能把年均拆成逐年序列，检验峰值年（1809年56.77m、1812年60.07m对1805年8.82m）的依赖度是否曾突破可吸收量级。11.5年为本轮选定的期间长度，取11或12年各值变动约±4\%。Domaine
Extraordinaire清册（萨克森贷款7.355m、威斯特法利亚1m、意大利股份3.97m）源出private
collection，无从核对。

\begin{center}\rule{0.5\linewidth}{0.5pt}\end{center}

\subsection{第五章　法国（五）社会与政权存续：征兵容忍、教会、保王派、舆论与继承}\label{ux7b2cux4e94ux7ae0-ux6cd5ux56fdux4e94ux793eux4f1aux4e0eux653fux6743ux5b58ux7eedux5f81ux5175ux5bb9ux5fcdux6559ux4f1aux4fddux738bux6d3eux8206ux8bbaux4e0eux7ee7ux627f}

\subsubsection{一、支持是条件机器}\label{ux4e00ux652fux6301ux662fux6761ux4ef6ux673aux5668}

法国社会对政权的支持不是忠诚，是一台条件机器：财产、和平、低征额、教俗均衡、面包。五个条件都满足时，它是帝国最稳的地基；缺一个，它开始计价。

【锚】征兵抵制率序列：réfractaires均值共和九年至十三年约27\%，1806--10年降至6\%，1813年起回升2--10\%；1810年官方在册逃兵与拒征存量164,770；1800年大赦申请者17.5万；1813年1月政府承认残余约五万。福雷斯特（Forrest，pp.91--92）：\literalquote{}villages
that had openly defied the government\textquotesingle s commands in 1793
or Year VII were by 1810 and 1811 filling their military quotas with
only the minimum of
protest\literalquote{}------动机是恐惧而非爱国；1810--11是例外好年份（经济危机把穷人推向军队），1813--14合规崩溃。

【锚】拿破仑1806年自认教士是当年征兵顺利的主因：\literalquote{}Les prêtres catholiques
sont d\textquotesingle un grand secours ; ils ont été cause que la
conscription de cette année a été beaucoup mieux que celle des années
précédentes\literalquote{}（转引Melchior-Bonnet，Napoléon et le
Pape，p.87）。教会合作是征兵合规的放大器；1809--13年的升级链条全部由罗马兼并这一政策选择触发（Hicks逐环）。若维持专约均衡，教会回到放大器的位置；若教皇被囚持续，继承时刻现成正统性否决者，各情景存续概率−5至10个百分点。

【推·中高】和平常备军的社会容忍上限：年征≤8--10万（占人口≤0.3\%）＋和平死亡率＋役期兑现＋教会合作→常备军55--70万（含外籍）可无限期维持；年征\textgreater20万或追征已豁免年级→12--18个月内合规崩溃（1813年实测）。锚点：1810--11合规峰、1813崩溃、1818年古维翁-圣西尔法四万／年平稳重启、亚眠后三十万常备军先例。这是裁量的阈值，不是硬阈值。

\subsubsection{二、精英倒戈的解剖}\label{ux4e8cux7cbeux82f1ux5012ux6208ux7684ux89e3ux5256}

【锚】1814年4月4日，内伊、勒费弗尔、麦克唐纳、蒙塞逼宫；4--5日马尔蒙第六军倒戈：后者使有条件退位（传子＋摄政）的谈判筹码消失------\literalquote{}Il
n\textquotesingle y aurait ni régence, ni succession
impériale\literalquote{}。元老院废位：塔列朗非法召集，在册九十名元老六十四人到会，六十六票通过废位；4月6日宪法第6条元老自保俸禄与sénatoreries（\literalquote{}les
sénateurs actuels\ldots sont maintenus\ldots La dotation actuelle du
Sénat et des sénatoreries leur
appartient\literalquote{}）；路易十八后给三万六千法郎年金，八十四人入贵族院。波尔多，1814年3月12日：市长林奇串通王党秘密网络，在贝雷斯福德仅六千人（部分葡军）抵近时撕勋章升白旗，国民卫队无首未抗；苏尔特判词\literalquote{}使全体人口蒙羞\literalquote{}，威灵顿3月29日公开切割：\literalquote{}C\textquotesingle est
contre mon avis et ma manière de voir\ldots que certaines personnes de
la ville de Bordeaux ont jugé convenable de proclamer Louis
XVIII\literalquote{}。精英倒戈＝少数密谋＋外军掩护＋民众被动，不是群众性改宗。

【推·中高】1814年精英倒戈是五个必要条件的合取：外军占首都、可信替代王朝、军事灾难、皇帝不在场、中枢家族撤离。S5和平继承时刻前三条不成立。falsifier（\literalquote{}1813--14文书显示精英已备另立秩序\literalquote{}）方向性成立但适用域受限：它证明显贵忠诚没有王朝存量、完全条件化，而不是任何继承都必倒戈。凡情景叠加对外战败＋外军逼近，按1814模式重判（存续概率×0.3--0.5）。

【锚】马莱1812年伪元老院令纲领（SHD dossier Malet转录）：\literalquote{}de traiter
immédiatement de la paix avec les Puissances belligérantes, de faire
cesser les malheurs de
l\textquotesingle Espagne\literalquote{}------废帝制、立即媾和、终结西班牙灾难、恢复荷兰与意大利独立、对逃兵大赦、迎教皇入巴黎、恢复出版自由。这是1812年反对派对社会最大公约数的暴露性文件，与1813年莱内报告、1814年迎阿图瓦的呼声三点重合：和平、减征、教俗和解。

【锚】公投：共和八年官方3,011,007票系吕西安最后时刻伪造（真实赞成少于1793年约三十五万，军队未投票）；1815年附加法仅1,552,942赞成／5,740反对，参与率约20\%（共和十年约45\%），军队票被官方膨胀至四十五万。正统性生产的上限由此可见。

保王派能力包络实测：无外军入境则非生存威胁。1804年卡杜达尔网络需英国运载资助；1815年旺代第四次起事名义约五万，被约两万正规军六周击溃（Lignereux：\literalquote{}远非民众保王主义\literalquote{}），但在滑铁卢当周钉住两万以上帝国机动兵力------其真实功能是外军入境时的\literalquote{}接收政权供应商\literalquote{}与机会成本制造者。

\subsubsection{三、寿限与继承}\label{ux4e09ux5bffux9650ux4e0eux7ee7ux627f}

【锚】INED-Bonneuil
1806年法国生命表：1803年34岁男性活到1821年的概率79.5\%，中位死亡年龄67.4岁（≈1837）；分年代：1820s死15.4\%、1830s死23.7\%、1840s死27.2\%。拿破仑实际死于51岁，属最差五分位，与胃癌病理一致。叠加病理三分支后裁量：皇帝1820s死40--50\%，1830s死30--35\%，1840s以后15--25\%。

【推·中】继承情景概率带（S5前提）：皇帝1820s死（幼主摄政）------交接平稳65--80\%，存续至1848
30--45\%；1830s死（成年新君）------交接80--90\%，存续45--60\%；1840s死------交接≥85\%，存续55--70\%；罗马王早逝旁支（史实1832年死于结核）使存续概率×0.6--0.75。加权后王朝无条件存续至1848≈40--55\%；若计入\literalquote{}行政国家存续而王朝更替\literalquote{}的形态，拿破仑制度秩序延续至1848的概率≥70\%。关键机制：显贵的制度忠诚不等于王朝忠诚（Petiteau：91\%的国家任职跨政权延续）；\textbf{摄政期必然开启\literalquote{}宪政化让渡市场\literalquote{}}（莱内清单→宪章→附加法，三代同款）。摄政的最优设计＝1813年2月5日法（母后名义＋不可撤免的摄政会）＋遗嘱指定联合班子；最大制度空白是1829年成年时摄政会解散程序在1813年法中没有规定。补偿池15--30m／年可覆盖显贵；sénatoreries三十六处；融资方式为独立轴。

未决：Duvillard原表未亲读（以Bonneuil重建替代，换口径只会使1820s死亡权重更高）；《Bulletin
des
lois》1813年第474／490号摄政法原刊未亲验；Woloch／Bergeron／Tulard／Bluche经书评转引；卡杜达尔行刑日（6-25
vs 6-28）、莱内表决数（223:31 vs
223:51）、三次\literalquote{}八万令\literalquote{}日期、公投反对票（5,657／5,740）四处口径冲突备案；1805--11年档案中显贵层是否存在\literalquote{}后拿破仑预案\literalquote{}未能检验------这是最重要的改判缺口。

\begin{center}\rule{0.5\linewidth}{0.5pt}\end{center}

\subsection{第六章　不列颠（一）政治：内阁、议会、王室、和平党与登陆后的决策}\label{ux7b2cux516dux7ae0-ux4e0dux5217ux98a0ux4e00ux653fux6cbbux5185ux9601ux8baeux4f1aux738bux5ba4ux548cux5e73ux515aux4e0eux767bux9646ux540eux7684ux51b3ux7b56}

\subsubsection{一、和平派多数与可持续执政不是同一门槛}\label{ux4e00ux548cux5e73ux6d3eux591aux6570ux4e0eux53efux6301ux7eedux6267ux653fux4e0dux662fux540cux4e00ux95e8ux69db}

【锚】1806年\literalquote{}全才内阁\literalquote{}（Talents）取得选举多数，仍因国王要求对天主教问题作终身不再提出的保证而倒阁：\literalquote{}When
the king demanded a pledge that the ministers, some of whom wished to
make further concessions, would not raise the issue again in his
lifetime, Grenville and his pro-Catholic colleagues
refused.\literalquote{}------王冠认可是稳固和平的必要一环。在位的保守政府也可以媾和；和平网络不自动转化为执政多数。稳固的和平需要王冠认可、盟友处置、海军安全与真实商业开放的组合，不需要惠特布雷德领一个纪律严密的\literalquote{}和平党\literalquote{}赢得大选。

【锚】1812年6月23日枢密令只对美国船舶及美国财产货物撤销指定命令，并保留合理通知后恢复权：\literalquote{}so
far as may regard American vessels, and their cargoes being American
property, from the 1st of August
next\literalquote{}------不等于英国全面放弃海权或对法战争。（Historical
Hansard网页以1812-04-24为标题却插载该令；本书按文书内日期引用。）

\subsubsection{二、英国能接受什么：六条件与两个窗口}\label{ux4e8cux82f1ux56fdux80fdux63a5ux53d7ux4ec0ux4e48ux516dux6761ux4ef6ux4e0eux4e24ux4e2aux7a97ux53e3}

【锚】亚眠条约第XIII--XIV条确有财产保护：\literalquote{}All the sequestrations laid on
either side on funds, revenues, and credits, of what nature soever they
may be, belonging to any of the contracting powers, or to their citizens
or subjects, shall be taken off immediately after the signature of this
definitive
treaty.\literalquote{}但这不等于法国保证战败英国的全部公债偿付、议会不受干预或不夺舰。1806年12月30日雅茅斯勋爵叙述塔列朗口头的\literalquote{}实际占有＋等价交换\literalquote{}基础；同场承认\literalquote{}there
were certainly no papers that could be produced in which the uti
possidetis was recognized formally by the French
government\literalquote{}。两者共同约束反事实：完整的保证包从未在桌上。

【推·中】若从1807年起法国不废西班牙王、不迫使葡萄牙沦为战场、不进攻俄国，而且英国连续看不到可投入战争的大陆伙伴，\literalquote{}1809--12年先出现削减无效贸易战与试探议和，1811--14年出现能够签署妥协和平的政府多数\literalquote{}。最可能的载体是保留王冠信任的在位政府调整政策，或吸收辉格、坎宁等中间人物的宽基内阁。指认某年、某位首相只能低置信。更早的1806--07不是虚构窗口：全才内阁已实际谈判；法国如果把领土交易写定、不反复加价，可将和平提前到这一窗口。反过来，法国继续吞并、索要英国交舰或废王，即使英国没有盟友，1812--16年可能换成应急内阁却仍不接受法国的最高要求。

【修正机制·中高】无盟友既降低陆战收益，也可能让英国把财力集中于海军，从而改善其拒绝坏和平的备选。1811--14和平窗\textbf{不能由\literalquote{}无大陆盟友\literalquote{}单独推出必然媾和}，必须依赖安全与贸易条件。1815年6月14日财政大臣万西塔特公开提出：法国若与大陆列强达成安排，\literalquote{}that
expense now incurred for our armies would cease, and the supplies at
present demanded for them could be applied to the service of our
navy\literalquote{}；同场蒂尔尼质疑欠款不会重生------\literalquote{}there seemno reason to doubt
that similar arrears would again
accumulate\literalquote{}。陆转海是当时桌面选项，不证明它已执行。

六条件的本质：承认法国的大陆第一权，不承认对英处置权，现在不组盟但不永远放弃。\literalquote{}吞并全大陆＋永久排英＋英国自愿永久中立\literalquote{}三者不能同时成立。1812年对西班牙约瑟夫王朝的接受障碍足以否定\literalquote{}只保葡萄牙和西西里就能和平\literalquote{}的低底价模型，但未证英国要求法国全面退至1792年边界。\textbf{红队指出六条件遗漏了安特卫普；第五十二章补论。}

\subsubsection{三、伦敦失守之后}\label{ux4e09ux4f26ux6566ux5931ux5b88ux4e4bux540e}

政府内迁与媾和并不互斥。在保存合法授权后进行条件议价，比伦敦失守即降或永久续战更符合证据包络。若法军持续施压、英方无迅速恢复路径、法国给有限条款（保王朝、议会、公债、主体舰队），1--3个月为政治议价的中心窗；若可信支付与基地仍在，3--6个月的首战季续战与之并列。海军驻区继承未执行的命令与非泰晤士基地------\literalquote{}the
new commander inherited the standing or unexecuted orders of his
predecessor\literalquote{}------支持短期失联仍可组织行动，但依赖合法指挥、基地交通、供应商信任与支付，不能推出首都失守后无限续战。撤粮令：福蒂斯丘所述1803年10月31日某版方案因摩尔与里士满公爵反对而被放弃；但奇尔科特（Chilcott
2006，印pp.242--243）引1803年11月19日索尔兹伯里郡会议命令，仍安排居民与牲畜的撤移路线------\literalquote{}the
greater part of this county is highly favourable for the removal, not
only of inhabitants, but of large flocks of sheep requiring space; they
must therefore take their route over the
downs\literalquote{}。计划存在不等于执行成功。

\subsubsection{四、1815--1848：蓝水转向而非退出欧洲}\label{ux56db18151848ux84ddux6c34ux8f6cux5411ux800cux975eux9000ux51faux6b27ux6d32}

最大行省化可以得到英国短期默认，但若同时永久排除贸易、集中沿海战争能力并威胁英国主权，稳定跨继承的英国中立较不可能；重新联盟须同时具备威胁、可战伙伴与国内供款。【锚】1823年坎宁--波利尼亚克备忘录把外国武力介入西属美洲视为新问题，\literalquote{}it
would consider any foreign interference, by force or by menace, in the
dispute between Spain and the colonies, as a motive for recognizing the
latter without
delay\literalquote{}------以贸易、承认与海军构成蓝水战略的机制，不自动复制史实拉美独立的时程。

未决：1803--05财政部／英格兰银行／海军的撤移与付款连续性原令未取得（\literalquote{}三十车黄金\literalquote{}等细节不列作已核能力）；1812年4月23日卡斯尔雷两信只有后世转述；克鲁泽专著正文、格洛弗、库克森目标章未获；非英语来源覆盖不足。

\begin{center}\rule{0.5\linewidth}{0.5pt}\end{center}

\subsection{第七章　不列颠（二）财政与信用：税、债、银行限制期、补贴与耐力年份}\label{ux7b2cux4e03ux7ae0-ux4e0dux5217ux98a0ux4e8cux8d22ux653fux4e0eux4fe1ux7528ux7a0eux503aux94f6ux884cux9650ux5236ux671fux8865ux8d34ux4e0eux8010ux529bux5e74ux4efd}

\subsubsection{一、1810--11：压力确有，信用未断}\label{ux4e00181011ux538bux529bux786eux6709ux4fe1ux7528ux672aux65ad}

【锚】1811年5月20日预算比较3\%类借款：\literalquote{}Last year the interest was 4l.
4s. 2d. per cent.; this year it was 4l. 14s.
11d.\literalquote{}------特定新增贷款的现金票息由4.208333\%升至4.745833\%，+53.75bp，融资仍成功；它不是全部公债的重定价，也不含分期缴款的完整IRR。英格兰银行《千年宏观数据》A31表：年均consol收益率1810年4.40854\%、1811年4.67417\%（+26.6bp）；1797、1798年为5.90660\%、5.93145\%------1811年的恶化低于1797--98已历的危机水平。A27表：1810年中央政府总收入£73.0m、1811年£71.0m，利息24.4m、24.6m；债息／收入33.4→34.6\%。Bordo--White（NBER
WP
3517，1990，表2印p.51）：1810、1811年净收入约68.39m、66.53m，债务负担23.40m、23.86m（Gayer--Rostow--Schwartz口径），同样不见断崖；其表3转引Mitchell
\& Deane（1962）p.442------本轮所得只是对校层级，未亲阅原表。

\subsubsection{二、封锁的时钟不是违约的时钟}\label{ux4e8cux5c01ux9501ux7684ux65f6ux949fux4e0dux662fux8fddux7ea6ux7684ux65f6ux949f}

【推·中】在1810基线票息24.4m、持续收入55--60m、非息支出48--50m、新借成本6--7\%且零名义税基增长的模型中，债息／收入60\%探针在1817--1823年触及，70\%探针在1820--1827年；起息±2扩展至1815--26与1818--30。它们是条件压力日期，不是预测违约日，也不是必然媾和日。加入首年维持史实G57.2、12--24个月阶梯减支的过渡后：S2至1848年均未越60\%探针，1830年债务代理£690--902m、1848年£645--1,080m------\textbf{不保证低于史实}（£840m／£831m）；S3适应线过渡后1844--47年亦越60\%。伦敦金融中枢被破坏才可能触发短时清算断裂------政府、账簿、税网保全时数周至两季度清算危机、6--18个月重建------1803年已在限制期，不能把再停兑当全新未用的保险。

\textbf{给对弈的结论}：英国不因财政被迫投降。战略败退不必等于财政崩溃；万西塔特1815年的陆转海是当时文书支持的政策反馈。未闭合：Mitchell--Deane／Sherwig原表；1811年商业救助实放与回收；1815债息32.2与Hume
30.450903、1811海军拨款19.822与20.276144差额未桥接；最低海军实物成本。

\begin{center}\rule{0.5\linewidth}{0.5pt}\end{center}

\subsection{第八章　不列颠（三）皇家海军与海军部：舰、人、木、坞与封锁的可持续性}\label{ux7b2cux516bux7ae0-ux4e0dux5217ux98a0ux4e09ux7687ux5bb6ux6d77ux519bux4e0eux6d77ux519bux90e8ux8230ux4ebaux6728ux575eux4e0eux5c01ux9501ux7684ux53efux6301ux7eedux6027}

\subsubsection{一、账面舰队不是可用舰队}\label{ux4e00ux8d26ux9762ux8230ux961fux4e0dux662fux53efux7528ux8230ux961f}

【锚】1808年初：战列舰海勤现役113、海勤预备13、合计126；港勤现役19、港勤预备44、在建或已订购48------\literalquote{}Line
113 197350 13 25089 126
222439\literalquote{}；账面大合计237不等于可用舰队。詹姆斯《Abstract》17--23（现代转录，未逐年影像对校）：1809--15年海勤现役113、108、107、102、102、99、47；海勤总127、124、124、120、124、118、109。1803--14年毛建成74艘战列舰，1808--12五年46艘------不是净增。1812年票定145,000＝seamen
113,600＋marines 31,400，不能称\literalquote{}十四万五千熟练海员\literalquote{}。

【锚】1813年11月议会：政府方以法国扩舰、美国与波罗的海任务及海员难以召回为由拒绝提前裁军------\literalquote{}If
we suddenly disbanded, it would not be so easy a task, on an emergency,
to recal our
seamen\literalquote{}；反对方主张把海军资源转陆战。1797年斯皮特黑德十六艘战列舰拒航，诉求为薪酬、给养与待遇；诺尔的扩大诉求与封锁泰晤士后遭镇压------人力的政治天花板真实存在。

\subsubsection{二、舰材：韧性上调，但不是充分替代}\label{ux4e8cux8230ux6750ux97e7ux6027ux4e0aux8c03ux4f46ux4e0dux662fux5145ux5206ux66ffux4ee3}

【锚】阿尔比恩（Albion，1926，p.356）：北美大桅1807→08年输出由4,442增至16,729------\literalquote{}the
exports from British North America jumped from 4,442 to 16,729 in a
single year\literalquote{}；橡木1810年逾17,000、1811年逾24,000
loads；p.368：1812年前赴英柚木仅一货船（非柚木战舰）。戴维（Davey，2009，博士论文，本轮全文入库）：1810年波罗的海给养改善，机制是替代与持续通商护航------不是大麻的充分替代；俄麻与熟练劳力仍是真瓶颈。麻铜实物余额、维修召回与国家总账未闭合。

\subsubsection{三、法胜大陆之后的皇家海军}\label{ux4e09ux6cd5ux80dcux5927ux9646ux4e4bux540eux7684ux7687ux5bb6ux6d77ux519b}

主判：法国的大陆胜利不导致英国自动重复1817年式深裁。持续海战可规划13--15万人、高档£18--22m（1812费用尺）；武装和平6--9万人。1830年端点（规划包络，非观测）：W态H130--155／A12
115--145／I 105--125；P态H130--155／A12 100--140／I
45--70。\textbf{法国六十艘若为当季真实敌对压力，英国的低和平档必须升档：至少78--97任务舰。}A12不可当作敌对压力F直接使用；1815年英国P停战后与法国停战前不直接比I；不把詹姆斯的burthen与法国排水量相除。海上优势收窄不消失；英国接受大陆妥协的门槛低于交舰队的门槛；至1848竞争共存。

\begin{center}\rule{0.5\linewidth}{0.5pt}\end{center}

\subsection{第九章　不列颠（四）陆军与本土防御：正规军、民兵、志愿兵与登陆战局}\label{ux7b2cux4e5dux7ae0-ux4e0dux5217ux98a0ux56dbux9646ux519bux4e0eux672cux571fux9632ux5fa1ux6b63ux89c4ux519bux6c11ux5175ux5fd7ux613fux5175ux4e0eux767bux9646ux6218ux5c40}

\subsubsection{一、三十八万不是三十八万}\label{ux4e00ux4e09ux5341ux516bux4e07ux4e0dux662fux4e09ux5341ux516bux4e07}

【锚】福蒂斯丘的1803年底\literalquote{}三十八万\literalquote{}口径含大不列颠志愿兵与自耕农骑兵全军衔；\literalquote{}三十四万\literalquote{}为自报effective
rank and
file。1804年1月GB志愿兵380,258全军衔：骑28,943／步341,011／炮10,304；1805年降为360,814------原表已看图。林奇（Linch）：1810年英格兰三十九个线列营24,764人，其中7,677病员，\literalquote{}leaving
only 17,087 men fit for
duty\literalquote{}；1811年UK名册55,938人仅五个营可增援威灵顿------单位骨架与体能而非总数决定可动员性。福蒂斯丘附表：1803至1813年普通征募共134,316、民兵转募99,755------后者是来源重分配，不可再加入本土总兵力。

\subsubsection{二、集中与运输}\label{ux4e8cux96c6ux4e2dux4e0eux8fd0ux8f93}

【锚】1804年已把登记车辆编入部队运输：\literalquote{}waggons should likewise be
marked, set apart, and provided with seats to convey the men from the
remoter districts to any threatened point\literalquote{}；车运计划估计\literalquote{}It was reckoned
that the waggons would carry the troops forty miles a day for three
days\literalquote{}，目的是向伦敦集结。奇尔科特p.244引WO 30/141
p.16：征用供给车辆须二十四小时内备好、计划日行25--30英里------作者认为可能乐观。\textbf{无1804年南岸四十八小时集中四万正规军的演练实绩。}英格兰南岸74座马特洛塔主要建于1805--1808，东岸29座为1809--1812：\literalquote{}Work
on the south coast towers began early in
1805\literalquote{}------1803--04不能赋予成熟塔链的阻滞能力。

\subsubsection{三、登陆战局的四十八格}\label{ux4e09ux767bux9646ux6218ux5c40ux7684ux56dbux5341ux516bux683c}

在\literalquote{}D+2完整合成军已上岸、英国海军D+2--7封断后援\literalquote{}的基准下：两万法军较易被遏制／争夺，四万是预警铰链，八万与十二万较可能取得初期局部优势------这不证明大规模卸载可达，也不是全国迫和的概率。运输：德布里埃尔（Desbrière，p.442）所刊战争部长致内伊函署1805年3月30日；pp.443--444将parc
de l\textquotesingle armée的107辆车、456匹挽马列为待另配份额------\literalquote{}Les
456 chevaux de trait pour les 107 voitures attachées au parc de
l\textquotesingle armée\literalquote{}------原1,616需求／1,391马位不能读成完整运输纵列已闭合。四组蒙特耶尔运载样板提供80,928战斗席但只有5,564马位，与八万人10\%马率假设相差2,436马位------两组假设的资源冲突，不是史实全军缺额。两潮（Saint-Haouen）与六潮（Glover）模型给出八万渡海约2--4日或4.5--6日，是串行时间压力测试，装卸量未校准。S5下英国常设陆军三档6.5--9／10--15／16--22万（all
ranks口径，排驻印、公司军、海军陆战队），低置信。取消半岛战争不可只给法国节省陆军而继续把同一英国远征军从本土扣除；英国的驻军分配与征募内生改变。

\begin{center}\rule{0.5\linewidth}{0.5pt}\end{center}

\subsection{第十章　不列颠（五）工业、商业与劳工：被排除二十五年能否存活}\label{ux7b2cux5341ux7ae0-ux4e0dux5217ux98a0ux4e94ux5de5ux4e1aux5546ux4e1aux4e0eux52b3ux5de5ux88abux6392ux9664ux4e8cux5341ux4e94ux5e74ux80fdux5426ux5b58ux6d3b}

\subsubsection{一、\literalquote{}出口不是贸易\literalquote{}}\label{ux4e00ux51faux53e3ux4e0dux662fux8d38ux6613}

【锚】1811年议会救济辩论直接质疑发货额可代表繁荣：\literalquote{}Exports were not
trade\literalquote{}------货物可滞留南美两三年；大陆对回运殖民品的封锁会反过来削弱海外客户的支付能力。1812年4月17日谢菲尔德请愿明确愿为英国独立承压，同时要求废除枢密令：\literalquote{}if
they regarded those Orders as necessary for supporting ... the just
rights and independence of the United Kingdom, they would willingly bear
the pressure without a murmur\literalquote{}------反贸易管制不等于接受法国霸权。

【锚】马歇尔《统计展示》（1833，印p.74／PDF
88）确认1811年美国官方出口£1,431,829，加加拿大1,909,689为3,341,518；英格兰银行A40表相应美国11.432m／北美13.342m比原表约多£10m------本轮保留原始值并另建校订全球总额£27.459m（Heckscher申报现价≠official）。

\subsubsection{二、替代有效，但危机年不足}\label{ux4e8cux66ffux4ee3ux6709ux6548ux4f46ux5371ux673aux5e74ux4e0dux8db3}

以1804--06平均为基期，旧官方价拉美＋亚洲总货物增量／北欧缺口在1811年为21.7\%，1808--12年净额合计87.3\%；纯拉美本产回款口径仍缺------证据同时支持有效替代与危机年替代不足。

【推·条件包络，非统计CI】1848年工业相对同年史实：S2
79--100\%；S3（排除至1848）54--79\%（排除1815--1840满二十五年为59--81\%，仍约1803年的1.72--2.36倍）；S5商业和平88--109\%。主线：持续的欧洲排除不会可靠地消灭英国工业或迫其交出政治独立；1810年代危机的机制是销量、回款、原料、粮价与阶级治理的叠加。\textbf{更稳定的法国霸权需要给英国商人可兑现的贸易，而非永久最大伤害。}保西王不自动节省葡萄牙战费。未决：纯拉美年度回款、工厂投资面板、克鲁泽原专著全文。

\begin{center}\rule{0.5\linewidth}{0.5pt}\end{center}

\subsection{第十一章　不列颠（六）海外帝国：西印度、加拿大、开普、地中海基地与帝国子午线}\label{ux7b2cux5341ux4e00ux7ae0-ux4e0dux5217ux98a0ux516dux6d77ux5916ux5e1dux56fdux897fux5370ux5ea6ux52a0ux62ffux5927ux5f00ux666eux5730ux4e2dux6d77ux57faux5730ux4e0eux5e1dux56fdux5b50ux5348ux7ebf}

【锚】1823年下院转述1820年海外军需官英国供款表约£1.629m；锡兰1817年地方收入£378,812而支出£450,816；1830年下院宣读1827年财政部信：五个新殖民地异期十年以上累计赤字£2.524m，英国供款£2.435m------\literalquote{}after
nearly the whole of it had been laid out there, they required large
remittances from this
country\literalquote{}。地方税、中央补贴与海军战略收益必须分账；殖民扩张短期更可能增加补贴和军费，不能单凭占领港口推断偿债能力增强。

【锚】1833年废奴法第LXIV条明确排除东印度公司领地、锡兰、圣赫勒拿：\literalquote{}or to
the Island of Ceylon, or to the Island of Saint
Helena\literalquote{}；LXV条开普延四月、毛里求斯延六月。殖民政府反驳莫尔斯沃思的军海支出分摊：无殖民地仍需舰队和邮船；少数白人自治可能进一步压迫西印度黑人多数------£4m不能直接当可避免成本。

贸易账（描述会计，非封锁因果）：以1806申报值为基年，1808年其余美洲本产出口增5.71m，北欧含法降5.41m，抵补1.055；对北欧加西葡降5.85m抵补0.976；1811年其余美洲仅增1.06m对北欧降6.07m，美国另降10.55m。以1805为基年，1808--10三年其余美洲增量连续超过北欧及北欧加西葡的下降------宽反证触发，蓝水韧性上调；但正式殖民地的份额未知（1808年若要由正式帝国部分覆盖北欧损失，至少需94.7\%增量属正式帝国，加西葡则需102.5\%，不成立）。亚眠条约第VI条明确开普与巴达维亚主权但各缔约国同税补给，第X条列马耳他中立、多国保证与接防条件------殖民主权与港口准入可分开交易，这是B8的真实当时工具锚点。

主线：法国大陆优势与英国海外续存可并存；英国保核心港链、差异化正式帝国与外国市场；S2比持续绞杀更可能诱使英国减少大陆干预。1848年非印度驻军35--50k、殖民中央支出£3.4--5.4m／年（含海军分摊，与B2／B3不可再相加），围绕1848年演说参照的裁量包络。NSW羊毛两来源口径冲突保留。

\begin{center}\rule{0.5\linewidth}{0.5pt}\end{center}

\subsection{第十二章　爱尔兰：联合王国的软肋还是幻想的第二战场}\label{ux7b2cux5341ux4e8cux7ae0-ux7231ux5c14ux5170ux8054ux5408ux738bux56fdux7684ux8f6fux808bux8fd8ux662fux5e7bux60f3ux7684ux7b2cux4e8cux6218ux573a}

\subsubsection{一、法方把爱尔兰动员写成条件，不是资源}\label{ux4e00ux6cd5ux65b9ux628aux7231ux5c14ux5170ux52a8ux5458ux5199ux6210ux6761ux4ef6ux4e0dux662fux8d44ux6e90}

【锚】1803年8月8日，拿破仑令德克雷向巴黎的联合爱尔兰人代理提出：\literalquote{}Decres
was told to offer to the agents of the United Irishmen then in Paris
25,000 men and 40,000 muskets, with artillery and ammunition, provided
that 20,000 of their countrymen joined the French army on its
landing.\literalquote{}1804年9月方案18,000＋25,000两波、目标都柏林。而1803年12月，应构成入侵军主体的爱尔兰军团\literalquote{}consisted
of but fortynine officers and thirteen
soldiers\literalquote{}（与Carles所记1804年1月20日四十三官一士官互证）。法方数字链全部经德布里埃尔的两条转引路径（Wheeler
\& Broadley 1908、Carles
1976），二者同源，不构成独立互证；W\&B将埃米特起义误系于1803年6月23日（应为7月23日），该书精度应降权。

【锚】1796--1798年驻爱部队中正规兵只占5--6\%（1793年68\%，1801--02年20\%）：\literalquote{}The
percentage of regular soldiers on the Irish establishment in the period
1793-1802 was as follows: 1793, 68 per cent; ... 1796, 6 per cent; 1797,
6 per cent; 1798, 5 per
cent\literalquote{}。1803年续战使爱尔兰义勇军由63,000增至约80,000；哈德威克拟设天主教连队遭阻（威克姆称其\literalquote{}unsafe\literalquote{}）。英国在爱尔兰的镇压主力是一支武装的、有地方情报的新教少数派民兵。1808--15年英军204个骑步兵团的十四万余名士兵样本中，爱尔兰出生者占32\%（英格兰53\%、苏格兰14\%）------爱尔兰同时是英国战争机器的人力支柱与法国设想的第二战场兵源。

【锚】1803年临时政府《宣言》：\literalquote{}2. From the same date, all transfers of
landed property are prohibited, each person, holding what he now
possesses, on paying his rent until the national government is
established\literalquote{}------革命派自己的文件切断了\literalquote{}法国胜利→土地革命→无饥荒\literalquote{}的推理链。

\subsubsection{二、动员上限与占领账}\label{ux4e8cux52a8ux5458ux4e0aux9650ux4e0eux5360ux9886ux8d26}

【推·中】动员上限四档：无法军300--3,000；1,000--2,000法军3,000--8,000；5,000--8,000法军10,000--25,000；18,000--25,000法军＋四万支枪30,000--60,000，都柏林失守概率0.55--0.70。参数：参与率以1798年韦克斯福德0.10--0.20与梅奥0.01--0.03为界，每法军带兵k≈4取自安贝尔实例，枪支为硬约束（本地无兵工）。爱尔兰的军事价值是法国投送量的函数，不是爱尔兰民意的函数；分水岭＝两万人＋四万支枪。若1804--05年驻爱英军实为四万以上（Kilmainham／WO
17逐月表未获），概率应降至0.30--0.45。

法国长期占领需33,000--66,000人（南三省3--6‰、阿尔斯特东北18--30‰），而爱尔兰财政汲取上限约£5.0--5.5m／年（1807年岁入£4.378m，1807--21加税反而减收）------\textbf{爱尔兰对法帝国是净消耗单位}，不得入A层，上限B级藩属。三方结构（天主教多数／八万武装新教义勇派／英国国家）：法军登陆只移除第三项，第二项退守阿尔斯特东北并得皇家海军支持------最接近的类比是西西里／加的斯式沿海阵地，不是1808年西班牙式全民反法。

三项未用潜力：可见的远征准备（锁住英军并阻止1811年英爱民兵互换法启用）；把爱尔兰军团扩为3,000--5,000实兵作登陆骨架；以巡航舰队袭扰科克--沃特福德的粮秣与腌肉出口（科克是英国封锁舰队与半岛军的给养基地）。

【推·中】S2／S5下天主教解放提前到1806--1813年，且不伴随1829年史实的四十先令自由持有人选举权剥夺（郡选民约216,000→37,000）------奥康奈尔式群众政治提前15--20年，议题直接跳至废除联合与土地。饥荒（疫病外生，所有情景发生）：S0
1.0--1.1M；法系＋大陆式征用管制＋开放外迁0.3--0.7M；法系＋战时封锁0.9--1.4M（可能更糟）；英治＋提前解放0.7--1.0M；决定变量排序：外迁通道＞救济制度类型＞土地产权＞政权国籍。置信低--中。

\begin{center}\rule{0.5\linewidth}{0.5pt}\end{center}

\subsection{第十三章　俄国（一）宫廷与外交精英：亚历山大能被什么收买、被什么推翻}\label{ux7b2cux5341ux4e09ux7ae0-ux4fc4ux56fdux4e00ux5babux5ef7ux4e0eux5916ux4ea4ux7cbeux82f1ux4e9aux5386ux5c71ux5927ux80fdux88abux4ec0ux4e48ux6536ux4e70ux88abux4ec0ux4e48ux63a8ux7ffb}

\subsubsection{一、弑君约束的贸易基座}\label{ux4e00ux5f11ux541bux7ea6ux675fux7684ux8d38ux6613ux57faux5ea7}

俄国不是一个可以被条约锁住的国家，因为它的沙皇不是一个可以被条约锁住的人------他被自己的贵族锁着。【锚】1801年政变前夜，祖博夫在晚宴演说中明列废黜保罗的理由之一：\literalquote{}the
insane act of a rupture with England was contrary to the essential
interests of the Russian nation, dried up the sources of its
wealth\literalquote{}（恰尔托雷斯基回忆录，卷一；回溯性，但兼采本尼希森亲述）。政变以亚历山大同意废黜（非谋杀）为前提，事后无法追究，因为\literalquote{}彼得堡社交界人人是共犯\literalquote{}。

【锚】皇太后对女儿婚姻持法定否决权：保罗一世遗诏以庄严敕令形式存于莫斯科圣母安息大教堂，\literalquote{}conférait
un véritable droit de veto dont elle ne manquerait pas
d\textquotesingle user\literalquote{}（Vandal，卷一，本轮亲读）；亚历山大虽可以君主意志压倒，但\literalquote{}对母亲以主上口吻说话于他不可忍受\literalquote{}。俄婚让步包须前置12--18个月管理这一制度性障碍------它没有被管理。

\subsubsection{二、亚历山大是条件性机会主义者}\label{ux4e8cux4e9aux5386ux5c71ux5927ux662fux6761ux4ef6ux6027ux673aux4f1aux4e3bux4e49ux8005}

【锚】1810年12月25日密信向恰尔托雷斯基提出俄版波兰复国方案（八问＋兵力表230,000对155,000）：\literalquote{}Such
a moment presents itself only once... The support on which the Poles can
rely is limited to the person of Napoleon, who cannot live for
ever.\literalquote{}1811年1月31日第二信给出全部担保：东界德维纳--别列津纳--第聂伯、自由宪法、纯个人联合，自报一线兵力240,500＋预备124,000；但矢言非波兰倒戈不先手：\literalquote{}Mais
tout change de face si les Polonais veulent se joindre à
moi\literalquote{}。1812年4月1日自述：\literalquote{}Your previous letters had left me too little
hope of success to authorise me to
act\literalquote{}------先发计划因波兰拒绝合作而放弃。

【锚】不媾和的决心1811年5月已成型（前引科兰古告别宣言）；入侵后经巴拉绍夫以荣誉担保重申\literalquote{}一兵在境即不议和\literalquote{}。奥尔登堡是可交易的面子议题：1812年2月亚历山大向瑞典特使勒文海尔姆指图愿以奥尔登堡补偿丹麦丢失挪威之损，\literalquote{}qu\textquotesingle il
sacrifierait volontiers malgré la parenté\literalquote{}。

【锚】1812年3月17日之夜斯佩兰斯基的流放是向贵族与舆论交出的\literalquote{}抵押\literalquote{}：阴谋由阿姆费尔特为首，联合警察大臣巴拉绍夫与阿拉克切耶夫；亚历山大对诺沃西尔采夫自白\literalquote{}L\textquotesingle ennemi
frappe à la porte de l\textquotesingle empire...Je joue gros
jeu\literalquote{}，私下重复\literalquote{}斯佩兰斯基不是叛徒\literalquote{}；彼得堡社交界欢庆此为\literalquote{}c\textquotesingle était
une première victoire sur les
Français\literalquote{}（席尔德转引）。任何让步包的窗口在斯佩兰斯基被献祭前（≤1811年底）有效得多。

\subsubsection{三、六个让步包×八个精英：唯一六列同击的是关税容忍}\label{ux4e09ux516dux4e2aux8ba9ux6b65ux5305ux516bux4e2aux7cbeux82f1ux552fux4e00ux516dux5217ux540cux51fbux7684ux662fux5173ux7a0eux5bb9ux5fcd}

把六个让步包（波兰保证／奥尔登堡补偿／关税容忍／多瑙公国与君士坦丁堡／共同印度远征／联姻）对八个精英集团（沙皇、鲁缅采夫、斯佩兰斯基、皇太后党、英派贸易贵族、军鹰、流亡主战网、教会）打分：\textbf{唯一能同时打动前六列的是关税容忍}（接受1810年敌制关税＋中立航运）------涅谢尔罗迭五项中唯一的硬红线即此；沙皇\literalquote{}大国自主\literalquote{}指令；斯佩兰斯基财政止血；皇太后党存量怨气释放；祖博夫演说的贸易条款。而它正是拿破仑为了S3所不能给的。\textbf{法俄同盟的死结是结构性的（封锁vs贵族生存），不是个人性格。}

【推·中高】俄国对法国霸权工程的可达上限是C层（被胁迫同盟），半衰期约四年------这里的\literalquote{}四年\literalquote{}是观察到的同盟政治寿命（1807-07→1811初→1812-04）加上政治安全裁量，\textbf{不是经济阈值}（第十四章否决\literalquote{}三至四年贵族破产链\literalquote{}）。不存在使俄国长期从属的已证路径；存在有文书锚点的\literalquote{}有条件冷和平\literalquote{}（涅谢尔罗迭1811年方案），代价是放弃排他封锁并容忍俄奥重新武装。A层无任何入口；B层无嫁接口------连沙皇自己的斯佩兰斯基法制化都因\literalquote{}太法国\literalquote{}而政治死亡。

B7三版：不战＋接受1810关税（v3）是\literalquote{}俄国受约束\literalquote{}意义上唯一有文书基础的达标支线；止于双河线的有限战争（v2）军事可停、政治不结果，造就\literalquote{}第二个西班牙\literalquote{}；史实版自毁。falsifier裁定：贸易主因说成立但需限定------贸易是敌意的放大器与政治化通道而非唯一源（沙龙敌意含当吉安公爵案等文化成分，无封锁的1803--05年俄廷已反法）；关税容忍买到的是\literalquote{}不战\literalquote{}而非\literalquote{}亲法\literalquote{}。\literalquote{}亚历山大自1808年起决意再战\literalquote{}不成立：1808--09年他在同盟内收取芬兰与两公国实利；成立的是弱版------自1810年末起视战争为大概率终局并系统备战，保留由对方开第一枪的决策纪律。

S5推演：法胜世界里神圣同盟无法作为欧洲秩序出现（其生成条件是击败拿破仑）；亚历山大的弥赛亚气质倒装为以俄为中心的东方正统集团；波兰宪政实验（1815宪法＋1818华沙演说，本轮核官方俄译转引链）证明其宪政主义是主权者的礼物语法，S5下实验场被删除，波兰从irredenta方向发酵。俄国在法制欧洲扮演史实中英国＋奥斯曼的混合角色：体系外武装平衡者，1848型总危机时最大的外部收割者。

未决：\literalquote{}1811年库拉金备忘录\literalquote{}未获原档（最近对应物为Kurakin
1808--09商务报告，Makarov 2024转引）；Lieven《Russia Against
Napoleon》与Rey《Alexander
I》专著经书评转引；1818华沙演说法文原本未获；塔列朗埃尔福特变节语为同源双传本，无第三方在场证据。

\begin{center}\rule{0.5\linewidth}{0.5pt}\end{center}

\subsection{第十四章　俄国（二）贵族、农奴经济与财政：大陆体系伤了谁}\label{ux7b2cux5341ux56dbux7ae0-ux4fc4ux56fdux4e8cux8d35ux65cfux519cux5974ux7ecfux6d4eux4e0eux8d22ux653fux5927ux9646ux4f53ux7cfbux4f24ux4e86ux8c01}

\subsubsection{一、减压始于1809，不是1811}\label{ux4e00ux51cfux538bux59cbux4e8e1809ux4e0dux662f1811}

【锚】克罗斯比（Crosby，America, Russia, Hemp, and
Napoleon，1965，本轮全文入库）：\literalquote{}In 1809, 376 vessels arrived at
Kronstadt and St. Petersburg as compared to sixty for
1808.\literalquote{}pp.209--213记1811年俄港殖民货过剩、麻价上升，美商俄绳货在美国与巴西亦滞销------反驳\literalquote{}封锁期贸易持续瘫痪\literalquote{}。p.33称莫斯科至波罗的海的雪橇运费低于等距海运------须区分季节性国内联运与向法国另辟整条陆路。特罗申（Troshin，2015）给彼得堡有限商品样本的贸易条件指数1808＝.43、1809＝.69、1810＝.74、1811近同------\textbf{0.74是价格贸易条件指数，不是出口量或贵族净收入的七成恢复}；其表1年份错排、口径不一，不作全国完整序列。

【锚】斯佩兰斯基1811年报告：\literalquote{}прекращение выпуска ассигнаций возвещено 2
февраля 1810 г., а восприяло начало свое только 1 генваря 1811
г.\literalquote{}------停发纸币1810年宣布、1811年初才实行，荷兰与汉堡的信用受法方行动冲击。戴维表29的1806年脂油目的地列加总79.88而印100、亚麻99.28而印100------缺格不可当零。

\subsubsection{二、损失的条件带}\label{ux4e8cux635fux5931ux7684ux6761ux4ef6ux5e26}

【推·参数包络】出口庄园直接毛现金损失：S2中立许可分支0.5--16.7\%（中心5.3\%）；S3真实加严7.2--41.3\%（中心22.8\%）；完全合法直贸的隔离贸易损失可为零。模型参数（e、b、n、r、p、d）公开，不是样本均值或统计置信区间。法国若要替英国吸收俄麻：以世纪之交六万作者吨麻出口与1806年目的地份额交叉，需直接市场×26.4或广义大陆×3.41------异年同价尺度，非已证容量（1806份额含对英转口中介故偏宽）。1762--82年英国进口大麻的95.9\%来自俄国（Davey
p.23：\literalquote{}between 1762 and 1782, from a total of 19,172 tons of hemp
imported by Britain, 18,392 came from
Russia\literalquote{}），但这不是拿破仑时期皇家海军的需求率。

【锚】涅谢尔罗迭1811年10月的和平建议仍以俄奥恢复财政与军队、制衡法国为目的------\literalquote{}celui
qu\textquotesingle au fond de leur pensée elles regardent comme leur
seul ennemi\literalquote{}（Vandal卷三）。不应只引和平语句当作永久接受法国霸权的承诺。

\textbf{主线}：法国以俄贸易／信用／关税及边境安全自主换有限克制，比封锁逼服更可信；S2的最强交换包＝中立贸易与关税自主＋汉萨信用＋奥尔登堡／普鲁士＋互撤兵担保；经济利益能支持合作，不能单独购买俄国永久服从。到1848年俄国仍是选择性工业化、改革迟滞的自主保守大国；\literalquote{}财政攻击\literalquote{}改称\literalquote{}负面财政外溢\literalquote{}。未闭合：全国海关实物量、分组庄园账、统一币种债务年度账、逐年乌拉尔成本。

\begin{center}\rule{0.5\linewidth}{0.5pt}\end{center}

\subsection{第十五章　俄国（三）军队与多线战争：能被击败吗、能被有限战争胁迫吗}\label{ux7b2cux5341ux4e94ux7ae0-ux4fc4ux56fdux4e09ux519bux961fux4e0eux591aux7ebfux6218ux4e89ux80fdux88abux51fbux8d25ux5417ux80fdux88abux6709ux9650ux6218ux4e89ux80c1ux8febux5417}

\subsubsection{一、纸面百万，在列减半，一线再减半}\label{ux4e00ux7eb8ux9762ux767eux4e07ux5728ux5217ux51cfux534aux4e00ux7ebfux518dux51cfux534a}

【锚】1812年6月俄军四层口径：在册近六十万（Lieven）／编制48万＋1,600炮（Bogdanovich）；一线三个西方面军含哥萨克仅24.2万；里恩（Riehn）：\literalquote{}the
grand total of all the forces in the first, second, and third lines was
488,000 men, of which about 428,000 gradually came into
action\literalquote{}；茹拉夫斯基财政口径131.8万为满编预算值，实际\literalquote{}在列\literalquote{}仅其现役栏目一半------\textbf{禁入总装}。Lieven一线242k与Riehn
190k禁互填。

【锚】征兵天花板（Keep，Soldiers of the
Tsar，本轮全文亲读摘录）：1705--1825共九十次levy；1801年前累计227万，亚历山大朝再征近二百万；被征占登记男丁比例彼得朝3.6\%→18世纪末7\%→19世纪初8\%；\literalquote{}The
greatest drain occurred in 1812, during Napoleon\textquotesingle s
invasion, when three levies were called, which took a total of 20 men
per unit\literalquote{}（每五百男丁二十人）。1825年军事殖民地\literalquote{}the Separate Corps of
Military Settlements, reportedly controlled 160,000
soldiers\literalquote{}（约全军三分之一）＋37.4万农民；抵抗序列（Bug 1817、Chuguyev
1819二千人被捕／275判死／52夹鞭25死、诺夫哥罗德1831）\literalquote{}episodic and
ineffectual\literalquote{}。

\subsubsection{二、B7
falsifier成立：不靠焦土与冬季}\label{ux4e8cb7-falsifierux6210ux7acbux4e0dux9760ux7126ux571fux4e0eux51acux5b63}

【推·中高】无焦土无冬季，俄军1812年底仍保正规野战军≥20万（Riehn自洽链488k／428k／正规≈33万；多瑙军与芬兰军的释出由外交决定）；克劳塞维茨\literalquote{}战争进程中几乎必然翻倍\literalquote{}的判断与数据吻合。追击也遭冬季重损------库图佐夫1812年12月19日自报：\literalquote{}Of
the 97,000 men whom Kutuzov had commanded at Tarutino before the
beginning of the campaign, 48,000 -- in other words almost half -- were
in hospital. Only 42,000 soldiers were still in the
ranks.\literalquote{}（Lieven）------但干部核心存活，1813年夏完全恢复：冬季是双刃。

B7有限战争四分支：A主线（中高）俄军不接受停线、逐年再生消耗；B（中）冬季有限反攻撞未受损法军为俄最坏枝；C（中低）接受割波兰和约------唯一杠杆是宫廷政变，法军无法军事制造；D（低）体系崩溃需四条件同时成立且与停线设定自相矛盾。\textbf{有限战争的最优期望＝一次休战。}复国波兰与停线胁迫在B7内互斥：停线利于淡化\literalquote{}侵略\literalquote{}定性以削弱俄本土轨融资，但\literalquote{}it
was in Napoleon\textquotesingle s power to re-establish a Polish state
of 15 million people on Russia\textquotesingle s borders and this would
be a huge threat to Russian
security\literalquote{}（Lieven引本尼希森文书）直接激活精英生存威胁认知。

【锚】战争财政双轨：\literalquote{}In 1812 Alexander had not requested a British
subsidy... fighting on his own territory he could finance the war
without great
difficulty\literalquote{}（Lieven）；本土轨（纸币＋实物捐输约一亿卢布＋常平仓）弹性大但不可携带；境外轨依赖英援（1813年£4.6m≈德意志战役全部预计支出）。法军止步边境省份＝把俄国留在其财政最强形态内作战。

【锚】库图佐夫对威尔逊的均势派证词：彻底摧毁拿破仑的遗产\literalquote{}his succession
would not fall to Russia or any other continental power, but to that
which commands the sea, and whose domination would then be
intolerable\literalquote{}------俄军最高指挥层存在不愿为英国火中取栗的流派，是B7谈判分支的人事基础。

\textbf{结论}：1812年之外击败俄国的军事条件是\literalquote{}更浅出\literalquote{}（边境省份合围歼灭干部核心，1812年6--8月窗口）或\literalquote{}不战而约束\literalquote{}（切英援＋关税，1811年已实质达成）。S5俄军＝纸面八十万缺额20\%，区域可畏（对奥斯曼、波斯、邻境镇压）、全欧不可持续（枪械＋外汇＋两千公里行军三重约束）；S5的稳定形态是欧洲两分而非俄国被约束；\textbf{俄军的存在封死S6东扩}；S5末期技术代差积累使俄军可能在边缘战区重演克里米亚式脆弱。未决：Shchepetil\textquotesingle nikov逐levy表；军事殖民地财务账；法方停线参谋文书（科兰古回忆为孤证）；1849匈牙利出兵规模无档案锚。

\begin{center}\rule{0.5\linewidth}{0.5pt}\end{center}

\subsection{第十六章　波兰-立陶宛：藩属过载、复国筹码与对俄杠杆}\label{ux7b2cux5341ux516dux7ae0-ux6ce2ux5170-ux7acbux9676ux5b9bux85e9ux5c5eux8fc7ux8f7dux590dux56fdux7b79ux7801ux4e0eux5bf9ux4fc4ux6760ux6746}

\subsubsection{一、忠诚度最高，榨取率也最高}\label{ux4e00ux5fe0ux8bdaux5ea6ux6700ux9ad8ux69a8ux53d6ux7387ux4e5fux6700ux9ad8}

【锚】1808年1月国王预算草案：\literalquote{}Dokument wskazywał 21 029 398 zł deficytu.
Przy przewidywanych dochodach na poziomie 21 997 700 zł wydatki miały
wynieść 43 027 098 zł, w tym 22 850 770 zł wydatków na
armię\literalquote{}------收入2,200万兹罗提、支出4,300万、军费2,285万、赤字2,100万（AE
CP
324,343转引）；1809年1月国王报拿破仑总收入2,900万、军费2,500万；1810年11月1日欠军需商1,500万。军费占比在Grab的\literalquote{}预算三分之二\literalquote{}与Kowalczyk线的\literalquote{}公共支出九成\literalquote{}之间，无年度决算系列可裁定。

法国给公国的从来不是现金补贴，而是减免与担保：1810年放弃印花4.35m法郎、盐款减1.65m、炮款减0.5m；承担八千波军军饷（巴约讷军事条款）；1811年担保Perregaux--Laffitte借款（实发7,795×1,000法郎）。反向流量为主：巴约讷2,000万法郎票据、驻法军供养。\textbf{净现金流方向是波兰→法国}------卡设falsifier不触发，但公国自身财政不能闭合仍成立，两命题须分开。

【锚】被占负担实测（Tokarz，1917）：\literalquote{}dostawy dla żołnierza rosyjskiego w
ciągu roku 1813 i czterech pierwszych miesięcy 1814 roku wyniosły sumę
157.928.176 złp.; od maja r. 1814 do maja 1815 skarb Księstwa wydał na
ten sam cel 89.942.911
złp\literalquote{}------同一土地在被占领状态下两年被提取近2.5亿兹罗提，约为公国正常年收入的八倍：可被短期榨取的物理上限远高于可持续税基，复国国家若沦为战场，其资产会被双向清空。

犹太人政策双轨：1807宪法与1808-05-09征兵令名义平等→1808年三敕令中止公民权十年并禁购地产→1812-01-29以年缴700,000兹罗提赎全体兵役、战时加倍至140万；会议王国1817--18原样恢复赎役制。

\subsubsection{二、动员上限与复国的疆界}\label{ux4e8cux52a8ux5458ux4e0aux9650ux4e0eux590dux56fdux7684ux7586ux754c}

【推·参数全注】B3-lite（3.2--4.5m人）常备35--55k／冲刺90--130k；B3-max（＋加利西亚，5.5--6m）45--65k／130--165k；胜俄后（＋立陶宛，6.5--8m）55--75k／150--200k。锚：公国宪法定额30k→1810--11纸面59,200（人口4.2--4.3m）→1812部长会议报74,722、Kukiel全口径108k（含法饷与国卫）；立陶宛编制17k基本填满、潜力估50k但被现金--军官--被装四瓶颈卡死；1831年会议王国（人口约4m）自41.8k扩至87k＋。稳态峰值≈人口2.2--2.6\%（含辅助），持续两年以上无外源现金则崩。财政闭合四条件：贸易出口开放／无巴约讷式抽取／驻军自域化／常备≤1.2\%人口。

【锚】亚历山大1811年1月31日的开价（恰尔托雷斯基回忆录附录）：\literalquote{}the
kingdom of Poland shall for ever be united to Russia, whose Emperor
shall in future bear the title of Emperor of Russia and King of
Poland\literalquote{}------复国＝旧波兰全部含俄占省份（白俄除外）至德维纳--别列津纳--第聂伯线，波兰族官员军队，自由宪法；sine
qua
non＝永属俄皇＋公国显要签名担保。1813年1月13日同人明拒米哈伊尔大公方案：立陶宛、波多利亚、沃里尼亚\literalquote{}无论如何说服不了俄国\literalquote{}交给别的君主。

\textbf{B3总判}：1807年\literalquote{}取消普鲁士＋复国波兰\literalquote{}不可一步到位------提尔西特的对俄和解目的与复国目的互斥（拿破仑1812年在维尔纽斯答维比茨基时自证镜像命题：对奥保证即不能全面复国）。可行形态是两段式：1807年做\literalquote{}公国＋但泽＋西普走廊\literalquote{}不称王国（俄可勉强接受，中置信），王国名号留作B7限制战争胜利后的战果；复国王国的最大可信疆界与亚历山大自开的双河线重合。缓冲净收益为正，且有同期直接实测（1811年先发计划因波兰拒绝合作而放弃；其兵力表以波军五万投俄为\literalquote{}不战推进到奥得河\literalquote{}的前提）。

\subsubsection{三、忠诚是双备函数}\label{ux4e09ux5fe0ux8bdaux662fux53ccux5907ux51fdux6570}

波兰忠诚是\literalquote{}复国期票可信度×贸易现金流\literalquote{}两变量皆备时近乎无价（1812年108k）、皆失时归零（1813年春）的非线性跳变。缓冲功能（要求军事化）与政体可持续性（要求宪制--财政平衡）结构性互斥；\literalquote{}产权＋贸易＋自域驻军\literalquote{}三件套是1811年报告已提出的缓解方案。S5：会议王国镜像实验证明\literalquote{}经济整合成功≠政治稳定\literalquote{}（卢贝茨基的盈余与1830年起义并行）；断裂四环：宪法空洞化／军队屈辱／外部革命行情／外派威胁；善治版（宪制＋贸易＋不外派）可过拿破仑寿限且波兰是继承危机中最后背叛的属国（替代选项＝被瓜分）；兵营版在15±7年内重演1830式起义（单样本类比外推）；\textbf{S6行省化是唯一能把最亲法单位推向抵抗的情景}。未决：Żółtowski决算系列、Kukiel《Wojna
1812》原页、Dundulis立陶宛法语系统。

\begin{center}\rule{0.5\linewidth}{0.5pt}\end{center}

\subsection{第十七章　奥地利：宫廷、财政破产、多民族结构与东南转向}\label{ux7b2cux5341ux4e03ux7ae0-ux5965ux5730ux5229ux5babux5ef7ux8d22ux653fux7834ux4ea7ux591aux6c11ux65cfux7ed3ux6784ux4e0eux4e1cux5357ux8f6cux5411}

\subsubsection{一、被证伪的\literalquote{}财政锁死\literalquote{}}\label{ux4e00ux88abux8bc1ux4f2aux7684ux8d22ux653fux9501ux6b7b}

【锚】1809年申布伦秘密条款II：\literalquote{}ne s\textquotesingle élève pas au-dessus
de 150 000 hommes pendant la durée de la guerre
maritime\literalquote{}------限军十五万（不是通说的五万）；第V条八千五百万法郎是一次性和约结算，不是年贡。Axtmann（1991）财政附录转Beer（1877）：1809年财政收入31.1、支出86.3、军费66.8；1810年25.0／76.1／52.2（百万银值盾CM，总支出不含偿本与纸券回收；页图已核）。1810年裁军保留军官、技术兵、休假兵与团骨架；1813年可扩至约29.8万动员编组／24.5万野战（经Messman转引Rothenberg，未亲核团级原账）。\textbf{\literalquote{}财政破产锁死至1820\literalquote{}的假说被证伪}：破产约束再战的成本与窗口，不永久剥夺能力。军费\literalquote{}250m→47m
florins\literalquote{}跨1811年换币未换成同银值，不据其算削减率。

【锚】1809年9月4日的真实照会要求波希米亚三圈割萨克森：\literalquote{}cède à la Saxe
les cercles de Leitmeritz, de Saatz et
d\textquotesingle Elnbogen\literalquote{}------这是谈判领土筹码，不是波希米亚独立方案。1809年的\literalquote{}肢解\literalquote{}是分层的：匈牙利独立的公开宣言、波希米亚三圈的谈判实案、伊利里亚的实际割让------不是一张已获批准的民族建国图。

【锚】1811年《普通民法典》（ABGB）可兼承认天赋人权、禁人身奴役与保留特别财政法规、分层土地权------不需法国占领才能法典化，也不是一次废除所有领主权。1812年法奥盟约第六条担保奥斯曼欧洲领土，故S5奥国东南转向的主线是多瑙交通、通商与政治斡旋，而非无条件吞并巴尔干（领土扩张另列有条件低置信分支）。

\subsubsection{二、跨一次继承的不平等伙伴}\label{ux4e8cux8de8ux4e00ux6b21ux7ee7ux627fux7684ux4e0dux5e73ux7b49ux4f19ux4f34}

【推·中高】主线：保王朝核心、有限援军、商业海口与明确边界，可使奥国成为跨一次继承的不平等合作伙伴；全哈布斯堡的A层并省不是联合满足英国中立、长期财政与一次继承的合理主线。但\literalquote{}留旧官员\literalquote{}本身不否定并省------需检验他们是否接受巴黎的最终税兵任命权。区分S5-A（早克制、避免1809再战，不必继承史实的割海口与1811年式破产）与S5-B（败后修复，以1812年盟约条款为当时可用工具）。占领兵力压力测试10--18万（低置信，非经验估计）；高度合作的减半检验5--9万。

【锚】1813年摄政决议修正1804年禁女性摄政条款：\literalquote{}l\textquotesingle Impératrice
mère réunit de droit à la garde de son fils mineur, la régence de
l\textquotesingle Empire.\literalquote{}若新君非其子，其摄政结束------奥国对外孙的亲缘利益不能自动转到法国旁支；联姻有利一次继承，不自动支持全意大利统一；罗马王早夭后奥方的保护动机减弱。

\textbf{红队修正}：A卡对奥地利的\literalquote{}不平等\literalquote{}来自史实的1805与1809失败。在没有普雷斯堡的世界里（路径B），法国对奥地利的\literalquote{}三项保证\literalquote{}是法国给奥地利的担保，不是约束奥地利的条约；1810年联姻的史实动机是瓦格拉姆后梅特涅的生存策略。第五十章按两条路径分别处理。统计帝国、海关区与反事实剩余领土不重合；1830年人口28.511m不直接移入1810年边界；A卡未冻结可比疆界的绝对人口，只给1810同疆指数100→1848约121--146。

\begin{center}\rule{0.5\linewidth}{0.5pt}\end{center}

\subsection{第十八章　普鲁士：残存国家、改革派与取消普鲁士的可行性}\label{ux7b2cux5341ux516bux7ae0-ux666eux9c81ux58ebux6b8bux5b58ux56fdux5bb6ux6539ux9769ux6d3eux4e0eux53d6ux6d88ux666eux9c81ux58ebux7684ux53efux884cux6027}

\subsubsection{一、条约口径的三处纠正}\label{ux4e00ux6761ux7ea6ux53e3ux5f84ux7684ux4e09ux5904ux7ea0ux6b63}

【锚】1808年9月8日巴黎条约：分立条款的42,000限军\textbf{有十年期限}（1809年1月1日起，期满\literalquote{}rentrer
dans le droit commun\literalquote{}），不得延至1848；正约第VII条三要塞法军驻军\literalquote{}Total
des trois garnisons 10.000
hommes.\literalquote{}（格洛高3,300、居斯特林2,800、什切青3,900），第VIII条饷由法方出，但住宿、口粮、草料、燃料灯火由普鲁士行政提供；第XIII条列七条军用道路------\literalquote{}撤军\literalquote{}不等于财政与主权脱身。分立条款对奥援军法文转录作\literalquote{}une
division de 10,000 hommes\literalquote{}、德文转录作\literalquote{}eine Division von 16.000
Mann\literalquote{}（两版均载1809年特例12,000），冲突未解，数据表只用42,000与12,000。1812年2月24日法普秘密军事协定第II条：\literalquote{}Se.
Majestät der König von Preußen wird ein Kontingent von 20.000 Mann
stellen\literalquote{}（步14,000／骑4,000／炮2,000、60炮）------不用克拉克转载片段的12,000。

\subsubsection{二、取消普鲁士：对俄有价，价格随时间上升}\label{ux4e8cux53d6ux6d88ux666eux9c81ux58ebux5bf9ux4fc4ux6709ux4ef7ux4ef7ux683cux968fux65f6ux95f4ux4e0aux5347}

【锚】Vandal刊述1807年11月12日给科兰古的训令：以俄取多瑙两公国交换法国处置西里西亚；1808年2月17日报告：\literalquote{}La
demande de Berlin effaroucherait peut-être
moins\literalquote{}------索柏林比索西里西亚更少令俄惊惧，因俄怕后者强化华沙。塔列朗回忆录记亚历山大对西里西亚试探答以\literalquote{}Ceci
est réellement une affaire d\textquotesingle honneur pour
moi\literalquote{}（回溯性自述）。到1811年10月，涅谢尔罗迭报告已把普鲁士存续写成最本质条款：\literalquote{}Je
regarde comme beaucoup plus important et même comme
l\textquotesingle objet le plus essentiel de
l\textquotesingle arrangement un article qui assurerait pour quelque
temps l\textquotesingle existence politique de la
Prusse.\literalquote{}并要求按偿款进度归还奥得要塞。【锚】施泰因曾同意使节提案：以三四万普军供法使用换减赔款和撤军------\literalquote{}ein
Truppenkorps von 30.000 bis 40.000 Mann Napoleon zur freien Verfügung zu
stellen\literalquote{}：当时可选的客户国交易存在，不需要假想全欧宪法。

\textbf{判定}：取消普鲁士不是不可能，而是在1807夏--1808年2月有可谈价格，1811年后几近等同选择对俄战争；先决条件是俄方缓冲得到满足且不同时推进大复波兰。取消王朝、将其领土并入德／波客户国、直接法省是三件事，第一项不自动节省后两者的驻军与行政成本：控制残普1.5--3万，取消后初期低／中／高摩擦2.5--4.5／4.5--7.5／8--12万，增量约1.5--6万，按Marion
600--900法郎／人年折合约900万--5,400万法郎／年。

【锚】接收国账：威斯特法利亚\literalquote{}Servicing the huge debt inherited from the
former rulers of this region and maintaining Jerome\textquotesingle s
extravagant court expenditures also burdened
Westphalia\textquotesingle s treasury.\literalquote{}（Grab
p.101）------1812年债务估1.4--2亿法郎，须养本国25,000人并支付驻境12,500法军，拿破仑又以免税赠地削减其岁入：\literalquote{}Das
Königreich Westphalen verlor durch die Abtretung der Dotationsdomänen
wichtige Einnahmequellen, was erheblich zum finanziellen Ruin des Landes
beitrug.\literalquote{}（Berding手稿）。行政延续与社会改革结构性互斥：留用容克与旧官降低接管成本，则土地／领主制改革被冻结。萨克森的显示偏好：\literalquote{}von
seiner dortigen Civilliste nie etwas bezogen, sondern aus seiner eigenen
Casse dem warschauer Staatsschatze nach und nach 30 Mill. Fl.
vorgeschossen, die erst 1828 von Rußland mit 450800 Thlr. vergütet
wurden\literalquote{}（ADB）------1778年拒以两劳西茨换领土，任华沙大公从未取当地civil
list；纯增量型接收可行，交换型将遭拖延。排序：可信有限客户国＞有俄方交易的取消分配＞不断再削残普＞全面复波＋强行法省化。

\subsubsection{三、改革能力与1813年动员分离}\label{ux4e09ux6539ux9769ux80fdux529bux4e0e1813ux5e74ux52a8ux5458ux5206ux79bb}

【锚】哈登贝格1807年纲领中的45,000步／25,000骑／80,000后备是建议非实有；柏林国民军：\literalquote{}volunteers
made up only a small part of the Berliner Landwehr, about 13\% according
to the records of the council and the
commission\literalquote{}------摊派4,000步兵与576骑兵，6月底仍缺800人，市长请求减额被国王拒绝；7月基金支出127,374塔勒而房租附加税仅收40,570、约10\%欠缴。\literalquote{}改革保留\literalquote{}不等于\literalquote{}解放战争必然重演\literalquote{}。国债口径冲突未解：Murau（2023）\literalquote{}At
the end of 1806, the Prussian debt had been 53.5 million thaler. By
1820, it amounted to 180 million thaler of interest-bearing
bonds\literalquote{}（1820年合计217.2m＝180有息＋11.2纸币＋25.9省债，占4.25年财政收入）vs
克拉克转载片段的1806年35m、1810年66m；财政危机不必然产出代议制------1820年敕令名义上要求新债须经等级会议同意，\literalquote{}the
Seehandlung extensively used the loophole that had been created and
provided off-balance-sheet government financing\literalquote{}。

\subsubsection{四、S5终局（史实前提下的残普）}\label{ux56dbs5ux7ec8ux5c40ux53f2ux5b9eux524dux63d0ux4e0bux7684ux6b8bux666e}

残普（提尔西特疆界＋西里西亚完整，无1815年西扩）1848年约620--810万人口，常备5.6--9.7万、危机总动员12.4--21.8万、各战区野战合计8.0--15.2万（显式假设算术）；普鲁士不自动取得史实1834年关税领导权------它依赖1815年西部领土与1831年黑森加入形成的通海关税屏障。A卡1815--30单一主战区危机野战9--14万与残普8.0--15.2万在单一战线重叠：\textbf{\literalquote{}奥强于普\literalquote{}的实质不在单场会战兵力，而在多战区并行能力与持续养兵的财政深度}。红队要求的\literalquote{}完整普鲁士＋北德主持权\literalquote{}情形见第六十六章。

\begin{center}\rule{0.5\linewidth}{0.5pt}\end{center}

\subsection{第十九章　南德诸邦：巴伐利亚、符腾堡、巴登------忠诚的价格}\label{ux7b2cux5341ux4e5dux7ae0-ux5357ux5fb7ux8bf8ux90a6ux5df4ux4f10ux5229ux4e9aux7b26ux817eux5821ux5df4ux767bux5fe0ux8bdaux7684ux4ef7ux683c}

\subsubsection{一、里德机制：忠诚变量是\literalquote{}谁担保我的吞并所得\literalquote{}}\label{ux4e00ux91ccux5fb7ux673aux5236ux5fe0ux8bdaux53d8ux91cfux662fux8c01ux62c5ux4fddux6211ux7684ux541eux5e76ux6240ux5f97}

【锚】里德条约（1813年10月8日）公开第4条：\literalquote{}Seine Majestät der Kaiser von
Österreich garantirt sowohl in Seinem Namen als im Namen Seines
Alliirten Seiner Majestät dem Könige von Bayern den freien und ruhigen
Besitz, so wie die volle Souveränetät über alle Staaten, Städte,
Domainen und Festungen, in deren Besitz Seine Majestät sich vor dem
Anfange der Feindseligkeiten befunden
hat.\literalquote{}秘密第5条：巴伐利亚野战至少36,000人（奥地利150,000）。忠诚终结的直接机制＝奥地利提供了替代性的领土担保。拿破仑有意让巴伐利亚的领土所得主要取自哈布斯堡，使其永久绑定法国（HdBG）；法国战运一转，巴伐利亚即暴露于奥地利复仇。

【锚】1808年5月巴伐利亚宪法第一条把邦联成员资格写进宪法：\literalquote{}§. 1. Das
Königreich Bayern bildet einen Theil der rheinischen
Föderation.\literalquote{}其动机（Weigand／HLB）一为抢先阻止拿破仑的邦联统一Fundamentalstatut、二为整合新旧领土；宪法许诺的国民代表机关从未召集。

血税实绩：巴伐利亚1812年约33,000人参加俄国远征（HdBG／Grab，可能同源）；符腾堡\literalquote{}an
dem Feldzug gegen Rußland 1812/13, aus dem von 15.800 Württembergern nur
300 zurückkehrten, bis zur Leipziger Völkerschlacht am 16. bis 18.
Oktober 1813. Erst im Vertrag von Fulda (1813) wechselte Friedrich zu
den Alliierten
über\literalquote{}------符腾堡直到莱比锡仍在法方作战，转向晚于巴伐利亚。1806年邦联配额：巴30,000／符12,000／巴登8,000。

【锚】Weis（1983）：\literalquote{}Die Gesamt-Nettoeinnahme des bayerischen Staates
durch die Säkularisation betrug 3,7 Millionen
fl.\literalquote{}------世俗化的财政净收益微小，经常性收支扣除养老金、堂区再拨与承接修道院债务后为赤字，唯一大额是1812年前变卖国有化不动产约二千万盾的一次性收入；对比1804年国债1.07亿盾、1811年1.18亿盾；购拜罗伊特1808年起须向拿破仑支付巨款（Weis\literalquote{}20
Millionen fl.\literalquote{} vs
HdBG\literalquote{}23,000,000法郎≈一年国入\literalquote{}，两口径并列）。蒂罗尔：1806--08税负+20\%、1808宪法废其特别地位并适用征兵、反教会措施→1809年四月全区起义，需法巴萨意联军约三万人秋季再征服；但HLB指出约瑟夫二世同类政策1780年代已几乎引发起义，1814年奥地利收回后并未废除巴伐利亚改革------这是现代国家与特权山地共同体的结构冲突，非法国体系独有。

【锚】Kopsidis--Bromley（2016，转引Fehrenbach 1983／Berding 1973）：\literalquote{}it
was Napoleon himself who pushed the \textquotesingle seigniorial counter
revolution\textquotesingle{} which preserved the seigniorial system in
southwest Germany until the revolution of
1848\literalquote{}------领主特权重定义为私产并入民法典；1815年邦联条例第14条追认之，且明定1807年巴伐利亚敕令为全德处理陪臣贵族的Basis
und
Norm。法国胜利情景下的南德社会结构＝法典外壳＋领主制内核＋陪臣贵族妥协，与史实同构。

\subsubsection{二、总判与四条件}\label{ux4e8cux603bux5224ux4e0eux56dbux6761ux4ef6}

【推·中高】第三德意志是法国霸权体系内成本效益最高的支柱：自带行政、自征税、自供兵，零法国驻军成本。持久的四个合取条件：担保独占（奥地利不得成为替代担保人）、血税≤阈值（年动员≈1\%人口且无灭军级损失）、主权底线（不吞并、不强推邦联宪法）、经济通道（封锁不毁农产与过境贸易）。S1／S2／S5下四条件可同时成立；\textbf{S6吞并南德为高置信负收益}：33k即时盟军换成自行承担60--90k驻军＋十年负净供兵。falsifier（\literalquote{}忠诚只系于法军胜利\literalquote{}）部分成立但不完整：1809年法国胜负未卜时三邦均出兵无动摇；1813年吕岑--包岑的胜利也未阻止滑出------真实触发器是\literalquote{}替代担保人出现∧血税超限（1812灭军）\literalquote{}的合取。B7宽俄版（避免1812式灭军）是支柱持久的关键前置。

1815--48：南德立宪与关税自组织是内生行动线（1808宪法、1815年弗里德里希主动召集等级会议、1828年巴符关税同盟均非外力强加）；法国胜利只改变其形态。若法国关税墙不降（史实对盟邦工业品高关税------贝格案例），南德在墙外自组织次级关税区；若普鲁士被弱化，该同盟以巴伐利亚为最大成员＝经济第三德意志，其对法离心力随工业化上升（中）。

\textbf{红队修正}：G1\literalquote{}担保独占／所得取自哈布斯堡\literalquote{}的条件在无普雷斯堡的路径B下只剩1803年教产部分；在S2本义下（普雷斯堡成立）直接适用。第五十章处理。

\begin{center}\rule{0.5\linewidth}{0.5pt}\end{center}

\subsection{第二十章　北中德意志：萨克森、威斯特法利亚、贝格、达尔贝格、汉萨与汉诺威}\label{ux7b2cux4e8cux5341ux7ae0-ux5317ux4e2dux5fb7ux610fux5fd7ux8428ux514bux68eeux5a01ux65afux7279ux6cd5ux5229ux4e9aux8d1dux683cux8fbeux5c14ux8d1dux683cux6c49ux8428ux4e0eux6c49ux8bfaux5a01}

\subsubsection{一、邦联宪法流产的主因是南德主权否决}\label{ux4e00ux90a6ux8054ux5baaux6cd5ux6d41ux4ea7ux7684ux4e3bux56e0ux662fux5357ux5fb7ux4e3bux6743ux5426ux51b3}

【锚】Karl Beck（1890，据Eberstein遗档；本轮只得开篇长段）：\literalquote{}dass aber
diese Bemühungen weniger an dem bösen Willen Napoleons scheiterten...
sondern vor allem an den Gegenbestrebungen der beiden mächtigsten
Fürsten des Rheinbunds, der Könige von Bayern und Würtemberg. Napoleon
musste aber auf diese Fürsten Rücksicht nehmen, weil sie ihm das
Haupt-Kontingent zu seinen Heeren
lieferten.\literalquote{}1806年10月16日筹备会因巴伐利亚与符腾堡使节缺席流产；邦联中央议事会从未开会------\literalquote{}A
central diet, which was supposed to regulate relations among members and
negotiate with Napoleon, never met.\literalquote{}（Grab
p.94）。法国的角色是\literalquote{}不强制\literalquote{}而非\literalquote{}否决\literalquote{}。制度化置信不上调，但约束内生于法国对南德供兵的依赖：S2／S5兵源充裕的世界存在\literalquote{}窄版Fundamentalstatut\literalquote{}（仲裁＋配额＋陆岚保障、无共同市场）路径，置信中低→中。

\subsubsection{二、并合承受度的排序}\label{ux4e8cux5e76ux5408ux627fux53d7ux5ea6ux7684ux6392ux5e8f}

【锚】贝格是北德唯一主动请求并入法国的单位（为得市场准入），被拒；否决力量含法国左岸保护主义利益集团------1810年9月16日科隆商会致Ladoucette书（AN
F12 549--550，Schmidt 1905附录E转录）：\literalquote{}nous sommes persuadés que tous
regarderont la réunion du grand-duché de Berg comme une des plus grandes
calamités qui puissent affliger leur
industrie.\literalquote{}贝格出口从封锁前六千万法郎跌至1811年一千一百万（Grab
pp.97--98）。

【锚】萨克森是大陆体系最大的工业赢家（Heckscher pp.302--306：\literalquote{}Of all
manufacturing countries on the Continent there is scarcely one which
developed so powerfully under the Continental System as
Saxony.\literalquote{}------开姆尼茨1801年双厂、莱比锡集市1810年库存65.5m法郎为当时估值）；而帝国驻邦联公使Bacher
1810年10月2日的报告本身已论证：强迫萨克森入关税警戒线＋征原棉税将使厄尔士山区赤贫化、工人外流奥地利、法国零收益。\textbf{S6／S3对萨克森的延伸在同时代官方内部已被否定。}

威斯特法利亚破产是结构性的（赠地半数领地＋养法军12,500＋债1.4--2亿；六年征兵七万＋志愿三万／二百万人口为藩属人均之冠；入俄2.2万生还约1,500）。汉堡1813：\literalquote{}Hamburg
wurde von den Franzosen mit einer Zwangskontribution über 48 Millionen
Françs (=25 Millionen Bancomark)
bestraft.\literalquote{}------6月8日帝国敕令宣布汉堡\literalquote{}法外\literalquote{}三月、三十富商为质、11月5日查封汉堡银行库存白银7,489,343马克·班科、年底逐出约两万无存粮居民、焚郊区约八千户；对照Marion
p.321汉萨三省1812年预期毛税仅38.791m------\textbf{罚没一次≈年税目标的1.3倍}，帝国在汉萨的边际财政形态已从征税退化为掠夺。

【锚】1809年多恩贝格起义的旗帜是旧帝国黑双头鹰（\literalquote{}Sieg oder Tod im Kampf
für das
Vaterland\literalquote{}），目标是复辟选侯旧制而非民族目标；威斯特法利亚自有军队（含多恩贝格本部两连猎兵）拒绝叛变并一击镇压之。北德抵抗定理：只要法军或藩属正规军在场且忠诚，一切暴动数日内被平；断裂条件外生（野战军失败或撤出）------1813年实测。萨克森国王弗里德里希·奥古斯特的反应函数（ADB
Flathe）：合法主义＋风险极小化；1813年即使在体系最弱时刻也只走到对奥秘约（4-20）＋托尔高中立令，吕岑一胜即在收到威胁信前主动回摆；莱比锡阵前士兵求离法而国王以\literalquote{}尔等之职责\literalquote{}答。

\textbf{排序}：法兰克福≈奥尔登堡≈梅克伦堡（政治可并、净值≈零或负、奥尔登堡有俄国外部成本实测）＜汉诺威（可并但兵源泄英）＜贝格（唯一并入优于藩属的单位，但被法国保护主义否决）＜汉萨（可并而财政自败、1813退出弹性极高）＜威斯特法利亚（并入＝自拆样板招牌）＜萨克森（不可并，帝国自身官僚已论证）。\textbf{S6在北德的可行边界即史实1811年边界。}北德是顺从带非支撑带；S5下形成\literalquote{}法制西北
vs 旧制东北\literalquote{}的双德意志结构轴；萨克森是体系健康度的领先指标。

\begin{center}\rule{0.5\linewidth}{0.5pt}\end{center}

\subsection{第二十一章　荷兰与瑞士：金融、殖民地、征兵与调停制样板}\label{ux7b2cux4e8cux5341ux4e00ux7ae0-ux8377ux5170ux4e0eux745eux58ebux91d1ux878dux6b96ux6c11ux5730ux5f81ux5175ux4e0eux8c03ux505cux5236ux6837ux677f}

\subsubsection{一、荷兰：可兼并，不可消化}\label{ux4e00ux8377ux5170ux53efux517cux5e76ux4e0dux53efux6d88ux5316}

【锚】Marion（IV，p.325，本轮亲核）：\literalquote{}la dette considérable, 78
millions, fut répudiée pour les deux tiers mais inscrite pour un tiers à
la charge du Trésor français, qui eut alors à payer outre 62 millions de
rentes françaises 26 de rentes
hollandaises\literalquote{}------年金78m否定三分之二，保留26m入法国国库，而法国自身永久年金仅62m：兼并荷兰使法国公债名目负担增四成，且烧毁了阿姆斯特丹的金融合作。拿破仑1810年7月25日致莫利安：\literalquote{}10
millions de florins pour la guerre et 6 millions de florins pour la
marine... la Hollande me rendit, pour payer la dette, 28 millions, ce
qui ferait une soixantaine de
millions\literalquote{}------六千万法郎／年是目标非实收；同信附言当晚召见戈赫尔：帝国需要荷兰技术官僚来消化荷兰。1810年7月致勒布伦：\literalquote{}désirant
conserver l\textquotesingle administration hollandaise, qui est plus
économique que la
nôtre\literalquote{}------兼并的行政理由不成立，动机纯为封锁、兵源与债务处置。

【锚】Joor（Soldiers, Citizens, Civilians，2009，本轮全文）：\literalquote{}the
research identified over 580 collective actions in the category of
\textquotesingle unrest\textquotesingle... included more than 80
revolts... the Emperor attempted to maintain a large army in the area,
averaging some 15,000
men\literalquote{}------1806--13年逾580起骚动、逾80起起义（≥20人暴力对抗）；动因：军事／征兵35\%、财税＋封锁26\%、反政权24\%、宗教仅4\%；起义的三分之二发生在并合后3.3年内；西荷兰常驻法军约15,000至1813年末；驻军一撤三周蒸发（1813年11月实测）。falsifier\literalquote{}骚乱主因征兵\literalquote{}成立，但\literalquote{}兼并可行性上调\literalquote{}仅在\literalquote{}兼并退化为藩属（不征兵不削债软海关）\literalquote{}时成立------那恰好取消兼并的全部财政军事目的；真正被上调的是B6b软兼并／藩属化路径。

\subsubsection{二、瑞士：不必并，也不能并}\label{ux4e8cux745eux58ebux4e0dux5fc5ux5e76ux4e5fux4e0dux80fdux5e76}

【锚】1803年9月27日法瑞军事协约：\literalquote{}La République Française entretiendra a
son Service seize mille hommes de Troupes Suisses... enrôlés librement
et
volontairement\literalquote{}------四团×4,000＋每团1,000补充营、限欧陆使用（Art.18）、瑞士本族军事司法（Art.19）、法国十日内回防承诺（Art.24）、有效期二十五年（Art.25）。瑞士的血税是条约化商品，不是征兵义务。《调停法》：\literalquote{}Sur
15.203 hommes, le contingent de Berne sera de
2292\literalquote{}------联邦兵额15,203、联邦经费仅490,507瑞郎／年、每邦常备≤200人、会期≤一月；新邦兵额占总额50.9\%；属地废除、法律平等、内部自由贸易三项1798年成果被联邦法锁定。调停模式以零驻军零财政成本换得十年顺从。

【锚】瓦莱兼并链条：1810年11月3日动机信（辛普隆公路花费一千五百万法郎不能为小国牺牲）------\literalquote{}mon
intention est de réunir le Valais à la
France\literalquote{}→11月12日敕令→11月14日贝尔蒂埃入锡永→税额定六万法郎≈每居民一法郎→拿破仑自认此省无钱可取（\literalquote{}on
ne tirerait point
d\textquotesingle argent\literalquote{}）；1813年12月26日省长弃逃。纯走廊型兼并、财政净负；提契诺占领、海关督察＝政策漂移侵蚀最优安排的实证。瑞士团入俄7,000--10,000／生还700--1,300（三源分歧）。

\textbf{总判}：荷兰＝可兼并不可消化（镇压永远可行、净收益永远为负、1813实测三周蒸发）；瑞士＝不必并也不能并（兼并即同时销毁产出并激活山地起义链）。行省化A层在此两单位的边界：到荷兰为止且得不偿失，到瑞士为止根本不应进入。正确的问法是\literalquote{}并了之后还能得到什么\literalquote{}。1815--48对照组：联合尼德兰王国（比方旧债约32m盾vs荷方约二十亿；340万vs200万人口席位55:55；1819--23语言宗教敕令）在优于法国的兼并条件下十五年即断；Sonderbund
1847（25--26天、93死510伤、战后无肢解无处决→1848宪法）以低烈度自愈。兼并模式的凝聚力随时间衰减，调停模式的凝聚力增长------这是P8的敏感轴。

\begin{center}\rule{0.5\linewidth}{0.5pt}\end{center}

\subsection{第二十二章　丹麦-挪威与瑞典：波罗的海门户与挪威价码}\label{ux7b2cux4e8cux5341ux4e8cux7ae0-ux4e39ux9ea6-ux632aux5a01ux4e0eux745eux5178ux6ce2ux7f57ux7684ux6d77ux95e8ux6237ux4e0eux632aux5a01ux4ef7ux7801}

\subsubsection{一、封锁的可执行上限是漏斗，不是闸门}\label{ux4e00ux5c01ux9501ux7684ux53efux6267ux884cux4e0aux9650ux662fux6f0fux6597ux4e0dux662fux95f8ux95e8}

【锚】戴维：\literalquote{}In 1805, 11,537 merchant ships passed through the Sound,
over half going to
Britain.\literalquote{}------护航单队400余艘，哈诺集结约500艘再启航；相对于丹挪炮艇单次最大截获47艘（1810年7月，丹麦炮艇vs挪威双桅舰两口径），拦截率上限在个位数百分比。升格所需的三个物理前提（占瑞典西南海岸／俄真诚闭港／海峡制海）从未同时成立。

【锚】1811年5月23日索马雷兹向海军部报告瑞军总司令塔瓦斯特在胜利号上的密谈：瑞典承诺不行任何敌对行动，\literalquote{}the
appearance of any hostile measures was only intended for Demonstration,
and in order to elude the Vigilance of French
Spies\literalquote{}------1810年11月对英宣战后的\literalquote{}假战争\literalquote{}是国家授权的双面政策。巴顿（Barton）：\literalquote{}Swedish
imports from Britain thus amounted in 1810 to nearly 7,000,000 riksdaler
in
value\literalquote{}；美旗船入俄峰值1811年约225艘。falsifier\literalquote{}哥德堡贸易已使封锁名存实亡\literalquote{}成立：S3的北欧环节应编码为漏斗，漏率对法国政策严格递增。

\subsubsection{二、S3的第一受害者是法国的盟友}\label{ux4e8cs3ux7684ux7b2cux4e00ux53d7ux5bb3ux8005ux662fux6cd5ux56fdux7684ux76dfux53cb}

【锚】丹麦财政寿命表：1807年参战→1812年券量1.27亿（五倍）→1813年1月5日敕令六比一重组（Rigsbank、4,600万上限、6\%不动产银税）------\literalquote{}kurantdalerens
kurs var faldet til kun 6 \% af dens pålydende værdi i
sølv\literalquote{}。挪威平年四分之一口粮靠进口，进口粮四分之三来自丹麦；1812--13年饥荒按欧洲标准仅马德里1811--12可比------副总督警报\literalquote{}a
quarter of the population would probably perish if things did not
improve rapidly\literalquote{}（Glenthøj \&
Ottosen，本轮全文）；瑞典战争党明文把饥荒当促统武器（阿德勒斯帕雷1811）。\textbf{S3越严，英瑞对挪威的绞索越紧，丹麦联盟越亏损。}

【锚】瑞典是可购买的，但价码在丹麦手里：1811年2月卡尔·约翰向拿破仑秘密求购------\literalquote{}Karl
Johan entered into secret negotiations with Napoleon for an alliance in
return for subsidies and increased territory, above all in Norway. The
emperor categorically opposed such an
arrangement.\literalquote{}（Barton）；1812年3月拿破仑改售芬兰只帮其抬价。法国买瑞典必卖丹麦的结构互斥无法用钱解开。基尔条约第4条以\literalquote{}与瑞典联合的王国\literalquote{}而非行省形式转让挪威，\literalquote{}Grønland,
Færøerne og Island ikke deri
indbegrebne\literalquote{}；第6条国债按人口与收入比例分担。芬兰模式（波尔沃保证＋精英赎买＋零汲取------\literalquote{}to
give this people a political existence in such a way that it does not
regard itself as conquered by Russia but on the contrary as attached to
it by the bonds of its evident
interests\literalquote{}）是北欧存在的真实招抚模板，但与法国并合省的汲取型模式财政不兼容。

\textbf{总判}：北欧是法国政策的镜子；最优是\literalquote{}少要求\literalquote{}（名义封锁＋许可证常态＋不惩瑞典→丹麦可维系）；S5挪威归属主线＝留丹麦联合王国联邦化（中），备选法版基尔／独立缓冲／斯堪的纳维亚联合（低），任何安排付\literalquote{}宪法折扣\literalquote{}；芬兰1809年安排在一切情景稳定（高）。未决：挪威分年超额死亡率、Heckscher原表、Saumarez
Papers原件、基尔\literalquote{}一百万\literalquote{}两口径（对丹军团补贴vs秘密补偿）。

\begin{center}\rule{0.5\linewidth}{0.5pt}\end{center}

\subsection{第二十三章　意大利王国：欧仁、米兰精英、意大利军与\literalquote{}第二法国\literalquote{}}\label{ux7b2cux4e8cux5341ux4e09ux7ae0-ux610fux5927ux5229ux738bux56fdux6b27ux4ec1ux7c73ux5170ux7cbeux82f1ux610fux5927ux5229ux519bux4e0eux7b2cux4e8cux6cd5ux56fd}

\subsubsection{一、王冠已在文书中腾出给亲生子}\label{ux4e00ux738bux51a0ux5df2ux5728ux6587ux4e66ux4e2dux817eux51faux7ed9ux4eb2ux751fux5b50}

【锚】1810年1月3日拿破仑致阿尔迪尼：\literalquote{}Par le statut du royaume, je
l\textquotesingle ai appelé au trône en cas que je n\textquotesingle aie
point d\textquotesingle enfant; mais, comme prince italien, il lui faut
un apanage\ldots{} cela formera l\textquotesingle apanage de la branche
du prince de Venise dans tous les
temps.\literalquote{}------欧仁由无子时继承人改设为带采邑的\literalquote{}威尼斯亲王支\literalquote{}。1810年2月17日元老院法令罗马并入法国，第7条皇太子称罗马王，第10条皇帝十年内罗马加冕；3月1日咨文定法兰克福归欧仁世袭，意大利仅留其\literalquote{}照料与行政\literalquote{}。\textbf{意大利王冠在同期文书群中腾出给未来的亲生子}------这有直核；那不勒斯的合并没有同期文书，仅存回溯层（《通信集》XXX印p.549的\literalquote{}次子\literalquote{}摄政方案）。

\subsubsection{二、汲取实绩与两个崩溃点}\label{ux4e8cux6c72ux53d6ux5b9eux7ee9ux4e0eux4e24ux4e2aux5d29ux6e83ux70b9}

【锚】Grab（2003，第十章）：\literalquote{}the government drafted 155,000 men between
1802 and 1814, and expanded the Italian army to 70,000 men by
1812\literalquote{}；俄役27,000人仅约1,000生还；驻王国25,000法军维持费＋意军开支通常占预算50\%；1802--11岁入近翻倍；峰值疆域89,600平方公里／670万人口。1809年双起义------威尼托（奥侵背景＋恢复威尼斯共和国诉求，数千农民威胁维琴察／特雷维索）与艾米利亚-罗马涅（反新磨坊税，盗匪与逃兵领导，一度威胁博洛尼亚）------\literalquote{}under
the influence of the revolt, Prina decided to revoke the milling
tax\literalquote{}。税负冲击×征兵冲击×合法性缺口的三元叠加才出系统性暴动：这是能力上限表的第一个崩溃点。

第二个崩溃点是1814年4月20日米兰的\literalquote{}雨伞日\literalquote{}：普里纳被私刑，\literalquote{}abilmente
alimentato da quanti -- in prevalenza ex nobili -- miravano a far
naufragare la candidatura di Eugenio di Beauharnais a sovrano di un
nuovo Stato cisalpino
indipendente\literalquote{}（DBI）------旧贵族刻意煽动，以使欧仁登基新独立国的候选破产；财政部档案与普里纳私人档案全毁。同期行政机器在库空与战败下仍效忠到底。崩的是王朝纽带与精英协调，不是国家机器------与F5的\literalquote{}继承时刻精英协调\literalquote{}命题同构。

【锚】经济附庸化（Heckscher
pp.296--301）：1810年三重规定------羊毛禁入、过境禁令、生丝只准出口里昂；\literalquote{}la
France avant
tout\literalquote{}信（1810年8月23日）；王国外贸约半对法；并入帝国的领土也不自动获法国市场待遇------荷兰与汉萨是\literalquote{}proselytes
of the gate\literalquote{}，\literalquote{}they were to be left without participation in the
advantages of the French
market\literalquote{}。\textbf{行省化分支的\literalquote{}经济平权\literalquote{}论据被同期实践证伪。}

\subsubsection{三、三分支}\label{ux4e09ux4e09ux5206ux652f}

{\def\LTcaptype{none} % do not increment counter
\begin{longtable}[]{@{}
  >{\raggedright\arraybackslash}p{(\linewidth - 4\tabcolsep) * \real{0.3333}}
  >{\raggedright\arraybackslash}p{(\linewidth - 4\tabcolsep) * \real{0.3333}}
  >{\raggedright\arraybackslash}p{(\linewidth - 4\tabcolsep) * \real{0.3333}}@{}}
\toprule\noalign{}
\begin{minipage}[b]{\linewidth}\raggedright
分支
\end{minipage} & \begin{minipage}[b]{\linewidth}\raggedright
在位期
\end{minipage} & \begin{minipage}[b]{\linewidth}\raggedright
继承后
\end{minipage} \\
\midrule\noalign{}
\endhead
\bottomrule\noalign{}
\endlastfoot
C 分治藩属（北意王国＋那不勒斯＋法国并合区） &
\textbf{0.55--0.65（主线）} & 0.30--0.40 \\
B 罗马王／次子统一王国（北＋中意先行，那不勒斯条件性并入） & 0.20--0.30
& \textbf{0.35--0.45（升为主线）}；含那不勒斯全统一0.10--0.18 \\
A 行省化\literalquote{}第二法国\literalquote{} & 0.10--0.20 & 0.10--0.15 \\
\end{longtable}
}

C是全部同期文书指向的唯一稳态，断点是继承时刻藩属王冠的再分配；B的通道是1805年法令第3条分冠的自动执行＋皇子1829年成年（拿破仑活到该日≈0.50--0.65）；A是惩罚态非设计态------逆全部同期制度承诺，抽走米兰\literalquote{}都城租金／官职垄断／本国军旗\literalquote{}三合流支柱。1815--48的校准器是奥治伦巴第-威尼西亚：\literalquote{}In
all, well over one-half of the Lombard and Venetian troops remained
loyal to
Austria\literalquote{}（Sondhaus）------三轮起义（1821／1831／1848）＋1848年30--35k伦-威籍奥军过半仍忠（约15,000开小差）＋行政机器全程运转；镇压需拉德茨基级七万驻军（Hilleprandt
1867转引）。征服者直辖\literalquote{}可运行、可镇压、不可赢得\literalquote{}。

\textbf{红队修正}：这些概率以史实王国（含威尼托1806、马尔凯1808、特伦蒂诺1810）为对象；1805年分冠条款的触发条件是\literalquote{}外军撤出那不勒斯／马耳他／科孚\literalquote{}。第六十五章按两条路径分别重估。未决：Rath
1941（付费墙）、Du
Casse十卷（E2：1813--14是否正式许欧仁世袭）、征兵155k与200k+两口径、磨坊税\literalquote{}撤销vs减半\literalquote{}、岁入82→144m里拉。

\begin{center}\rule{0.5\linewidth}{0.5pt}\end{center}

\subsection{第二十四章　那不勒斯与西西里：缪拉、卡拉布里亚与英国的岛}\label{ux7b2cux4e8cux5341ux56dbux7ae0-ux90a3ux4e0dux52d2ux65afux4e0eux897fux897fux91ccux7f2aux62c9ux5361ux62c9ux5e03ux91ccux4e9aux4e0eux82f1ux56fdux7684ux5c9b}

\subsubsection{一、藩属的压力测试极限}\label{ux4e00ux85e9ux5c5eux7684ux538bux529bux6d4bux8bd5ux6781ux9650}

【锚】拜永条约（1808年3月30日）对缪拉的义务：\literalquote{}obliged the king of Naples
in time of war to contribute 18,000 infantry, 3,000 cavalry, and a train
of twenty-five pieces of artillery to France, the costs to be met
entirely by the Neapolitan treasury wherever they
served\literalquote{}（Davis；同书p.252另作16,000／2,500／20，两处并列）；代养驻境内法军；造并六艘战列舰；执行封锁；秘密条款以缪拉在法全部财产作担保。税收：\literalquote{}from
57,000,000 lire in 1808 to 61,000,000 in 1809, and 73,000,000 in
1812\literalquote{}；1812年军费占总预算80\%；1813年战争预算57m里拉≈1806年全部岁入；已无领地可卖。缪拉军（不含境内法军）：\literalquote{}From
29,000 men, exclusive of French troops in the kingdom, which it numbered
in July 1812, it had risen to 35,000 in October, and to 46,000 a year
later\literalquote{}（Johnston）；1815年野战约35,000＋5,000骑＋60炮（御前会议自称八万；Castellano账面17,000与Johnston两口径并存）；托伦蒂诺败于军需崩溃。

【锚】下卡拉布里亚一省：\literalquote{}between November 1806 and March 1811... 5,421
bandits had surrendered or been captured, killed, or executed in the
province of Lower Calabria
alone\literalquote{}（王家宪兵队1812年报告）；Manhès自报1810--11冬杀／处决900人。镇压峰值成本≈3.0--4.5m
ducats／年（900法郎／兵年×兵力区间的推断值，非档案实额），约为王国税收（≈13m
ducats）的23--35\%------严格版falsifier（\literalquote{}超过税收\literalquote{}）不成立；修订版成立：卡拉布里亚两省实收≤0.45m
ducats／年而吸收≥3m／年，\textbf{净流出为自身产出的6--8倍}；全国军事总负担1808年已达税收三分之二、1812年80\%。A层并合的边际不可持续；B层藩属在低需索档位可持续。

\subsubsection{二、英国的岛}\label{ux4e8cux82f1ux56fdux7684ux5c9b}

【锚】英国对西西里补贴：1805年密约£170,000一次性＋年£150,000（参战升£300,000）；\literalquote{}to
pay a subsidy of £300,000 annually, backdated to September
1805\literalquote{}（1807年3月23日条约）；后增至£400,000／年，实际流向流亡者年金与王后间谍网；1812年9月条约7,300西西里师归英指挥、费用从补贴扣（Gregory）。1812年西西里宪法：两院制、议会独占立法、司法独立、陪审制、废封建权利与限嗣继承（有补偿）------\literalquote{}only
the literate had the right to vote and the right to sit in parliament
(which meant about 15,000 men out of a population of a
million-and-a-half)\literalquote{}；西西里宣布独立、王不得离岛。英国吞并西西里的选项：代理人层反复提出（Moore
1807、Drummond 1806、Howick 1807未发训令、Bentinck 1814\literalquote{}Sogno
filosofico\literalquote{}），内阁层每次否决------\literalquote{}measures which would have in fact
rendered Sicily a province of Great
Britain\literalquote{}（Rosselli）；英国模式是基地＋商业特权非领土主权（1816年商约10\%关税减让即变现）。以西西里为跳板：防御近绝对，进攻上限师级袭击＋牵制------1810年一次威胁集结锁死法-那四万人一夏；西西里＝不沉的航母，但弹药库在别处。

\subsubsection{三、1820--21三定理}\label{ux4e09182021ux4e09ux5b9aux7406}

烧炭党：1810年那不勒斯首个会所，1813年卡拉布里亚首个会所几可肯定由省督Briot创建；1815年估≥30,000（纯属猜测），1816年保王将军Nunziante称仅卡拉布里亚就有逾五万（严重夸大）；1820年革命由萨莱诺烧炭党1818年议会预先策划、以西班牙宪法为纲领。由此三定理供S5直接采用：\textbf{藩属立宪周期律}（法式行政国家＋过大军官团＋财政紧缩→霸权松动即政变，周期5--7年）；\textbf{外军平叛依赖律}（法国将扮演1821年奥地利的角色）；\textbf{西西里分离律}（大陆每次立宪，岛内同步分裂）。

\textbf{主线}：南意在法国霸权工程中唯一可持续的位置是B层法典化藩属＋海军义务承包＋封锁豁免区；其史实失败是需索强度的失败而非制度设计的失败；对帝国净贡献≈0，价值在位置与否决权。未决：Finley
1994未获（以Davis ch.11＋Johnston＋宪兵1812年报替代）；Le
Brethon刊本止于1810年8月；镇压成本无档案实额。

\begin{center}\rule{0.5\linewidth}{0.5pt}\end{center}

\subsection{第二十五章　教廷与欧洲诸教会：正统性的跨国否决者}\label{ux7b2cux4e8cux5341ux4e94ux7ae0-ux6559ux5ef7ux4e0eux6b27ux6d32ux8bf8ux6559ux4f1aux6b63ux7edfux6027ux7684ux8de8ux56fdux5426ux51b3ux8005}

\subsubsection{一、否决权的三层}\label{ux4e00ux5426ux51b3ux6743ux7684ux4e09ux5c42}

【锚】1813年枫丹白露专约签署数周后，教皇撤回：\literalquote{}nous regardons comme nuls
les deux écrits, tant celui qui a été fait à Savone que celui du 25
janvier\literalquote{}（Pacca回忆录，1833年法文版，本轮亲核）------撤回函援引帕斯卡尔二世1111年先例。教皇个人可被隔离与疲劳战术压弯（1811年萨沃纳笔记、1813年1月25日两次实证），但黑衣枢机网络两次均于数周内促成撤回；撤回的动力是教会法先例逻辑（防各国宫廷援引）而非信众反弹------无信众反弹是因为帝国封锁了消息。\textbf{否决权的实质是使被迫签署丧失终局性，而不是拒签的能力。}三层：教皇个人可压服＜枢机-教会法机构不可压服＜信众约束弱且需与征兵／占领合取。

【锚】拿破仑亲述的备选方案：1801年谈判破裂时当面威胁以亨利八世方式改宗------\literalquote{}Si
Henri VIII, qui n\textquotesingle avait pas la vingtième partie de ma
puissance, a su changer la Religion de son
pays\literalquote{}（Consalvi回忆录，1866）；圣赫勒拿回溯自述曾在天主教与新教间择一。国教分裂选项在决策空间内存在，但被其自身否决。

【锚】1808年春教廷已正式书面提出：\literalquote{}s\textquotesingle est déjà déclarée
prête à fermer, dans la guerre actuelle, ses ports aux
Anglais\literalquote{}------本次战争内关闭对英港口＋武装守海岸，但拒绝加入进攻性联盟。可谈判空间：封锁的事实执行可换，盟约形式与世俗主权不可换。

\subsubsection{二、比利时的双实证与西班牙的反转}\label{ux4e8cux6bd4ux5229ux65f6ux7684ux53ccux5b9eux8bc1ux4e0eux897fux73edux7259ux7684ux53cdux8f6c}

比利时是专约和解可行性与脆弱性的双实证：1801--09年八年有效（多数教士宣誓、征兵合规改善）；1809年绝罚＋囚教皇后逆转------\literalquote{}a
growing segment of the clergy refused to pray for him following the
papal
excommunication\literalquote{}（Grab）；根特、图尔奈司铎集体拒认新主教，数百司铎被捕。专约体制的忠诚是转授忠诚，抵押品是教皇本人的合作。西班牙教会动员的双重口径：Grab称教士在起义与洪塔中骨干；Esdaile反证\literalquote{}most
of the guerrilla bands that were formed in Catalonia were wholly
secular\literalquote{}------教会的角色在合法性供给而非军事指挥。\textbf{B4反转判定（中高）}：保留波旁＋联姻版下西班牙教会可挤压可合作------教会已在波旁治下接受Consolidación卖教产（1798--1808教皇简函合作，\literalquote{}Godoy
also alienated the Church by selling extensive ecclesiastical
properties\literalquote{}）；无王朝真空则无洪塔，无洪塔则教士无合法性罢工的载体。

\subsubsection{三、S5／S6三形态}\label{ux4e09s5s6ux4e09ux5f62ux6001}

7.1\literalquote{}第二个1801\literalquote{}和解均衡（中高）------构件全部历史存在过：教廷1808年已书面愿意事实关港；1813年撤回函的公开可让渡域；枫丹白露第5--9条含给教皇实权补偿条款。7.2武装加利刚（中）------可运转十年，二十年尺度官吏教会与地下教会双轨化（衰减速率无面板）。7.3教廷部门化（低）------罗马1809--14试点实证：行政存在，忠诚为零，\literalquote{}They
obeyed Pius VII\textquotesingle s orders to refrain from swearing
allegiance to the new
ruler\literalquote{}。教会是A层可扩性的第一证伪器；S6下不是引爆器而是腐蚀剂------不推翻帝国，但使天主教抵抗区的驻军账永久上浮。\textbf{和解是\literalquote{}长久\literalquote{}判据下性价比最高的单项投资}：教会只要罗马与授职权，比显贵年金、军队胜利、市场信用都便宜。最大行省化与教会和解不可兼得------与海权同构。1815--48复兴三引擎与极端教权主义四模板的移植判定：《Mirari
Vos》全文亲核------教廷自己会谴责\literalquote{}上帝与自由\literalquote{}路线。未决：7.1窗口概率依赖F1克制先验；Chadwick／Hales／Latreille／Schwarzfuchs／Broers
2002／Caiani
2021八专著未获（AA会话403），以d\textquotesingle Haussonville谱系＋Pacca／Consalvi一手回忆＋Grab全书＋Caiani
2019论文替代；西班牙教士总量两口径（20万vs14.8--17万）并列。

\begin{center}\rule{0.5\linewidth}{0.5pt}\end{center}

\subsection{第二十六章　西班牙（一）宫廷与国家：1807--08桌面上的全部选项与政体终局}\label{ux7b2cux4e8cux5341ux516dux7ae0-ux897fux73edux7259ux4e00ux5babux5ef7ux4e0eux56fdux5bb6180708ux684cux9762ux4e0aux7684ux5168ux90e8ux9009ux9879ux4e0eux653fux4f53ux7ec8ux5c40}

\subsubsection{一、三重超载的旧制机器}\label{ux4e00ux4e09ux91cdux8d85ux8f7dux7684ux65e7ux5236ux673aux5668}

【锚】财政硬数（Herr，Rural Change and Royal Finances）：\literalquote{}Vales in
existence in May 1808 \ldots{} 1,893,422,682\literalquote{}------vales
reales存量18.93亿rv（Consolidación基金报告，ANP AF IV
1608B，法方自审数），1808年3月贴水63\%；岁入1803--07仅505--605M
rv（不含美洲）对支出均800M；美洲汇入从1802--04年均193M跌至1805--07年28M；对法补贴每月六百万法郎经Ouvrard--Hope墨西哥汇票机器支付（rv月额Marichal
16M vs Herr
24M系法郎换算差，总装以6M法郎／月为准）；教产变现本土约1,300--1,650M
rv＋美洲15.4M
pesos------强制赎回毁墨西哥信贷体系，是1810年起义前紧张的根源。Marichal：\literalquote{}in
1809--1811, total income was 1,000 million reales, of which 60\% came
from the
Americas\literalquote{}；1812--14跌至5\%；无墨西哥白银则摄政府与加的斯议会撑不过1809--11。\textbf{任何西班牙政权的财政生死系于美洲忠诚×制海通道×加的斯--韦拉克鲁斯商人网三件套。}

\subsubsection{二、1807--08桌面上的十六项选项}\label{ux4e8c180708ux684cux9762ux4e0aux7684ux5341ux516dux9879ux9009ux9879}

【锚】枫丹白露条约（1807年10月27日＋29日军事附约）：葡萄牙三分------\literalquote{}el
sur, el «Principado de los Algarves», comprensivo de esta provincia y
del Alentejo, quedaría bajo el gobierno de Godoy, con el título de
«príncipe de los Algarves»\literalquote{}（La
Parra，戈多伊传）；北部予埃特鲁里亚太后，中部留置可与英换直布罗陀／特立尼达；西王称三国protector并将获\literalquote{}两美洲皇帝\literalquote{}衔；各出兵25,000步＋3,000骑。拿破仑拒公布、拒设界委使其执行空心化。

【锚】联姻回应链（La
Parra，费尔南多七世传）：费尔南多1807年10月11日亲笔求婚信→拿破仑11月4日收信无评论→11月13日致卡洛斯四世否认收信（ignoro
si
existe）→卡洛斯11月18日自提联姻→\textbf{1808年1月10日拿破仑书面\literalquote{}Consiento
de buen grado a este enlace. Pero V.M. debe comprender que ningún hombre
de honor querría aliarse con un hijo deshonrado por su
Declaración\literalquote{}}（西廷读为婉拒）→2月3日条件化为换埃布罗割地筹码→5月巴约讷九点方案第9点降为退位甜头（娶皇帝侄女）。联姻始终是工具非政策，每次\literalquote{}同意\literalquote{}都附当下不可执行条件且价码递增。\textbf{波旁立宪终局只在拿破仑改行B4b的前提下成立。}

【锚】cuestiones
proponibles（十八条，杜罗克／塔列朗口授伊斯基耶多，1808年3月5日达马德里）：占领边境诸省至总和平后一年→以葡萄牙全境换比利牛斯--埃布罗间诸省；实现后允联姻但若亲王不驯可废储另立一子；自认西班牙存在第三党\literalquote{}anarquistas,
cuyos designios se alargaban hasta el extremo de aspirar a una reforma
capital de la monarquía española, con semejanza, según unos, a la
Constitución inglesa y, según otros, a la Constitución
americana\literalquote{}。埃布罗方案拟由约瑟夫治理割取区，因约瑟夫拒换那不勒斯而阻。B4c埃布罗兼并＝史实1810--12已试验并失败的方案（埃布罗以北事实剥离1810年2月至1812年2月），坍缩为B4a。

\subsubsection{三、军事价值是二值函数}\label{ux4e09ux519bux4e8bux4ef7ux503cux662fux4e8cux503cux51fdux6570}

【推·中高】同盟态＝25,000远征军＋港口＋船厂＋补贴净流入；占领态＝−29至36万驻军（Oman两切片，第二章）；二值之间无中间态------1808--13年约瑟夫的半合法态收获两头最劣。\textbf{A层南界＝比利牛斯。}

\subsubsection{四、终局概率带（1848视点）}\label{ux56dbux7ec8ux5c40ux6982ux7387ux5e261848ux89c6ux70b9}

{\def\LTcaptype{none} % do not increment counter
\begin{longtable}[]{@{}
  >{\raggedright\arraybackslash}p{(\linewidth - 4\tabcolsep) * \real{0.3333}}
  >{\raggedright\arraybackslash}p{(\linewidth - 4\tabcolsep) * \real{0.3333}}
  >{\raggedright\arraybackslash}p{(\linewidth - 4\tabcolsep) * \real{0.3333}}@{}}
\toprule\noalign{}
\begin{minipage}[b]{\linewidth}\raggedright
终局
\end{minipage} & \begin{minipage}[b]{\linewidth}\raggedright
S2／S5（法胜＋B4b）
\end{minipage} & \begin{minipage}[b]{\linewidth}\raggedright
法在半岛军事全胜
\end{minipage} \\
\midrule\noalign{}
\endhead
\bottomrule\noalign{}
\endlastfoot
波旁绝对主义-大臣改革制（1815--30型延长） & \textbf{0.40--0.55} & --- \\
波旁立宪（钦定宪章／复召Cortes） & 0.30--0.45 &
波旁复辟（含立宪）0.30--0.45 \\
波拿巴替换／波拿巴立宪 & 0.05--0.15 & \textbf{0.35--0.55} \\
共和国 & ≤0.02 & 0.03--0.08 \\
割据长期化 & --- & 0.10--0.20 \\
\end{longtable}
}

共和硬上界≤0.08：需三条件合取（双王朝失信＋无王朝焦点的战争延长＋美洲共和样板回传），史实1868--73才凑齐；1848年前任何情景至多凑两项------这是迟发共和，不是不可能的共和。falsifier
1：1808年前本土自由派共和文书本轮与上轮均未发现；拿破仑方情报所称\literalquote{}美宪派\literalquote{}实指宪政限权派。若PARES／BNE出土1808年前本土共和纲领则改判。未决：热罗姆候选无同期文书；Fugier两卷未直读；巴斯克三省＋纳瓦拉战前税基缺；巴约讷洪塔首会65／75两说。

\begin{center}\rule{0.5\linewidth}{0.5pt}\end{center}

\subsection{第二十七章　西班牙（二）社会与抵抗：游击、洪塔、加的斯与政治文化}\label{ux7b2cux4e8cux5341ux4e03ux7ae0-ux897fux73edux7259ux4e8cux793eux4f1aux4e0eux62b5ux6297ux6e38ux51fbux6d2aux5854ux52a0ux7684ux65afux4e0eux653fux6cbbux6587ux5316}

\subsubsection{一、总起义有开关}\label{ux4e00ux603bux8d77ux4e49ux6709ux5f00ux5173}

【锚】\literalquote{}The Spanish uprising drew its character from the decapitation of
political authority brought about by Napoleon\textquotesingle s
sequestration of the Spanish
monarchy\literalquote{}------1808年2月法军夺庞普洛纳、圣塞巴斯蒂安、菲格拉斯、巴塞罗那四要塞后并无总起义；起义各城延迟数日、由显贵与下级军官的阴谋组引导。反向自然实验：1823年，\literalquote{}unlike
Napoleon\textquotesingle s invasion in 1808, the French met virtually no
popular resistance. The networks of Spanish elites who had mobilised the
people in 1808 could not (or would not) do the same in
1823\literalquote{}（Lawrence）------法军56,000＋西班牙保王军约35,000入境几乎未遇抵抗，尽管1822年2月国库收入88\%已用于军费。1796--1808法西同盟十二年（含特拉法加、月贡六百万法郎、拉罗马纳1.3万人赴丹麦）无任何反法起义；1793--95反法圣战限于边境且随和平终止。

\textbf{反应函数总律}：西班牙社会对法国政策的敏感轴只有三根------王朝符号、税与征发、教会；三轴齐触（B4a）则总起义，触一轴（B4c、S3）则区域／阶层化抵抗，不触（B4b及B2／B3／B5等）则无感。二日事件的三条件＝王位处置公开化＋首都占领军开枪＋等待外壳的失意中层，缺一即退化为可平息的骚乱。B4b的falsifier判负（中高）：保留波旁下不存在同等的起义触发条件；风险不在反法起义而在王朝内爆（阿兰胡埃斯式宫廷政变、继承争端）与财政破产。

\subsubsection{二、纳瓦拉：在最贵的地方买最少的服从}\label{ux4e8cux7eb3ux74e6ux62c9ux5728ux6700ux8d35ux7684ux5730ux65b9ux4e70ux6700ux5c11ux7684ux670dux4ece}

【锚】Tone（The Fatal Knot，转Joseba de la Torre）：\literalquote{}the French regime
collected over 152 million reales in taxes and requisitions during the
period 1808 to 1814... over 40 percent of Navarre\textquotesingle s
normal five-year agricultural and industrial
output\literalquote{}；成熟游击下征收崩溃（1811--12实收\textless29\%，1812--13仅10\%）；镇压兵力8,000→13,000仍无法控制乡村。B4c的兼并带（纳瓦拉、巴斯克、加泰）恰是抵抗结构最强的区域------小农＋分散聚落＋村社民主＋司铎密度＋走私枪支＋foral动员机器；1717--22年关税线北推已失败一次；阿桑萨1810年专程赴巴黎警告设省必战。成本带（模型）：常驻镇压＋边防4.5--8万（纳瓦拉1.5--2.5万／巴斯克1--1.5万／加泰2--4万）；兼并区财政净贡献第一个十年为负；转正条件＝一代人（15--25年）＋低税＋市场整合＋不碰教会；加泰是唯一可能的\literalquote{}合作化兼并区\literalquote{}（棉业市场利益＋反马德里积怨，低中置信）。

【锚】初始洪塔的社会构成（Fraser）：\literalquote{}SPAIN: 29 juntas, 585 members...
Municipal authorities 140 (23.9)... Clergy 138 (23.6)... Military
officers 107
(18.3)\literalquote{}；王室地方当局14.7\%；有衔贵族仅3.2\%、商人2.2\%、劳动阶层3.2\%------洪塔是旧制度中层的接管，不是民众民主或大贵族复辟。人口账（Fraser九地区三十万人样本加权）：\literalquote{}pre-war
350,000 to 510,000 people (4.0 to 5.8 percent...) can be estimated as
missing compared to 215,000 to 375,000 (2.4 to 4.2 percent) in the
war\literalquote{}------战前灾难（饥荒／黄热病）损失大于战争；游击腹地（纳瓦拉+6.5\%、巴斯克+4.3\%、加利西亚+12.9\%）战时自然增长为正，代价集中在会战、围城与饥荒走廊（埃斯特雷马杜拉−14.1\%、加泰−8.6\%）。Esdaile\literalquote{}可能超百万\literalquote{}为宽口径叙述值，两口径并列。

游击史学的两极（Tone结构说vs
Esdaile人造说）在反事实关键点上收敛：游击不是西班牙人的固有属性，而是占领＋国家解体＋榨取的函数；Esdaile：除纳瓦拉、巴斯克外多数partida士兵系枪口强征，1808年夏后到处是逃役与反征兵骚乱；Tone：纳瓦拉志愿率仅受查人口3--4\%，教士不成军。

\subsubsection{三、无占领的宪政}\label{ux4e09ux65e0ux5360ux9886ux7684ux5baaux653f}

加的斯式激进单院宪法以王位空缺为成立条件；但财政破产（岁入350M vs
支出850M＋、债11,000M
reales）、臃肿军官团、美洲战争厌战三项结构压力在无占领世界仍存其二。【推·中高】宪政会出现但晚到（1815--30窗口替代1810--12）、温和化（传统Cortes／两院宪章形态优先，1812式单院宪法低置信）、更血腥地内战化（缺反法黏合剂，自由派与保王大众直接对撞）；1820年革命的动力学全套不需要法国占领变量（欠饷＋美洲战争＋共济会＋商港）。【锚】卡洛斯战争证大众保王主义是地区结构变量而非全国常数：\literalquote{}The
poorest areas of Spain, which were in Extremadura and Andalucía, were
lost causes for Don
Carlos\literalquote{}（Lawrence）------北方三区可养万人级反自由派军队数年，政府需二十万以上军队＋约18,000外国志愿者才能压平，但动员强制多于志愿。法胜世界1815--48三断层恒在------土地／señorío清算（1813年反封建暴动已点燃，1814年瓦伦西亚贵族因此投向政变）、教会财产与修会（1798先例→1834
matanza结构条件）、中央化vs
foral；任何政权只改变爆发顺序与镇压者身份；管理的单点＝教会协约（堂区司铎薪俸国家化则大众保王主义失去中介神经）与军饷。

\begin{center}\rule{0.5\linewidth}{0.5pt}\end{center}

\subsection{第二十八章　葡萄牙与巴西宫廷：英国的港口、布拉干萨的大西洋}\label{ux7b2cux4e8cux5341ux516bux7ae0-ux8461ux8404ux7259ux4e0eux5df4ux897fux5babux5ef7ux82f1ux56fdux7684ux6e2fux53e3ux5e03ux62c9ux5e72ux8428ux7684ux5927ux897fux6d0b}

\subsubsection{一、占得住壳，拿不到瓤}\label{ux4e00ux5360ux5f97ux4f4fux58f3ux62ffux4e0dux5230ux74e4}

【锚】1807年10月22日英葡秘密公约：\literalquote{}England pledged an escort in case the
royal family fled to Brazil, and guaranteed never to recognize as king
of Portugal any prince who was not the legitimate heir of the house of
Bragança\literalquote{}（Manchester，本轮全文）------英国获权占马德拉，摄政王承诺舰队不落法手；在法军入境前五周，帝国的海洋部分已被预划入英国保护圈。坎宁1807年11月7日训令：\literalquote{}In
case the prince regent failed to go to Brazil, Strangford and Smith were
empowered to repeat the Copenhagen affair in Lisbon (F. O., 63/56,
Canning to Strangford, Nos. 11, 13,
14)\literalquote{}------葡舰队在任何分支下都不会完整落入法国之手。出走规模：8,000--15,000人、36舵、流通货币之半、八千万克鲁扎多金银、八艘战列舰；\literalquote{}A
delay of two hours in weighing anchor would have resulted in the
detention of the fleet within the harbor as the wind shifted immediately
after the narrows had been
passed\literalquote{}------窗口真实存在，但只差数日至数小时。

【锚】1810年商约：\literalquote{}the obligations and conditions expressed or implied
in it shall be perpetual and immutable; and they shall not be changed or
affected in any manner in case His Royal Highness the Prince Regent of
Portugal, His Heirs or Successors, should again establish the seat of
the Portuguese Monarchy within the European Dominions of the
Crown\literalquote{}------英货15\%低于葡船1808年的16\%，宗主国待遇不如外邦，直到1810年10月18日才拉平。Pedreira（HAHR
2000）：\literalquote{}In 1808--1813, exports of national commodities to Brazil had
declined to 22.4 percent of the value reached in 1796--1806 \ldots{}
Reexports of Brazilian goods to Europe sunk to 11.6 percent of the
prewar yearly average\literalquote{}；1825--30再出口跌至战前17\%。Maxwell：\literalquote{}for most
of that century 60 percent of the state\textquotesingle s revenues
derived from Brazil\literalquote{}------失去巴西即失去财政基座。

【锚】卡洛塔的拉普拉塔主张：英国内阁是否决人（坎宁两道训令）、西德尼·史密斯是越权推手、若昂借壳谋取拉普拉塔；布宜诺斯艾利斯存在真实的Carlotismo秘密会社------\literalquote{}A
secret society counting such illustrious names as Manuel Belgrano,
Martin Pueyrredon, and Mariano Moreno, was formed \ldots{} to proclaim
D. Carlota Joaquina regent of all Spain\literalquote{}。

\subsubsection{二、里斯本军港的价值是制海权的函数}\label{ux4e8cux91ccux65afux672cux519bux6e2fux7684ux4ef7ux503cux662fux5236ux6d77ux6743ux7684ux51fdux6570}

【锚】1808--09年塞尼亚文的俄国分舰队：法国占领里斯本九个月仍无法使用港内友军舰队；朱诺1808年7月3日函请合作被拒；最终俄英公约把舰队交英国存押------\literalquote{}The
ships of war of the Emperor of Russia, now in the Tagus \ldots{} shall
be delivered up to Admiral Sir Charles Cotton, immediately \ldots{} to
be sent to England, and there held as a
deposit\literalquote{}。制海权不在手时，里斯本军港的价值≈零。Beresford葡军：\literalquote{}The
field-strength of the Portuguese regular forces should have been,
according to its establishments, 56,000 men. In September 1809 there
were only 42,000 men with the colours\literalquote{}（Oman
III）；1812年初33,000正规军支撑威灵顿；英国直接承饷其半，补贴定为维持三万人之费；摄政委员会岁入仅及所需之半。

\textbf{主线}：葡萄牙是拿破仑体系里\literalquote{}占得住壳拿不到瓤\literalquote{}的极端案例------帝国的瓤（王朝、舰队、八千万克鲁扎多、巴西、贸易网）1807年11月29日一日出海；法国最优是维持1803--06年的付费中立（月一百万里弗实测）而非入侵------半岛战争的葡萄牙段是拿破仑送给英国的战略礼物。枫丹白露可执行性：地面四周内可执行（三支西班牙师实占指定区），目标不可达（标的物提前四十八小时出海），结构自败（法国自己三个月后单方废约改索埃布罗以北）。\textbf{对行省化的教训：领土瓜分对首都可搬到另一半球的海洋帝国天然打不中要害。}葡本土上限B级藩属且为负净值（−55～−70m法郎／年）。S5下巴西入英圈成美洲第二大国，葡语世界永久分流；法制里斯本的\literalquote{}1820\literalquote{}更像1817年戈梅斯·弗雷雷的放大版（革命者失去\literalquote{}合法且可迎回的王\literalquote{}）。未决：朱诺一亿法郎贡赋实征额无逐项账（估名目一两成）；Sherwig对葡逐年补贴表；1810条约与1815敕令为Manchester转引FO档。

\begin{center}\rule{0.5\linewidth}{0.5pt}\end{center}

\subsection{第二十九章　奥斯曼帝国：改革、政变、藩镇与海峡}\label{ux7b2cux4e8cux5341ux4e5dux7ae0-ux5965ux65afux66fcux5e1dux56fdux6539ux9769ux653fux53d8ux85e9ux9547ux4e0eux6d77ux5ce1}

\subsubsection{一、可租用，不可征用}\label{ux4e00ux53efux79dfux7528ux4e0dux53efux5f81ux7528}

【锚】1808年3月圣彼得堡瓜分谈判死锁于达达尼尔：\literalquote{}Mon opinion est que nous
ne pouvons céder ni Constantinople ni les
Dardanelles.\literalquote{}------俄方（鲁缅采夫与亚历山大）坚持君士坦丁堡与达达尼尔皆归俄；法方君堡可议、达达尼尔绝不让；科兰古按训令提出双锁分持被拒（Vandal，卷一）。埃尔福特公约（1808年10月12日）未完成瓜分：只托九至十二月间认两公国归俄（第8条俄方改稿）；\literalquote{}Les
hautes parties contractantes s\textquotesingle engagent
d\textquotesingle ailleurs, disait l\textquotesingle article 16, à
maintenir l\textquotesingle intégrité des autres possessions de
l\textquotesingle empire ottoman\literalquote{}------Vandal定性为共管扣押（mise sous
séquestre）而非真保证；第12条把瓜分悬置到\literalquote{}一年内再会\literalquote{}。falsifier成立，但精确断点在达达尼尔而非君士坦丁堡。

【锚】1812年3月拿破仑请马哈茂德二世率八万人入波兰的邀盟（随法奥同盟与保证奥斯曼完整的消息一同抵京）被实质无视：\literalquote{}Mahmud
II and his grand vizier still recalled what was generally perceived as
the treachery of Tilsit, and appear never to have seriously considered
such an
alliance.\literalquote{}（Aksan）。1807年春奥斯曼同盟期内，波斯尼亚穆斯林贝伊两度拒绝25,000法军借道去多瑙前线（只允少量军官），Bayraktar本人也反对法军到多瑙------\literalquote{}the
governor informed him that neither he nor the beys would ever permit the
French to march across their territory to the
Danube\literalquote{}（Shaw）。\textbf{盟友体系内部存在对大规模异教徒部队过境的多重否决。}

\subsubsection{二、反直觉：法国对俄威胁越大，奥斯曼的和约越好}\label{ux4e8cux53cdux76f4ux89c9ux6cd5ux56fdux5bf9ux4fc4ux5a01ux80c1ux8d8aux5927ux5965ux65afux66fcux7684ux548cux7ea6ux8d8aux597d}

【锚】俄因对法备战，1811--12年从要求全部两公国退到Seret线再退到Prut线，并放弃两个秘密条款与同盟要求：\literalquote{}The
idea of an alliance, and the two secret articles, never ratified by the
Ottomans, were dropped from Russian
demands.\literalquote{}（Aksan）。若B7＝不战，俄无退让压力，奥斯曼的条款更苛。这条推断（中高）供P6／P8。

【锚】新军（Nizam-ı Cedid）1806年底实册\literalquote{}22,685 men and 1,590 officers
enrolled in the Nizam-i Jedid army, of whom approximately half were
stationed in
Anatolia\literalquote{}；在鲁米利亚复制征兵完全失败，欧洲部分军力几乎完全依赖独立地方显贵------这是1806年埃迪尔内事件与1807年政变的结构背景。耶尼切里三层账：\literalquote{}some
135,000--145,000 names were on the rolls, while not more than 10,000
were on duty in
Istanbul\literalquote{}（1826年）；1809年出征点名4,000人但坚持领14,000份饷；esame已证券化、全社会持有------废除耶尼切里＝大规模既得利益清算，故需三年准备（换教长、倍增炮兵、进口五万支列日火枪）。

\textbf{主线（中高）}：奥斯曼是每次被强制就掉头的自主机器；S4的最优形态是低成本牵制器（海峡工事＋军械＋现金换对俄牵制与海峡关闭）；大军过境与合营瓜分均被实测证据否决；1812年法奥第六条证明法国官方立场回摆保全。S5下1815--48三大内生进程（集权×耶尼切里×埃及）方向不变；希腊更可能被镇压或成法国单边保护的小自治体（无纳瓦里诺类比，依赖法国国内philhellenism强度，未量化）；埃及保叙利亚更久（无1840伦敦公约类比）；奥斯曼比史实更完整但更依赖法国单一保护。未决：1800--15奥斯曼中央逐年决算不存在于公开英语文献（财政判断基于1784基线＋1826／1839点值）；İrad-ı
Cedid年入Juchereau数（50--75M
kuruş）与中央国库口径（14.7M）相差数量级；BOA拒盟原档经İsmail转引；1807年5月政变若新军被召能否镇压无当日兵力独立记录。

\begin{center}\rule{0.5\linewidth}{0.5pt}\end{center}

\subsection{第三十章　埃及、穆罕默德·阿里与红海通道}\label{ux7b2cux4e09ux5341ux7ae0-ux57c3ux53caux7a46ux7f55ux9ed8ux5fb7ux963fux91ccux4e0eux7ea2ux6d77ux901aux9053}

\subsubsection{一、等距的结构：商业向英、技术向法、政治两拒}\label{ux4e00ux7b49ux8dddux7684ux7ed3ux6784ux5546ux4e1aux5411ux82f1ux6280ux672fux5411ux6cd5ux653fux6cbbux4e24ux62d2}

【锚】Driault（Mohamed Aly et Napoléon，1925，本轮全文）：\literalquote{}Il eût
préféré s\textquotesingle entendre là-dessus avec la France\ldots{}
Napoléon ne songeait en ce sens qu\textquotesingle à une expédition
militaire : une entente politique sans doute eût mieux
valu.\literalquote{}------阿里1810年拒英国\literalquote{}代办独立\literalquote{}提议并更愿与法国达成政治协约，但拿破仑只筹划军事远征。拿破仑1810年9月17日问德克雷：\literalquote{}Quand
aurai-je les moyens de porter en Égypte\ldots{} total 40.000
hommes\ldots{} C\textquotesingle est à mon plan de campagne pour
1812.\literalquote{}；1811年3月8日令土伦备齐埃及远征所需并研究尼罗河船型；同期派Boutin／Nerciat入埃叙侦察要塞。全部停在准备段，闸门是制海权。塞巴斯蒂亚尼1802年使命与1803年1月30日《Moniteur》报告\literalquote{}Six
mille Français suffiraient aujourd\textquotesingle hui pour conquérir
l\textquotesingle Égypte\literalquote{}由Driault 1904
pp.25--26转印本补锚（非原影，出处Testa II, 46--61）。

【锚】1810--13年埃及粮食经马耳他comptoir卖往西班牙与葡萄牙：\literalquote{}the wali
decided to open a warehouse or comptoir in Malta which would sell grain
in Spain and Portugal and buy whatever Egypt needed from
Europe\literalquote{}（Marsot）------实质喂养英国的地中海--半岛战争体系；法国领事Drovetti逐月记录赴马耳他粮船并估粮贸利润45--50m
piastres。\textbf{S3封锁越严此线越值钱，阿里越不可能配合封锁。}埃粮1810--13实证可养5--8万东地中海军队。

红海舰队的天花板：\literalquote{}eighteen ships varying in weight between 100 and 150
tons were built in the Bulaq arsenal and carried in pieces on camel back
to
Suez\literalquote{}（Marsot），另向马斯喀特苏丹租船------苏伊士无舰材无大厂，S4经埃及的军团级海运不成立。红海运兵分支1803--15关闭的充分原因在法国侧：无制海权＞只备征服不备协约＞运力天花板＞阿里态度（仅列第四）。

\subsubsection{二、天花板由阿里的目标序决定}\label{ux4e8cux5929ux82b1ux677fux7531ux963fux91ccux7684ux76eeux6807ux5e8fux51b3ux5b9a}

【锚】1828--30年法国提议阿里代取的黎波里、突尼斯、阿尔及尔：\literalquote{}he was
approached by the French to undertake an expedition to capture the
Barbary states of Tripoli, Tunis and Algiers\ldots{} the expedition was
too costly to merit taking the risks involved and that it would divert
him from the provinces he strongly
desired\literalquote{}（Fahmy）------法国随后1830年自征。埃及作为法国工具的天花板在阿里的目标函数（独立世袭＞叙利亚＞舰材基地＞商路），不在法国的支付能力。军力三口径不得互填：Fahmy档案口径\literalquote{}by
the mid-1830s the number of conscripts had already reached
130,000\literalquote{}（人口约五百万的2.6\%）；Scott
1837实地口径：额定90,550、实额≤80,000埃籍正规，叙募＋非正规可近倍；逃兵累计陆军60,000＋海军20,000（每二征一逃）、欠饷15--20个月------动员上限已越可持续线。阿里致Najib
Efendi自评：\literalquote{}although we are men of war (ehlli harbdan), yet we are
still in the A B of that art (alif ba), whereas the Europeans are way
ahead of us\literalquote{}；Scott 1837判埃海军\literalquote{}not powerful enough to cope with the
fleets of the great European
powers\literalquote{}------埃及舰队不改变英法海军差距的主曲线。

【锚】科尼亚（1832年12月21日）：\literalquote{}The Ottoman forces lost 92 cannons,
3,000 casualties and 10,000 prisoners. The Egyptians, by contrast,
suffered 262 casualities and 530
wounded.\literalquote{}（Aksan）------奥斯曼五万对埃军一万五千；马尔蒙评若即刻进军可在一月初抵于斯屈达尔引发君士坦丁堡革命。1830年代埃新军＞奥斯曼新军，欧洲式陆军可在东地中海被本地政权复制。1835年疫死约五十万；1841年债八千万法郎。

\textbf{判定}：埃及是不可收编的乘数：法国可买到粮仓、威慑存在（英国为此付亚丁、前沿使团、波斯湾保险费）、技术代理陆权，买不到主权通道；A层行省化＝净消耗单位（1798--1801三年实测＋阿尔及利亚十万驻军曲线类比），B级法典化藩属是天花板；1815--48分岔主变量＝欧洲能否协调------法胜情景无纳瓦里诺、无1840四国干预→阿里保叙利亚成\literalquote{}阿拉伯--奥斯曼副帝国\literalquote{}，但内部汲取极限与继承结构不变→1848过载回调。巴尔巴里：阿尔及尔在册奴隶峰值仅1,645（1813）；突尼斯私掠收益占beylik岁入15.7\%（活动高峰年）；马赛自马格里布年到港船148→28（1793--99）------百万法郎级租金国家，私掠曲线随法国海权真空涨落；1830阿尔及尔类比不可外推埃及。未决：al-Jabartī原文、Ghorbal、Douin文书卷、Bowring原卷未获；1830远征军37,000／600舰船为通说。

\begin{center}\rule{0.5\linewidth}{0.5pt}\end{center}

\subsection{第三十一章　波斯、阿富汗与旁遮普：陆路通印的走廊}\label{ux7b2cux4e09ux5341ux4e00ux7ae0-ux6ce2ux65afux963fux5bccux6c57ux4e0eux65c1ux906eux666eux9646ux8defux901aux5370ux7684ux8d70ux5eca}

\subsubsection{一、加尔达纳使团的自我勘测}\label{ux4e00ux52a0ux5c14ux8fbeux7eb3ux4f7fux56e2ux7684ux81eaux6211ux52d8ux6d4b}

【锚】加尔达纳（Gardane）本人拟定的对印计划（Driault
1904，本轮全文）：阿勒颇线七至八月、特拉布宗线五至七月抵印度河，赫拉特为集结点，俄可在阿斯特拉巴德设仓，\literalquote{}il
faudra toujours un cheval ou une autre bête de somme pour deux hommes...
Il faudra beaucoup de paille hachée, il n\textquotesingle y a pas
d\textquotesingle autre
fourrage\literalquote{}，居民将逃亡。参谋勘测分工有名有实：Bianchi
d\textquotesingle Adda（斯库塔里--德黑兰）、Trézel（巴格达--伊斯法罕--设拉子--波斯湾）、Fabvier（德黑兰--伊斯法罕）、Truilhier（向马什哈德）。1808年1月26日对波斯军评估：骑兵144,000（纪律差但骑术佳）、步兵60,000（无价值）、炮手2,500而炮兵实际不存在（旧炮牛拖、炮弹出膛即炸）；1808年秋实测：\literalquote{}2
ou 300 cavaliers russes ont suffi à chasser tous les Persans de la rive
gauche de l\textquotesingle Araxe\literalquote{}。

【锚】芬肯施泰因条约：第12条法军过境须另订专约（路线、给养、运输、辅助军），第13条一切供应按本国人同价由法军自付，第14条通行权\literalquote{}elles
ne pourront être étendues, par des traités postérieurs, ni à
l\textquotesingle Angleterre, ni à la
Russie\literalquote{}------法俄联军过波斯需重新谈判；商约的Kharg岛让与以格鲁吉亚归还（第3--4条兑现）为前提。四项政治前提两两互斥：芬肯施泰因第14条排俄×俄要海峡×波斯要格鲁吉亚×奥斯曼要完整。

\subsubsection{二、四个明码标价的安全采购者}\label{ux4e8cux56dbux4e2aux660eux7801ux6807ux4ef7ux7684ux5b89ux5168ux91c7ux8d2dux8005}

【锚】卡扎尔中央财政的现金层极薄：沙对琼斯1809年自述省入不超省支、贡品为实物、现钱仅够养兵与宫廷、珠宝无法变现；英国价格序列：1809年预备约£100k／年＋训练1.6万步兵，\literalquote{}in
the definitive treaty the amount was fixed at 200,000 tomauns (or about
150,000l.)\literalquote{}（Kaye，1812年3月14日定约），1814年德黑兰条约保留但加非挑衅条件。杜兰尼（Elphinstone，1815）：名义岁入近三crore卢比、实收远低于二crore（约半为采邑）、王室现金仅0.9
crore；\literalquote{}Shauh-Shuja had only five {[}guns{]} when he took the field at
Peshawer in 1809\literalquote{}。锡克（Murray／Prinsep 1831--33）：\literalquote{}Making a total
Khalsa Revenue of 1,48,81,500... not less than ten crores of
rupees\literalquote{}（Govindgarh窖藏≥£10m）；总兵力82,014（正规27,752），火炮376门＋骆驼炮370；Prinsep同期预言其政府无制度化基础、一人死则机器停。

【锚】1809年英国三条约栅栏（Aitchison卷VII，全文已核）：阿姆利则4月25日（英不干预萨特莱杰河以北／兰吉特限河左岸驻军）、喀布尔6月17日（阻法波过境、英付费义务限法波联合入侵、禁法国人入境）、信德8月22日第四条：\literalquote{}The
government of Sindh will not allow the establishment of the tribe of the
French in
Sindh\literalquote{}；波斯1809年3月12日预备约第3条废止对欧一切条约＋不许任何欧洲军队过境赴印。明托四使团＋Karrack案＋三条约十个月落成------\textbf{法国东方压力最确凿的产出，是英属印度前沿体系提前二十年成形。}

\subsubsection{三、后勤包络的裁定}\label{ux4e09ux540eux52e4ux5305ux7edcux7684ux88c1ux5b9a}

【推】单季无预置无里海航运的≥2万法军波斯段行军关闭（高置信；实测锚：贝尔德1801年数千人分队过120英里即需5,000骆驼；英印1838年孟加拉纵队限9,500战兵配30,000骆驼）；分梯队＋四仓预置＋阿斯特拉巴德俄运＋首年1,500--3,000万法郎白银可使2--3万人1.5--2个战役季抵赫拉特（工程边缘可行，中低）。走廊四政权的主要威胁（波斯→格鲁吉亚交付；阿富汗诸派→内战胜利；兰吉特→两河间霸权承认；信德→双重挤压豁免）法国1807--09均无交付能力。S4的理性形态是\literalquote{}永久的加尔达纳\literalquote{}------常设军事使团＋对俄担保＋二十万toman级白银津贴；1815--48均衡：法得波斯C层盟国与牵制筹码，英得提前二十年的前沿帝国；锡克与信德终局不变只改时刻表；赫拉特是S4从威胁升级为战争的门槛（1837--38年英国反应链完整观测：先Kharg后战争）。未决：Trézel／Truilhier原勘测数据（可改判中档包络）；Amini／Atkin／Yapp／Ingram／Kelly经Iranica转引；波斯人口无同期可信总数（500--1,000万）；1810年俄方\literalquote{}归还占地\literalquote{}和约提案系波斯编年史孤证。

\begin{center}\rule{0.5\linewidth}{0.5pt}\end{center}

\subsection{第三十二章　东印度公司：财政、军队、贸易与承压能力}\label{ux7b2cux4e09ux5341ux4e8cux7ae0-ux4e1cux5370ux5ea6ux516cux53f8ux8d22ux653fux519bux961fux8d38ux6613ux4e0eux627fux538bux80fdux529b}

\subsubsection{一、\literalquote{}贸易养帝国\literalquote{}不成立}\label{ux4e00ux8d38ux6613ux517bux5e1dux56fdux4e0dux6210ux7acb}

【锚】公司请愿书（本轮亲读）：\literalquote{}all such sums of money, together with the
said sum of 6,289,405l., which has arisen from the surplus profits of
the trade carried on by the petitioners, as herein-before mentioned,
have been absorbed by payment of debts and expences incurred in respect
of the territorial acquisitions in
India\literalquote{}------1793--1812累计贸易利润£6,289,405（已扣商业费用、伦敦债券利息与股息），年均≈£0.32m；同期领地账年缺口£1.1--2.7m。领地账1803--06连续赤字：1803-04实绩岁入£13,273,044、支出£13,214,114、债息£1,534,758，扣明古连补给与进口货销售后净缺口£1,124,403；1805-06预算超支£2,651,939；结构上孟加拉盈余、马德拉斯与孟买赤字。董事会1812年正式函件：\literalquote{}wars
with the Native Powers have been repeatedly carried on, to the vast
accumulation of the Indian debt, now advanced from eight million
sterling, at which it stood in 1793, to about thirty-two
millions.\literalquote{}1807年起对伦敦开出汇票£13,104,924。Cartwright
1812年3月4日总账：经营所需资本£51,182,137；股本£7,780,000、伦敦债券£7,000,000、浮动与无息负债£7,787,953，差额约£28.6m须以印度计息债务维持（OCR印38,614,174，按51,182,127−22,567,953还原为28,614,174）。

\subsubsection{二、三问三答}\label{ux4e8cux4e09ux95eeux4e09ux7b54}

【锚】董事会1812年立场：\literalquote{}from the avowed design of France to invade our
Indian possessions with great armies by land, it may be necessary still
largely to augment our European force in that quarter. For a war of this
description the Indian revenues, if unincumbered with debt, would be
very inadequate. It would be an European war for European
objects\literalquote{}------同时承认在印王军\literalquote{}above twenty thousand
men\literalquote{}且为最昂贵兵种。1813年请愿自述：若不大幅削减印度军费就不会有足够盈余可汇，伦敦义务的履行\literalquote{}most
essentially\literalquote{}靠贸易。\textbf{英国被逐出欧洲后印度能否自养}：养政府与军队\literalquote{}能\literalquote{}（高）；再向伦敦汇出利息与政治开支\literalquote{}不能\literalquote{}（高）；负担欧洲级对法战争\literalquote{}不能\literalquote{}（高）。关键推论：欧洲排除\textbf{放松}而非收紧印度体系的约束------切断的是欧洲市场，而偿付能力来自中国与印度洋；真正的断点在伦敦信用而非印度税收。

【锚】茶--中国贸易：英国国库从中获得的关税与消费税1792-3至1808-9共£32,290,599（年均约£1.9m）；Milburn宽口径\literalquote{}forming
a total of £32,290,599, which, added to the above £30,739,208, forms a
beneficial total to the nation, of £63,029,807, on an average of 17
years, £3,707,636 per
annum\literalquote{}；该税基在大陆体系物理可及范围之外。孟加拉鸦片对华输出1795-6仅Sicca
Rs
2,51,450，1805-6已达32,94,570；鸦片专卖收入1804-5为£725,895、\literalquote{}1807-8\_
801,467\literalquote{}------战时唯一高速增长且不依赖欧洲市场的税源。汇款机器：\literalquote{}They
reserve to themselves two thirds of the chartered tonnage of their ships
destined from Bombay to China ... the freight, together with the
proceeds of the goods, shall be paid into the Canton treasury, for the
latter of which bills are granted on England at the current rate of
exchange\literalquote{}；孟买对华棉花出口峰年八万包约三千万磅。

【锚】1812年4月8日帕斯利、斯旺西、伯明翰请愿：\literalquote{}the Petitioners feel, at
this time, when they are nearly precluded from any trade with the
continent of Europe, that it becomes essentially necessary that the
merchants and manufacturers of this kingdom should look to new sources
for a vent for their respective
commodities\literalquote{}------\textbf{大陆体系是1813年特许状开放的有文书因果}。1806年后至1854年下院未再提交年度印度预算：\literalquote{}though
it is now only two years short of half a century since an Indian budget
was last presented to the House of Commons\literalquote{}（Vernon Smith
1854）------公司得以长期展期赤字的制度条件。1851--52年印度登记债利息£1,967,359、在英支付领地开支£2,506,377，余£531,265------home
charges是单调上升的棘轮。

\subsubsection{三、双侧脆弱性定理}\label{ux4e09ux53ccux4fa7ux8106ux5f31ux6027ux5b9aux7406}

【锚】1796→1808公司自有欧籍步兵由六团降至三团，土著步兵26团增至59团、土著骑兵8团增至16团：\literalquote{}the
European infantry, before confifting of 12 battalions, was converted
into fix regiments; the native infantry ... was formed into 26 regiments
of two battalions each\literalquote{}。1809年马德拉斯军官哗变的结构原因：\literalquote{}the higher
allowances granted to the officers of the Bengal army, and the undue
proportion of commands which had been recently bestowed upon the
officers of the Royal army\literalquote{}（Wilson）；1806年韦洛尔兵变：\literalquote{}About 850
sepoys were killed, the rest escaped by the sally-port on the southern
face\literalquote{}，第69团阵亡100人加15人伤重死亡。\textbf{增派国王军→公司欧籍军官晋升被挤压（1809实测）；减欧籍骨干扩土著兵→兵变风险上升（1806实测）。任何可信的法国威胁都把英国推向第一侧，且这次有外敌在场。}明托1807年：\literalquote{}20,000
French troops commanded by French officers, supplied with all the means
which attend such armies, would not, I am firmly convinced, subdue
us---but would create a struggle which must be attended with great
present
calamity\literalquote{}------运载瓶颈不在船而在逐次集结与落地后的印度盟友。反应函数三段：先要钱（1808
£2.46m、1813
£2,294,426索款）→再要人（更多国王军团）→只要伦敦解禁且海军配合就主动清场（1810波旁-毛里求斯、1811爪哇为实测）。公司无独立和战权。承压分档：L1可覆盖／L2年缺1--2.6m／L3需国库3--5m。未决：1803--15逐年印度岁入支出连续序列（HC
1855 Return，Dutt
1906第XXIII章转引，本地OCR拆散）；现代专著（Bowen／Peers／Philips／Webster）因工具限制不可用。

\begin{center}\rule{0.5\linewidth}{0.5pt}\end{center}

\subsection{第三十三章　印度诸邦与法属印度洋：马拉塔、锡克、迈索尔遗产与德卡恩}\label{ux7b2cux4e09ux5341ux4e09ux7ae0-ux5370ux5ea6ux8bf8ux90a6ux4e0eux6cd5ux5c5eux5370ux5ea6ux6d0bux9a6cux62c9ux5854ux9521ux514bux8fc8ux7d22ux5c14ux9057ux4ea7ux4e0eux5fb7ux5361ux6069}

\subsubsection{一、联络确存，失败在时序、中介与截获}\label{ux4e00ux8054ux7edcux786eux5b58ux5931ux8d25ux5728ux65f6ux5e8fux4e2dux4ecbux4e0eux622aux83b7}

【锚】Sen（French in India，本轮全文）：\literalquote{}he also received a Brahmin
agent from Sindhia, who came to solicit the French to land an
expeditionary force on the Malabar Coast, between Bombay and Goa, with a
promise of assistance from the
Marathas\literalquote{}------Binot在本地治里投降前数日接待辛迪亚使者；德卡恩1803年11月派Courson／Durhène携致Perron、辛迪亚、霍尔卡三方信件保证武装援助与同盟；二人1804年1月登陆后至浦那被马拉塔当局交给英国人。拿破仑1804年9月向德卡恩特使Lefebvre细问：\literalquote{}Napoleon
questioned him closely about the chances of success in India with so few
men as asked for by Decaen, and expressed grave doubts about the
sincerity of the Marathas\literalquote{}。Gordon：\literalquote{}it was characteristic of the
Maratha polity that the inner circle of power changed with each
generation... those who left did not remain some sort of loyal
opposition\literalquote{}------马拉塔政体无可签约的集体主体；1803年12月诸约禁止欧洲人在任何马拉塔军队服役，\literalquote{}largely
because Shinde\textquotesingle s European officers (mainly French)
deserted and disbanded their troops\literalquote{}。

\subsubsection{二、四个计划窗口全被欧洲关闭}\label{ux4e8cux56dbux4e2aux8ba1ux5212ux7a97ux53e3ux5168ux88abux6b27ux6d32ux5173ux95ed}

【锚】1803年德卡恩自请六战列舰＋3,000步兵＋500骑；1805年1月16日致德克雷信（Corr.
X 8279）三分舰队20,000法军＋3,000西军；1808年5月13日（Corr. XVII
13877）洛里昂20舰4,000兵＋布雷斯特20舰13,000兵；\literalquote{}On the 10th June
Napoleon gave precise instructions: 12 ships were to start from Nantes,
13 from Lorient and 31 from Brest, carrying a total of 19,600 soldiers
and 10,600 sailors\literalquote{}（Corr. XVII
14078）------西班牙起义后永久放弃。十年实送≈1,200人＋五艘巡航舰。英方同期完成两次同量级跨洋远征（1810毛里求斯6,848／11,306／16,000三口径，1811爪哇≈12,000）证明屏障是制海、中继与优先级，不是距离与季风。1810年法兰西岛防御：陆军约900＋陆战队炮兵300＋国民自卫队800≈两千（含忠诚不可恃的爱尔兰战俘）；\literalquote{}Decaen
issued a proclamation in August 1809, calling upon slave-holders to
place a certain number of slaves at his disposal for recruitment in the
army. But this order had to be withdrawn because of the opposition of
the European population, who feared a general rising of the
slaves\literalquote{}------基地社会结构对动员的硬约束。缅甸轴线：1809年爪哇Daendels遣de
la Houssaye赴缅谈判\literalquote{}establishing an alliance with the Burmese Government
for the subversion of the British Empire in India\literalquote{}，抵达前被英截获。

【锚】路径C（军官--军械输出）已被历史验证：Allard／Ventura
1822年经波斯入拉合尔建Fauj-i-Khas（1825年5--6千人，法语口令、三色旗配鹰徽）；英国反制以对波斯补助金为条件迫其解雇全部法国军官。

\textbf{四闸门模型}：出港×中继×季风窗×岁入基地，1803--15从未同开。分层带：骚扰层已达上限自供养；千人级军官--军械层1803--08可行；万人级需S1／S2且四闸门同开；十万级任何情景不可行。终局谱系：非武装归还（1814／15实局）→武装商站（S2可争）→桥头堡（需1830年前不可得的印度洋海军均势）→军官市场国家化（拉合尔实证）→征服（否定）。\textbf{印度对法国霸权的最优用法是威胁在场消耗英国}（1810--11英为拔巢动用2.4万人两舰队）。B4伊比利亚是最大单点改判杆：1808年6月10日的方案在西班牙不变时可以出港。未决：Prentout经双转引；德卡恩档案（Caen
145卷）与本地治里P.A.手稿未亲阅；Caroline／Manche／Vénus三舰派遣日期两说；Pulo
Aura货值疑OCR损。

\begin{center}\rule{0.5\linewidth}{0.5pt}\end{center}

\subsection{第三十四章　美国：中立承运人、禁运与1812}\label{ux7b2cux4e09ux5341ux56dbux7ae0-ux7f8eux56fdux4e2dux7acbux627fux8fd0ux4ebaux7981ux8fd0ux4e0e1812}

\subsubsection{一、可购买的服务供应商，不是可结盟的权力}\label{ux4e00ux53efux8d2dux4e70ux7684ux670dux52a1ux4f9bux5e94ux5546ux4e0dux662fux53efux7ed3ux76dfux7684ux6743ux529b}

【锚】Pitkin《统计展示》（1835）：美国转口贸易1806年60,283,234美元、1807年本国产品出口48,699,592美元；1808禁运年分别塌落至12,997,414与9,433,546（＝1807年的21.8\%与19.4\%）；1814年转口仅145,169美元。Perkins（Prologue
to War，1961）：\literalquote{}Britain scarcely benefited at all, for less than 5 per
cent of all reexported commodities went to England and her colonies.
\ldots{} It cast the protection of a neutral flag over supplies Napoleon
woefully
needed\literalquote{}------转口货几乎不进入英帝国市场，英国的敌意有物质基础，而这条贸易实为法国殖民地体系唯一规模化的出口通道。美国只卖三样：中立旗承运、对英政治与商业摩擦、廉价私掠战；不卖陆军、战列舰队与可靠同盟义务。它的价格是\literalquote{}不扣押美国船\literalquote{}，而拿破仑从未支付：1810年3月23日朗布依埃敕令下令扣押一切进入法国港口的美船。\textbf{最严封锁＝最快失去美国管道}------与北欧、俄国、法国海军、英国工业四卡同构。

\subsubsection{二、军事潜力被政体封顶}\label{ux4e8cux519bux4e8bux6f5cux529bux88abux653fux4f53ux5c01ux9876}

【锚】\literalquote{}the House rejected naval building on January 27 by a vote of
sixty-two to
fifty-nine\literalquote{}（1812年1月27日，十二艘战列舰＋二十艘巡航舰案）；盐税先以57:60被否；1--5月正规军只招到\literalquote{}few
more than 1,000
men\literalquote{}。财政为\literalquote{}只有海关的国家\literalquote{}（Pitkin：常态收入几乎全部来自进口与吨位税及公地出售，内税仅战时开征），禁运与战争自毁税基（海关收入1808年16,363,550→1809年7,296,020美元）。\textbf{美国在1848年前不构成第三海军变量。}

【锚】falsifier（对法扣船怨恨足以触发对法战争）前半成立、推论不成立：\literalquote{}by
the narrow margin of two votes, defeated Michael Leib\textquotesingle s
proposal to issue letters of marque and reprisal, the same to take
effect against France as well as
England\literalquote{}（1812年6月15日参议院）；肯塔基州议会与查尔斯顿跨党集会同时谴责两国；麦迪逊拒绝双重宣战的自述理由是党派损益（\literalquote{}arms
the federalists with new
matter\literalquote{}）。单独对法战争缺三项必要条件（法国不强征美国水手、北美无可侵法属领土、反法立场被编码为联邦党立场）。美国杠杆应\literalquote{}重新定价\literalquote{}而非整体下调：高注释、短系结、可被法国自身行为切断；计价单位是转口额与私掠船数。战争概率（四必要条件逐情景检查）：P（英美全面战争1812--15）＝S1
0.05--0.12／S2 0.30--0.45／S3 0.55--0.75；P（美法武装冲突）＝S2≤0.08／S3
0.25--0.40；P（对两国同时报复）＝S3
0.35--0.55。强征水手（门罗称6,257人／联邦党称\textless800，三口径已登记）与枢密令是英国单方供给的开战理由，与法国政策无关，故S2不能归零。

\subsubsection{三、白银的物理载体}\label{ux4e09ux767dux94f6ux7684ux7269ux7406ux8f7dux4f53}

【锚】Marichal（Bankruptcy of Empire）：\literalquote{}they were also authorized to
cash the drafts issued by the Consolidation Fund on the royal treasuries
of Mexico and to transport the silver pesos on U.S. (neutral) merchant
ships to
Europe\literalquote{}------1805--08年墨西哥白银入欧（含法国账户）的物理载体是美国中立商船：Hope
\& Co雇David
Parish组织船队，节点为费城（Parish／Craig）、巴尔的摩（Robert
Oliver，1806--08至少38航次）、纽约（Gracie）、新奥尔良（Nolte，约20＋20艘），每航次随船5--10万比索，担保为十张Consolidación汇票共8,484,375
pesos。1808年10月至1811年2月西属美洲官方运银抵加的斯\literalquote{}total
29,408,015\literalquote{}pesos（墨西哥2,400万，约总铸币量一半），多批由英国军舰承运------同一条管道在1808年后换旗与换受益人；谁控制管道取决于西班牙的政治归属，而非美国政策。到1833年广州贸易付款结构已从银转为汇票：\literalquote{}Bills
on England, £1,013,988, . . . \$4,772,516 \textbar{} Dollars and
bullion, . . .
682,519\literalquote{}；同年美国本国棉织品出口2,532,567美元中逾九十万销墨西哥------美国对拉美的角色从\literalquote{}运银者\literalquote{}变为\literalquote{}卖布者\literalquote{}。

\subsubsection{四、S5：门罗式原则提前，对象改为法国}\label{ux56dbs5ux95e8ux7f57ux5f0fux539fux5219ux63d0ux524dux5bf9ux8c61ux6539ux4e3aux6cd5ux56fd}

【锚】1823年文本两原则并列：\literalquote{}we should consider any attempt on their
part to extend their system to any portion of this hemisphere as
dangerous to our peace and safety. With the existing colonies or
dependencies of any European power we have not interfered and shall not
interfere.\literalquote{}------只要法国不帮西班牙重建统治，美国就不介入；但宣言无执行能力，故三分支：英美事实协作0.50--0.65／自建关税--海军0.20--0.30（1827--33海关收入已达1,971--2,903万美元）／向法国购买式妥协0.10--0.20（1803购地与1819年法使海德·德·纳维尔调停亚当斯--奥尼斯为先例）。扩张由人口决定：\literalquote{}1810
.... 7,239,903 \textbar{} 1820 9,638,131 \textbar{} 1830 ....
12,866,020\literalquote{}（约24年翻番）------法国即使欧洲全胜也无法给西班牙或其他客户国提供能对抗该曲线的移民供给，只能改变时点与代价：佛罗里达可推迟至1825--35，1848年前抵太平洋概率0.60--0.80。S6在美国维度不可观测，仅经\literalquote{}西班牙归属\literalquote{}间接传导。未决：1812--15军事行动与战费序列依赖Perkins单一研究传统；1810--12对法索赔的官方量化汇总（若量级近或超英国扣船损失，美法武装冲突概率上调至0.40--0.55）；美旗输银总额（400--800万比索仅为压力尺度）。

\begin{center}\rule{0.5\linewidth}{0.5pt}\end{center}

\subsection{第三十五章　海地与加勒比奴隶社会：法国的殖民底线}\label{ux7b2cux4e09ux5341ux4e94ux7ae0-ux6d77ux5730ux4e0eux52a0ux52d2ux6bd4ux5974ux96b6ux793eux4f1aux6cd5ux56fdux7684ux6b96ux6c11ux5e95ux7ebf}

\subsubsection{一、黄热病是非免疫人口流量的函数}\label{ux4e00ux9ec4ux70edux75c5ux662fux975eux514dux75abux4ebaux53e3ux6d41ux91cfux7684ux51fdux6570}

【锚】Geggus（1979）：英军1794年10\%／月→1797年2.5--4.2\%／月；牙买加驻军1803--36年均13\%（含1808／1819／1822／1825／1827五个疫年），任一年≥6\%；疫年17--30\%仅作压力测试------三个数字不得合成（本卡曾自造\literalquote{}期望年损耗7.6--15.4\%\literalquote{}的伪校准，已撤回：13\%是全期年均已含疫年，再与17--30\%加权即重复计入）。免疫来自\literalquote{}在流行区成长已有既往感染的人群\literalquote{}，不是族裔天然免疫。1802--03年的损失\literalquote{}八千毛／六千超额\literalquote{}不可相加；海员损耗大部在1803年前已沉没，不入反事实回收。

\subsubsection{二、层级带与命运的地点}\label{ux4e8cux5c42ux7ea7ux5e26ux4e0eux547dux8fd0ux7684ux5730ux70b9}

【推·主观条件概率】A再征服复奴\textless0.08；B法典化藩属0.12--0.30；C不常驻登陆的承认／关税交易0.55--0.80（促成签字的概率，非净收益）；D总带0.30--0.55------D直接领有0.08--0.18，D′归西班牙主权取条约权益0.25--0.45。\textbf{海地的命运在马德里决定（B4）。}B被C支配：法国能给客户政权的海地现政权已拥有------土地已由1805年宪法总则第12条收归国有并由佩蒂翁分给军官士兵，教产亦无可卖（帝制宪法50--52条不设国教、国家不供养任何教士）；1814年马卢埃方案给克里斯托夫的是托尔蒂岛、给佩蒂翁一派的是\literalquote{}例外承认与白人平等的公民权\literalquote{}，两项都低于它们已经拥有的。C不支配D：D的特有产出（土地与劳动力容量）是C给不了的；在\literalquote{}承认＋商业准入＋不输出革命\literalquote{}的目标函数下C最优，若目标改为\literalquote{}获得可撷取领土\literalquote{}则必须走D并支付代价。D层应读作\literalquote{}一个10--15年的撷取窗口\literalquote{}而非永久边疆：西班牙1809年7月重接管圣多明各东部，被同时代称为\literalquote{}España
Boba\literalquote{}，1821年短命独立、1822年被海地并吞------存续十二年；D′把财政负担转给了一个已被实测不能承担它的伙伴。

【锚】糖：法属殖民地糖出口1789年141,000,000法磅→1818年归零，咖啡保四成；甜菜糖1813年实产1,100,000公斤（Heckscher印p.293，见第四十章），对1788年法属殖民地糖95,045吨为约1.2\%（原用展望值3,500吨得3.7\%，已更正）。\textbf{蔗糖回流杀死甜菜工业→本土无反殖民地利益集团→法国废奴推迟至1855--1875}；甜菜替代能力越小，机制越强。奴隶贸易增量有界+15--40\%。1825年索赔一亿五千万不能当现金实收（1838年实际债权60m／90m，付款延至1883年，款入存款局补偿旧殖民者私人索赔，国库净额未知）。

\subsubsection{三、与陆军的接口}\label{ux4e09ux4e0eux9646ux519bux7684ux63a5ux53e3}

以在场5,000--8,000（q\_C
0.65--0.80→在册6,250--12,308；欧洲在场率0.85→调驻损5,313--10,462在场）：仅调驻，S2余额−11,027至−16,176，S5＋33,824至＋38,973（占原余5.5万的9.7--19.0\%）。固定300m陆军预算下按F2公式E＝0.85×{[}(300m−c\_C×G\_C)/700＋200,000{]}：有限据点S2缺13,304--26,637、S5余23,363--36,696；D层直接领有S2缺28,482--71,099、S5从＋21,518转−21,099------D层上沿可吃尽S5全部余量并转赤。从欧洲调驻者原已按700拨费，只新增热带溢价1.875--8.615m。c\_C（1,000--1,400法郎／兵年）与q\_C是未校准假设；Marion的那不勒斯900不是热带系数；年征已含死亡，不再追加480--1,040补员；人力与财政耦合，不再说\literalquote{}真正约束只是财政\literalquote{}。静态财政缺口不证明\literalquote{}10--15年窗口\literalquote{}，该期限是另行类比。HT／IT2的协作与HT×F2的联算均不是独立互证。未决：Dubois／Girard／Blackburn／Drescher／Ferrer／Régent／Popkin未获；Persée《Outre-Mers》2003海地专号只得题录；1809--21东圣多明各逐年财政与驻军数。

\begin{center}\rule{0.5\linewidth}{0.5pt}\end{center}

\subsection{第三十六章　新西班牙与古巴：白银、克里奥尔精英与Consolidación}\label{ux7b2cux4e09ux5341ux516dux7ae0-ux65b0ux897fux73edux7259ux4e0eux53e4ux5df4ux767dux94f6ux514bux91ccux5965ux5c14ux7cbeux82f1ux4e0econsolidaciuxf3n}

\subsubsection{一、控制越强，实得越少}\label{ux4e00ux63a7ux5236ux8d8aux5f3aux5b9eux5f97ux8d8aux5c11}

新西班牙对法国的汲取价值与法国对西班牙的控制强度\textbf{反向移动}。汲取需要三样东西同时在场------被本地官僚承认的合法委托人、肯兑付的总督、英方默许的航线；行省化摧毁第一样，摧毁第二样的动机，并因必然伴随排他封锁而摧毁第三样。这是\literalquote{}最大行省化≠最大有效海权\literalquote{}（F3）、\literalquote{}S3与S6在海岸带相互抵消\literalquote{}（F4）之后同一结构定理的第三次出现。

【锚】Marichal：\literalquote{}The functionaries later calculated that they had
obtained approximately 20 million pesos (400 million reales) between
1804 and 1808, although a substantial part of these funds never arrived
to Spain but were paid out in Holland to the Crown bankers and/or to
French Treasury
officials\literalquote{}------1804年8月2日至1808年2月8日马德里对美洲皇家国库开出的汇票约二千万比索，仅新西班牙一地；这是\literalquote{}法国能从墨西哥白银抽取多少\literalquote{}的史实上限锚。\literalquote{}It
can be estimated that Hope had around ten million pesos in bills of
exchange in their power (previously sent by Ouvrard) with the obligation
to transfer to the French Treasury at 75\% of their face value,
discounting expenditures and
commissions\literalquote{}（p.187）。Buist原文证实合约价：\literalquote{}Baring Bros. \& Co. were
committed to pay 3 francs 75 centimes for
each\literalquote{}piastre；1806年5月10日Mollien--Izquierdo协定：西班牙实际承担仅Compagnie的实际赤字六千万法郎，以一千万比索西属美洲汇票（议定3.75法郎／比索）加2,400万法郎Espinosa汇票支付；1806--07年Hope--Baring双合同均为买方、按每比索3.75法郎于巴黎支付、每月上限200万法郎；1807年7月扩约时最新汇票3.85法郎。\textbf{3.75法郎相对5法郎名义＝75\%，相对巴黎5.2--5.4≈69--72\%，还须扣费用------不可笼统\literalquote{}净得75\%\literalquote{}。}

【锚】\literalquote{}British war frigates arrived in Veracruz on six opportunities on
behalf of the Gordon \& Murphy firm, loading a total of 8,651,641 silver
pesos which were transported to
England.\literalquote{}（1806年12月至1808年8月），另Hope合约的Diana一船3,829,835
pesos------战时敌国军舰替西班牙（乃至间接替法国债权链）运银，\literalquote{}封锁\literalquote{}在白银环节从来不是二元开关。

\subsubsection{二、三情形与阿尔马登的刀}\label{ux4e8cux4e09ux60c5ux5f62ux4e0eux963fux5c14ux9a6cux767bux7684ux5200}

【推·四因子连乘（汲取额×送达率×法国分成×交付折扣）】S-A波旁续任亲法：法国年实得90--330万pesos（中枢190万；1804--08已发生的峰值口径210--240万）＝450--1,650万法郎＝F4陆军预算的1.3--5.5\%；S-B波拿巴国王：0--60万------而同一套机器在1808-10至1811-02给了加的斯2,400万；S-C藩属化：0--30万，英国经贸易取4--10百万。对照1803年补贴条约年要求960--1,920万pesos（Marichal两处reales／peso换算不一致，故给区间）：覆盖率可持续口径13\%、全区间5--34\%、峰值口径11--25\%。\textbf{墨西哥白银在任何情景下都填不满补贴条约的三分之一。}运费佣金再扣0.85--0.95无来源（假设）；法国分成0.25--0.45来自单点推断。

【锚】水银是硬约束：年产250万marcs需水银17,000--25,000
quintales，阿尔马登峰值仅18,000--20,500；Brading \&
Cross记17世纪\literalquote{}运量腰斩→五年后产出崩落\literalquote{}；洪堡1802年亲见瓜纳华托\literalquote{}les
produits se vendoient à des prix très-bas, parce que le manque de
mercure entravoit l\textquotesingle amalgamation, et que toutes les
mines étoient encombrées de
minerais\literalquote{}。法国控制西班牙即控制阿尔马登，可在3--7年内把墨西哥年铸币从2,000--2,400万砍到600--1,200万pesos------但受害者排序是英国（英格兰银行储备、东印度公司对华白银、对盟国现金补贴）＞西班牙＞墨西哥矿业社会，法国自身无所得。这是破坏性工具而非吸管。洪堡：\literalquote{}la
valeur de l\textquotesingle or et de l\textquotesingle argent des mines
du Mexique est presque d\textquotesingle un quart plus petite que la
valeur du produit
territorial.\literalquote{}------新西班牙农业总产值（约2,900万piastres）大于金银产值约三分之一；把新西班牙当纯白银单位是结构性误判。

\subsubsection{三、政变的因果链与起义的硬边界}\label{ux4e09ux653fux53d8ux7684ux56e0ux679cux94feux4e0eux8d77ux4e49ux7684ux786cux8fb9ux754c}

【锚】1808年9月15日政变首领Gabriel de
Yermo当日的法律身份是一名Consolidación债务人：\literalquote{}El 13 de enero de 1806,
solamente se le exigieron 131,200. Al negarse a pagar, vino la orden de
embargo de la más valiosa de sus propiedades, la hacienda de
Temixco\ldots{} En septiembre de 1808, había liquidado 86,000 pesos.
Estaba a punto de resolverse su asunto y sabía que tenía todo en su
contra\literalquote{}（INEHRM，印pp.54--59亲读）。补贴条约（1803）→Consolidación推广美洲（1804）→总督追索大商人与教会债权（1805--08）→被追索者在王位危机中夺权（1808-09-15）：四环之间不需要\literalquote{}民族感情\literalquote{}中介。falsifier正面检验分裂裁决：1799年\literalquote{}弯刀阴谋\literalquote{}确为独立纲领，总督阿桑萨1799年11月30日密报自承\literalquote{}existe
en esta América una antigua división y arraigada enemistad entre
europeos y criollos, enemistad capaz de producir las más funestas
resultas\literalquote{}------断层先于危机（高）；但该纲领属边缘小组（1,000
pesos、两支手枪、五十把弯刀），1809年巴利亚多利德精英阴谋的目标仍是\literalquote{}以费尔南多七世之名执政的洪塔\literalquote{}------触发论应重写为\literalquote{}1808把已存在二十年的断层从边缘阴谋提升为合法机构的公开议题\literalquote{}。

【锚】莫雷洛斯1813年《Sentimientos de la Nación》：\literalquote{}1o. Que la América
es libre e independiente de España y de toda otra Nación, Gobierno o
Monarquía\ldots{} 20o. Que las tropas extranjeras o de otro reino no
pisen nuestro suelo, y si fuere en ayuda, no estarán donde la Suprema
Junta.\literalquote{}第9条官职只给美洲人，第16条限定港口与10\%关税------起义纲领在文本层面先验排除\literalquote{}法国客户国式墨西哥\literalquote{}。

\subsubsection{四、古巴：三重否决的均衡}\label{ux56dbux53e4ux5df4ux4e09ux91cdux5426ux51b3ux7684ux5747ux8861}

古巴的\literalquote{}永远忠诚\literalquote{}是军官团（1808年7月27日洪塔案仅73签名、46为半岛人，被蒙塔尔沃一句话压下）、奴隶主（恐惧圣多明各重演）、糖商（需要合法旗号保住美国市场）三重否决的均衡；后两重对法国政策高度敏感------法国是1802年奴隶制的复辟者，其大陆体系与\literalquote{}给古巴糖留美国出口\literalquote{}直接冲突。【锚】杰斐逊禁运是\literalquote{}大陆体系式贸易断绝施于加勒比种植园经济\literalquote{}的唯一实测：\literalquote{}descendió
a la mitad el comercio con Norteamérica, lo que triplicó los precios de
los productos de importación y no pudieron venderse dos tercios de la
cosecha de azúcar de
1808\literalquote{}------一个收获季内对美贸易腰斩、进口物价三倍、三分之二糖收成滞销。1848年归属：西0.40--0.55／美0.20--0.35／英0.05--0.12／独立0.03--0.08／法系≤0.05；任何法国主导安排下古巴脱离概率上升，去向主线是美国。未决：Buist
pp.295--356的\literalquote{}75\%\literalquote{}中枢（已由X1全文核为3.75法郎合约价）；von
Wobeser附录、Romero Sotelo汞与矿产1810--21、TePaske \&
Klein国库序列、Hamnett／Anna／Brading专著、AGI Ultramar 833、ANF AF IV
1082／1608；半岛人口7,906（Navarro y Noriega）vs
70,000--80,000（洪堡）近十倍之差直接决定\literalquote{}统治层厚度\literalquote{}。

\begin{center}\rule{0.5\linewidth}{0.5pt}\end{center}

\subsection{第三十七章　南美：加拉加斯、拉普拉塔、秘鲁、新格拉纳达与英国渗透}\label{ux7b2cux4e09ux5341ux4e03ux7ae0-ux5357ux7f8eux52a0ux62c9ux52a0ux65afux62c9ux666eux62c9ux5854ux79d8ux9c81ux65b0ux683cux62c9ux7eb3ux8fbeux4e0eux82f1ux56fdux6e17ux900f}

\subsubsection{一、英国自判征服无望}\label{ux4e00ux82f1ux56fdux81eaux5224ux5f81ux670dux65e0ux671b}

【锚】卡斯尔雷1807年5月1日内阁备忘录（Castlereagh Correspondence
vol.VII，本轮全文）：军事征服\literalquote{}relieve us from the hopeless task of
conquering this extensive country, against the temper of its
population\literalquote{}，转向以\literalquote{}辅助者与保护者\literalquote{}身份、用当地武装促成独立；自设收益仅两项------剥夺敌人资源、为英国制造品打开市场。首选政体不是共和国而是波旁系王子的立宪君主国：\literalquote{}it
ought to be our policy to support the pretensions of some member of the
Bourbon family, as naturally connected with that Monarchy, rather than
any Prince more immediately connected with our own
Crown\literalquote{}（考虑奥尔良公爵，由米兰达与迪穆里埃随行）。对价见1806年6月9日科克伦--米兰达条款：排他通商＋他国10\%差别关税＋债权执行＋领事。波帕姆1803年备忘录自承：\literalquote{}if
the enemy, either from France or Spain, is able to land a European force
on the coast, and by any immediate sacrifices or concessions to the
natives obtain their favour, the enterprise from this country may have
the greatest difficulties to
encounter\literalquote{}------同期文书对\literalquote{}法国以自治／贸易让步先发制人\literalquote{}可行性的直接支持。1806年布宜诺斯艾利斯的翻脸机制由西班牙方内部报告独立证实：\literalquote{}The
English had been invited under the persuasion that they would declare
the independence of Buenos Ayres; but those who had so calculated were
mistaken.\literalquote{}英国1807年自设南美兵力上限13,200（对美战争之虑消失可再抽3,000），而奥克穆蒂认为仅有效占领拉普拉塔一区即需15,000；扶植方案的物资核心是\literalquote{}A
stand of arms for 20,000 infantry and 10,000 cavalry\literalquote{}。

\subsubsection{二、法国的投送实测：两名专使与一条双桅船}\label{ux4e8cux6cd5ux56fdux7684ux6295ux9001ux5b9eux6d4bux4e24ux540dux4e13ux4f7fux4e0eux4e00ux6761ux53ccux6845ux8239}

【锚】Torrente（1829年保王派回溯叙述）：1808年7月15日一艘英国战舰与一艘携带两名专使的法国双桅船同日抵加拉加斯，\literalquote{}por
un bergantín francés que traía á su bordo dos Comisionados con
instrucciones para hacer reconocer á José Napoleón por Rei de España i
de las Indias\ldots{} se resolvió la jura del Soberano legítimo, i el
desprecio de las órdenes del intruso\literalquote{}；拿破仑使者萨塞奈\literalquote{}la llegada en el
día 13 de Mr. Sassenay, emisario de Napoleón para hacer reconocer por
Rei de España i de las Indias á su hermano
José\literalquote{}（1808年8月13日蒙得维的亚）反而使对费尔南多七世的宣誓提前完成，利涅尔斯的法国血统成为埃利奥与蒙得维的亚组织保王洪塔的理由。保王军主体是美洲castas：\literalquote{}Hubo
época en que el Comandante Boves reunió 12.000 de ellos, entre los
cuales habría apenas 200 europeos; i los Vireyes Abascal, Pezuela, i La
Serna llegaron á tener de 15 á 20.000 hombres sobre las armas, no
entrando á veces por mil los peninsulares\literalquote{}。米兰达1806年远征在科罗被\literalquote{}una
fuerza igual á la de los espedicionarios, compuesta la mayor parte de
Indios\literalquote{}驱逐------证伪英方情报\literalquote{}印第安人是政府天敌\literalquote{}。（不应把史实投送值当所有反事实的绝对物理上限；萨塞奈1892年专著被Gallica验证页拦截，法方训令未阅。）

【锚】莫雷诺1809年《Representación》：\literalquote{}más de cuatro millones fueron
introducidos, y entre confiscaciones y derechos apenas recogió la aduana
noventa y seis mil pesos, debiendo haber entrado en ella millón y
medio\literalquote{}------禁令执行率6.4\%；蒙得维的亚洪塔开放贸易后以二十笔生意入账逾七十万pesos并按时发饷，布宜诺斯艾利斯欠饷五个月以上。秘鲁（Hamnett
IEP
2000）：正规军1809年1,500→1813年2月8,000＋民兵四万；1816年债务一千一百万pesos；Morán--Yarango
2021给AGN利马主财库逐年战费表及其\literalquote{}不含军饷\literalquote{}口径警告；Echeverri：波帕扬矿奴保王，1781年普查24\%印第安／17\%奴隶。Torrente革命前贵金属年产表总计36,063,272
pesos（秘鲁栏把1780--89十年合计标为年产，编者仅加\literalquote{}So in
MS.\literalquote{}，不可直接用作1806年产量）。

\subsubsection{三、两个变量决定南美}\label{ux4e09ux4e24ux4e2aux53d8ux91cfux51b3ux5b9aux5357ux7f8e}

【推·中高】南美1803--48的走向由两个变量决定------是否存在被本地精英承认的合法国王、出海通道是否开放。法国只能操作前者（借西班牙王位名义），英国只能操作后者（航运与同盟资格），二者都无法武力占领任何总督区。法国\literalquote{}助西班牙保美洲\literalquote{}的唯一可行版本是B4b（保波旁）＋B6b（制度化中立船／自由港），收益上界400--800万pesos／年白银（毛口径，未乘送达率与法国分成------MX四因子90--330万为承重数，红队I项）并剥夺英国\literalquote{}唯一供货者\literalquote{}地位；派兵版本全部不可行；\textbf{S6在美洲的代价是确定的}------合法性归零、白银与市场同时丧失，P2的行省化净收益须扣除美洲现值；S3在南美的净效果是给英国送钱。1848年地图概率带：S2主权实体中位4（3--7）、\literalquote{}帝国存续但联邦化\literalquote{}0.40--0.55（含新西班牙------与MX冲突，P8裁定全西属美洲保留≤0.10）；S1中位7；S0／S3中位9；S6中位10（9--12）且法国收益为零。跨情景不变量：巴西不回欧洲法制秩序、上秘鲁归属恒为争端、蒙得维的亚---布城港口之争持续、巴拉圭趋于孤立自治。未决：Lynch／Halperín
Donghi／Adelman／Anna／Street／Kaufmann／Waddell／Colombeia未获（工具限制）；卡斯尔雷\literalquote{}Old
Spain is, in fact, a French
province\literalquote{}段的页眉OCR\literalquote{}180M\literalquote{}须确认属1807年5月而非1808年件。

\begin{center}\rule{0.5\linewidth}{0.5pt}\end{center}

\subsection{第三十八章　国际金融家：Hope、Baring、Rothschild、Ouvrard与白银合同}\label{ux7b2cux4e09ux5341ux516bux7ae0-ux56fdux9645ux91d1ux878dux5bb6hopebaringrothschildouvrardux4e0eux767dux94f6ux5408ux540c}

\subsubsection{一、反既有债权人，亲新设债权人}\label{ux4e00ux53cdux65e2ux6709ux503aux6743ux4ebaux4eb2ux65b0ux8bbeux503aux6743ux4eba}

拿破仑体制不是\literalquote{}反信用\literalquote{}，而是对被并金融中心适用\literalquote{}三分之一统一债\literalquote{}规则（1797年法国本土；1805年热那亚1.44亿里拉→1,850万法郎、1.5\%、一举否认近1.22亿，且至帝国崩溃仅一半获登记；1810年荷兰tiercering
78m→26m），对自设藩属公债则给预算优先权＋法定利息上限（1805年意大利王国财政总法第13条）------而该债务在法国战败后仍被维也纳体系承认、1816--20由外交公约分摊、付息至1866年。\textbf{法国霸权在金融上的失败是政策性而非结构性的。}

【锚】falsifier（1810--13法国无法在阿姆斯特丹／汉堡以可比利率融资）结果成立但因果须反转：法国在两地一笔未借，因为它选择没收而摧毁了市场------兼并汉堡后课4,800万法郎罚金并搬空银行7,506,956
Mark Banco（准备率实测100.235\%）。决定性反例：\literalquote{}Laffitte：La France ne
trouverait pas un petit écu à emprunter sur aucune place de
l\textquotesingle Europe.\literalquote{}（1816）；次年一个战败被占领的法国以发行价55募到3.15亿法郎，钱主要是法国人自己的。应下调的不是\literalquote{}法国信用国家化\literalquote{}的可行性，而是它在拿破仑执政下的发生概率（估0.15--0.30）。

\subsubsection{二、白银链的物理实态：封锁封的是船，封不住债权}\label{ux4e8cux767dux94f6ux94feux7684ux7269ux7406ux5b9eux6001ux5c01ux9501ux5c01ux7684ux662fux8239ux5c01ux4e0dux4f4fux503aux6743}

【锚】Buist（At Spes Non
Fracta，1974，本轮全文）：经Parish代理自韦拉克鲁斯出口约645万比索，而Hope分类账显示直运欧洲仅两批共60,000比索；主渠道是以比索向北美商行预付运欧货款、商行再加年息6\%对欧洲收货人开票。\literalquote{}as
a general rule Hamburg was preferred for remittances to the Continent
because it was on the exchange market there that the deficit in
Britain\textquotesingle s trade with Russia was
adjusted.\literalquote{}------为避免压阿姆斯特丹--巴黎汇率，汉堡Matthiessen \&
Sillem、鹿特丹Littledale \& Dixon、安特卫普David Parish \&
Co.与法兰克福Bethmann
Brothers对Hope开票并自行汇巴黎。摩擦成本可量化：\literalquote{}The total percentage of
dues gradually increased, from 11 \% in May 1806 to 33 \% in the second
half of
1807\literalquote{}------跨阵营白银转移的摩擦成本1806年5月宽松期≤23.5\%（税费11\%＋保险12½\%），1807年下半年加严期≤62.5\%（税费33\%＋保险12½\%＋所有权洗白折扣17\%）；加严把摩擦推高约2.7倍但未推到100\%；要到100\%需同时控制墨西哥、美国与英国海军。

【锚】法国承包商体系的年净动员上限被1804--05实测为约6,000万--1亿法郎：1804年4月对西3,200万汇票＋6月1.5亿要求（先扣4,800万欠货款）→\literalquote{}At
the end of September a mere 1,200,000 francs
remained.\literalquote{}（1805年法兰西银行）→11月挤兑、Compagnie披露赤字7,000万（莫利安复查后真实赤字6,000万、负债1.42亿）。在统合公债实施前，法国任一年的超预算军事支出受此硬顶约束------\textbf{奥斯特里茨在财政上被\literalquote{}必须在十二月前结束\literalquote{}约束}。

\subsubsection{三、金融中心的四种职能与罗斯柴尔德}\label{ux4e09ux91d1ux878dux4e2dux5fc3ux7684ux56dbux79cdux804cux80fdux4e0eux7f57ux65afux67f4ux5c14ux5fb7}

【锚】1801--06年法兰克福资本市场的边际买盘是黑森选帝侯的英国补贴回流（Ullmann，本轮全文）：\literalquote{}Mit
jährlichen Neuinvestitionen von knapp 1 Mill. fl. sowie der Wiederanlage
zurückgezahlter Kapitalien in Höhe von fast 1,5 Mill. fl. war der Fürst
zwischen 1801 und 1806 der mit Abstand größte
Anleihezeichner.\literalquote{}------英国的补贴总额（£4,200万，B2）同时是欧陆资本市场的流动性来源；S1会在关闭战争的同时关闭这条管道，\literalquote{}法国胜利→法兰克福崛起\literalquote{}不成立（其最大发行客户奥地利1797年已退出，最大商行Bethmann
1810年被法国罚逾36万法郎）；罗斯柴尔德1806年承接的正是这条管道的托管权。把\literalquote{}中心\literalquote{}拆为发行／清算／贵金属／对外债权四职能：S2／S5下1848年是\literalquote{}发行在巴黎、清算在伦敦、白银在汉堡\literalquote{}的三分天下；伦敦因英镑汇票、英格兰银行与美洲贸易融资三项结构理由在任何情景不丢清算职能。

【锚】Ferguson（据家族账簿）：\literalquote{}In 1815 the combined capital of the
Rothschild houses in Frankfurt and London was at most £500,000. In 1818
the figure was £1,772,000; in 1825 £4,082,000; and in 1828
£4,330,333.\literalquote{}（Baring对应£374,365／429,318／452,654／309,803）；1836年五家合计略逾£600万；\literalquote{}1818--44年平均年利润上法兰克福与巴黎两家均高于伦敦家\literalquote{}------削弱\literalquote{}没有伦敦就没有罗斯柴尔德\literalquote{}的强命题。波拿巴秩序下形态转为以巴黎为中心的帝国财政承包家族，1836年当量期望约£230万（R-A
0.40--0.55／R-B
0.15--0.25／R-C被行政没收0.20--0.35）；法制欧洲给它更早的法律平等与更差的财产安全。

【锚】Ziegler：1818年艾克斯会议修约------\literalquote{}the principals among them -
Metternich, Nesselrode and Hardenberg - had all participated personally
in the loan and would have lost heavily if the terms had not been
amended.\literalquote{}------当政治家个人持有主权债时，债务重组的政治成本上升、违约概率下降。可迁移至1815--48秩序稳定性机制。

\subsubsection{四、给他卡的硬参数}\label{ux56dbux7ed9ux4ed6ux5361ux7684ux786cux53c2ux6570}

并省金融代价＝名义1/3×登记率1/2≈面值1/6；跨阵营白银摩擦≤23.5\%至≤62.5\%、从未到100\%；承包商体系年净动员上限6,000万--1亿法郎；欧洲外汇市场深度£10万（官方单点）--£70万（多节点私人网络）；1814年1月全欧两月内可搜集的法国硬币仅£15万→英国1813--15的约束是结算而非财政；1815--48资本市场年均新发行包络（模型非观测）S0
£15--26m、S2／S5 £19--34m、S3 £9--17m、S6
£6--13m；法国rentes均衡收益S5下4.0--6.5\%（做统合公债）或7.5--9.0\%（不做）。最高优先低成本改判项：关闭加来--敦刻尔克--格拉沃利讷的金几尼换伦敦汇票通道（1812年4月样本£27,300金→£65,798面值汇票，比值2.41）------该通道因莫利安\literalquote{}英国金属外流是其虚弱之象\literalquote{}的重商主义误判被容忍并发许可；须与F3（保护清算节点）的内在冲突同时评估。未闭合：汉堡agio月度序列（Soetbeer
1855／Preis-Courant／Staatsarchiv
Hamburg）------可直接判定封锁伤金属还是只伤信用；Herries回忆录只读索引；Gille正文、Lévy-Leboyer、Bouvier未取；Buist的Account
X重建含作者自标假设（53船货2.9m美元、Parish 721,000美元）不可当实测。

\begin{center}\rule{0.5\linewidth}{0.5pt}\end{center}

\subsection{第三十九章　大陆体系作为执行机器与地下经济}\label{ux7b2cux4e09ux5341ux4e5dux7ae0-ux5927ux9646ux4f53ux7cfbux4f5cux4e3aux6267ux884cux673aux5668ux4e0eux5730ux4e0bux7ecfux6d4e}

\subsubsection{一、走私成本是执行强度的价格}\label{ux4e00ux8d70ux79c1ux6210ux672cux662fux6267ux884cux5f3aux5ea6ux7684ux4ef7ux683c}

【锚】莱茵边界走私保险费率的完整序列（Aaslestad \&
Joor编，Rowe章，印p.212）：\literalquote{}The rate for crossing the Rhine frontier, in
contrast, ran at about 30 per cent in July 1809, and would rise to about
50 per cent at the end of the year, a rate that was maintained into
1811. This compares to only 6 per cent for the Rhine in 1800, 10 per
cent from 1801 to 1802, and 26 per cent from October 1806 to June
1808.\literalquote{}------1812年编制增至35,000后保费未再升：边际执行产出在≈50\%处饱和。Heckscher（印p.210）：\literalquote{}In
Strassburg there were \textquotesingle insurers\textquotesingle{} of
different grades, the chief of which charged a commission of from
4{[}0{]} to 50 per cent.; in 1809 it was considered that the expenses of
passing the frontier of France were, as a rule, 30 per cent., while the
above-mentioned new customs line between Rees and Bremen could be broken
through for 6 or 8 per
cent.\literalquote{}------同一时点不同边界段的执行强度差可达5--8倍；新建关线的边际执行力不足老线的七分之一。

【锚】机器的规模：\literalquote{}in 1812 the total number of customs agents in the
French Empire peaked at approximately 35,000, including four thousand
agents de bureaux, and 30,750 organized into
\textquotesingle brigades\textquotesingle{} that patrolled the
frontiers. Of the latter 3200 deployed to the départements réunis
bordering the
Rhine\literalquote{}------莱茵并合四省每十万人口241名海关人员（相当于今日英格兰威尔士警察密度），但仅约3\%为本地人、逾80\%来自旧法国：\textbf{外来占领式而非本地嵌入式}。财政自毁（模型，可复算）：1809年法国关税净收11.6m，而同年27,200名海关人员仅按最低级préposé年薪500法郎计的纯工资即13.6m------工资／关税＝117\%（海关自1806年3月16日兼征盐税，但盐税不需要三万名边境旅，方向不逆转）。

【锚】缉获率的两个不相容锚点：Rist述汉堡海关\literalquote{}每三四个背货人抓一个\literalquote{}（人次口径，被抓者多为诱饵）；Marzagalli据档案\literalquote{}In
August 1810, customs agents seized 9300 kilograms of colonial goods at
the gate between Hamburg and Altona. Moreover, this quantity seems to
have represented at most five percent of the quantity
smuggled.\literalquote{}（重量口径）。两渠道定理：批发渠道可定价压制，零售渗透不可关闭。价格楔子：巴黎糖8→12法郎／公斤
vs
伦敦1.35--2；特里亚农原糖税从价等价150--222\%；残差租金占巴黎价37.5--63.8\%；殖民品／粮食物价比恶化9.2--17.5倍（梅梅尔谷价−60至−80\%、不来梅麦价−62\%）。

\subsubsection{二、特里亚农的官方经济学：\literalquote{}把走私成本国有化\literalquote{}}\label{ux4e8cux7279ux91ccux4e9aux519cux7684ux5b98ux65b9ux7ecfux6d4eux5b66ux628aux8d70ux79c1ux6210ux672cux56fdux6709ux5316}

【锚】Heckscher（印p.214）：海关总监Collin de
Sussy与Montalivet提议国家与走私者竞争，\literalquote{}the importation of the hitherto
forbidden goods was to be permitted, but only on payment of a duty that
exactly corresponded to an amount which, as we have seen, the smuggling
business had previously cost. In that case no more goods would come into
the country than had been the case beforehand, but the profit would fall
to the state instead of to the
smugglers.\literalquote{}------关税水平内生于执行失败程度，而非保护目标。1810年特里亚农--枫丹白露的二分法（禁工业品／重税殖民品）不是意志薄弱，而是机器触顶后的理性重整：\textbf{执行机器能把走私的从价成本推到AVE
50--62\%（硬上限55±8\%），足以排除有大陆替代品的英国工业品（成本劣势35--60\%），远不足以排除无替代品的殖民品。}Ellis（Alsace，1981，印p.286，本轮亲读）：\literalquote{}What
was envisaged was not a Continental Zollverein in which all the
resources of the Empire, of the subject and allied states, could be
harnessed through reciprocal favours and equitable sacrifices to the
common cause of defeating Britain, but rather a vast
\textquotesingle Uncommon Market\textquotesingle{} geared to French
interests.\literalquote{}

结构后果：封锁越严，法国财政越依赖它所禁止的贸易；抽租收入主要来自并地式存量清算（1810--11一次性池≈3.07亿法郎：关税105.9＋拍卖≈150＋荷尔斯泰因42.5＋法兰克福9）而非常态流量，财政需求持续生产新的兼并动机，每次兼并又延长不可守的边界------这条螺旋是内生的，不需要俄国冬天。执行失败→兼并的螺旋有完整史实链（1795默兹线→莱茵商人求并→1798得偿→\literalquote{}the
problems of policing the new frontier without controlling the right bank
beyond provoked calls after 1798 for intervention on the right
bank.\literalquote{}→Kehl
1808、荷兰与汉萨1810、提契诺1810）；刹车盘也是执行地理学：贝格1811年4月四千人签名请愿求并入被拒，理由之一是莱茵河作为关税边界远优于\literalquote{}毫无主要天然障碍\literalquote{}的贝格东界------\textbf{为行省化提供一个非政治的、技术性的边界判据。}

\subsubsection{三、Ellis法则与保护版}\label{ux4e09ellisux6cd5ux5219ux4e0eux4fddux62a4ux7248}

\literalquote{}the efficacy of the Blockade stood in inverse ratio to the extent of
Napoleon\textquotesingle s military
involvements\literalquote{}（Ellis印p.202）------S3与大战不相容，要最大化封锁必须先停战。对英伤害的峰值需要\literalquote{}大陆严格控制∧英美断裂\literalquote{}，法国只掌握一半变量。S3严格版年净账为负（关税15--35
vs 执行成本28--38.5 vs
税基损伤−60至−110）；许可证版财政可持续但目的自毁；本卡提出史实未走的第三条\literalquote{}保护版\literalquote{}：中率财政关税40--60\%（恰压在保费上限之上）＋普遍准入取代逐船许可＋不再为封锁而并地，估法国关税可稳在60--90m／年、边境旅可缩至1.8--2.2万人。\textbf{红队指出：这一方案中\literalquote{}关税＝走私成本国有化\literalquote{}有当时提案（特里亚农本身），\literalquote{}普遍准入\literalquote{}是英国商人的诉求，X2自标MODEL------两半的证据等级不同，第五十一章按此标注。}S5关税同盟化：A法国优先分层准入0.50--0.60（1830）／0.35--0.45（1848）；B有限关税同盟0.15--0.25／0.25--0.35；C解体0.15--0.25／0.20--0.30；D逆向同盟0.05--0.10／0.05--0.15------推动力是财政摩擦（Heckscher印p.227\literalquote{}一次完税后自由流通\literalquote{}已实施），阻力是\literalquote{}la
France avant tout\literalquote{}与皇帝寿限。未决：Marzagalli
1999全书、Crouzet两卷、Dufraisse原文、Tarlé俄意专著均转引；三套法国关税收入口径（Heckscher／Marion／AHAD：1807＝77m、1812＝109.5m）量级不合并列；海关总署分年支出账、分港缉获量值表、许可证分年发放数为改判证据的关闭条件。

\begin{center}\rule{0.5\linewidth}{0.5pt}\end{center}

\subsection{第四十章　科技与工业能力：欧洲技术前沿在法制秩序下会落在哪}\label{ux7b2cux56dbux5341ux7ae0-ux79d1ux6280ux4e0eux5de5ux4e1aux80fdux529bux6b27ux6d32ux6280ux672fux524dux6cbfux5728ux6cd5ux5236ux79e9ux5e8fux4e0bux4f1aux843dux5728ux54ea}

\subsubsection{一、核账}\label{ux4e00ux6838ux8d26}

【锚】蒸汽：英国1800年35,000马力、1830年160,000（Kanefsky系列经Crafts
2004表3）；法国1810年约200台（通行口径）、1839年才33,000马力、1848年65,000（NBER
a01172a源Annuaire 1922）；\literalquote{}En 1815 la France redécouvrit
l\textquotesingle industrie anglaise et s\textquotesingle aperçut alors
qu\textquotesingle en matière de machines à vapeur elle était en retard
de trente ans\literalquote{}；蒸汽纺纱厂空窗二十五年（1787→1812 Mulhouse
Dollfus-Mieg）。纱锭：Montalivet 1813年2月25日Exposé，\literalquote{}le département de
la Seine comptait 133 448 broches en activité \ldots{} 10 p. 100 du
nombre total de broches en activité dans
l\textquotesingle Empire\literalquote{}------帝国总纱锭≈133万（根特Escaut省134,000锭为帝国第一省）vs
英国≈460--500万，比值≈0.28；帝国自产原棉仅覆盖需求12\%。生铁：法国困于木炭（Creusot焦炭孤例1809--12年产2,300--3,000吨），而1780年代法国曾是欧洲最大木炭铁国；比利时以法国六分之一人口产近同量煤（1840--44：4.1M
vs 3.5M吨），1847年焦炭生铁率89.5\% vs
法国42.6\%。法国领先项：化学（Leblanc碱价80--100→10法郎／100kg）、Jacquard、Chappe电报网（1813年已达阿姆斯特丹、威尼斯、美因茨；峰值534站／5,000余公里）、互换性制造。

【锚】互换性制造的发明与遗忘（Alder）：Blanc
1790年11月20日荣军院演示互换枪机（Vincennes造约1,000件）；Vandermonde工场1794--95；帝国回归工匠分散生产；至1850年代\literalquote{}the
French military had never heard of the method---let alone of Honoré
Blanc\literalquote{}。这是本轮识别的最大未用潜力：S5重启Blanc路线技术无障碍、政治被否（低置信被采纳）。

\subsubsection{二、封锁是发明之乳母，不是发明之母}\label{ux4e8cux5c01ux9501ux662fux53d1ux660eux4e4bux4e73ux6bcdux4e0dux662fux53d1ux660eux4e4bux6bcd}

Heckscher的机制操作化：\literalquote{}necessity is the nurse of
invention\literalquote{}：封锁加速配方型技术（Leblanc碱、Koechlin染料1810--11超英），冻结载体型技术（动力织机帝国仅根特／Sennheim零星；滚筒印花晚英二十年；焦炭孤例）；1814／1817和平冲击测试------未完成转换者全部倒退（法棉业1814崩溃、萨克森印花业1817被击穿）。falsifier（大陆1806--15蒸汽年增量≥英国30\%）不成立：至多5--10\%（萨克森棉纺几无蒸汽、瑞士走水力、Cockerill
1817才造蒸汽机）→扩散加速分支不上调，载体型技术差距25--30年主线维持。英国技术封锁时间轴（Jeremy
1977）：禁技工移民1719起、禁机器出口1750起（清单不含蒸汽机与机床）；1824废移民禁、1825机器出口改许可、1843年8月全废；贸易委员会1819年自认对技工\literalquote{}从来无效\literalquote{}------禁令只抬高转移成本不能归零，1815年和平＋废禁才是洪峰。

【锚】甜菜糖口径（供全书）：1813年公开展望≈350万公斤／334厂；内政部后报\literalquote{}they
had only got 1,100,000 kgs. of
sugar\literalquote{}，334许可证仅158份实际使用（Heckscher印p.293，上轮页图核）。

\subsubsection{三、1848年前沿位置与铁路}\label{ux4e091848ux5e74ux524dux6cbfux4f4dux7f6eux4e0eux94c1ux8def}

【推】发明前沿（蒸汽／机床／铁路／廉价铁）一切情景留在英国（高置信）；法国持化学、丝机、光报、爆破弹舰炮四利基；采用差距S5收窄（法达史实的1.1--1.4倍）、S3双输（英国＝史实54--79\%且大陆更糟）、S6最慢（财政透支＋资本外逃＋治安挤占工程师团）；S6唯一技术正项＝小步兼并贝格。\textbf{技术论据支持\literalquote{}小行省化\literalquote{}否决\literalquote{}大行省化\literalquote{}。}铁路：帝国版\literalquote{}比利时模式\literalquote{}（国家干线＋桥路／矿务团执行＋特许营运＋战略线优先）比史实法国1832--42议会十年争吵更可能；但资本瓶颈（tiercement毁阿姆斯特丹资本市场＋预算竞争）使1848帝国（含比-莱茵）网仅2,500--5,000公里（中心3,500）≈史实法＋比之和；S3下英国投入品断流→铁路时代后移5--15年；S6\textless1,000公里。Chappe网S5下马德里／汉堡／罗马线1815--25补齐为大概率，加速的是中枢--边疆控制回路（政令可印1--2日）而非市场信息------F1的环路时延8--12天仅部分缓解，S6治理密度约束仍在。比-莱茵-萨克森工业带真实存在（Mons--Liège--Aachen煤铁＋根特棉＋Verviers毛＋萨克森棉）但被帝国关税几何主动切割；Beugnot／Bacher的关税同盟倡议被\literalquote{}la
France avant
tout\literalquote{}压倒；任何情景下1848年此带都是欧陆重工业的心脏。未决：大陆各国逐年新装蒸汽机统计；Landes／Pollard／Henderson／Horn／Chassagne／Daumas／Bret／Belhoste全文；Jacquard
1812年11,000台与Shrapnel史述为同源传统。

\begin{center}\rule{0.5\linewidth}{0.5pt}\end{center}

\subsection{第四十一章　思潮、民族主义、宗教复兴与社会运动}\label{ux7b2cux56dbux5341ux4e00ux7ae0-ux601dux6f6eux6c11ux65cfux4e3bux4e49ux5b97ux6559ux590dux5174ux4e0eux793eux4f1aux8fd0ux52a8}

\subsubsection{一、1813年的祛魅}\label{ux4e001813ux5e74ux7684ux795bux9b45}

【锚】Fichte《对德意志民族的演讲》接受史（Breazeale导论检索摘录）：\literalquote{}Nearly
six hundred copies of the Addresses were sold within a
month\literalquote{}；1813--14年解放战争语境中其名与文本\literalquote{}很少被援引\literalquote{}；1824年再版被普鲁士审查官拒绝；1862年百年诞辰才被特赖奇克封为民族先驱；法占当局对讲演无动作。Tugendbund：\literalquote{}die
maximal 700 Mitglieder besaß und Ende 1809 vom preußischen König
förmlich aufgehoben
wurde\literalquote{}------1809年现存名录23个分会、在册约620人，全部在城镇且与占领／围城经历地理相关，乡村零分会；法奥警探高估其重要性。1809年奥方号召无响应。1813年普鲁士动员构成：\literalquote{}Ende
Mai waren 7000 f. I. zu Fuß und 3000 zu Pferd
aufgestellt\literalquote{}（志愿猎兵约一万）；总动员280,720（Hagemann转引）≈男性人口3\%------志愿者约占动员额十分之一，\literalquote{}民族起义\literalquote{}的人力主体是国家强制征召。Planert（2007，经两篇专业书评亲读转引）：\literalquote{}inherited
loyalties, localisms, and confessional attitudes determined the world
view and actions of most
people\literalquote{}（p.659）；1813年\literalquote{}全德民族起义\literalquote{}在普鲁士之外不成立（p.489）；\literalquote{}德意志民族战争\literalquote{}是事后建构（p.620）。

falsifier（\literalquote{}1806--12占领区民族话语已达1813水平\literalquote{}）不成立；但祛魅双向：1813基准本身被同源证据削低，故反事实的减速幅度比神话口径小；且占领磨损真实存在（Tugendbund地理＝占领地理）。不得读成\literalquote{}占领无代价\literalquote{}。

【锚】挪威案例（Glenthøj \& Ottosen pp.186--190亲读）：\literalquote{}the mood in
Norway is such that the conclusion of peace with Britain is regarded as
the only means of saving
Norway\literalquote{}------英国封锁＋1812饥荒＋哥本哈根调度失灵→事实自治（1807政府委员会）＋差异化受苦→分离主义萌芽，无需任何\literalquote{}解放战争\literalquote{}或民族主义宣传。\textbf{此机制在法制欧洲照常运转，产物是藩属分离主义而非泛德意志主义。}

\subsubsection{二、三大主义在法制欧洲的形态与1848}\label{ux4e8cux4e09ux5927ux4e3bux4e49ux5728ux6cd5ux5236ux6b27ux6d32ux7684ux5f62ux6001ux4e0e1848}

自由主义分体制内法典--宪政请愿主流（莱内／附加法模板）与科佩流亡支（贡斯当对统一化的批判为同时代机制预言，ch.12--13亲读）；意大利民族主义快于德意志（统一的制度容器）；社会主义偏向圣西门技术官僚＋密谋两极，议会改良支最弱。【锚】宪章请愿量级（Hansard亲读）：1842年\literalquote{}very
nearly 3,500,000\literalquote{}签名；1848年委员会十三名文书员十七小时点验\literalquote{}the number
of signatures has been ascertained to be
1,975,496\literalquote{}（宣称5,706,000）------百万级群众签名政治需合法请愿渠道＋廉价印刷＋工业人口三者齐备，S5欧陆缺第一项，同等压力必改道表达。卢德运动三因结构（Jarrige
pp.1--3）：\literalquote{}sur 9000 ouvriers des filatures, 3000 seulement étaient
employés\literalquote{}（1811年6月曼彻斯特）×1809年废除保护性旧规×1799--1800结社法入地------市场崩溃项直接受大陆体系政策分支调制。烧炭党三口径（1821年回忆录，反烧炭党官方视角，无独立互证）：\literalquote{}They
amounted to from twenty-four to thirty thousand from the very beginning
of their
establishment\literalquote{}；卡拉布里亚一地50,000--60,000（Nunziante）；1820年前后传闻200,000＋；士兵服从会社意志而非军纪。

【锚】Broers直读（印pp.179--182）：拿破仑改革在意大利\literalquote{}教化群众\literalquote{}目标彻底失败，民众虔诚反因迫害强化------\literalquote{}emerged,
if anything, strengthened by Napoleonic
persecution\literalquote{}；经拿破仑时代，教会不再可能从属俗权且\literalquote{}比其自由派对手更能自我重塑\literalquote{}。S5若维持教俗决裂，等于把十九世纪最强的自我重塑型反对机器养在体制外。

\textbf{主线（中高）}：无1813年，德意志政治民族主义降档慢行15--25年（文化民族主义不减速），载体从战争记忆换为关税--征兵申诉；群众化在任何情景都需国家级动员事件点火。法制欧洲意识形态死穴排序：\textbf{天主教正统诉讼＞继承合法性真空＞民族主义＞自由主义＞社会主义}。1848：体系级复合危机落在1845--52窗口概率≈0.6--0.8（歉收外生恒在，承E1）；取欧陆同步革命波形态≈0.35--0.55；主形态＝巴黎辐射型立宪--面包--继承复合危机，诉求语言是宪章／法典／反征兵而非民族统一；帝国使巴黎成为单一点火点，同步机制强于史实1848。思想／继承／教会风险在P8聚合时勿双计。未决：Planert／Hagemann原书、Rath、Thompson、Chase全文；约二十项【权重记忆】数字（报刊发行量、识字率、大学生数、卢德镇压兵力）未亲核。

\begin{center}\rule{0.5\linewidth}{0.5pt}\end{center}

\subsection{第四十二章　生态、人口与资源储量：欧洲能养多少军队与舰队}\label{ux7b2cux56dbux5341ux4e8cux7ae0-ux751fux6001ux4ebaux53e3ux4e0eux8d44ux6e90ux50a8ux91cfux6b27ux6d32ux80fdux517bux591aux5c11ux519bux961fux4e0eux8230ux961f}

\textbf{储量不等于到厂可用流量}------这是本章对全书的一句话贡献。185条资源行分十六维、8个可复算假设场景。

\begin{itemize}
\tightlist
\item
  木材：1810--12年七组船厂逐年7／9／8艘下水推翻\literalquote{}年六舰物理不可能\literalquote{}；但长期熟材与修船账未明。1669年法国森林令对公社／教会林划保留四分之一；Graham抽样1775年Clérans\literalquote{}Between
  15,000 and 20,000 trees ... only 228 trees ... considered suitable for
  their
  purposes.\literalquote{}------18世纪私林的适材率不能直接当1810年帝国公林；法国1805年橡木丰富与局部船港稀缺可并存（安特卫普流域近62\%大橡木；Caffarelli供给估算含主观折损且英国封锁出海）。
\item
  大麻：1762--82年英国进口俄麻95.9\%（Davey
  p.23）不是拿破仑时期皇家海军需求率；英麻库存与实耗未核。
\item
  硝石：1794年法国817万公斤峰值是刮墙含硝土的应急，不是人工硝床稳态；1828年约50万公斤；东印度公司500吨是合约供额非印度总产；1828产量不能倒推1813年战场短缺。
\item
  马：1813年法国马群350万，其中83.8万未满四岁；征用上限档是设值；各国马群同口径缺。
\item
  粮：敖德萨单港小麦出口1815年171,708公吨、1816年188,259公吨（原表Φύλλο1!C9／C10，C7＝1813缺值，网站将缺失单元渲为零，不可当观测零）；俄国原刊1812--48旧单位与缺报覆盖未亲阅。挪威、爱尔兰的粮食依赖见相关章。
\item
  人口：英国人口OWID／Gapminder序列早段GB／后段含爱尔兰有边界断点------Porter
  1801年GB 10.943m、1811年12.591m、1821年GB
  14.392m、UK约21.283m；不得直接拿GBR长序列做动员比。法国1800年约2,900万（quality＝LOW，不得把现代疆界序列当政体人口）。
\item
  粮秣模型：1.5公斤粮／人日、10公斤料／马日；E1的S6-max马十万与F2接口的12.5万不同，差250吨／日由马群解释。
\end{itemize}

【推】S2／S5较S3／S6的生态后勤条件更稳，但所有运输区间均为透明情景而非史实。改判证据：法舰熟材到厂结存与修舰挤出、英海军俄麻实耗与替代库存、历史帝国疆域人口与各国马群同口径、俄粮原表------这些可改变资源净约束与P2／P4数值。

\begin{center}\rule{0.5\linewidth}{0.5pt}\end{center}

\subsection{第四十三章　并合省内圈：比利时、莱茵、皮埃蒙特、利古里亚、日内瓦的整合实绩}\label{ux7b2cux56dbux5341ux4e09ux7ae0-ux5e76ux5408ux7701ux5185ux5708ux6bd4ux5229ux65f6ux83b1ux8335ux76aeux57c3ux8499ux7279ux5229ux53e4ux91ccux4e9aux65e5ux5185ux74e6ux7684ux6574ux5408ux5b9eux7ee9}

\subsubsection{一、整合的实义是\literalquote{}条件性行政合流\literalquote{}}\label{ux4e00ux6574ux5408ux7684ux5b9eux4e49ux662fux6761ux4ef6ux6027ux884cux653fux5408ux6d41}

内圈五区（比利时、莱茵左岸、皮埃蒙特、利古里亚、日内瓦／莱芒，合计≈785万人）在\textbf{国家机器层}七至十二年达标：行政≤3年、治安3--8年、征兵照额8--11年、司法即时高接受、税制形式≤3年。但\textbf{认同层}除莱芒外无一交割------1813--14年的压力测试里，无人为帝国而战，也无反法起义，制度留而主权送走。整合的实义是一桩交易：本地显贵以服从换产权与职位，法国以Rechtsstaat程序换税与兵。Woolf的判语：\literalquote{}the
policy of ralliement was successful...the French regime had won
acceptance---at least as long as its presence seemed solid and
permanent\literalquote{}（印p.203）。

【锚】省长的法籍化梯度：\literalquote{}None of the eighteen prefects who served in the
Rhenish departments originated
locally\literalquote{}（Rowe印p.87）------萨伏依＋日内瓦七任全法籍；莱茵十六任中十五法籍；比利时二十八任中仅一比利时人＋两意大利人；皮埃蒙特二十二任中九意大利人，且1802年首任六省全部本地人、1805年为并利古里亚而全部撤换------先本地化开局、后法国化正常化。副省长与司法反向：莱茵1800年259名治安法官与公证人中仅12名非本地，Roer初审法院本地：法籍＝4:1。

【锚】皮埃蒙特征兵全序列（共和七年至1810年十一个年级）：在册166,556、réformés
53,740（≈三分之一）、实得38,242（＝在册23\%）；共和十一年仅500／4,000出发，拿破仑亲笔\literalquote{}le
pays ne sera français que par la
conscription\literalquote{}；共和十三年起各省良好（Sesia除外）；Frasca总评皮埃蒙特按征兵／人口比在并合国中出类拔萃。莱茵：insoumission
30\%（共和九至十三年）→34\%（1806--10），始终高于帝国均值（21\%→31\%）且为阿尔萨斯（8\%→17\%）约两倍；desertion
13\%→3\%为帝国最优。收敛的是入伍后服从与照额交付（1811年\literalquote{}征前抵抗几乎消失\literalquote{}），拒不报到从未收敛------靠1811／1813年机动纵队压平（Roer在逃者2,000→670）。

【锚】显贵归附曲线2--9年：科隆1800年拒市长职→1803--04旧Bürgermeister三人接受任命；皮埃蒙特106名炮兵军官1798年整建制转法军、1810年\literalquote{}三十最高纳税人\literalquote{}150人中114为旧贵族；天花板明确------Roer选举团1804年拒选贝尔纳多特、1810年改选被革职法官；荣誉军团1814年38,000人中平民仅1,500（\textless4\%）；拿破仑须以\literalquote{}剥夺继承权与没收财产\literalquote{}威胁强征皮埃蒙特世家子弟入圣西尔。

【锚】税：Marion
p.321（经F4图核）1812年并合地预计毛额------比利时83M（人均24.8法郎）、莱茵37.5M（20.8）、皮埃蒙特33M（17.7）、利古里亚16M（≈25--29，人口分母带旗标）；莱芒不在表中，因它自1798年按老法国计税本身即整合标志。唯一省级实账：Roer
1803年实纳11,138,406法郎 vs
旧制约6.25M（+78\%）------名义差，非兼并净因果；Rowe判莱茵可能为财政净流入区（基建＋国家机器支出）。

\subsubsection{二、教会维度的倒U形}\label{ux4e8cux6559ux4f1aux7ef4ux5ea6ux7684ux5012uux5f62}

1801--08专约红利（比利时1803年多数教士宣誓、讲坛劝征兵；莱茵三教区隶梅赫伦、Berdolet赢得信任）→1809绝罚后系统性拉断：根特／图尔奈司铎集体拒认新主教、亚琛继任主教永未获授职、Roer
1812年360堂区空缺且司铎1,006/2,696为前修士（专约教会无再生产能力）、韦尔切利主教证词\literalquote{}三分之一信众不领圣体\literalquote{}。\textbf{教会项是系统变量而非地方变量，S6可行性直接依赖对罗马的政策分支。}

\subsubsection{三、S6的吞吐量}\label{ux4e09s6ux7684ux541eux5410ux91cf}

【推·中】每并合五百万人需12--15省长＋40--60副省长＋成建制双语司法财税中层（auditeur管道1813年仅452人）、常年宪兵370--470／百万人＋窗口期数千野战军机动预备、7--12年机器期（藩属孵化期可折抵2--7年）。并合节奏≈每代人（20--25年）一圈：1803年起算到1848年最多完成\literalquote{}内圈＋荷兰--西北德--北意\literalquote{}的两代半径。\textbf{S6与S3不兼容}（封锁摧毁海洋省份的税基与合法经济），只能搭配S2型对英共存。falsifier\literalquote{}莱茵／比利时征兵税收五年内趋同本土\literalquote{}：征兵维不成立（比利时前五年恰为最坏窗口；莱茵并入后五年insoumission不降反升；真实收敛8--11年且仅desertion与照额交付收敛），税收维部分成立（形式≤3年，人均实收与间接税深度1812年仍单列册）。可扩展性维持\literalquote{}中速有条件\literalquote{}；真约束在干部池、双语带、教会均衡三个存量瓶颈。未决：利古里亚聚合征兵序列与三省宪兵数；Almanach
impérial 1810／1812逐省官员名录；逐年逐省配额--实得表（须AN AF IV
1123--24或Vallée纸本）。

\begin{center}\rule{0.5\linewidth}{0.5pt}\end{center}

\subsection{第四十四章　并合省外圈：托斯卡纳、罗马、荷兰、汉萨、伊利里亚、加泰罗尼亚的断裂点}\label{ux7b2cux56dbux5341ux56dbux7ae0-ux5e76ux5408ux7701ux5916ux5708ux6258ux65afux5361ux7eb3ux7f57ux9a6cux8377ux5170ux6c49ux8428ux4f0aux5229ux91ccux4e9aux52a0ux6cf0ux7f57ux5c3cux4e9aux7684ux65adux88c2ux70b9}

\subsubsection{一、纸面机器十八个月建成，社会咬合由三个接口决定}\label{ux4e00ux7eb8ux9762ux673aux5668ux5341ux516bux4e2aux6708ux5efaux6210ux793eux4f1aux54acux5408ux7531ux4e09ux4e2aux63a5ux53e3ux51b3ux5b9a}

行省化的边界是政治--经济边界而非行政边界。八个单位除瓦莱外至少一个接口（征兵／教会／海关）失败，且失败模式六型互异------\textbf{S6没有单一修补方案}。封锁与兼并互为拆台；抵抗强度与现代性不同构（荷兰前提最好抵抗最密、瓦莱最差零抵抗）；断裂变量是政策交换的缺失------外圈连市场准入这一第一补偿都没给。

【锚】罗马两省1810年首次征兵（和平期、强驻军）：\literalquote{}un tiers de
réfractaires avant, une moitié de déserteurs après\literalquote{}（Madelin，La Rome de
Napoléon，1906，脚注引AN AF IV
1509警务公报1810-08-17／22，转引谱系）------官方仍评价\literalquote{}比预期顺利\literalquote{}。1811年级：\literalquote{}sur
536 conscrits du Trasimène, 80 seulement parviennent à
s\textquotesingle échapper\literalquote{}（三倍押送下仍逃80人）；罗马60名补入第17线列者军医汰除27名重病（54\%）；1812年教士宁烧洗礼登记册不交名单，特拉西梅诺诸canton开始起事；1813年法国档案自陈农村匪患已具\literalquote{}un
caractère insurrectionnel
grave\literalquote{}（Madelin印p.375）------外圈负例有保质期。

【锚】汉堡（Aaslestad章）：\literalquote{}1475 vessels arrived in Hamburg without
impediment. British goods estimated at 590,000 tons were sold openly
without any
seizures\literalquote{}（柏林敕令后六个月）；汉堡付150万法郎贿款；阿尔托纳路线日均6,000--10,000人次走私、法方没收≤5\%；并合后435家糖厂剩约40家、印花布27场→1811年4场、1811--12失业制造业15,000＋航运翻倍、救济17,000＋；1813年2月24日总爆发连锁四城。补偿的缺失：\literalquote{}To
their bitter disappointment, tariff barriers remained...
Napoleon\textquotesingle s licence system privileged French ports at the
expense of annexed ports.\literalquote{}------这是外圈与内圈结果分岔的关键机制变量。

【锚】伊利里亚：\literalquote{}In November 1810, the Emperor ordered the conscription
of 18,000 soldiers, but during the next three years authorities called
up only 3000 men each
year.\literalquote{}（Grab）；反征兵骚乱多处、数百逃兵入匪帮；封锁无法执行、\literalquote{}走私成国民工业\literalquote{}（Bundy语经Grab）；1812年支出13.5m其中军海9.9m、收入长期不敷。荷兰：见第二十一章（80余起起义的67\%在并合后）。

【锚】加泰罗尼亚的法理自供状：1812年1月26日四省敕令从未刊入帝国法律公报、无元老院令追认（事实并合而不敢法律并合）；1813年3月15日下加泰民政撤销复归军政＝行政自认失败；征兵从未实施。语言让步实验（Mercader
Riba）：奥热罗1810年3月使加泰语官方化，9月即逆转；1812年4月绍夫兰咨询的Audiència委员会结论\literalquote{}la
insurrecció no s\textquotesingle afeblí pas gens amb
això\literalquote{}------起义丝毫未因此减弱，只使文书更难，建议法典译卡斯蒂利亚语。\textbf{全西班牙的A层试点在最大强制＋最大让步下双败。}瓦莱＝外圈唯一零抵抗单位：1810年11月26日要员与教会宣誓后驻军降至约一千；两副省长皆本地要人；但财政净负------\literalquote{}能并的不值钱，值钱的不能并\literalquote{}。

\subsubsection{二、外圈的账}\label{ux4e8cux5916ux5708ux7684ux8d26}

驻军依赖的榨取区：6--10万驻军≈S6-lite全部富余（2.9万）的2倍以上；野战军一败即刻归零。外圈财政贡献区间{[}−10,
+40{]}m／年，替换纸面144m。S6外推主线（中高）：北意技术可行政治亏本（没收自治溢价、征兵实得向罗马折扣滑动）；德西分叉于关税墙（贝格型工业区在\literalquote{}并合＋真准入\literalquote{}下可内圈化，海岸带必重演汉萨自败）；全西班牙已证伪。falsifier（\literalquote{}罗马／荷兰／汉萨兼并后三年无系统性抵抗→外圈上调\literalquote{}）不成立：三地均有系统性抵抗；但需精确化------抵抗在驻军成本下均可管制（除加泰），西班牙式战争未在任何新并合省重演：\textbf{外圈可管、不可省化}------可占有并榨取（负净值），不能转化为供兵供税的正常省。未决：荷／汉／罗年度征兵实得面板（若实得率\textgreater80\%则总判下调）；托斯卡纳／帕尔马为最薄弱单位；Pivec-Stelè／Bundy／Mercader
Riba原书；1813年汉堡法军撤出日期两说。

\begin{center}\rule{0.5\linewidth}{0.5pt}\end{center}

\subsection{第四十五章　反抗的决定因素：三十五个地区---时段的比较}\label{ux7b2cux56dbux5341ux4e94ux7ae0-ux53cdux6297ux7684ux51b3ux5b9aux56e0ux7d20ux4e09ux5341ux4e94ux4e2aux5730ux533aux65f6ux6bb5ux7684ux6bd4ux8f83}

\subsubsection{一、拆开\literalquote{}反抗\literalquote{}}\label{ux4e00ux62c6ux5f00ux53cdux6297}

三十五个地区---时段单元、十六个地区簇，把\literalquote{}反抗\literalquote{}拆成互不混算的三个结局：大规模起义22／35、跨季持续武装抵抗5／35、败局倒戈5。四条件（地形／宗教政策冲突／精英吸纳／外部军事机会）作crisp-set配置比较；留一行与留一地区簇交叉检验；十三个先行冻结规格的敏感性。

【推】样本内五个持续抵抗正例（威尼托1809、那不勒斯1799、卡拉布里亚1806、蒂罗尔1809、西班牙1808）全部同时具备宗教政策冲突与外部军事机会；精英吸纳仅1/5出现，地形4/5（威尼托为平原反例）。外部军事机会O＝0时，无论政策挑衅0--3项，持续抵抗率均为0；O＝1且挑衅≥2项时十三个单元全部发生大规模起义。\textbf{推论：行省化A层的地理边界是英奥俄军事投送的边界，不是语言或文化的边界。}（n＝5且为目的抽样，不构成总体必要性证明；O与结局有双向因果------西西里对卡拉布里亚的资助即由预期反抗强度决定。）

【锚】Woolf独立支持外援条件：\literalquote{}Apart from Spain, only two revolts
seriously challenged French authority: the endemic war in Calabria and
the Tyrol rising. Both were dependent on outside support from a major
power.\literalquote{}（印p.243）；西西里的西德尼·史密斯持续资助卡拉布里亚匪帮\literalquote{}as
otherwise they would have
dissolved\literalquote{}。推翻\literalquote{}天主教＝抵抗\literalquote{}：Woolf明确西班牙\literalquote{}both episcopate and
secular clergy preached obedience to Joseph\textquotesingle s
government\literalquote{}（印p.245）------起义源于统治阶级全面的信用崩溃而非教会动员；荷兰档案分类独立佐证：\literalquote{}only
4 per cent can be considered religious in
origin.\literalquote{}（Joor印p.202）却仍大规模抗争。故条件R定义为\literalquote{}宗教政策冲突\literalquote{}而非天主教人口。

【锚】比利时须按时段分列：\literalquote{}Except for the peasant revolt in 1798,
Belgium never challenged the French rule in any significant
way.\literalquote{}（Grab印p.83）------1798年正例（数百拒誓教士被捕、35人流放圭亚那、9月引入征兵后10月起义、12月初镇压）；1800--13年负例。同一人口与宗教构成，结局相反：身份--地形固定决定论不成立。两个异常案例：莱茵1796--99以约十分之一人口的驻军（136,000--187,000）压住最高风险配置------\literalquote{}Armed
resistance was out of the question, given French military
power.\literalquote{}（Rowe印p.90，比较论证而非识别出的驻军因果效应）；苏黎世1804年Bockenkrieg在\literalquote{}保精英＋保自治\literalquote{}的调停制下仍发生农民起义且非反法------安抚精英与安抚民众可能互相冲突（恢复旧精英的地租什一税赎金正是点火原因）。征兵冲击是时间递减变量：\literalquote{}conscription
seems to have been accepted increasingly with resignation, although it
took time before it could acquire in each annexed region the quality of
an unavoidable
curse\literalquote{}（Woolf印p.232）------解释比利时1802与荷兰1811同配置而结局相反；四条件缺一个时间维度（并合年资×征兵制度化）。

\subsubsection{二、反证未通过，降级为定性比较}\label{ux4e8cux53cdux8bc1ux672aux901aux8fc7ux964dux7ea7ux4e3aux5b9aux6027ux6bd4ux8f83}

契约反证未通过；失败模式是\textbf{弃权而非判错}：十三个规格中模型从未把比利时判成正例或西班牙判成负例；比利时1798为配置奇点（T0R1E0O0全样本仅此一例），比利时1802落在与荷兰1811同格的矛盾格（一致性0.333）。检验脚本初版曾把\literalquote{}比利时1798应为正例\literalquote{}错误施加于两个结局，已自我更正。八区预测（三政策档下的条件裁量区间）：非统计概率、不得相乘；德西H、巴伐利亚H、加泰P、那不勒斯A四行落在样本中不存在的配置（1111），证据等级低。单编码者设计，不声称任何编码信度指标达标；教士密度全样本缺测；驻军／人口率仅零星可得，\literalquote{}镇压可替代政策让步\literalquote{}仍是假说；Lignereux并合省rébellions专题（PUR
2013，DOI
10.4000/books.pur.134235）被Anubis阻断------若其显示\literalquote{}反叛\literalquote{}在并合省是常态化协商形式，22/35为低估；波兰与伊利里亚缺地方研究，结局保留{[}0,1{]}。

\begin{center}\rule{0.5\linewidth}{0.5pt}\end{center}

\subsection{第四十六章　整合绩效面板：并合省、藩属与旧法国的量化比较}\label{ux7b2cux56dbux5341ux516dux7ae0-ux6574ux5408ux7ee9ux6548ux9762ux677fux5e76ux5408ux7701ux85e9ux5c5eux4e0eux65e7ux6cd5ux56fdux7684ux91cfux5316ux6bd4ux8f83}

\textbf{行政常规≠净兵源≠财政自付}------三者不能混同；A层并省的收益必须相对于原B层藩属已提供的兵与钱计算增量，并同时扣除新增驻军与承接的财政义务。

【锚】Rowe表3同源对照：1806--10莱茵征前逃避33.65\%高于阿尔萨斯17.12\%，入伍后逃亡3.00\%却低于13.23\%------同一地区是否\literalquote{}整合成功\literalquote{}随结果变量反转；单一整合指数会掩盖选择、替身与阶段差异。复算：莱茵逃兵下降10.25个百分点、阿尔萨斯下降9.66，两者变化差仅−0.59------不可把共同下降全归兼并整合。比利时被征率216,111/3,028,705＝7.135\%（不是6.12\%，也不是去重的人口被征概率）；全期累计÷年数是期均，不是成熟年度流量；皮埃蒙特\literalquote{}我不知道500已出发\literalquote{}非精确500到营；Roer名义税+78.21\%不是兼并净因果；Marion
1812全为预期非实收；Rouanet--Piano
2023正文的逃役反事实增量3,256与附录2,812冲突（27\%/49\%更接近2,812的分母；3,256实为额外送军，2,812/10,499≈26.78\%）------保留冲突，不冻结\literalquote{}净增兵＝0\literalquote{}。

【推·参数包络】给定五年含训练、年留存0.90、可部署率0.80（Q2口径），每百万人第10年净可用兵：L
−38至+5,318；M −7,102至+2,190；H
−26,729至−6,034。\textbf{它不是统计预测或概率}；同时有效受训服役批次为四批（不是\literalquote{}十年只征四批\literalquote{}）；全口径驻军成本g是最承重假设，尚需逐地年账替换；L/M/H映射是假说。D1的\literalquote{}7--12年\literalquote{}与Q2的\literalquote{}3--6/6--12年\literalquote{}不能拼成已识别的统一收敛期。116行来源／缺口、8项刊本估计、46复算、3,840条模型路径、108敏感性行；未取得ICPSR微观复制包，未建1800--13完整省年面板，未做新回归或统计LOOCV------缺数据不写成检验失败。

\begin{center}\rule{0.5\linewidth}{0.5pt}\end{center}

\subsection{第四十七章　市场的证词：公债、汇率与保费}\label{ux7b2cux56dbux5341ux4e03ux7ae0-ux5e02ux573aux7684ux8bc1ux8bcdux516cux503aux6c47ux7387ux4e0eux4fddux8d39}

\subsubsection{一、债价不是王朝赔率}\label{ux4e00ux503aux4ef7ux4e0dux662fux738bux671dux8d54ux7387}

【锚】英格兰银行／Neal月末3\% consol与Vaslin→Monnet月度法国5\%
rente（\literalquote{}Jusqu\textquotesingle à la Seconde Guerre mondiale, nous prenons
le taux de marché de la rente perpétuelle émise par
l\textquotesingle État français, à 5 \% avant 1825 et à 3 \%
ensuite\literalquote{}；\literalquote{}Ces taux ont été recueillis par Jacques-Marie Vaslin,
Alphonse Courtois et David Le Bris, et compilés dans la thèse de David
Le Bris\literalquote{}），各168月，17事件四种宽窗（描述性，非事件因果）：

{\def\LTcaptype{none} % do not increment counter
\begin{longtable}[]{@{}
  >{\raggedright\arraybackslash}p{(\linewidth - 4\tabcolsep) * \real{0.3333}}
  >{\raggedright\arraybackslash}p{(\linewidth - 4\tabcolsep) * \real{0.3333}}
  >{\raggedright\arraybackslash}p{(\linewidth - 4\tabcolsep) * \real{0.3333}}@{}}
\toprule\noalign{}
\begin{minipage}[b]{\linewidth}\raggedright
事件
\end{minipage} & \begin{minipage}[b]{\linewidth}\raggedright
英债理论价
\end{minipage} & \begin{minipage}[b]{\linewidth}\raggedright
法债理论价
\end{minipage} \\
\midrule\noalign{}
\endhead
\bottomrule\noalign{}
\endlastfoot
提尔西特（1807-06→08） & −2.15\%（安慰剂第47.4百分位） &
\textbf{+12.38\%}（约95.9百分位） \\
埃尔福特（1808） & +0.57\% & +0.60\% \\
侵俄（1812-05→07） & −7.55\%（94.8百分位，混美英战争与英国军费冲击） &
--- \\
莱比锡 & --- & −25.16\% \\
巴黎陷落／退位（1814-03→05） & +7.57\% & \textbf{+17.71\%} \\
滑铁卢 & −4.42\% & +15.61\% \\
\end{longtable}
}

提尔西特：大陆胜利大幅改善法国债的估值，而未迫使英国即时信用崩溃。1814年：拿破仑退位之际法国国家债券涨价------直接削弱\literalquote{}债价＝王朝存续赔率\literalquote{}的强映射（报价同时定价未来法国财政与王朝替代）。1813年1月伦敦对拿破仑死亡或被俘的保险3几尼／£100＝3.15\%毛费，不是王朝存续赔率。

【锚】Antipa（2016，JEH）：\literalquote{}For the period spanning the years from 1805
to 1808 the agio\textquotesingle s stability simply reflects the lack of
observations.\literalquote{}------1805--08金agio\literalquote{}平稳\literalquote{}是缺测；\literalquote{}it rose above 10
percent after 1808 and reached a peak of 45 percent
mid-1813\literalquote{}反映英镑对黄金实际贬值而非政治征服赔率。特拉法加的消息11月6日传伦敦，11月14日断点后英国3\%
consol收益率区制均值5.10→4.95\%，不是单日变化。Antipa \&
Chamley（2017）：\literalquote{}the exchange rate fluctuated around 34 schillings
(Flemish banco) per pound until the end of
1808\literalquote{}，此后渐贬（约24--28为图像粗读），论文归因伊比利亚战费与盟友补贴------价格不唯一识别主权终局。海险：1809年伦敦→加的斯60\%为两条高尾航线非伦敦平均；1812年英港→布宜诺斯艾利斯12.5--18.8\%经Llorca转引；Danson
1894《Our Next War\ldots{} premiums paid at Lloyd\textquotesingle s from
1805 to 1816》未获（Hathi 403）。

\subsubsection{二、有条件概率压力，否定伪精度}\label{ux4e8cux6709ux6761ux4ef6ux6982ux7387ux538bux529bux5426ux5b9aux4f2aux7cbeux5ea6}

提尔西特\literalquote{}英国坏状态\literalquote{}二态后验10.9--40.8\%，仅在主观先验10--30\%、支付差20--60\%、事件归因25--100\%的显式模型下成立；莱比锡法公债减损20.5--92.9\%不能等于王朝崩溃。最终安慰剂表34行、敏感性306行、概率压力918行。未决：Neal日度底本、DFIH日价、伦敦--巴黎／汉堡／阿姆斯特丹逐月报价、同合同Lloyd\textquotesingle s保费；月度取价日与Vaslin底本未核清；1806年断点的CI印成1805年视为原文错误不擅改。

\begin{center}\rule{0.5\linewidth}{0.5pt}\end{center}

\section{第二部　如何赢：路径的对弈}\label{ux7b2cux4e8cux90e8-ux5982ux4f55ux8d62ux8defux5f84ux7684ux5bf9ux5f08}

\subsection{第四十八章　对弈规则、奇迹的计价与三条互斥定理}\label{ux7b2cux56dbux5341ux516bux7ae0-ux5bf9ux5f08ux89c4ux5219ux5947ux8ff9ux7684ux8ba1ux4ef7ux4e0eux4e09ux6761ux4e92ux65a5ux5b9aux7406}

\subsubsection{一、规则}\label{ux4e00ux89c4ux5219}

对弈引擎（P6，\texttt{p6\_wargame.py}，无依赖、确定性输出）不是预测器，是一致性记账器：每年一回合------法国行动→十个单位按反应函数响应（英国、俄国、奥地利、普鲁士、西班牙、南德、教廷、美国、北欧、奥斯曼）→人力／财政／战区结算→\literalquote{}同一机动军不得两处\literalquote{}硬检查。参数全部取自单位卡：法国在册400--550k、盟军150--250k、在场率0.85；姿态需求（千人在场）S2
540、S5 490、S6-lite 566.25、S6-max 922.25、S6-full
1,102.25；边际增量------西班牙占领+290至+360、英伦登陆+80至+130、征俄+250（法籍主力份额）、有限对俄+180、外圈治安+60至+100、卡拉布里亚+22、取消普鲁士+60、东方远征+30；财政------高收入＋可兑盟约案陆海合计495m，无盟约中心案陆军余量230m；700法郎／兵年（外圈900）；关税------严格版15--35m、执行成本28--38.5m、税基损伤−60至−110m；保护版60--90m、执行14--24m。

\textbf{三栏计价}是方法的核心：M类奇迹＝外生、法国不可购买、概率≤0.25；C类选择＝法国自己的决策，按F1先验计价（工具性0.45--0.70、非工具性0.05--0.25）；L类命运＝寿限与继承。只有M类计入契约的奇迹计数------否则任何反事实策略推演在定义上都不可能存在。红队指出这一区分的代价：路径B有三个C类节点（C3对英写定交易且不再加价0.15--0.25；C7不并荷汉0.15--0.25；C8统合公债0.15--0.30）恰在奇迹阈值上，只因被归入\literalquote{}选择栏\literalquote{}而不计。本书接受这一批评（第五十三章）：三栏分开计价是正当的，但\textbf{\literalquote{}零奇迹\literalquote{}不等于\literalquote{}高实现概率\literalquote{}}，标题不得偷换。

引擎的弱有效性检验：以史实政策向量运行（S0对照线），人力账自1808年西班牙废立起转为NEGATIVE，1812年需求112.8--119.8万对供给38.75--56万，且当年触发\literalquote{}伊比利亚＋俄境\literalquote{}两战区硬冲突------与F2独立给出的\literalquote{}占领态吞噬29--36万＝常态包络60--80\%\literalquote{}同向。这只弱检验人力／财政的符号次序；K1--K3（控制／俄受约束／英接受）是计分卡手填，不是引擎输出------红队P项，本书标明。

\subsubsection{二、三条互斥定理}\label{ux4e8cux4e09ux6761ux4e92ux65a5ux5b9aux7406}

把四十四张卡的反应函数放在同一张桌上，出现三组在单张卡上看不见的互斥：

\textbf{T1 封锁互斥。}
最严封锁（S3）同时与（i）大规模战争（Ellis法则：\literalquote{}the efficacy of the
Blockade stood in inverse ratio to the extent of
Napoleon\textquotesingle s military
involvements\literalquote{}）、（ii）俄国不战（关税容忍是唯一六列同击的让步，制度上等于放弃S3）、（iii）北欧盟友的财政生存（丹麦财政寿命5--6年；挪威粮食杠杆倒挂）、（iv）美国中立管道（价格＝不扣船）、（v）西属美洲白银（三件同在才能汲取）、（vi）法国自己的关税收入（严格版15--35m
vs
执行28--38.5m）互斥。六条分别由X2、R1、SC、US、MX、F4给出，方向一致------但红队N项提醒：R1与R2是同一国家的两张卡，US／MX／X1共用白银管道材料，独立束数应下调；三条定理各自至少有两束独立支撑，不是\literalquote{}十几张卡各自独立证实\literalquote{}。\textbf{放弃S3是一切可行路径的必要条件。}

\textbf{T2 行省化互斥。}
A层扩张越过内圈之后，兵力账（F2）、财政账（F4）、治安账（D2）、盟军账（G1）、整合账（Q2）同时转负，且它们不是独立的账而是同一笔账的多个面：外圈是封锁执法的对象，而封锁摧毁的正是外圈的税基。\textbf{\literalquote{}多吞＝多兵\literalquote{}在第十年前恒为负。}

\textbf{T3 目标函数互斥。}
最大陆权⊥最大海权（F3自陈）；最强控西⊥最大白银（MX）；最严封锁⊥最大对英伤害（X2：法国只掌握合取式的一半）；B5独立得出\literalquote{}霸权耐久性最优点≠对英伤害最大点\literalquote{}。\textbf{法国的最优不是最大化任何单项，而是选一个有界的目标函数。}

\begin{center}\rule{0.5\linewidth}{0.5pt}\end{center}

\subsection{第四十九章　被否决的三条路：入侵、绞杀、东方}\label{ux7b2cux56dbux5341ux4e5dux7ae0-ux88abux5426ux51b3ux7684ux4e09ux6761ux8defux5165ux4fb5ux7edeux6740ux4e1cux65b9}

\subsubsection{一、路径A：S1入侵线------最快的赌博，不是最快的胜利}\label{ux4e00ux8defux5f84as1ux5165ux4fb5ux7ebfux6700ux5febux7684ux8d4cux535aux4e0dux662fux6700ux5febux7684ux80dcux5229}

1803年布洛涅集结，全部资源压海峡，不动西班牙不动意大利；1804年继续集结与舰队集中计划，外交安抚奥俄；1805年海峡窗口开------这是第一个奇迹节点。

M-A1（0.10--0.25，奇迹）：八万人＋炮马弹药D+2完整上岸且运输队未被截击。B4的证据：无四十八小时集中四万正规军的实绩，无D+2完整卸载的校准，蒙特耶尔样板马位缺2,436，B3无固定七十二小时预警但责任区编制在位。M-A2（0.10--0.20，奇迹）：登陆后持续跨海补给------两潮／六潮是不同的时间模型，英舰队未失效即切断。M-A3（0.25--0.45，非奇迹）：伦敦失守后英国在1--3个月内签有限和约而非迁政府续战------B1a两分支并列，续战以基地＋可信支付为条件。M-A4（0.15--0.30，非奇迹）：1805--06奥俄不利用法军渡海组织第三次同盟------史实1805年同盟已成型，法军主力离陆是最强诱因。

计分：四项外生事件，严判据（≤0.25）两件奇迹，宽判据（≤0.30）三件------主线资格在判据敏感区；\textbf{不能作主线，但不构成逻辑排除}。1806年K1\literalquote{}部分\literalquote{}、K2\literalquote{}否\literalquote{}、K3\literalquote{}条件性是\literalquote{}；两个硬奇迹未兑现即回落到史实型多线战争，1815年全否。B4的判定给这条路一个精确的形状：两万孤军可封控、四万是铰链、八万完整炮马D+2上岸有初期局部优势、十二万不增长期征服力；八万渡海约2--4日或4.5--6日的串行压力测试要求连续同长的制海权；即使伦敦失守，B1a只给1--3个月的有条件议价窗，且前提是法方保王朝、议会、公债与主体舰队------法国要的越多，英国续战的窗口越长。

\subsubsection{二、路径C：S3封锁绞杀------四个奇迹}\label{ux4e8cux8defux5f84cs3ux5c01ux9501ux7edeux6740ux56dbux4e2aux5947ux8ff9}

1806年柏林敕令并为执法而并地；1807年提尔西特把俄国拖入封锁；1808年为闭合伊比利亚海岸废立西班牙；1810年特里亚农--枫丹白露、并荷兰汉萨；1812年为封锁闭环征俄------Ellis法则的自毁点。所需奇迹：M-C1（0.05--0.15）俄国自愿长期执行排他封锁而不索关税容忍；M-C2（0.10--0.20）美国在扣船政策下仍维持中立运输管道（1810年朗布依埃敕令扣押一切进法港美船）；M-C3（0.05--0.15）英国工业在二十五年排除下崩解到被迫接受法国最高条件（B5：S3下1848工业为史实54--79\%，不是崩解）；M-C4（0.05--0.15）严格封锁与大规模战争同时维持效力。四件奇迹，K1--K4全否。引擎的财政结算：严格封锁年关税净−133.5至−53m（含税基损伤），可得401.8m对需求438m以上，逐年NEGATIVE。\textbf{这条路不是\literalquote{}更狠一点就赢\literalquote{}，而是自己吃掉自己的税基、盟友与白银。}

\subsubsection{三、路径D：S4东方转向------三个奇迹}\label{ux4e09ux8defux5f84ds4ux4e1cux65b9ux8f6cux5411ux4e09ux4e2aux5947ux8ff9}

1807年芬肯施泰因＋加尔达纳使团；1808年圣彼得堡瓜分谈判死锁于达达尼尔；1809--10年波斯走廊远征＋伊比利亚未了。M-D1（0.10--0.25）≥两万法军经波斯走廊两个战役季内抵赫拉特并保持补给（PS：单季无预置＝关闭；分梯队＋四仓＋里海航运为边缘可行中低）；M-D2（0.05--0.15）奥斯曼允许大军过境或合营瓜分不破裂（1807年连25,000援军过境波斯尼亚都被地方贝伊否决）；M-D3（0.05--0.15）万人级远征的四闸门同开（1803--15从未同开）。三件奇迹，且走廊四政权都是明码标价的安全采购者，法国交付不了他们的主要威胁。S4的理性形态从来不是远征，而是\literalquote{}永久的加尔达纳\literalquote{}------一个常设军事使团、对俄担保、二十万toman级的白银津贴；其产出是英属印度前沿体系提前二十年成形。

\subsubsection{四、三条被否决的路的共同教训}\label{ux56dbux4e09ux6761ux88abux5426ux51b3ux7684ux8defux7684ux5171ux540cux6559ux8bad}

三条路各自最大化一个单项------渡海、封锁、东方------而T3已经证明单项最大化在体系层面互斥。它们的失败不是运气，而是目标函数：入侵线把胜负押在两个物理奇迹上；绞杀线要求四个国家违背自己的生存利益；东方线要求四个走廊政权接受法国交付不了的安全。\textbf{剩下的只有一条路：一个有界的目标函数。}下一章把它拆成两个版本。

\begin{center}\rule{0.5\linewidth}{0.5pt}\end{center}

\subsection{第五十章　主线的重构：S2本义（史实至提尔西特＋两个开关）与路径B（1803年即切换）}\label{ux7b2cux4e94ux5341ux7ae0-ux4e3bux7ebfux7684ux91cdux6784s2ux672cux4e49ux53f2ux5b9eux81f3ux63d0ux5c14ux897fux7279ux4e24ux4e2aux5f00ux5173ux4e0eux8defux5f84b1803ux5e74ux5373ux5207ux6362}

\subsubsection{一、P6交付的路径B，及其被红队指出的自洽问题}\label{ux4e00p6ux4ea4ux4ed8ux7684ux8defux5f84bux53caux5176ux88abux7ea2ux961fux6307ux51faux7684ux81eaux6d3dux95eeux9898}

P6的最优构造线（称\literalquote{}S2′\literalquote{}）是这样一条世界线：1803年不为马耳他重开战、把争端交第三方仲裁；1804年称帝并对英写定领土交易清单且承诺不再加价；1805年不进德意志，以意大利王冠换奥地利的商业海口与边界保证，第三次同盟不成型；1806年全才内阁早窗谈判，不建大邦联宪法只做仲裁＋配额；1807年对俄给关税容忍、不称波兰王国；1808年不废西王（B4b）；1809年教廷和解；1810年不并荷兰汉萨，改分层准入，奥地利联姻；1811年统合公债；1812年以\literalquote{}欧洲总和约\literalquote{}签英法条约；1813--15降为S5和平姿态。引擎结算：人力全程TIGHT（需求540对供给467.5--680）、财政全程正、无战区冲突；所需外生事件三项（M-B1英国签约0.40--0.60，M-B2俄国十年不战0.35--0.50，M-B3奥地利不再入场0.45--0.60），无一≤0.25。P6据此判\literalquote{}存在论成立\literalquote{}。

红队的批评击中要害。\textbf{这条路径把产生\literalquote{}控制\literalquote{}的原因移除了，却把\literalquote{}控制\literalquote{}的结果留在了账上。}
P6用来判定K1\literalquote{}大陆四层控制\literalquote{}的版图表把奥地利列为C层（依据1809年申布伦限军十五万）、普鲁士列为C层（依据1808年巴黎公约限军四万二千）、\literalquote{}莱茵邦联（1808）三十九邦一千四百万人\literalquote{}列为B层、华沙公国列为B层------四项都是路径B明文移除的1805、1806、1807、1809四场战争的和约产物。引擎里，路径B的奥地利＝ALLY\_PASSIVE、普鲁士＝NEUTRAL、俄国＝NONBELLIGERENT：三个未被击败、疆域完整的大国的\literalquote{}不动\literalquote{}，不是\literalquote{}受约束\literalquote{}。P7为路径B裁定的十四项前提冲突（普鲁士完整、无华沙公国、无普雷斯堡、无申布伦\ldots\ldots）恰是\literalquote{}路径B离开了单位卡反应函数的估计域\literalquote{}的症状------但P7只清理了普奥、意、葡、北欧几处，没有清理K1本身。同样，路径B的英法战争状态在P6（\literalquote{}英国继续封锁与补贴外交\literalquote{}、\literalquote{}汉诺威归还乔治三世\literalquote{}、引擎britain＝HOSTILE\_ARMED→NEGOTIATING且无一年签约）与P7（\literalquote{}1803年起战争未重开\literalquote{}、\literalquote{}无枢密令\literalquote{}、\literalquote{}训练竞争替代交战\literalquote{}）之间相互矛盾；M-B1的0.40--0.60借自B1b以战争持续为前提的1811--14窗口。还有七处\literalquote{}保留史实果实而移除其原因\literalquote{}的条目：无普雷斯堡而有莱茵邦联与巴符王号，且神圣罗马帝国的命运未交代；无1809年而有1810年奥地利\literalquote{}主动求取\literalquote{}联姻；无枫丹白露、无法军入西而阿兰胡埃斯政变\literalquote{}照常发生\literalquote{}；威尼托留奥而威尼斯船厂持续下水、欧仁有威尼斯支；华沙公国三说并存；完整普鲁士得北德主持权又\literalquote{}无领导核\literalquote{}。

本书的处置：\textbf{不抛弃路径B，但重新命名它所达到的东西，并把被P6跳过的竞争路径补算出来。}

\subsubsection{二、S2本义：史实打到提尔西特，再切换}\label{ux4e8cs2ux672cux4e49ux53f2ux5b9eux6253ux5230ux63d0ux5c14ux897fux7279ux518dux5207ux6362}

PROJECT.md对S2的定义------\literalquote{}不废西王、不打俄、对英武装共存\literalquote{}------三个动词的宾语（1808年的西王、1812年的俄国、1803年起的英法战争）都以提尔西特后的处境为前提；四十四张单位卡的S2列也都在这个处境上估计。P6却用\literalquote{}不打任何仗\literalquote{}的S2′替换了它。本书把被替换的路径叫做\textbf{S2本义}（引擎标号E）：

\begin{itemize}
\tightlist
\item
  \textbf{1803--1807年按史实运行}：马耳他决裂、布洛涅、乌尔姆--奥斯特里茨--普雷斯堡（1805）、耶拿--奥尔施泰特与柏林敕令、莱茵邦联（1806）、弗里德兰--提尔西特（1807年7月）。这些胜利是史实事件，不入奇迹栏；它们的果实------普雷斯堡对奥的割地与赔款、提尔西特对普的削减与华沙公国、俄国入体系、莱茵邦联、巴符王号、约瑟夫在那不勒斯------在这条路径上是\textbf{真实的条约产物}。
\item
  \textbf{自1807年秋／1808年起切换}：不签米兰敕令，把特里亚农型的保护性财政关税提前到1808年（关税＝走私成本的国有化，这是海关总监的当时提案）并对俄给关税容忍（接受俄国的中立航运与1810年式关税，即B7″v3），不为封锁而并地；不废西王------联姻去条件化，收西班牙的25,000远征军与加的斯--费罗尔船厂（B4b）；对奥改保证（普雷斯堡之后不再羞辱，A卡S5-A）；1809年教廷和解、不并罗马；1810年不并荷兰汉萨，分层准入；奥地利联姻（在普雷斯堡之后的奥地利有史实的动机）；1811年统合公债，不吞奥尔登堡；1812年不征俄，在B1b的1811--14窗口谈欧洲总和约；1813年起降为S5姿态。
\end{itemize}

本书作者用P6引擎（不改原卡任何文件，新脚本与输出存于工作区\texttt{final/r\_5b1357a9c6-calc/}）跑了这条路径，两个变体：E\_literal按引擎原状态机（1806--07年的柏林敕令一旦置\literalquote{}血税超限\literalquote{}即永久，南德盟邦DEFECT\_RISK、供给−80至−120k）；E\_reset在1808年切换为保护性关税且无对俄战争后重置该状态。【机】重置的依据是G1的里德机制：南德忠诚的终结是\literalquote{}替代担保人出现∧血税超限（1812年灭军）\literalquote{}的合取（1809年法国胜负未卜时三邦均出兵无动摇；符腾堡直到莱比锡仍在法方作战），单凭1806--07年的封锁负担不触发；适用条件是无奥地利的替代担保（无里德式条约）且无一次灭军级损失------S2本义不征俄恰满足后者。E\_literal则是引擎原状态机的保守读法，两者并列。结果：

{\def\LTcaptype{none} % do not increment counter
\begin{longtable}[]{@{}
  >{\raggedright\arraybackslash}p{(\linewidth - 12\tabcolsep) * \real{0.1429}}
  >{\raggedright\arraybackslash}p{(\linewidth - 12\tabcolsep) * \real{0.1429}}
  >{\raggedright\arraybackslash}p{(\linewidth - 12\tabcolsep) * \real{0.1429}}
  >{\raggedright\arraybackslash}p{(\linewidth - 12\tabcolsep) * \real{0.1429}}
  >{\raggedright\arraybackslash}p{(\linewidth - 12\tabcolsep) * \real{0.1429}}
  >{\raggedright\arraybackslash}p{(\linewidth - 12\tabcolsep) * \real{0.1429}}
  >{\raggedright\arraybackslash}p{(\linewidth - 12\tabcolsep) * \real{0.1429}}@{}}
\toprule\noalign{}
\begin{minipage}[b]{\linewidth}\raggedright
年
\end{minipage} & \begin{minipage}[b]{\linewidth}\raggedright
姿态／封锁
\end{minipage} & \begin{minipage}[b]{\linewidth}\raggedright
人力需求（千在场）
\end{minipage} & \begin{minipage}[b]{\linewidth}\raggedright
供给（E\_literal／E\_reset）
\end{minipage} & \begin{minipage}[b]{\linewidth}\raggedright
财政可得
\end{minipage} & \begin{minipage}[b]{\linewidth}\raggedright
财政需求
\end{minipage} & \begin{minipage}[b]{\linewidth}\raggedright
状态
\end{minipage} \\
\midrule\noalign{}
\endhead
\bottomrule\noalign{}
\endlastfoot
1805 & S2／无 & 540 & 467.5--680 & 540.0 & 498 &
人力TIGHT，财政OK，多瑙战区 \\
1806 & S2／严格 & 540 & 387.5--560／同 & 401.8 & 438 &
\textbf{财政NEGATIVE}（关税净−133.5至−53） \\
1807 & S2／严格 & 540 & 387.5--560／同 & 401.8 & 438 &
\textbf{财政NEGATIVE} \\
1808 & S2／保护 & 540 & 387.5--560／467.5--680 & 551.0 & 488 &
财政OK（关税净36--76） \\
1809--1811 & S2／保护 & 540 & 同上 & 551.0 & 498--528 & OK \\
1812 & S2／保护 & 540 & 同上 & 551.0 & 548 & OK（余3m，紧） \\
1813--1815 & S5／保护 & 490 & 同上 & 551.0 & 523--538 & OK \\
\end{longtable}
}

无战区冲突。所需外生事件与路径B相同（M-B1、M-B2、M-B3），奇迹计数为零；C类节点只剩C4--C10（无C1马耳他、C2
1805不进德意志、C3对英写定交易），十节点独立乘积的算术上下界从路径B的2.9e-6--2.9e-4升到9.6e-5--3.0e-3（10--33倍）------这个乘积按P6纪律不得当路径概率，只作序数比较：\textbf{S2本义需要更少的克制节点。}引擎的两处提示要一起读：1806--07的史实封锁年财政为负------这就是史实，柏林敕令时期的法国财政本来就靠贡赋填（F4）；1812年财政余额只有3m，它说明\literalquote{}签约后必须立刻兑现和平红利\literalquote{}是预算恒等式（S5姿态后才有13--28m余量）。

\subsubsection{三、逐国重填K1：两条路径的对照}\label{ux4e09ux9010ux56fdux91cdux586bk1ux4e24ux6761ux8defux5f84ux7684ux5bf9ux7167}

{\def\LTcaptype{none} % do not increment counter
\begin{longtable}[]{@{}
  >{\raggedright\arraybackslash}p{(\linewidth - 4\tabcolsep) * \real{0.3333}}
  >{\raggedright\arraybackslash}p{(\linewidth - 4\tabcolsep) * \real{0.3333}}
  >{\raggedright\arraybackslash}p{(\linewidth - 4\tabcolsep) * \real{0.3333}}@{}}
\toprule\noalign{}
\begin{minipage}[b]{\linewidth}\raggedright
大陆强权
\end{minipage} & \begin{minipage}[b]{\linewidth}\raggedright
S2本义（E）：约束形态与文书
\end{minipage} & \begin{minipage}[b]{\linewidth}\raggedright
路径B：约束形态与文书
\end{minipage} \\
\midrule\noalign{}
\endhead
\bottomrule\noalign{}
\endlastfoot
奥地利 &
C层。普雷斯堡（1805）割威尼西亚／蒂罗尔＋赔款为真实条约；1808年后法国改给保证（A卡S5-A：保王朝核心、有限援军、海口与边界）；无西班牙泥潭则1809年窗口不开（A卡\literalquote{}S5-A早克制避免1809\literalquote{}）；1810年联姻有史实动机。M-B3
0.45--0.60------但红队正确指出：若A卡判1805年受辱后1809年再战概率显著上升，须下调；本书取0.40--0.60并标此不确定。
& \textbf{无条约。}
奥地利未败、疆域完整（含威尼西亚、蒂罗尔）；法国的\literalquote{}三项保证\literalquote{}是给奥地利的担保而非约束奥地利的条约；1810年联姻的动机无锚。只能落L4\literalquote{}受约束中立\literalquote{}且无文书------K1\literalquote{}部分\literalquote{}。 \\
普鲁士 &
C层。提尔西特（1807）＋1808年9月8日巴黎公约（限军42,000十年期、三要塞法军10,000由普方供宿粮、七条军用道路）为真实条约。
& \textbf{无条约。}
普鲁士完整、未败、未改革；P7给它\literalquote{}北德邦联式安排的主持权\literalquote{}------这是1806年普鲁士自己想要、1866年才得到的结构，无锚点。K1\literalquote{}否\literalquote{}。 \\
俄国 &
C层（被胁迫同盟）。提尔西特同盟为真实条约；关税容忍（B7″v3）使R1的\literalquote{}有条件冷和平\literalquote{}半衰期由四年升至8--12年；不称波兰王国、不吞奥尔登堡。M-B2
0.35--0.50。 &
D层（受约束中立），无同盟条约；R1的让步包分析以提尔西特后的处境为前提，需移植。M-B2同带但证据更弱。 \\
西班牙 &
C层（同盟）。1796年圣伊尔德丰索同盟为真实条约；B4b＝保留波旁＋联姻去条件化（1808年1月10日\literalquote{}Consiento
de buen grado\literalquote{}为真实文书）。 &
同左（1803--08年的法西同盟在两条路径上都是真实的）。 \\
那不勒斯 &
B层。1806年征服与约瑟夫／缪拉体制为真实（1806年是路径E的史实段）；拜永条约（1808）真实。
& 波旁两西西里存续（P7
PRE-04）------1805年它是第三次同盟成员、英国的地中海陆上支点；它是否构成\literalquote{}可充当联盟锚的敌对大陆强权\literalquote{}，P6／P7未讨论。 \\
莱茵邦联 &
B层。1806年7月12日邦联条约、1808年三十九邦一千四百万人、承诺126,000人------真实；G1\literalquote{}担保独占／所得取自哈布斯堡\literalquote{}四条件直接适用。
&
1806年\literalquote{}窄版邦联\literalquote{}------但无普雷斯堡则无奥皇退位、诸邦退出神圣罗马帝国的法律通道与皇帝授王号的依据均未交代；G1条件只剩1803年教产部分。 \\
华沙公国 &
B层。真实（1807）；PL卡\literalquote{}两段式\literalquote{}（公国不称王国）直接适用；波兰忠诚的\literalquote{}复国期票×贸易现金流\literalquote{}双备函数可运行。
&
P6三说并存（不创设／只保公国／B层320--450万）。本书取\literalquote{}不创设\literalquote{}（无提尔西特则无公国），波兰维持1795年瓜分结构（P7
PRE-03）。 \\
荷兰 &
B层。路易的荷兰王国存续（1810年不并）；1810年3月路易条约（已签）为当时文书。
& 同左。 \\
教廷 & D层＋专约（1809年不并罗马；1808年教廷已书面愿意事实关港）。 &
同左。 \\
葡萄牙 & D层付费中立（无枫丹白露）。 & 同左。 \\
英国 &
K3：1811--14窗口签有限和约（B1b直接适用，因为战争自1803年持续、大陆盟友消失）。
& K3：无战世界------须按1802--03亚眠检验重新定价（第五十二章）。 \\
\end{longtable}
}

\textbf{判定。}
在S2本义下，\literalquote{}大陆四层控制\literalquote{}的每一层都有这条路径内部真实存在的条约或当时提案支撑，K1可填\literalquote{}是\literalquote{}；\literalquote{}俄国受约束\literalquote{}以提尔西特同盟＋关税容忍的形态成立；\literalquote{}英国不败而接受\literalquote{}以B1b窗口的形态成立。在路径B下，K1只能填\literalquote{}部分\literalquote{}：它达到的是\textbf{1801--03年吕内维尔--亚眠式西欧霸权的固化（primacy）}------A层≈内圈、B层＝意大利王国与窄版邦联、C层＝西班牙同盟------而对三个完整的东欧大国没有任何约束文书；它的可持续控制范围明显小于史实1810年帝国。

\textbf{因此本书把\literalquote{}控制版\literalquote{}的结论改写为两行：} -
\textbf{S2本义}：控制版存在（有条约支撑），置信中；两个开关的时点是1807／08年；世界线的1803--1807年即史实，1808年后按单位卡的S2列装配（第四部）。
-
\textbf{路径B}：以primacy形态存在、以control形态未证（红队所建议的措辞）；它的价值在于证明\literalquote{}有界目标函数\literalquote{}本身能自洽，但它离开了单位卡的估计域，第五十八章把它的相关年表条目降为低置信。

一个必须写明的推论：S2本义的\literalquote{}控制\literalquote{}是\textbf{用1805--1807年三场大战的胜利买来的}。它们在史实中发生了，所以不算奇迹；但它们也带来了史实的代价------特拉法加、莱茵邦联的血税义务、俄国贵族的封锁怨恨、奥地利的复仇动机------这些代价在1808年后的开关设计中都要偿付。S2本义并不比路径B\literalquote{}便宜\literalquote{}，它只是把账单从\literalquote{}克制\literalquote{}栏转到了\literalquote{}战争\literalquote{}栏。

\subsubsection{四、S2本义的版图（1812年固化态）}\label{ux56dbs2ux672cux4e49ux7684ux7248ux56fe1812ux5e74ux56faux5316ux6001}

A层并合省：旧法国（≈2,900万）＋比利时、莱茵左岸、皮埃蒙特、利古里亚（1805年，史实）、日内瓦、帕尔马（1808年，史实）------≈37--39百万人口代理；托斯卡纳有条件（在B4b世界里埃特鲁里亚可作为对西班牙波旁的补偿保留为B层，或按史实1808年并入------P2列为条件直辖）；不并荷兰、汉萨、罗马、伊利里亚、加泰罗尼亚。B层法典化藩属：意大利王国（含威尼托）、那不勒斯（缪拉）、莱茵邦联三十九邦、华沙公国、荷兰王国、瑞士（调停制）、威斯特法利亚。C层条约同盟：西班牙（波旁）、普鲁士（残普）、俄国（被胁迫同盟）、奥地利（不平等伙伴）、丹麦。D层受约束中立：教廷、葡萄牙、瑞典、汉萨（不并）、奥斯曼。\textbf{与P2的边界版（约39.5--40.5m）一致；与史实1810年帝国（≈42m
A层＋外圈）的差别正是外圈。}

\begin{center}\rule{0.5\linewidth}{0.5pt}\end{center}

\subsection{第五十一章　两个开关的解剖}\label{ux7b2cux4e94ux5341ux4e00ux7ae0-ux4e24ux4e2aux5f00ux5173ux7684ux89e3ux5256}

\subsubsection{一、开关一：大陆体系由\literalquote{}排他封锁\literalquote{}改为\literalquote{}保护性财政关税\literalquote{}}\label{ux4e00ux5f00ux5173ux4e00ux5927ux9646ux4f53ux7cfbux7531ux6392ux4ed6ux5c01ux9501ux6539ux4e3aux4fddux62a4ux6027ux8d22ux653fux5173ux7a0e}

这个开关的本质是把封锁从一种战争行为改为一种税制。它之所以承重，不是因为它\literalquote{}温和\literalquote{}，而是因为它同时满足六个互不相关单位的反应函数：

\begin{enumerate}
\def\labelenumi{\arabic{enumi}.}
\tightlist
\item
  \textbf{执行机器（X2）}：走私成本的硬上限≈55±8\%
  AVE，中率关税40--60\%恰压在其上------国家取代走私者收租，\literalquote{}the profit
  would fall to the state instead of to the
  smugglers\literalquote{}（Heckscher印p.214）。
\item
  \textbf{俄国宫廷（R1）}：关税容忍是六个让步包中唯一六列同击的一项；它买到\literalquote{}不战\literalquote{}，把冷和平半衰期由四年抬到8--12年。
\item
  \textbf{俄国经济（R2）}：出口庄园毛损S2许可分支中心5.3\%，S3中心22.8\%。
\item
  \textbf{北欧（SC）}：S3的第一受害者是法国盟友；名义封锁＋许可证常态＋不惩瑞典→丹麦可维系。
\item
  \textbf{美国（US）}：价格＝不扣美国船；S2下英美战争概率0.30--0.45（S3
  0.55--0.75），而美法冲突≤0.08。
\item
  \textbf{法国财政（F4）}：严格S3总关税15--35m对执行成本28--38.5m与税基损伤−60至−110m；保护版关税净36--76m与史实和平年净额同量级（1806--07实测51--61m）。
  此外B5：英国商人需要可兑现的利润出口；B1b：真实商业开放是英国接受的四要素之一。
\end{enumerate}

\textbf{证据等级的标注（红队O项）}：这一开关的两半证据等级不同。\literalquote{}关税＝走私成本国有化\literalquote{}有当时提案------它就是特里亚农（1810）本身，Collin
de
Sussy与Montalivet的方案；\literalquote{}对俄关税容忍\literalquote{}有当时的对手菜单------1810年12月俄国关税、沙普塔尔所记俄方许可证配额提议（回忆录孤证）、1812年4月8日俄方\literalquote{}先撤军\literalquote{}照会；\literalquote{}不再为封锁而并地\literalquote{}有反例锚（路易的荷兰条约已签、贝格请愿被拒）。但\literalquote{}公布税则的普遍准入取代逐船许可证\literalquote{}是英国商人的诉求（1812年请愿），X2自标MODEL------\textbf{它在法方没有当时提案。}B8规则\literalquote{}仅采当时提过的方案\literalquote{}因此只覆盖开关一的四分之三；最后四分之一是模型。这直接连到第五十二章的最大未校准项：英国对中率关税的反应。

\subsubsection{二、开关二：西班牙由\literalquote{}废立\literalquote{}改为\literalquote{}保留波旁并联姻\literalquote{}（B4b）}\label{ux4e8cux5f00ux5173ux4e8cux897fux73edux7259ux7531ux5e9fux7acbux6539ux4e3aux4fddux7559ux6ce2ux65c1ux5e76ux8054ux59fbb4b}

它同时满足十一个单位：SP1（B4b是全菜单性价比最高的单项------同盟＋2.5万军＋六百万法郎／月
vs
占领−29至36万军；A层南界＝比利牛斯）；SP2（保留波旁不面对1808式总起义；三轴不触则无感）；F2（西班牙增量22.5万占S6-max缺口63\%）；F4（西班牙350／350／0：占领地榨取全被就地消耗）；MX（控制越强白银越少：S-A
90--330万 vs S-C
0--30万）；SA（法国\literalquote{}助西班牙保美洲\literalquote{}的唯一可行版本＝B4b＋B6b）；IN2（1808年6月10日的五十六舰方案在西班牙不变时可以出港）；B5（保西王不自动节省葡萄牙战费------但取消半岛战场释放英国陆军）；B6；PT（维持1803--06付费中立）；HT（海地的命运在马德里决定）；IT3（保留波旁下西班牙教会可挤压可合作）。

它的当时文书全在桌上：费尔南多1807年10月11日求婚信；拿破仑1808年1月10日的书面同意（附荣誉条件）；枫丹白露条约的同盟与出兵条款；cuestiones
proponibles第五条的\literalquote{}实现后允联姻\literalquote{}。B4b要做的只是\textbf{去条件化}------把联姻从筹码变成政策。F1给这一节点的先验是工具性0.35--0.55（王朝收益＋25,000军＋白银管道）。

红队D-3项：在路径B（无枫丹白露、无法军入西）的世界里，阿兰胡埃斯政变（1808年3月）的史实触发------缪拉军逼近与宫廷南逃计划------不存在，B4b\literalquote{}以调停人入场\literalquote{}所需的宫廷危机如何生成，未交代。在S2本义下，枫丹白露与法军入西（1807年11月至1808年2月）是史实段，危机照常生成；开关在1808年3--5月启动：不召巴约讷废立，而以调停人身份承认费尔南多七世、完成联姻、撤出四要塞以外的驻军、维持同盟。这是两条路径又一处证据等级的差别。

\subsubsection{三、其他开关}\label{ux4e09ux5176ux4ed6ux5f00ux5173}

\begin{itemize}
\tightlist
\item
  教廷（C6，工具性0.40--0.60）：1808年教廷已书面愿意事实关港；枫丹白露第5--9条含实权补偿；\literalquote{}第二个1801\literalquote{}的构件全部存在过。
\item
  荷兰汉萨（C7，非工具性0.15--0.25）：这是最贵的克制节点之一------它要求放弃直接支配；路易1810年3月的条约已签而被作废，是它\literalquote{}曾在桌上\literalquote{}的证据。
\item
  统合公债（C8，非工具性0.15--0.30）：X1判金融技术可用、拿破仑采纳概率0.15--0.30。
\item
  不征俄（C9，工具性0.35--0.55）：1811年4月致函愿先行裁军回到一年前状态（通信17579）是姿态；1812年4月8日俄方照会是最后出口；F1给不战版0.20（0.10--0.30），1811年4月19日机器启动后≈0.10------P6的0.35--0.55偏高，本书取0.20--0.45并标此分歧。
\item
  胜利后十八个月（C10，0.25--0.40）：F1的0.7冲动反向。
\end{itemize}

\begin{center}\rule{0.5\linewidth}{0.5pt}\end{center}

\subsection{第五十二章　英国如何\literalquote{}不败而接受\literalquote{}：条件、窗口、安特卫普}\label{ux7b2cux4e94ux5341ux4e8cux7ae0-ux82f1ux56fdux5982ux4f55ux4e0dux8d25ux800cux63a5ux53d7ux6761ux4ef6ux7a97ux53e3ux5b89ux7279ux536bux666e}

\subsubsection{一、两种世界，两种事件}\label{ux4e00ux4e24ux79cdux4e16ux754cux4e24ux79cdux4e8bux4ef6}

M-B1（英国在未被击败时签署承认法国大陆优势的条约，0.40--0.60）在S2本义下定价的是B1b描述的那个事件：战争自1803年持续，1807年后大陆伙伴接连消失（奥地利1805年败后受约束、普鲁士削减、俄国入体系、西班牙保持同盟、葡萄牙付费中立），英国\literalquote{}连续看不到可投入战争的大陆伙伴\literalquote{}，1809--12年先出现削减无效贸易战与试探议和，1811--14年出现能签署妥协和平的政府多数------载体是在位保守政府的政策调整或宽基内阁。1811年的信用压力（票息+53.75bp）、1812年的枢密令部分撤销（对美）、卢德与请愿，都是这个窗口的史实侧证。M-B1
0.40--0.60在此可用；B1b自己的限定一并携带：不能由\literalquote{}无大陆盟友\literalquote{}单独推出必然媾和，须有安全与贸易条件；指认某年某相只能低置信。

在路径B的无战世界里，事件是另一个：\literalquote{}亚眠和平在意大利加冕、邦联、巴符王号、法西联姻、奥婚之后能否维持九年\literalquote{}。B1b明言\literalquote{}承认现状不等于同意拿破仑今后每一项扩张都合法\literalquote{}，从未给这个事件定价；F1对1802--03窗口的克制先验只有0.05--0.15。\textbf{这两个事件不可并存，也不可互借概率。}本书对路径B的K3取\literalquote{}未定价\literalquote{}，对S2本义的K3取B1b带。

\subsubsection{二、六条件与被遗漏的第七条：安特卫普}\label{ux4e8cux516dux6761ux4ef6ux4e0eux88abux9057ux6f0fux7684ux7b2cux4e03ux6761ux5b89ux7279ux536bux666e}

B1b的六条件：保留独立海权、王朝与议会、公债、可执行的商业准入、盟友处置、王冠认可。红队指出工作区里躺着一条没有任何卡使用的同期锚点------1814年5月30日巴黎和约第XV条：\literalquote{}Antwerp
shall for the future be solely a commercial port.\literalquote{}（法文本：\literalquote{}Dorénavant
le port d\textquotesingle Anvers sera uniquement un port de
commerce.\literalquote{}）同条前文把荷兰的舰船与兵工厂（\literalquote{}especially the fleet in the
Texel\literalquote{}）排除在法国可保留的份额之外。这是英国在\textbf{战胜}时写进和约的价格；它告诉我们英国的战略心脏在哪里：1793年斯海尔德开放是开战理由，1809年瓦尔赫伦远征的目标就是安特卫普的舰队与船厂（F3已录），1815年后的荷兰屏障与斯海尔德政策一脉相承。F3的判断：\literalquote{}这一结果解释法国为什么重视该港，也解释英国为什么对其存在特别敏感。\literalquote{}

【机】英国对低地国家的安全关切是一个跨越1689--1914的长时段常量（斯海尔德开放为1793年开战理由、瓦尔赫伦远征、荷兰屏障、比利时中立），其适用条件是英国仍以海峡安全为第一安全利益------在任何法国不败的世界线里都成立。问题在于路径B与S2本义都要求两件事同时成立：(i)
英国在武装和平下\literalquote{}接受\literalquote{}（K3）；(ii)
法国以安特卫普为主厂（1807--13年下水五十二艘中安特卫普十八艘）把1830年的战列舰效能比推到.876（分问③）。B3的升档触发器以\literalquote{}可出海舰数≤40--60\literalquote{}表述，没有把安特卫普本身当触发器。\textbf{这两件事在同一世界线内很可能互斥。}本书给出两条终局带（第六十三章详）：

\begin{itemize}
\tightlist
\item
  \textbf{终局α（接受＋差距不缩）}：英国接受的价格含安特卫普非军事化或限舰。法国海军的起点减少十二舰、年下水减三------P4／F3的\literalquote{}失安特卫普\literalquote{}分支：1830年A12主线降至50（W）／71（P）；P4的W基准下吨比.444、效能比.267（\literalquote{}单一冲击，英国反应保持W基准\literalquote{}）；在P基准下按F3的71舰对主线111舰线性折算，吨比≈.58、效能比≈.56（本书作者的比例近似，非P4模型重跑）。这是\literalquote{}英国不败而接受\literalquote{}的\textbf{最可能形态}：英国用它在1814年真的写进和约的那一条来换承认。
\item
  \textbf{终局β（不接受＋追赶）}：法国保留安特卫普军港并扩建，英国把它当升档触发器，保持W战备------P4的\literalquote{}法P、英保W\literalquote{}：吨比.902、效能比.731；K3不成立，武装共存代替和约（P6自认最可能的终局）。
\end{itemize}

在B1b的窗口里，英国的要价会落在α；法国接受α的先验是F1的非工具性克制（放弃直接支配的一部分海军资源），0.15--0.25。\textbf{因此M-B1的0.40--0.60应读作\literalquote{}英国愿意签α型和约的概率\literalquote{}，而法国愿意接受α型条款的概率是另一个乘子。}这是本书对P6三个外生事件的最重要修正。

\subsubsection{三、K3的可核方法}\label{ux4e09k3ux7684ux53efux6838ux65b9ux6cd5}

引擎从未产生\literalquote{}英国签约\literalquote{}状态，K3是手填的。可观测的代理：王冠认可（乔治三世的天主教立场、摄政后的变化）、国内多数（1811--12年枢密令辩论的票数、1812年谢菲尔德与伯明翰请愿的措辞）、海军安全（安特卫普与斯海尔德条款、舰队交出的拒绝）、商业准入（许可证数量与关税水平）。改判证据：1806--07枢密令辩论与1812年撤令条件；Board
of
Trade与珀西瓦尔／坎宁的立场文书；1806年英法谈判的法方书面条款；1813--14卡斯尔雷对低地国家与安特卫普的训令。

\begin{center}\rule{0.5\linewidth}{0.5pt}\end{center}

\subsection{第五十三章　概率：从1803年出发，这条路有多现实}\label{ux7b2cux4e94ux5341ux4e09ux7ae0-ux6982ux7387ux4ece1803ux5e74ux51faux53d1ux8fd9ux6761ux8defux6709ux591aux73b0ux5b9e}

\subsubsection{一、两个综合卡的口径}\label{ux4e00ux4e24ux4e2aux7efcux5408ux5361ux7684ux53e3ux5f84}

P6（共因模型）：条件于\literalquote{}法国采纳并维持有界目标框架\literalquote{}，三国默许由同一个潜变量（法国可信自限）驱动且正相关，联合0.30--0.50；无条件（乘以F1的采纳先验）0.018--0.090；再乘秩序存续0.008--0.059（条件版0.135--0.325）。P8（λ模型）：F采纳.15--.30，W持续.40--.60，E／R／A＝.40--.60／.35--.50／.45--.60，J＝(1−λ)ERA＋λ·min(E,R,A)，λ＝0--.70→C1815
0.38--7.27\%（中点2.49\%），O1848
0.17--4.44\%（中点1.27\%）；若坚持P6的强共因（λ＝.70--1）则条件J
.264--.50、无条件C 1.8--9.0\%、O
0.81--5.85\%。P8的方法论异议：既已条件于法国采纳克制，再把\literalquote{}共同法国政策\literalquote{}当英俄奥的强共因可能重复计入。Tornado：λ
1.259pp、F采纳.847pp、W .508pp、R .448pp、红线份额r
.416pp（r＝.30时中心O由1.27\%降至≈0.97\%）、E .217pp、A
.119pp。本书并列两口径不裁决。

\subsubsection{二、把标题改回来}\label{ux4e8cux628aux6807ux9898ux6539ux56deux6765}

红队F项：§1.1把\literalquote{}存在，且不需要任何外生奇迹\literalquote{}作为标题，而正文的无条件概率是0.4--9\%，三个C类节点恰在奇迹阈值，F1对1802--03窗口克制给0.05--0.15，P6自己认为最可能的终局是\literalquote{}英国始终不签约的武装共存\literalquote{}。本书接受：\textbf{\literalquote{}逻辑可构造\literalquote{}不等于\literalquote{}现实可行\literalquote{}。}结论表改为：

{\def\LTcaptype{none} % do not increment counter
\begin{longtable}[]{@{}
  >{\raggedright\arraybackslash}p{(\linewidth - 2\tabcolsep) * \real{0.5000}}
  >{\raggedright\arraybackslash}p{(\linewidth - 2\tabcolsep) * \real{0.5000}}@{}}
\toprule\noalign{}
\begin{minipage}[b]{\linewidth}\raggedright
命题
\end{minipage} & \begin{minipage}[b]{\linewidth}\raggedright
判决
\end{minipage} \\
\midrule\noalign{}
\endhead
\bottomrule\noalign{}
\endlastfoot
存在一条内部自洽、其全部构件在当时都被提出过的路径，达到本项目操作化的\literalquote{}控制整个欧洲\literalquote{}
& 是------以S2本义的形态（有条约支撑）；以路径B的形态只达到primacy \\
条件于法国采纳并维持有界目标框架，该路径走通的概率 & P6 0.30--0.50；P8 J
0.063--0.404（强共因0.264--0.50） \\
从1803年出发的无条件概率 & P6 0.018--0.090；P8 0.0038--0.0727 \\
最可能的实际终局（P6备选一） &
大陆控制、俄国不战、英国不签约的武装共存------不满足K3 \\
它要求拿破仑做的事 &
在三个节点（对英写定交易且不再加价／不并荷汉／自缚财政）做他在1802--03与1810--11两个史实和平窗口里都没做的事；S2本义少一个（无C3），但多两场已打过的仗的代价 \\
\end{longtable}
}

\subsubsection{三、C1的调和与C9的分歧}\label{ux4e09c1ux7684ux8c03ux548cux4e0ec9ux7684ux5206ux6b67}

P6把1803年马耳他节点（C1）标为工具性0.45--0.60，而F1对同一窗口的显示偏好判定是0.05--0.15（B1的(i)\literalquote{}以马耳他换和平\literalquote{}≈0.15）。两者相差三到四倍。本书的处置：C1只属于路径B；S2本义没有C1。这个分歧因此不影响主线，但它是路径B概率带的一个未解口径差。C9（不征俄）P6给0.35--0.55，F1给不战版0.20（0.10--0.30）；本书取0.20--0.45，并注明它是S2本义最贵的一个节点------1812年是\literalquote{}S2的生死关\literalquote{}，只有把\literalquote{}接受1810年关税＋配额\literalquote{}包装成拿破仑的胜利才能压住开战（F1）。

\subsubsection{四、这条路的\literalquote{}现实\literalquote{}是什么样的}\label{ux56dbux8fd9ux6761ux8defux7684ux73b0ux5b9eux662fux4ec0ux4e48ux6837ux7684}

从1803年的布洛涅出发，一个最可能的、由本书证据支撑的版本是这样的：拿破仑照史实打赢1805、1806、1807三场战争，然后------在1807年秋或1808年春------做出三个他在史实中没有做的选择：把封锁变成关税，把西班牙变成联姻的盟友，把教皇留在罗马。俄国因此不战，奥地利因此不再入场，英国因此在1811--14年间失去继续战争的大陆理由，并在安特卫普问题上开出价码。到1812年，法国控制了大陆的四层，俄国被条约与关税约束，英国要么签α型和约（差距不缩），要么保持武装共存（差距在P档追赶而无和约）。这条路从1803年出发的概率是个位数百分比，因为它要求一个\literalquote{}靠自己爬上来的军人\literalquote{}在胜利之后停下来；它是\literalquote{}合理\literalquote{}的，因为每一个构件都在当时的桌上；它是\literalquote{}可行\literalquote{}的，因为账算得平；它不是\literalquote{}现实\literalquote{}的，如果\literalquote{}现实\literalquote{}的意思是\literalquote{}拿破仑会那样做\literalquote{}。

\begin{center}\rule{0.5\linewidth}{0.5pt}\end{center}

\section{第三部　行省化：最重要的问题}\label{ux7b2cux4e09ux90e8-ux884cux7701ux5316ux6700ux91cdux8981ux7684ux95eeux9898}

\subsection{第五十四章　三层回答与五本账}\label{ux7b2cux4e94ux5341ux56dbux7ae0-ux4e09ux5c42ux56deux7b54ux4e0eux4e94ux672cux8d26}

\subsubsection{一、三层回答}\label{ux4e00ux4e09ux5c42ux56deux7b54}

\literalquote{}法国把欧洲行省化，是否真有合理、现实、可行的路径？\literalquote{}------答案取决于\literalquote{}行省化\literalquote{}指什么。

\textbf{字面版（A层覆盖大陆）：不存在这样的路径。不是\literalquote{}很难\literalquote{}，是账算不平。}
置信高。这是本书全部证据里最稳固的一层：F2的S6-max中心缺12.2万在场兵、S6-full缺30.2万；D2的既有外圈治安即吃掉6--10万＝S6-lite全部人力富余（2.9万）的2倍以上；F4的S6每新增五百万人在贫弱高摩擦区每百万人净−11m；G1的吞南德＝33k盟军换60--90k驻军；R3的俄军存在封死S6东扩；P2的丙层（再并南德、那不勒斯、全西班牙至77.3m人口）中心需求99.04万在场军对供给51.06万、年军费缺532.6m，低驻军高税余的角点仍缺228.7m，连主动构造的联合乐观案也缺215,514在场兵与174.98m／年；P8判全欧A层在既定资源包络内结构为零，模型外0--0.5\%只是分析师的失模容差。这些账目彼此独立到足以互证，参数扰动不足以翻转。

\textbf{边界版（A层限于内圈＋托斯卡纳／帕尔马型＋条件性贝格型）：存在，且已被史实部分实现。}
置信中高（方向）、中（边界与时点）。约39.5--40.5百万人口代理的直辖核心（P2）------旧法国＋比利时、左莱茵、皮埃蒙特、利古里亚、莱芒；帕尔马另列；托斯卡纳、瓦莱条件直辖；贝格须获真实的法国市场准入------1806--1810巩固，1825--1835制度稳态。P6的\literalquote{}已核上限≈3,700万\literalquote{}（旧法国≈2,900万＋内圈785万＋托斯卡纳／帕尔马型未取得独立人口）与P2的差是口径差（基年、纳入单位），不是判定分歧。

\textbf{控制版：以S2本义的形态存在（有条约支撑），以路径B的形态只达到primacy。}
见第五十章。

\subsubsection{二、阶段地图：哪里并、哪里不并}\label{ux4e8cux9636ux6bb5ux5730ux56feux54eaux91ccux5e76ux54eaux91ccux4e0dux5e76}

P2的阶段地图给每个单位三个时期的处置（甲1806--1810／乙1810--1820主线／丙1820--1840及1848审查）：比利时、左莱茵、皮埃蒙特、利古里亚、莱芒------既有A，固定税兵与地方职位，继承期保产权与教会协议；帕尔马留A；托斯卡纳条件过渡，低征募与港口贸易下留A，若持续赤字则转B（\literalquote{}不以省名强保\literalquote{}）；瓦莱按战略走廊考核，补贴型A；贝格------市场平等且本地官僚留任才转A，无平等准入则取消A；威斯特法利亚与其余德西------保军政机器，\literalquote{}条约协作而非为抹掉国界而兼并\literalquote{}；意大利王国------保米兰财政、军旗及两冠分离承诺，罗马王的意大利王位安排须另谈，不等于法国省；荷兰------保B并改善通商，\literalquote{}不以港口控制为由吞并\literalquote{}；汉萨------商业条约C，不保法国专属许可证歧视；罗马------主权归教廷，法律、财产与宗教权利分项协约；伊利里亚------自治边防或和约归还，\literalquote{}不将军事克罗地亚改成普通征兵省\literalquote{}；加泰与其余西班牙------保波旁主权，\literalquote{}停止另造正统性真空\literalquote{}；南德------君主B／C，继承时重订盟约；那不勒斯本土------降战争支出、保本地军队，西西里另案。\textbf{主线边界有意小于1812年的纸面最大直辖疆界：不是法国再也无法占领荷兰或罗马，而是为稳态减去那些可以通过条约取得主要战略服务、却要用兼并承担全部债务和政治责任的地区。}

\subsubsection{三、联合账：最宽的人力案尚且不等于能付钱}\label{ux4e09ux8054ux5408ux8d26ux6700ux5bbdux7684ux4ebaux529bux6848ux5c1aux4e14ux4e0dux7b49ux4e8eux80fdux4ed8ux94b1}

固定假设：F2的S2任务基线54万在场（内圈7.85m×4‰＝3.14万已含）；原北意25,000、威国12,500、南意20,000法军任务在新增驻军中扣除；人力偏宽基线＝法国50万在册＋盟军20万在册、在场率0.85；并省失去的独立军先扣除、再加入成熟的新法国兵；法国需要支付的在册人数＝(总在场需求−剩余自养盟军在场)/0.85；年军费缺口＝所需法国在册×700法郎−陆军预算300m−新地区可汇现金；退并省使旧法国收入减少的净额另作0／30／60m压力。

{\def\LTcaptype{none} % do not increment counter
\begin{longtable}[]{@{}
  >{\raggedright\arraybackslash}p{(\linewidth - 10\tabcolsep) * \real{0.1304}}
  >{\raggedleft\arraybackslash}p{(\linewidth - 10\tabcolsep) * \real{0.1739}}
  >{\raggedleft\arraybackslash}p{(\linewidth - 10\tabcolsep) * \real{0.1739}}
  >{\raggedleft\arraybackslash}p{(\linewidth - 10\tabcolsep) * \real{0.1739}}
  >{\raggedleft\arraybackslash}p{(\linewidth - 10\tabcolsep) * \real{0.1739}}
  >{\raggedleft\arraybackslash}p{(\linewidth - 10\tabcolsep) * \real{0.1739}}@{}}
\toprule\noalign{}
\begin{minipage}[b]{\linewidth}\raggedright
情景
\end{minipage} & \begin{minipage}[b]{\linewidth}\raggedleft
直辖人口
\end{minipage} & \begin{minipage}[b]{\linewidth}\raggedleft
军事需求在场
\end{minipage} & \begin{minipage}[b]{\linewidth}\raggedleft
人力余额
\end{minipage} & \begin{minipage}[b]{\linewidth}\raggedleft
法国需支付在册
\end{minipage} & \begin{minipage}[b]{\linewidth}\raggedleft
年军费缺口
\end{minipage} \\
\midrule\noalign{}
\endhead
\bottomrule\noalign{}
\endlastfoot
甲 内圈＋帕尔马 & 38.25m & 544.6k & +50.4k & 440.7k & 8.5m \\
乙筛选无贝格 & 39.48m & 551.1k & +43.9k & 448.4k & 13.8m \\
乙筛选含贝格 & 40.48m & 558.6k & +33.7k & 464.2k & 20.9m \\
乙四地全并 & 51.35m & 622.4k & −70.0k & 628.2k & 135.7m \\
乙四地＋历史其他外圈 & 56.68m & 662.3k & −116.7k & 683.1k & 174.2m \\
丙 南德、南意、西班牙全并 & 77.32m & 990.4k & −479.8k & 1,135.1k &
532.6m \\
历史直辖地图量级对照 & 43.58m & 584.5k & +3.7k & 495.7k & 47.0m \\
\end{longtable}
}

征募与自治军机会成本：乙四地＋外圈年名义征额31,001、年抵营假设20,688、十年后成熟新增在场38,936，而失去原B军在册104,000；丙分别为59,170／31,912／60,106／170,000。即使100\%保留原B军，乙四地＋外圈中值仍缺约6.70万在场、120.4m军费；丙仍缺39.23万、430.6m。\literalquote{}若保留本地军权、预算权、王朝与司法否决点才取得合作，其实际制度已接近B而不是均一法国省。\literalquote{}

\textbf{有限版也不免费。}
按陆军300m、20万自养盟军、700法郎／兵年，含贝格仍缺20.9m。可闭合的示例：盟军在册25万，陆军330、海军165，退并省减少旧中央净税30------可支付法国434,286在册；并贝格失7,000原盟军后余243,000自养盟军；合计575,693在场对558,600需求余17,093（无贝格方案余25,686）。改判阈值：含贝格在该组合下多损失旧净税约14.1m或多增约1.71万在场任务即耗尽剩余空间；退并省净税损为60m而非30、或同时增加3万在场任务，方案转负。\textbf{成立条件是盟约兑现和海军节奏让步，而不是凭空得到和平红利。}

\subsubsection{四、五本账为什么是同一笔账}\label{ux56dbux4e94ux672cux8d26ux4e3aux4ec0ux4e48ux662fux540cux4e00ux7b14ux8d26}

【机】占领型帝国的边际账：每一块新领土的净贡献＝（实收税＋净供兵）−（驻军＋行政＋承接债务＋失去的盟军与市场）；史实类比是奥治伦巴第--威尼西亚（1815--48：可运行、可镇压、不可赢得，驻军七万）与法属阿尔及利亚（1846／47十万驻军）；适用条件是新领土有自己的正统性焦点（王朝、教会、foral）且外部军事机会不为零。兵力账（F2）、财政账（F4）、治安账（D2）、盟军账（G1）、整合账（Q2）在A层越过内圈之后同时转负，而它们不是五笔独立的账：外圈是封锁执法的对象，封锁摧毁的正是外圈的税基；并省失去的是自养盟军，得到的是第十年才转正（且只在hazard6口径下）的新省兵；驻军密度越高，教会与征兵两个接口越容易断；断了就要更多驻军。Q2的净兵模型（第10年每百万人L
−0.04至+5.32千、M −7.10至+2.19千、H
−26.73至−6.03千）和F2的役期桥（hazard6 +0.790‰ vs
五年含训练−0.463‰）说的是同一件事：\textbf{新省的净供兵在最有利的假设下才勉强为正，在中等安全负担下已为负。}

\subsubsection{五、按路径疆域再基经常收入（红队i\_cbd9c97817，本书作者重算）}\label{ux4e94ux6309ux8defux5f84ux7586ux57dfux518dux57faux7ecfux5e38ux6536ux5165ux7ea2ux961fi_cbd9c97817ux672cux4e66ux4f5cux8005ux91cdux7b97}

F4的T＝750m锚于1810年6月普通账已纳入的疆域与税目（含皮埃蒙特、利古里亚、托斯卡纳、帕尔马、罗马、比利时、莱茵；不含7月并入的荷兰与12月并入的汉萨）；F4自给宽下行650、高收入案800。P6的1813年恒等式用的是高收入案：800净经常收入＋40可兑现盟约＝300旧债／非军＋25新补偿＋20储备＋300陆军＋195海军。两条路径的A层都比1810年6月疆域小。用Marion
p.321的1812年预期毛税（净毛比226.39/342.26≈0.66）近似扣除：

{\def\LTcaptype{none} % do not increment counter
\begin{longtable}[]{@{}
  >{\raggedright\arraybackslash}p{(\linewidth - 8\tabcolsep) * \real{0.1667}}
  >{\raggedright\arraybackslash}p{(\linewidth - 8\tabcolsep) * \real{0.1667}}
  >{\raggedleft\arraybackslash}p{(\linewidth - 8\tabcolsep) * \real{0.2222}}
  >{\raggedleft\arraybackslash}p{(\linewidth - 8\tabcolsep) * \real{0.2222}}
  >{\raggedleft\arraybackslash}p{(\linewidth - 8\tabcolsep) * \real{0.2222}}@{}}
\toprule\noalign{}
\begin{minipage}[b]{\linewidth}\raggedright
路径
\end{minipage} & \begin{minipage}[b]{\linewidth}\raggedright
扣除项（毛）
\end{minipage} & \begin{minipage}[b]{\linewidth}\raggedleft
净扣除
\end{minipage} & \begin{minipage}[b]{\linewidth}\raggedleft
T中心
\end{minipage} & \begin{minipage}[b]{\linewidth}\raggedleft
T高
\end{minipage} \\
\midrule\noalign{}
\endhead
\bottomrule\noalign{}
\endlastfoot
S2本义（保利古里亚，不并罗马；托斯卡纳条件） &
罗马16.5［＋托斯卡纳22.5］ & 10.9--25.8 & \textbf{724--739} &
774--789 \\
路径B（另不并利古里亚） & 罗马16.5＋托斯卡纳22.5＋利古里亚16 & 36.4 &
\textbf{714} & 764 \\
\end{longtable}
}

代入恒等式（陆海合计＝T＋盟约40−300−25−20）：

{\def\LTcaptype{none} % do not increment counter
\begin{longtable}[]{@{}lrrr@{}}
\toprule\noalign{}
情形 & 陆海合计 & 海军195时陆军余量 & 对陆军线300的缺口 \\
\midrule\noalign{}
\endhead
\bottomrule\noalign{}
\endlastfoot
S2本义·高收入＋盟约 & 469--484 & 274--289 & \textbf{11--26m} \\
S2本义·中心＋盟约 & 419--434 & 224--239 & 61--76m \\
S2本义·中心无盟约 & 379--394 & 184--199 & 101--116m \\
路径B·中心（盟约存疑） & 369（＋40） & 174（214） & 86--126m \\
F4宽下行650 & 305（＋40） & 110（150） & 150--190m \\
\end{longtable}
}

结论：\textbf{\literalquote{}陆300＋海195\literalquote{}只在高收入案＋可兑现盟约下接近闭合（缺11--26m）；在中心案下缺60--115m。}MAN-02／FIS-02的\literalquote{}闭合-条件\literalquote{}须改写为\literalquote{}紧闭合，条件更苛\literalquote{}：或海军降至130--165m（P2闭合案的165、F3的低档），或法国在册降至40--43万并以25万自养盟军补足（这正是F2／P2的主线），或盟约转移达到约100m／年------这个量级恰好是史实的外部现金依赖（F4：123.1m／年）。\textbf{S2本义比路径B多出的盟约转移（普鲁士贡赋、邦联配额折现、意大利王国分摊）是它在财政上更可行的又一个理由}；而在路径B里，\literalquote{}盟约40m\literalquote{}的来源本身存疑。Marion数字是预期而非实收、F4的T是构造值------本表是量级校核，不是决算。

\begin{center}\rule{0.5\linewidth}{0.5pt}\end{center}

\subsection{第五十五章　为什么外圈必断}\label{ux7b2cux4e94ux5341ux4e94ux7ae0-ux4e3aux4ec0ux4e48ux5916ux5708ux5fc5ux65ad}

外圈的证据是本书里独立来源最多的一组（第二十一、二十、十二、二十八、二十四、二十六、二十七、四十四、四十五章），此处只列每个单位的一句话与它的锚：

\begin{itemize}
\tightlist
\item
  \textbf{荷兰}：可兼并不可消化------年金78→26m烧毁资本市场；西荷兰常驻15,000法军至1813；580余起骚动、80余起起义、三分之二在并合后；驻军一撤三周蒸发。拿破仑自认荷兰行政比法国省钱，兼并动机纯为封锁、兵源与债务处置。
\item
  \textbf{汉萨}：可并而财政自败------1813年罚金4,800万≈汉萨年税目标38.791m的1.3倍；并合连市场准入这一第一补偿都没给。
\item
  \textbf{贝格}：唯一主动求并的单位，被法国左岸保护主义否决（科隆商会1810年9月16日书）------\literalquote{}能并的不值钱，值钱的不能并\literalquote{}的反面：值钱且愿意并的，被法国自己的利益集团拒绝。
\item
  \textbf{萨克森}：不可并------帝国自己的驻邦联公使Bacher
  1810年10月2日已论证延伸关税＝厄尔士山赤贫化＋法国零收益。
\item
  \textbf{威斯特法利亚}：并入＝自拆样板招牌；接收国账证明分给客户王朝≠法国省钱。
\item
  \textbf{罗马}：1810年首征三分之一拒征、一半开小差；1813年\literalquote{}un caractère
  insurrectionnel
  grave\literalquote{}；教会的忠诚为零------教廷是A层可扩性的第一证伪器。
\item
  \textbf{伊利里亚}：令征18,000实得3,000／年；\literalquote{}走私成国民工业\literalquote{}。
\item
  \textbf{加泰罗尼亚}：1812年敕令不敢刊公报、1813年民政自撤------全西班牙的A层试点在最大强制＋最大让步下双败；纳瓦拉榨取1.52亿reales仍征收崩溃。
\item
  \textbf{爱尔兰}：法国长期占领需33--66k而汲取上限≈£5--5.5m→净消耗单位；上限B级。
\item
  \textbf{葡萄牙}：占得住壳拿不到瓤------B级且−55至−70m／年。
\item
  \textbf{那不勒斯}：A层净负资产------卡拉布里亚两省净流出为自身产出的6--8倍，1812年军费占预算80\%。
\item
  \textbf{瑞士}：不必并也不能并------调停模式零成本换16,000条约兵。
\item
  \textbf{波兰}：S6行省化是唯一能把最亲法单位推向抵抗的情景。
\item
  \textbf{西班牙}：A层南界＝比利牛斯；350／350／0。
\end{itemize}

Q1给出这张地图的结构判据：外部军事机会O＝0时持续抵抗率为0，O＝1且挑衅≥2项时十三个单元全部大起义------\textbf{A层的地理边界是英奥俄军事投送的边界，不是文化边界}；D2的六型失败模式（渐进磨损／先行榨取／合法性罢工／城市暴动反复／滞后爆发／战争吞噬＋荒地治理）互不相同，故S6扩张不存在单一的修补方案；断裂点＝封锁伤害×精英补偿的交点；外圈失败是政策组合选定而非社会禀赋注定（罗马与加泰除外：对政策不敏感）。X2补了技术判据：莱茵河可守，贝格东界不可守。\textbf{\literalquote{}外圈可管、不可省化\literalquote{}}------可占有并榨取（负净值），不能转化为供兵供税的正常省。

\begin{center}\rule{0.5\linewidth}{0.5pt}\end{center}

\subsection{第五十六章　乙层：技术可闭合、政治未证}\label{ux7b2cux4e94ux5341ux516dux7ae0-ux4e59ux5c42ux6280ux672fux53efux95edux5408ux653fux6cbbux672aux8bc1}

乙层（北意王国、德西、伊利里亚、加泰新增13.1m；合并历史外圈≈56.676m）在P2的普通中心账下不闭合（缺11.67万在场、174.2m／年）。但P2主动构造的联合乐观反例------低驻军＋100\%保留原B层军队建制与自养税源并扣除新A兵＋陆军350m／700法郎＋其他任务减30,000在场＋高税余------得需求571,650在场、物理供给余23,598、军费余19.23m；同案丙仍缺215,514与174.98m。\textbf{因此乙全并不能以\literalquote{}纯技术不可能\literalquote{}结案。}

F2二轮的核对：P2的成熟函数本已是四批受训／年留存0.9的五年含训练口径，不是F2的hazard6；乙行余额中新增省兵只贡献247.6人，基础(500k＋200k)×0.85＝595k−571,650＝23,350仍正------役期口径桥不自动推翻乙反例。真实的承重点因此转移到两处：(a)
\textbf{法国50万在册＋原自养军20万能否在承诺役期、补充率与350m下长期维持}------F2的维持式C＝500k×(1/6＋0.025)/0.9≈10.6万／年，高于F5和平年征≤8--10万的社会容忍带（两卡隐含役期与退出率口径未合）；(b)
\textbf{\literalquote{}保留军队建制≠保留王朝主权\literalquote{}}------若军队真正改受法国直接指挥、地方税经中央授权持续自养，它仍可能是A，但这要IT1／G1／SP1的政治证据判定：北意分冠（1805年art.3）、加泰主权、南德四条件、地方自养与低驻军能否同时实现------证据不足。P8的压力算：乙至1848给定O为0.12--2.52\%（选择吞并.15--.30×资源角.20--.40×精英归顺.10--.30×继承存续.40--.70，纯P8压力计算，非P2估计）。

\textbf{裁定}：乙层作为1830--40年代的低置信政治协约支线并列保留，不写成主线，也不抹平。这是PROJECT.md预案的分叉触发点（\literalquote{}P2判技术可行而单位卡判政治不可行\literalquote{}）；用户裁定不开分叉，本书把它挂为未决。改判证据：P2乙案逐地法军／盟军不重复计算的人力与自养税册；IT1／G1合法转隶文书；1812年同口径在岗文官与宪兵实员；荷／汉／罗年度征兵面板。反过来的改判也要写明：这些角点\literalquote{}证明当前包络下全并没有自动自付，不是逻辑上排除一切未来税基翻番的世界\literalquote{}；若要借1840年代的工业增长救1830年的全并，必须另证增长不会被十多年驻军、信用破坏和战争反复吃掉------不能直接把P5／X3的和平增长移植到S6。

\begin{center}\rule{0.5\linewidth}{0.5pt}\end{center}

\subsection{第五十七章　官员、宪兵、语言与教会：如何真正实施}\label{ux7b2cux4e94ux5341ux4e03ux7ae0-ux5b98ux5458ux5baaux5175ux8bedux8a00ux4e0eux6559ux4f1aux5982ux4f55ux771fux6b63ux5b9eux65bd}

\subsubsection{一、人才存量不能从历任名单倒推}\label{ux4e00ux4ebaux624dux5b58ux91cfux4e0dux80fdux4eceux5386ux4efbux540dux5355ux5012ux63a8}

【锚】1809年11月12日陆军部长报告：\literalquote{}À l\textquotesingle aune de ce
document, l\textquotesingle Arme est forte de 15 474 hommes ... Jamais
les brigades n\textquotesingle auraient pu supporter le départ de quatre
mille de leurs membres hors des frontières de
l\textquotesingle Empire.\literalquote{}------全国宪兵15,474；西班牙的4,000宪兵一半从内地、一半从正规军仓底招募，作者认为国内不能承受全额抽出。1811年制度条令26,000宪兵\literalquote{}Encore
ne s\textquotesingle agit-il ici que de chiffres prévus par les
différents règlements, le nombre de gendarmes étant certainement bien
inférieur sur le terrain.\literalquote{}------编制非实员。1813年30,600\literalquote{}À côté de cette
force « territoriale » ou prévôtale ... plusieurs forces spécialisées
... passent de 15 689 hommes en 1801 à 30 600 hommes en
1813.\literalquote{}------含军中、海事、精锐、巴黎、西班牙等专种，不能当1812年新省治安可自由拨付的库。Lignereux书评：\literalquote{}Ouverte
aux indigènes sans atteindre les quotas fixés, la gendarmerie recrute
majoritairement parmi les anciens Français, au contraire des
commissaires de police d\textquotesingle origine essentiellement
locale.\literalquote{}------累计7,571名原并省宪兵、620名警察专员，前者多数原法国出身；本地宪兵比例≤0.5仍属有利假设。

新增岗位（相对甲阶段，中值）：乙筛选含贝格------8个首长席位、1,280专业行政骨架（外调272）、1,341宪兵（法国来源804）；乙四地＋历史外圈------68／10,880（3,032）／15,191（9,783）；丙------138／22,080（7,952）／50,693（36,205）。乙全并十年外调专业骨架约303人／年、法国来源宪兵约978人／年，另加本地招募与全部替补；丙二十年平摊仍需每年外调骨架约398、法国来源宪兵约1,810。auditeur管道1813年仅452人。\literalquote{}人口更大可以多招兵，却不自动多出愿为法国征兵的村长与书记。\literalquote{}

\subsubsection{二、1.5倍检验：条件否决，不是已证否决}\label{ux4e8c1.5ux500dux68c0ux9a8cux6761ux4ef6ux5426ux51b3ux4e0dux662fux5df2ux8bc1ux5426ux51b3}

省长首长位：史实1811年130省量级对照，乙四地160、含外圈178都未超195；丙248超过------但只能比较最高职位数量。宪兵：乙四地中值24,110，若1812年同口径实际供给小于16,073且无本地替代则超过1.5倍；乙加外圈29,031对19,354；丙64,533对43,022。\textbf{1812年同口径在岗文官与宪兵实员未得，故行政硬阈值暂为条件否决。}否决全欧并省的承重证据是联合账与政治机会成本，不是伪造一个行政人员硬天花板。

\subsubsection{三、语言与教会}\label{ux4e09ux8bedux8a00ux4e0eux6559ux4f1a}

【锚】Woolf记录意大利语与德语人员需求的差别及1811--12年为德语职位调整人事；Rowe说明法籍莱茵宪兵须会德语。可执行设计：中央统一总账与基层可理解的程序分离------甲阶段保留地方诉讼语言与双语告示；乙给新任干部一至两年语言、地方法与账务见习，配本地副手；民事法律翻译先于征募加额；德语区不以法语考试一刀切清洗地方行政；丙不以撤除自治学校或教会追求一代人语言同化，\literalquote{}否则培训节省会转为驻军支出\literalquote{}。加泰的语言让步实验（1810年奥热罗）已证明：让步不能替代主权。教会：D1的倒U形与IT3的三层否决权说明，教会项是系统变量------\literalquote{}第二个1801\literalquote{}（专约＋教皇有限主权）是A层能扩到内圈之外的必要条件，而不是充分条件。

\begin{center}\rule{0.5\linewidth}{0.5pt}\end{center}

\section{第四部　法胜欧洲的运作，1815--1848}\label{ux7b2cux56dbux90e8-ux6cd5ux80dcux6b27ux6d32ux7684ux8fd0ux4f5c18151848}

\subsection{第五十八章　骨架年表（修正后）}\label{ux7b2cux4e94ux5341ux516bux7ae0-ux9aa8ux67b6ux5e74ux8868ux4feeux6b63ux540e}

\subsubsection{一、说明}\label{ux4e00ux8bf4ux660e}

P7交付的338条四栏年表（\texttt{P7\_chronicle.csv}）以路径B为骨架，其\literalquote{}无提尔西特级联\literalquote{}（无华沙公国、无哥本哈根1807、无芬兰战争、无1806--12俄土战争、无芬肯施泰因、普鲁士完整无改革、无1809、无并合浪潮、加勒比无易手、毛里求斯--爪哇不失）与十四项前提裁定（PRE-01至14）只对路径B成立。本书的主线是S2本义（第五十章），因此1803--1807年即史实，上述级联多数\textbf{不发生}------这也是S2本义比路径B更贴近单位卡估计域的原因。下表是S2本义的骨架年表：每条给锚点或机制引与置信；【B】标出路径B在同一年的分岔；不重复P7的338条，只列改写世界结构的条目。

\subsubsection{二、1803--1807：史实段}\label{ux4e8c18031807ux53f2ux5b9eux6bb5}

{\def\LTcaptype{none} % do not increment counter
\begin{longtable}[]{@{}
  >{\raggedright\arraybackslash}p{(\linewidth - 6\tabcolsep) * \real{0.1000}}
  >{\raggedright\arraybackslash}p{(\linewidth - 6\tabcolsep) * \real{0.6100}}
  >{\raggedright\arraybackslash}p{(\linewidth - 6\tabcolsep) * \real{0.2100}}
  >{\raggedright\arraybackslash}p{(\linewidth - 6\tabcolsep) * \real{0.0800}}@{}}
\toprule\noalign{}
\begin{minipage}[b]{\linewidth}\raggedright
年
\end{minipage} & \begin{minipage}[b]{\linewidth}\raggedright
条目
\end{minipage} & \begin{minipage}[b]{\linewidth}\raggedright
锚／机
\end{minipage} & \begin{minipage}[b]{\linewidth}\raggedright
置信
\end{minipage} \\
\midrule\noalign{}
\endhead
\bottomrule\noalign{}
\endlastfoot
1803 &
5月英法决裂；布洛涅集结；德克雷向爱尔兰人开价25,000人＋40,000枪（条件二万爱尔兰人加入）；9月27日法瑞军事协约16,000条约兵；法国把汉诺威作为对英筹码
& 史实；IE、NL卡 & 史实 \\
1804 &
称帝；当吉安案使俄廷反法（无封锁的1803--05年俄廷已反法）；德卡恩接待辛迪亚使者；Ouvrard--Hope白银合同机器启动
& 史实；R1、IN2、X1 & 史实 \\
1805 &
特拉法加；乌尔姆--奥斯特里茨；普雷斯堡（威尼西亚入意大利王国、蒂罗尔入巴伐利亚、巴符王号）；塔列朗斯特拉斯堡备忘录未采；橡木普查显示安特卫普流域62\%大橡木
& 史实；F1、F3 & 史实 \\
1806 &
莱茵邦联（7月12日）；神圣罗马帝国终结；耶拿；柏林敕令；意大利币制法令；那不勒斯约瑟夫；1806年10月16日邦联筹备会因巴符缺席流产
& 史实；G1、G2、P5 & 史实 \\
1807 &
弗里德兰、提尔西特（7月）：华沙公国、残普、俄入体系；芬肯施泰因（5月）；巴约讷条款对波军；枫丹白露对葡（10月）；朱诺入里斯本（11月30日）、布拉干萨出海（11月29日）；哥本哈根（9月）；米兰敕令（12月17日）------\textbf{S2本义的最后一个史实事件}
& 史实 & 史实 \\
\end{longtable}
}

【B】路径B：1803年马耳他仲裁，1805年不进德意志，1806年\literalquote{}窄版邦联\literalquote{}，1807年无提尔西特------奥普俄三国完整、无华沙、丹麦保舰队、芬兰留瑞典。

\subsubsection{三、1808--1812：开关期}\label{ux4e0918081812ux5f00ux5173ux671f}

{\def\LTcaptype{none} % do not increment counter
\begin{longtable}[]{@{}
  >{\raggedright\arraybackslash}p{(\linewidth - 6\tabcolsep) * \real{0.1000}}
  >{\raggedright\arraybackslash}p{(\linewidth - 6\tabcolsep) * \real{0.6100}}
  >{\raggedright\arraybackslash}p{(\linewidth - 6\tabcolsep) * \real{0.2100}}
  >{\raggedright\arraybackslash}p{(\linewidth - 6\tabcolsep) * \real{0.0800}}@{}}
\toprule\noalign{}
\begin{minipage}[b]{\linewidth}\raggedright
年
\end{minipage} & \begin{minipage}[b]{\linewidth}\raggedright
条目
\end{minipage} & \begin{minipage}[b]{\linewidth}\raggedright
锚／机
\end{minipage} & \begin{minipage}[b]{\linewidth}\raggedright
置信
\end{minipage} \\
\midrule\noalign{}
\endhead
\bottomrule\noalign{}
\endlastfoot
1808 &
\textbf{开关一}：不签枫丹白露关税以外的新兼并令；1月起以特里亚农型分档关税替代排他封锁（Collin
de
Sussy／Montalivet的\literalquote{}走私成本国有化\literalquote{}提前两年）；对俄给中立航运与关税容忍（R1唯一六列同击项）。\textbf{开关二}：3月阿兰胡埃斯政变照史实发生（法军在境）；4--5月不召巴约讷：以调停人承认费尔南多七世、完成联姻（1月10日\literalquote{}Consiento
de buen
grado\literalquote{}去条件化）、撤四要塞以外驻军、保留同盟。无马德里五月二日、无总起义（SP2三条件缺一）。9月巴黎公约对普（42,000限军十年期）照签；埃尔福特（10月）以两公国归俄换整体保全奥斯曼（第16条）；西属美洲无1808主权断裂，Yermo政变的Consolidación触发链仍在------墨西哥1808年9月政变可能仍发生但对象是总督而非王朝焦点（MX：断层先于危机）
& F1、X2、R1、SP1、SP2、PR、OT、MX & 中 \\
1809 &
无西班牙泥潭→奥地利1809年再战窗口不开（A卡S5-A：需法军深陷伊比利亚为触发）；\textbf{但}红队要求标注：奥地利在1805年受辱后的复仇动机仍在，若判其再战概率显著上升则M-B3下调------本书取奥地利再战0.25--0.40，若发生则按史实瓦格拉姆--申布伦（含伊利里亚、限军十五万）装配。教廷：不并罗马；\literalquote{}第二个1801\literalquote{}和解------教廷已书面愿意事实关港（1808年春），换取教皇领地主权与授职权。蒂罗尔起义不发生（1809年的引信是奥军入境）。
& A、IT3、G1 & 中 \\
1810 &
不并荷兰（路易条约3月已签，保留）、不并汉萨、不并瓦莱；分层准入＋许可证常态化；奥地利联姻（在普雷斯堡与不受辱的背景下奥方动机为财政与安全，梅特涅的生存策略弱化------本书标联姻概率0.45--0.60，若不成则俄婚受皇太后否决≈不成，罗马王链断，见第六十五章）；西班牙联姻（费尔南多娶皇帝侄女）；美国：无朗布依埃扣船令→马孔第二法案后美法关系不恶化；1810年12月俄国关税照史实颁布，法国不抗议（这就是关税容忍）；瑞典：卡尔·约翰当选（史实由瑞典内政决定）但其向拿破仑求购挪威被拒照旧
& F1、NL、A、US、R1、SC & 中 \\
1811 &
统合公债与法定利率上限（C8，非工具性0.15--0.30------本书标为主线最弱一环，若不做则rentes
7.5--9\%）；不吞奥尔登堡；4月致俄函愿先行裁军（通信17579）真实发出且被兑现；涅谢尔罗迭10月方案（保普鲁士存续、归还奥得要塞）被接受为对俄安排的基础；斯佩兰斯基不被献祭（其被流放的触发是备战）；英国：枢密令1811年辩论、信用压力+53.75bp、卢德------B1b窗口打开
& X1、R1、PR、B2、B5 & 中 \\
1812 & \textbf{不征俄}；C9
0.20--0.45。俄土战争照1806--12史实进程但布加勒斯特和约条款比史实苛刻（无法国威胁则俄不退让------OT）；美英战争：S2下0.30--0.45（强征水手与枢密令是英国单方开战理由）；若发生，法国不介入。英法：欧洲总和约谈判------安特卫普成价码（第五十二章α／β）。法财政余额仅3m（引擎）：签约后必须立即兑现和平红利。
& R1、R3、OT、US、B1b、F4 & 中 \\
\end{longtable}
}

\subsubsection{四、1813--1820：固化期}\label{ux56db18131820ux56faux5316ux671f}

{\def\LTcaptype{none} % do not increment counter
\begin{longtable}[]{@{}
  >{\raggedright\arraybackslash}p{(\linewidth - 6\tabcolsep) * \real{0.1000}}
  >{\raggedright\arraybackslash}p{(\linewidth - 6\tabcolsep) * \real{0.6100}}
  >{\raggedright\arraybackslash}p{(\linewidth - 6\tabcolsep) * \real{0.2100}}
  >{\raggedright\arraybackslash}p{(\linewidth - 6\tabcolsep) * \real{0.0800}}@{}}
\toprule\noalign{}
\begin{minipage}[b]{\linewidth}\raggedright
年
\end{minipage} & \begin{minipage}[b]{\linewidth}\raggedright
条目
\end{minipage} & \begin{minipage}[b]{\linewidth}\raggedright
锚／机
\end{minipage} & \begin{minipage}[b]{\linewidth}\raggedright
置信
\end{minipage} \\
\midrule\noalign{}
\endhead
\bottomrule\noalign{}
\endlastfoot
1813 &
姿态降为S5（需求490）；摄政法（2月5日）照史实颁布（它是1812年危机的产物；在无征俄世界里其触发条件是皇帝亲征西班牙／英伦的可能性，本书标为1813--15年间颁布，中）；英法和约α（安特卫普商业港化＋法保主体舰队＋殖民地互换＋商业准入）或β（无和约、武装共存）；若α，B3的P态深裁开始（6--9万人）；若β，英保W战备
& F1、F5、B1b、B3 & 中 \\
1814 &
无巴黎和约、无维也纳会议、无神圣同盟（R1：其生成条件是击败拿破仑）；莱茵航行自由＋废堆栈权谈判启动（P5：1815--25完成，比史实1831提前6--16年，条件是皇帝肯放弃把河流当武器）；安特卫普港池续建不受阻断
& R1、P5 & 中 \\
1815 &
藩属立宪周期律第一轮（IT2：5--7年）：那不勒斯、威斯特法利亚、荷兰要求代表制与预算权，以Lainé式请愿为形态；法国国内莱内式请愿（体制内）；西班牙：1815--30窗口的宪政压力（SP2：晚到、温和化）；法郎区四国货币公约提前版谈判（P5
M1 0.50--0.65） & IT2、F5、SP2、P5 & 中 \\
1816--17 &
无坦博拉级歉收缓冲；1816年无雨年（外生恒在）------帝国粮政（最高限价与征发传统）应对；1817年英国若α已深裁则海军人力6--9万；瑞典--挪威：挪威留丹麦联合王国联邦化（SC主线中）；葡萄牙：巴西宫廷不回欧洲（PT：巴西入英圈）
& E1、B3、SC、PT & 中低 \\
1818--20 &
藩属立宪周期律第二轮；那不勒斯1820年式军官政变（IT2）由法国扮演史实奥地利的平叛角色；西班牙1820年式革命动力学（欠饷＋美洲战争＋共济会）在无占领世界仍有其二------法西联合宪兵压制；德意志：1818--19南德宪法照史实（G1：内生）；关税：法国许可证行政下放开始（X2三条件之一）
& IT2、SP2、G1、X2 & 中 \\
\end{longtable}
}

\subsubsection{五、1820--1848：继承与秩序}\label{ux4e9418201848ux7ee7ux627fux4e0eux79e9ux5e8f}

{\def\LTcaptype{none} % do not increment counter
\begin{longtable}[]{@{}
  >{\raggedright\arraybackslash}p{(\linewidth - 6\tabcolsep) * \real{0.1000}}
  >{\raggedright\arraybackslash}p{(\linewidth - 6\tabcolsep) * \real{0.6100}}
  >{\raggedright\arraybackslash}p{(\linewidth - 6\tabcolsep) * \real{0.2100}}
  >{\raggedright\arraybackslash}p{(\linewidth - 6\tabcolsep) * \real{0.0800}}@{}}
\toprule\noalign{}
\begin{minipage}[b]{\linewidth}\raggedright
年
\end{minipage} & \begin{minipage}[b]{\linewidth}\raggedright
条目
\end{minipage} & \begin{minipage}[b]{\linewidth}\raggedright
锚／机
\end{minipage} & \begin{minipage}[b]{\linewidth}\raggedright
置信
\end{minipage} \\
\midrule\noalign{}
\endhead
\bottomrule\noalign{}
\endlastfoot
1820s &
皇帝死亡0.40--0.50（F5加权）；若死→幼主摄政（1813年法：母后名义＋不可撤免摄政会）→\literalquote{}宪政化让渡市场\literalquote{}开启（莱内清单→宪章）；交接平稳65--80\%；奥地利对外孙的保护动机启动（A卡：若新君非其子则摄政终止）；X1：继承时刻的债务与补偿池15--30m／年；MX：墨西哥自治王国存续至1820s
0.30--0.45 & F5、F1、A、X1、MX & 中 \\
1823 &
无史实法军入西（B4b世界里西班牙无1820--23自由派三年）；坎宁--波利尼亚克型备忘录以\literalquote{}法国不帮西班牙武力重建美洲统治\literalquote{}为条件的英国不承认（US／B6：门罗式原则提前且对象改为法国）
& SP2、B6、US & 中 \\
1825--30 & 有限版A层制度稳态（P2）；法郎区浅铸币同盟（P5
M1）；莱茵自由航行完成；铁路帝国干线（比利时模式：国家干线＋桥路团＋战略线）启动；X3：Blanc互换性重启的政治窗口（低）；英国：蓝水转向，殖民中央支出£3.4--5.4m（B6）；埃及：阿里对叙利亚（无1840伦敦公约类比→保叙利亚更久）
& P2、P5、X3、B6、EG & 中 \\
1829 &
罗马王成年（3月20日）；拿破仑活到该日≈0.50--0.65；若在世→1805年art.3分冠触发条件（外军撤出那不勒斯／马耳他／科孚------在α和约下可成立）；北中意王国交予皇子的通道打开（IT1
B分支）；1813年法未规定摄政会解散程序------制度空白 & F1、F5、IT1 &
中 \\
1830 &
皇帝1830s死30--35\%；成年新君交接80--90\%；波兰：兵营版会议王国式起义15±7年周期的第一个到期点（PL），善治版不发；那不勒斯／西西里第三轮立宪；法国国内\literalquote{}带盔甲的七月王朝\literalquote{}形态（F5≥70\%行政国家存续）
& F5、PL、IT2 & 中 \\
1832--35 &
萨克森、汉诺威照史实最晚废农奴制（P5反比规律：主权盟邦最浅最晚）；有限关税同盟化（X2
B分支1848视点0.25--0.35）；法兰克福／美因茨清算机构（若B分支）；埃及科尼亚（1832）后无欧洲协调干预→阿里保叙利亚
& P5、X2、EG & 中 \\
1840s &
皇帝1840s后死15--25\%；英属印度1848财政£20--25m、公司纯行政化（IN1）；旁遮普吞并进程照史实（IN2）；美国抵太平洋0.60--0.80、佛罗里达推迟至1825--35（US）；古巴归属西0.40--0.55；铁路帝国网2,500--5,000km（中值3,500）
& IN1、IN2、US、MX、X3 & 中 \\
1845--48 &
歉收外生恒在；体系级复合危机0.6--0.8；同步革命波0.35--0.55；主形态＝巴黎辐射型立宪--面包--继承复合危机；俄国是唯一完整的反体系军事极；秩序存续0.45--0.65、王朝在位0.35--0.50
& X4、E1、R3、F5、P8 & 中 \\
\end{longtable}
}

【B】路径B在1808年后与S2本义合流于大部分开关，但差别恒在：无华沙（波兰1795线）、普鲁士完整无改革（第六十六章）、波旁那不勒斯存续、威尼托留奥、无芬肯施泰因（波斯早入英轨）、无俄土战争、丹麦保舰队与挪威、法属印度洋两岛不失（IN2威胁在场存续）、东圣多明各法属（HT）、无奥婚动机（罗马王链须另证）。这些条目在P7的338条里有装配，但按本书判定属于\textbf{估计域之外的低置信推演}。

\begin{center}\rule{0.5\linewidth}{0.5pt}\end{center}

\subsection{第五十九章　经济秩序：分层准入、白银区、节点基建与三层农村}\label{ux7b2cux4e94ux5341ux4e5dux7ae0-ux7ecfux6d4eux79e9ux5e8fux5206ux5c42ux51c6ux5165ux767dux94f6ux533aux8282ux70b9ux57faux5efaux4e0eux4e09ux5c42ux519cux6751}

\subsubsection{一、法国霸权的商业价值可以称量}\label{ux4e00ux6cd5ux56fdux9738ux6743ux7684ux5546ux4e1aux4ef7ux503cux53efux4ee5ux79f0ux91cf}

【锚】Ellis（Napoleon\textquotesingle s Continental Blockade: The Case
of Alsace，1981，附录G、I，印pp.285--289，源A.N. F12
251）：1810年法国出口377m官方值，其中德意志诸邦143.4m＋意大利王国51.6m＋其他意邦20.5m＋瑞士21.2m＋荷兰44.6m＝281.3m（\textbf{74.6\%}）流向体系内大陆市场；1806年为58.0\%。1817年------复辟、体系解体、和平------这一比重跌到\textbf{33.8\%}（156.7/464m）；德意志诸邦一项143.4→56.2m（−61\%），意大利王国51.6→15.1m（−71\%）。\literalquote{}欧洲行省化\literalquote{}在法国资产负债表上不是抽象荣耀，而是每年约一亿二千万法郎官方值的制度性需求------约1806年出口总额的27\%、约Chaptal所计制造业厂主利润总额182m的三分之二。\textbf{这是全书对\literalquote{}法国为什么要控制欧洲\literalquote{}的最硬的经济回答。}

同一份材料立即给出反面：1817年出口总额464m几乎回到1806年的465m------它是靠把货卖到别处（美洲、黎凡特、附录I十五个目的地以外的30.5\%）补回来的。\textbf{帝国市场与海洋市场对法国是替代品。}于是全部情景归约为一个结构：S3只给法国帝国市场；S1／S2给帝国市场＋部分海洋市场；S5两者兼得；S6把帝国市场扩到纸面最大而把海洋市场与税基一并烧掉。

【锚】Chaptal《De l\textquotesingle industrie
françoise》（1819）制造业毛产值1,820,102,409法郎的完整分解：本国原料416m＋进口原料186m＋工资844m＋一般费用192m＋厂主利润182,005,221------利润份额10.00\%，固定资本消耗＋利息10.55\%，工资46.37\%，进口原料10.22\%（马克思式口径s/v≈21.6\%、s/(c+v)≈11.1\%、c/v≈0.94）。\textbf{保护关税能抬高那10\%，抬不出一个能自筹铁路与焦炭高炉的资本集团。}

\subsubsection{\texorpdfstring{二、1848年的账（条件包络，\texttt{P5\_growth.csv}可复算）}{二、1848年的账（条件包络，P5\_growth.csv可复算）}}\label{ux4e8c1848ux5e74ux7684ux8d26ux6761ux4ef6ux5305ux7edcp5_growth.csvux53efux590dux7b97}

1848年人均工业化指数（Bairoch，UK1900＝100）：

{\def\LTcaptype{none} % do not increment counter
\begin{longtable}[]{@{}lrrrrr@{}}
\toprule\noalign{}
区块 & S0 & S2 & S3 & S5 & S6 \\
\midrule\noalign{}
\endhead
\bottomrule\noalign{}
\endlastfoot
旧法国 & 16.8 & 19.8--28.1 & 12.1--17.9 & 18.6--28.1 & 13.3--18.2 \\
内圈并合省 & 22.4 & 23.6--30.9 & 18.4--23.1 & 23.6--29.4 & 20.6--25.1 \\
德意志 & 12.6 & 11.3--15.5 & 10.9--14.1 & 12.0--15.0 & 10.2--12.4 \\
意大利 & 9.2 & 8.9--11.5 & 7.9--9.8 & 8.9--11.2 & 7.7--9.3 \\
英国（承B5） & 48.4 & 38.2--48.4 & 26.1--38.2 & 42.6--52.8 &
36.3--48.4 \\
\end{longtable}
}

大陆六区工业总量与英国之比：S0 1.17；S2 1.55--1.67；S3
1.59--1.74（大陆仅80--107\%史实）；S5
1.37--1.51（大陆103--140\%、英国88--109\%）；S6
1.22--1.27（大陆81--104\%）。\textbf{S3拿到最高的比值却拿到最低的总量}------它靠打断英国取胜；\textbf{S5是唯一大陆与英国同时高于史实的格}（正和霸权）；\textbf{S6在每一本经济账上都被S5支配}，且区域离散度最高（2.69）。制造业利润份额S3
12--15\%＝把毛产值2--5个百分点（约3,600--9,100万法郎／年）由消费者转给受保护厂商------这是\literalquote{}封锁红利\literalquote{}的真实量级与真实收款人。

\subsubsection{三、货币：一个白银区，命运绑在马德里}\label{ux4e09ux8d27ux5e01ux4e00ux4e2aux767dux94f6ux533aux547dux8fd0ux7ed1ux5728ux9a6cux5fb7ux91cc}

【锚】1803年法国法定金银比价1:15.5；1806年3月21日法令把法国币制延伸至意大利王国（5g、900‰、百分制）；威斯特法利亚1808年1月11日折价1塔勒＝3.885法郎，排他采用推迟到邻邦一并采用（从未实现）；皮埃蒙特1816、比利时1832、瑞士1850、意大利1862单方采用法国10法郎金币为法偿；\textbf{1865年12月23日四国货币公约}------法国外交部1867年照会自述它\literalquote{}重建了事实上曾存在于这四国之间\literalquote{}的货币同盟，涉及人口七千万（FRUS
1867
doc.291）。LMU的边界：仅金币与5法郎银币＋辅币，不含铜币、纸币与存款，无共同机构、无共同央行。

【推·中高】法郎区的形成不需要法国军事胜利；胜利改变时点（1865→1815--30）、边界（扩到德意志西南，北德塔勒区显著低）而\textbf{改不了深度}------拿破仑因assignat记忆拒绝一切统合公债与纸币（Buist印pp.289--290），中枢容量3,000--4,200封／年。三形态：M1浅铸币同盟0.50--0.65（1830）／0.45--0.60（1848）；M2含发行机构的深同盟0.05--0.12／0.10--0.20（继承后由财政官僚推动）；M3三分裂0.25--0.40。双金属制的Gresham效应使1803--48法郎区在流通上以白银为主------\textbf{S5的货币史与西班牙--美洲史是同一条线}：B4b＋美洲不脱离→白银流量维持2,000--3,000万pesos／年量级、通缩压力小于史实；废立＋美洲独立→1820--40年代白银短缺同构或更糟。\textbf{命题P5-M：S5若同时选\literalquote{}兼并西班牙\literalquote{}与\literalquote{}法郎区双金属制\literalquote{}，就在1820年代给自己制造通货紧缩}------\literalquote{}最大行省化≠最优秩序\literalquote{}在货币账上的第四次复现。

\subsubsection{四、基建：瓶颈解除事件，不是里程}\label{ux56dbux57faux5efaux74f6ux9888ux89e3ux9664ux4e8bux4ef6ux4e0dux662fux91ccux7a0b}

【锚】1804年8月15日巴黎《莱茵河航行octroi公约》132条（第2条：莱茵河在航行与商业上永远被视为两帝国的共同河流）；1808年10月27日敕令设Magistrat
du
Rhin（CCNR直系祖先）；\textbf{1810年拿破仑削减了自己的作品------关闭奈梅亨以上航行并保留给法国船舶}；科隆商会捍卫堆栈权；史实堆栈权至1831年美因茨公约才取消。Ellis附录D（斯特拉斯堡商会1812年11月10日统计表第10号）：斯特拉斯堡全年河运1809年26,051公吨、1810年23,584、1811年26,557------日均65--73吨、约两三条驳船；对照敖德萨单港小麦出口1815年171,708吨------\textbf{一个黑海港的单一品种年出口量是莱茵河上阿尔萨斯首府全部河运贸易的六倍强}。帝国内陆水运是毛细血管不是大动脉；这同时解释海上贸易为何不可替代，以及X2的55\%走私保费为何足以排除英国工业品（需被排除的实物量本来就不大）。

【锚】安特卫普：Bonapartedok 1807开挖、1811年1月1日进船；Willemdok
1812通航；其余岸壁1815--37才建；1839年比荷条约第9条设斯海尔德航行税，1863年比利时以17,141,640荷兰盾赎买。【推·中高】S5下1839年式河口税法理上不可能，港池续建不受阻断→安特卫普1815--48成大陆最大北海商港，替代阿姆斯特丹--鹿特丹提前30--50年；\textbf{硬约束}：其商港价值是制海权与对英贸易的函数（1808--13主权在手港池建成而海运照样锐减）；\textbf{利益冲突}：要安特卫普的港还是要阿姆斯特丹的资本市场------帝国会选港，代价是大陆最深的私人资本池已在1810年被自己摧毁，此后铁路与重工业融资只能靠国家预算。

【锚】辛普朗1801--05（全年通行所需隧道被巴黎否决，收容院1831年才成）；圣康坦运河1810年投用（北方船与瓦兹河船尺度不同须换船，补水渠1826）；1811年12月16日第7644号法令把帝国道路分三等229条，目的是\literalquote{}réduire
la part du
Trésor\literalquote{}。\textbf{帝国工程成果是节点级的}，网络效应被船型不统一、配套延后一代人、维护费推给地方三件事截断；S5的基建红利只能按\literalquote{}瓶颈解除事件\literalquote{}（河流法权、船闸标准、河口主权、堆栈权）计算。

【锚】铁路：比利时1834年5月1日法令建国有铁路系统（1843年556km）；法国至1841年底仅319km营运；1842年6月11日大干线法确立国家下部＋公司上部的二分。【推·中】S5帝国铁路1848年2,500--5,000km（中3,500）≈史实法＋比之和；按25--40万法郎／km（量级假设），3,500km需8.75--14亿法郎分十五年＝年5,800--9,300万------恰在F4中心案可腾挪与不可腾挪的边界上，\textbf{与海军扩建竞争同一笔钱}。形态更接近\literalquote{}国家干线＋国家经营\literalquote{}的比利时模型：建设速度由财政部限定，铁路首先是军事--行政资产（战略线优先，路线＝Chappe光报线重演），大陆不会出现史实1840年代由铁路股票催生的大众投机与股份公司法演进。

\subsubsection{五、三层农村：改革深度与军事价值成反比}\label{ux4e94ux4e09ux5c42ux519cux6751ux6539ux9769ux6df1ux5ea6ux4e0eux519bux4e8bux4ef7ux503cux6210ux53cdux6bd4}

【锚】德意志各邦废除农奴制年份：奥地利1781/82；巴登1783；普鲁士1807；巴伐利亚、威斯特法利亚1808；符腾堡1817；\textbf{萨克森1832、汉诺威1833}（最后两个）。普鲁士1811年9月14日《调整敕令》§12补偿须割让三分之一耕地、§37非世袭佃农领主可并入一半；1816年5月29日《声明》把适用范围限于具备全部特征的\literalquote{}农民持有地\literalquote{}------\literalquote{}对农民不利地改变了1811年条款\literalquote{}（GHDI编者按）。

【推·中高，命题P5-A】法国直接统治的地方（比利时、莱茵左岸、皮埃蒙特）封建义务无补偿废除、农民保留全部土地；法国需要其军队与财政的主权盟邦（萨克森、巴伐利亚）改革最浅最晚------因为提供兵员、现金与过境权的只能是既存统治阶级，触动其土地权就是触动供给方（与G2的邦联宪法流产主因、G1的主权底线、D1的显贵交易、Kopsidis--Bromley的\literalquote{}拿破仑本人推动的领主制反革命\literalquote{}独立同向）。\textbf{S5到1848年的德意志农村是三层并存：莱茵--比利时自耕农带、客户王国\literalquote{}法令自由＋实体领主制\literalquote{}带、主权盟邦\literalquote{}旧领主制＋延后赎买\literalquote{}带。}\literalquote{}法国胜利＝欧洲农民得地\literalquote{}不成立。无封建义务的农民占大陆农业人口：S0
30--45\%，S5 40--55\%（裁量区间）。

\textbf{贡赋→税收是一次阶级转移（命题P5-C）}：战时负担落在被占领地的可没收资产，和平期必须落在本国与并合省的土地与消费；四条红线------加直接税（触莱内红线）／加droits
réunis（触大众红线）／削军（触军队红线）／借债（触拿破仑本人的assignat红线）------F4的495m平衡案实际选了\literalquote{}削军＋不借债\literalquote{}。S5秩序的阶级基础：显贵得低直接税与产权安全；农民得无封建义务但承担征兵与间接税；工业资产阶级得保护关税租金；军队得缩编后的职位与赠产；\textbf{断裂点是继承时刻}。

\textbf{马克思式读法（P5，中）}：S5的大陆经济是\literalquote{}两种剥削形式的拼接体\literalquote{}------核心区已完成向市场中介的剩余占有（利润率仅11\%），外围区保留超经济强制的剩余占有；帝国的政治功能正是维持这一拼接，而外围社会关系被冻结正是外围市场狭小、法国工业品出口天花板偏低的原因。\textbf{这是S5秩序内生的增长约束，且不能被\literalquote{}更多的行省化\literalquote{}解决------帝国的兵源结构与它的市场结构互为对方的上限。}

\subsubsection{六、关税同盟：有一个，但不叫关税同盟}\label{ux516dux5173ux7a0eux540cux76dfux6709ux4e00ux4e2aux4f46ux4e0dux53ebux5173ux7a0eux540cux76df}

【锚】Beugnot两三次提出把莱茵邦联发展为关税同盟，Bacher支持，巴黎并未照此行事（Heckscher印pp.295--296）；Keller--Shiue（2014）史实关税同盟使城市对小麦价差下降5.5\%；Heckscher印p.227\literalquote{}一次完税后自由流通\literalquote{}已实施。【推·中高】三条差别：发动机不同（史实由普鲁士1818年关税法与小邦财源需求推动；S5由法国许可证行政下放与继承后政权需要藩属精英财政合作推动------\literalquote{}推动它的是财政摩擦而非理念，阻止它的是一个会死的人\literalquote{}）；中心不同（史实柏林且排除奥地利；S5只能是巴黎或以美因茨／法兰克福为清算地的邦联机构------残普不自动取得1834式领导权）；内容不同（史实对外统一税率＋内部自由；S5的B分支是消费税化＋一次完税自由流通，对外保护留给法国单方决定）。X2形态带：A法国优先分层准入0.50--0.60（1830）／0.35--0.45（1848）；B有限关税同盟0.15--0.25／0.25--0.35。

\subsubsection{七、英国领先的含义变了}\label{ux4e03ux82f1ux56fdux9886ux5148ux7684ux542bux4e49ux53d8ux4e86}

人均层：英国1848人均工业化42.6--52.8（S5）对法国18.6--28.1、内圈23.6--29.4------领先倍数从史实2.88收窄到1.5--2.8，从未消失。总量层：大陆六区总量在史实中已是英国的1.17倍------\literalquote{}大陆总和赶上英国\literalquote{}在史实中已发生，不是法国胜利的成果。前沿层：不变（X3：发明前沿一切情景留英）。\textbf{工业追赶不自动转化为海军追赶}（F3：约束在船厂地理、木材通道、合格海员与预算，不在钢铁产量）。1830--48年经济政策的最可能执行者不是一个波拿巴，而是显贵--高级官僚--银行家构成的、有预算权议会的行政国家------\literalquote{}带盔甲的七月王朝\literalquote{}；改判证据：若F1的0.7冲动在和平期同样成立，则皇帝在世期间政策波动远高于本推演，S5前半段不可按七月王朝写。

\begin{center}\rule{0.5\linewidth}{0.5pt}\end{center}

\subsection{第六十章　社会、思想、教会与继承}\label{ux7b2cux516dux5341ux7ae0-ux793eux4f1aux601dux60f3ux6559ux4f1aux4e0eux7ee7ux627f}

\subsubsection{一、支持是条件机器，和平年它运转得最好}\label{ux4e00ux652fux6301ux662fux6761ux4ef6ux673aux5668ux548cux5e73ux5e74ux5b83ux8fd0ux8f6cux5f97ux6700ux597d}

F5的五个条件（财产、和平、低征额、教俗均衡、面包）在S5前四条同时成立的年份多于史实任何一年。年征≤8--10万（人口≤0.3\%）＋役期兑现＋教会合作→常备军55--70万（含外籍）可无限期维持（锚：1810--11合规峰、1818年古维翁-圣西尔法四万／年、亚眠后三十万常备军先例）。\textbf{但}F2／P2的主线需要法国在册40--45万＋盟军20--25万，若法国在册升到50万且役期六年、退出2.5\%，维持式补充≈10.6万／年------高于F5容忍带上沿。两卡的口径未合，本书以此为S5人力账的第一个未闭合项。

【锚】马莱1812年伪元老院令：\literalquote{}de traiter immédiatement de la paix avec
les Puissances belligérantes, de faire cesser les malheurs de
l\textquotesingle Espagne\literalquote{}------反对派的最大公约数是和平、减征、教俗和解。S5给了前两项；第三项取决于教廷开关。

\subsubsection{二、无胜利帝国的合法性经济（红队j项）}\label{ux4e8cux65e0ux80dcux5229ux5e1dux56fdux7684ux5408ux6cd5ux6027ux7ecfux6d4eux7ea2ux961fjux9879}

红队指出：耐久模型（P6／P8）只计\literalquote{}和平红利\literalquote{}，未计一个不再打胜仗的帝国对军官团、显贵与王朝合法性的侵蚀。本书把它作为独立侵蚀项讨论，给方向与幅度区间，不改概率主表。

【锚】\literalquote{}moi, je ne le puis pas, parce que je suis un soldat
parvenu\literalquote{}（1813年6月26日，梅特涅追述，回溯性）；1814年4月的倒戈是五条件合取（外军占首都＋可信替代王朝＋军事灾难＋皇帝不在场＋中枢家族撤离），S5和平继承时刻前三条不成立（F5）；显贵91\%的国家任职跨政权延续（Petiteau）------制度忠诚不等于王朝忠诚；赠产（Domaine）1810年初年收益\textgreater1,800万、受赠4,035人，1814年登记≈2,473万（F4）；军官团在无战争时失去晋升、战利品与赠产的三条通道。

【机】类比（权重知识，未在单位卡内核）：史实1815--30复辟军官团（半薪军官demi-soldes）的密谋与1830年七月革命中军队的中立；史实1830--48七月王朝以阿尔及利亚为军官团的出口；普鲁士1815--48和平期的军官团保守化而非政变化。适用条件：军官团的出口（殖民战争／海军扩建／宪兵与行政职位）是否存在。

【推·中】在S5，军官团合法性侵蚀的形态是\textbf{F1的0.7冲动的社会基础}------它不是皇帝的个人性格，而是一个军人帝国对新光荣的结构性需求；出口是殖民（阿尔及尔式远征）、海军扩建（吸收陆军军官进海军行政）、宪兵与省制扩编。若三个出口都被财政关闭（F4中心案陆海合计425m、无余量），侵蚀幅度按史实demi-soldes计：1815--20年密谋频率上升、但无外军配合则不致命（F5保王派实测）。对S5存续概率的影响：−3至−8个百分点（本书裁量，非模型），集中在1815--25年而非继承时刻；继承时刻的真正风险仍是F5的\literalquote{}宪政化让渡市场\literalquote{}与旁支问题。改判证据：1810--13年军官团对和平前景的书面态度（警务日报中对军官言论的记录）；1802--03亚眠和平期军官团的密谋率。

\subsubsection{三、教会}\label{ux4e09ux6559ux4f1a}

IT3的三形态：7.1\literalquote{}第二个1801\literalquote{}和解均衡（中高）------构件全部历史存在过；7.2武装加利刚（中）------可运转十年，二十年尺度双轨化；7.3教廷部门化（低）------行政存在、忠诚为零。教会是A层可扩性的第一证伪器，S6下不是引爆器而是腐蚀剂。\textbf{和解是\literalquote{}长久\literalquote{}判据下性价比最高的单项投资}------教会只要罗马与授职权，比显贵年金、军队胜利、市场信用都便宜。S5与教会和解的合取使比利时的专约红利（1801--09）成为常态：多数教士宣誓、讲坛劝征兵。1815--48天主教复兴的三引擎在S5照常运转但教廷自己会谴责\literalquote{}上帝与自由\literalquote{}路线（《Mirari
Vos》全文亲核）；Broers：民众虔诚经拿破仑迫害反而强化------S5若维持教俗决裂，等于把十九世纪最强的自我重塑型反对机器养在体制外。

\subsubsection{四、思想的死穴排序}\label{ux56dbux601dux60f3ux7684ux6b7bux7a74ux6392ux5e8f}

X4：法制欧洲意识形态死穴排序＝\textbf{天主教正统诉讼＞继承合法性真空＞民族主义＞自由主义＞社会主义}。无1813年，德意志政治民族主义降档慢行15--25年（文化民族主义不减速），载体从战争记忆换为关税--征兵申诉；群众化在任何情景都需国家级动员事件点火；挪威机制（封锁＋饥荒＋调度失灵→事实自治→分离主义萌芽，无需民族主义宣传）在法制欧洲照常运转，产物是藩属分离主义。自由主义分体制内法典--宪政请愿主流（莱内／附加法模板）与科佩流亡支；意大利民族主义快于德意志（统一的制度容器）；社会主义偏向圣西门技术官僚＋密谋两极。百万级群众签名政治需合法请愿渠道＋廉价印刷＋工业人口三者齐备，S5欧陆缺第一项，同等压力必改道表达。

\subsubsection{五、继承}\label{ux4e94ux7ee7ux627f}

F5的继承树（S5前提）：皇帝1820s死（幼主摄政）------交接平稳65--80\%，存续至1848
30--45\%；1830s死（成年新君）------80--90\%／45--60\%；1840s死------≥85\%／55--70\%；罗马王早逝旁支使存续×0.6--0.75。加权无条件王朝存续至1848≈40--55\%；\literalquote{}行政国家存续而王朝更替\literalquote{}≥70\%。关键机制：摄政期必然开启\literalquote{}宪政化让渡市场\literalquote{}（莱内清单→宪章→附加法，三代同款）；最优设计＝1813年2月5日法（母后名义＋不可撤免的摄政会）＋遗嘱指定联合班子；最大制度空白＝1829年成年时摄政会解散程序在1813年法中无规定。补偿池15--30m／年可覆盖显贵；融资方式（境外资产内地化／本土直接税／droits
réunis加档）决定政治后果。A卡：若新君非皇后之子，她的摄政终止------奥地利对外孙的亲缘利益不能自动转到法国旁支。

\textbf{罗马王的前提}（红队D-2项）：史实1811年3月20日出生链的前提是1810年奥婚，奥婚的史实动机是1809年瓦格拉姆后的奥地利生存策略。S2本义下，奥地利在普雷斯堡之后、无1809再战：联姻动机为财政与安全（不受辱的奥地利仍在提尔西特体系外寻找位置），本书标0.45--0.60；若不成，俄婚受皇太后法定否决（保罗遗诏）且需前置12--18个月管理，≈不成；则拿破仑的继承人链条转向：约瑟芬不育已定，离婚另娶（萨克森、巴伐利亚、符腾堡公主均在1809--10年名单上）------继承人出生年份1811--13，成年1829--31。\textbf{\literalquote{}罗马王\literalquote{}作为同一人物只在奥婚支成立；作为\literalquote{}皇太子\literalquote{}的功能在其他支成立但奥地利的保护动机消失。}

\begin{center}\rule{0.5\linewidth}{0.5pt}\end{center}

\subsection{第六十一章　各国在秩序中的位置}\label{ux7b2cux516dux5341ux4e00ux7ae0-ux5404ux56fdux5728ux79e9ux5e8fux4e2dux7684ux4f4dux7f6e}

一句话给每个单位在S5（S2本义固化后）的位置、它的上限、以及它最可能的断裂方式；置信与改判证据回第一部各章。

{\def\LTcaptype{none} % do not increment counter
\begin{longtable}[]{@{}
  >{\raggedright\arraybackslash}p{(\linewidth - 6\tabcolsep) * \real{0.2500}}
  >{\raggedright\arraybackslash}p{(\linewidth - 6\tabcolsep) * \real{0.2500}}
  >{\raggedright\arraybackslash}p{(\linewidth - 6\tabcolsep) * \real{0.2500}}
  >{\raggedright\arraybackslash}p{(\linewidth - 6\tabcolsep) * \real{0.2500}}@{}}
\toprule\noalign{}
\begin{minipage}[b]{\linewidth}\raggedright
单位
\end{minipage} & \begin{minipage}[b]{\linewidth}\raggedright
位置
\end{minipage} & \begin{minipage}[b]{\linewidth}\raggedright
上限
\end{minipage} & \begin{minipage}[b]{\linewidth}\raggedright
断裂方式
\end{minipage} \\
\midrule\noalign{}
\endhead
\bottomrule\noalign{}
\endlastfoot
法国 &
核心；税收国家＋缩编军＋分层准入租金；\literalquote{}带盔甲的七月王朝\literalquote{}（1830后） &
陆海合计425--495m；治理密度不足以支撑S6 &
继承时刻的宪政化让渡；0.7冲动；军官团出口关闭 \\
内圈并合省 & A层；体系内唯一英式煤铁区位；征兵照额8--11年 &
认同不交割，制度留而主权可送走 & 教会拉断（若罗马开关关闭）；外军入境 \\
意大利王国 & B层分治藩属（在位期）→继承后北中意统一王国（0.35--0.45） &
经济附庸化（\literalquote{}la France avant tout\literalquote{}）；军旗与都城租金 &
继承时刻藩属王冠再分配；雨伞日式精英四裂 \\
那不勒斯 & B层＋海军义务承包＋封锁豁免区；净贡献≈0 & 低需索档位可持续 &
藩属立宪周期律（5--7年）；法国扮1821奥地利 \\
教廷 & D层＋专约（第二个1801） & 事实关港可换；主权与盟约形式不可换 &
继承时刻的正统性否决（若被囚持续） \\
南德三邦 & B／C层支柱，零驻军成本 &
四条件（担保独占、血税≤阈值、主权底线、经济通道） &
替代担保人出现∧血税超限；关税墙不降→墙外自组织 \\
北中德 & 顺从带非支撑带；萨克森＝体系健康指标 & 法制西北 vs
旧制东北双结构 & 野战军失败或撤出（外生） \\
荷兰 & B层王国（不并），改善通商 & 可镇压、不可消化 &
征兵＋削债＋软海关三者任一 \\
瑞士 & 调停制，16,000条约兵 & 不必并也不能并 &
兼并（唯一断裂方式）；提契诺式政策漂移 \\
华沙公国 & B层，两段式（不称王国） &
稳态峰值人口2.2--2.6\%；无外源现金两年即崩 &
复国期票失信×贸易现金流断→归零；兵营版15±7年起义 \\
普鲁士（残普） & C层；42,000限军十年期后\literalquote{}rentre dans le droit commun\literalquote{} &
1848年620--810万人、常备5.6--9.7万 &
俄国保护价上升；1818年起限军到期后的再武装 \\
奥地利 & C层不平等伙伴；东南转向（多瑙交通、商业海口） &
常年16--23万军；多战区并行能力 &
1805受辱的复仇动机；继承时刻若新君非其外孙 \\
俄国 & C层被胁迫同盟→有条件冷和平（半衰期8--12年） &
自主、受约束、区域可畏、全欧不可持续 &
关税容忍撤回；波兰王国名号；奥尔登堡式吞并 \\
西班牙 & C层同盟（波旁＋联姻） &
25,000远征军＋船厂＋六百万法郎／月；美洲合法性焦点 &
王朝内爆（宫廷政变、继承争端）；财政破产；教会协约 \\
葡萄牙 & D层付费中立（月一百万里弗） & 壳可占瓤不可得 & 入侵（唯一） \\
丹麦--挪威 & C层；挪威留联合王国联邦化 & 财政寿命5--6年（若严封锁） &
S3加严；瑞典购买 \\
瑞典 & D层\literalquote{}假战争\literalquote{}双面 & 可购买，价码在丹麦手里 & 无 \\
英国 & α和约或β武装共存；蓝水转向 & 海上优势收窄不消失；至1848竞争共存 &
法国同时满足复战三门槛（威胁／伙伴／供款） \\
爱尔兰 & 英治＋提前解放（1806--13） & 奥康奈尔式群众政治提前15--20年 &
饥荒外生（0.7--1.0M） \\
奥斯曼 & C层条约国，缩小未瓜分；单一保护 & 每次被强制就掉头 &
集权×耶尼切里×埃及三进程 \\
埃及 & 不可收编的乘数；B级藩属天花板 & 保叙利亚更久（无1840类比） &
内部汲取极限与继承结构→1848过载回调 \\
波斯走廊 & \literalquote{}永久的加尔达纳\literalquote{}（若S4）或英轨（若不付费） &
C层盟国＋牵制筹码 & 赫拉特门槛 \\
东印度公司 & 照常扩张，1848财政£20--25m，公司纯行政化 &
欧洲排除放松而非收紧印度约束 & 伦敦信用（不是印度税收） \\
法属印度洋 & 威胁在场收租（英国付永久前沿--舰队溢价） & 军官--军械输出层
& 四闸门永不同开 \\
美国 & 可购买的服务供应商；门罗式原则提前且对象改为法国 &
1848前非第三海军变量 & 扣船（法国自伤） \\
海地 & C层承认／关税交易（0.55--0.80促成签字） & D层＝10--15年撷取窗口 &
马德里决定 \\
新西班牙 & 自治王国存续至1820s 0.30--0.45→1830--45转独立≥0.70 &
白银90--330万pesos／年（S-A） & Consolidación式追索；阿尔马登刀 \\
南美 & 帝国联邦化0.40--0.55（与MX冲突，P8裁≤0.10全保留） &
两变量：合法国王×出海通道 & 英国航运与同盟资格 \\
金融家 & 发行在巴黎、清算在伦敦、白银在汉堡 & rentes
4.0--6.5\%（做统合公债）或7.5--9.0\% & 没收（法国自伤） \\
\end{longtable}
}

\begin{center}\rule{0.5\linewidth}{0.5pt}\end{center}

\subsection{第六十二章　1848}\label{ux7b2cux516dux5341ux4e8cux7ae0-1848}

\subsubsection{一、危机的形态}\label{ux4e00ux5371ux673aux7684ux5f62ux6001}

【锚】歉收外生恒在（E1）；史实1845--47中欧歉收；X4的宪章请愿量级（1842年近350万签名、1848年委员会点验1,975,496）；Jarrige的卢德三因结构；Broers的民众虔诚强化。【机】1848的史实同步机制是巴黎点火＋铁路／电报／报刊传导＋各邦立宪诉求的同型化；适用条件：一个可辐射的中心、一组同型的诉求语言、一个共同的粮食--信贷冲击。

【推·中】S5下体系级复合危机落在1845--52窗口的概率≈0.6--0.8；取欧陆同步革命波形态≈0.35--0.55；主形态＝\textbf{巴黎辐射型的立宪--面包--继承复合危机}，诉求语言是宪章／法典／反征兵而非民族统一；帝国使巴黎成为单一点火点，同步机制强于史实1848。三个内生时钟------皇帝寿限、藩属立宪周期5--7年、单边冲动0.7／18个月------取代外敌成为主要风险。

\subsubsection{二、各处的形态}\label{ux4e8cux5404ux5904ux7684ux5f62ux6001}

\begin{itemize}
\tightlist
\item
  \textbf{法国}：若皇帝已死（1848年前累计≈0.75--0.85），危机落在继承政权头上------\literalquote{}带盔甲的七月王朝\literalquote{}面对1848的方式与史实七月王朝相似：秩序存续0.45--0.65、王朝在位0.35--0.50、\literalquote{}行政国家存续而王朝更替\literalquote{}≥70\%。若皇帝仍在世（79岁，累计概率≈0.15--0.25），F1的0.7冲动与体能折减0.6--0.8并存------本书不给出在世版的1848形态，因其在估计域之外。
\item
  \textbf{英国}：第三次宪章请愿被拒、预期骚乱未爆------以\literalquote{}稳定君主议会资本主义＋强工人压力并存\literalquote{}进入1848；无革命；财政D≈645--1,080m、I/R\textless60\%；若α和约已使海军深裁，1848年的英国对大陆危机的介入能力低于史实（无1848年英国斡旋的余地）。
\item
  \textbf{德意志}：各邦立宪化浪潮＋反霸权税兵总清算运动；\literalquote{}德意志统一\literalquote{}议程因无普鲁士领导核而更弱、因反法情绪而更泛德意志化------两力相消；奥地利有限宪制让步＋匈牙利宪制危机；能否升级为对法战争取决于法国继承政权的凝聚力。
\item
  \textbf{意大利}：分治版＝要宪章＋意大利关税同盟；统一王国版＝要求完全主权\literalquote{}毕业\literalquote{}------均非反法独立战争（前提：王冠已落地＋教会已和解）；若两前提缺失则1814式精英四裂重演。
\item
  \textbf{西班牙}：城市立宪骚动＋北方教权动员双向脉冲被法西联合宪兵压制；终态＝法国轨道内的行政--立宪君主国。
\item
  \textbf{波兰}：善治版无起义；兵营版第一个到期点1830±7已过，1848为第二个。
\item
  \textbf{俄国}：唯一完整的反体系军事极（纸面80万／实有50--65万＋紧急增至110万带）------最大外部收割者或干预者；1849年匈牙利式出兵的对象在S5是法国体系的东缘（华沙公国／加利西亚），条件是法国继承政权当时的凝聚力。
\item
  \textbf{东方与美洲}：奥斯曼是避难所与旁观者；埃及过载回调；美国抵太平洋进程中（0.60--0.80）、得州照失；西属美洲＝自治王国群或已部分独立；巴西仍葡属（无1808开港加速器，PT）；海地独立且付费；法属安的列斯蓄奴（HT：废奴推迟至1855--75）。
\end{itemize}

\subsubsection{三、1848是不是霸权的终点}\label{ux4e091848ux662fux4e0dux662fux9738ux6743ux7684ux7ec8ux70b9}

【推·中】不是自动的终点。史实1848没有推翻任何一个大国的国家机器，只推翻了七月王朝与梅特涅个人；S5的1848更可能推翻的是继承政权的\textbf{形态}（摄政会→宪章议会）而不是法国的大陆优势------除非它与俄国的干预、英国的复战三门槛、藩属的同步分离\textbf{合取}。P8的O1848（秩序在1848仍在）0.17--4.44\%（中点1.27\%）是从1803年算的无条件带；条件于1815年秩序已固化的D\textbar C为0.444--0.611------即\textbf{一个已经固化的法制秩序活过1848的概率约一半}，低于P7的0.45--0.65上沿，两口径并列。改判证据：1815--48任何一年法国继承政权的凝聚力指标（军队对摄政会的服从记录、显贵请愿的规模）；俄国对华沙的态度序列。

\begin{center}\rule{0.5\linewidth}{0.5pt}\end{center}

\section{第五部　分问}\label{ux7b2cux4e94ux90e8-ux5206ux95ee}

\subsection{第六十三章　英法海军差距缩到什么程度}\label{ux7b2cux516dux5341ux4e09ux7ae0-ux82f1ux6cd5ux6d77ux519bux5deeux8dddux7f29ux5230ux4ec0ux4e48ux7a0bux5ea6}

\subsubsection{一、答案的形状}\label{ux4e00ux7b54ux6848ux7684ux5f62ux72b6}

差距会缩小，不会消失；缩小的幅度由四个变量决定------战和状态（W持续海战／P武装和平）、法国海军预算层级（165／195／220--260m）、安特卫普的军港地位、英国是否因法国压力升档。\textbf{其中第三与第四个变量在旧综合里被漏掉了，而它们互相绑定。}

\subsubsection{\texorpdfstring{二、1830年端点（模型非观测，\texttt{P4\_annual.csv}／\texttt{verify\_p4.py}可复算；\texttt{verify\_p4.py}自陈\literalquote{}no
historical calibration
performed\literalquote{}）}{二、1830年端点（模型非观测，P4\_annual.csv／verify\_p4.py可复算；verify\_p4.py自陈\literalquote{}no historical calibration performed\literalquote{}）}}\label{ux4e8c1830ux5e74ux7aefux70b9ux6a21ux578bux975eux89c2ux6d4bp4_annual.csvverify_p4.pyux53efux590dux7b97verify_p4.pyux81eaux9648no-historical-calibration-performed}

{\def\LTcaptype{none} % do not increment counter
\begin{longtable}[]{@{}
  >{\raggedright\arraybackslash}p{(\linewidth - 8\tabcolsep) * \real{0.1579}}
  >{\raggedleft\arraybackslash}p{(\linewidth - 8\tabcolsep) * \real{0.2105}}
  >{\raggedleft\arraybackslash}p{(\linewidth - 8\tabcolsep) * \real{0.2105}}
  >{\raggedleft\arraybackslash}p{(\linewidth - 8\tabcolsep) * \real{0.2105}}
  >{\raggedleft\arraybackslash}p{(\linewidth - 8\tabcolsep) * \real{0.2105}}@{}}
\toprule\noalign{}
\begin{minipage}[b]{\linewidth}\raggedright
条件
\end{minipage} & \begin{minipage}[b]{\linewidth}\raggedleft
法A12
\end{minipage} & \begin{minipage}[b]{\linewidth}\raggedleft
英A12
\end{minipage} & \begin{minipage}[b]{\linewidth}\raggedleft
舰体吨代理比
\end{minipage} & \begin{minipage}[b]{\linewidth}\raggedleft
A12×训练效能比
\end{minipage} \\
\midrule\noalign{}
\endhead
\bottomrule\noalign{}
\endlastfoot
W持续海战，中档 & 84 & 140 & .742 & .449 \\
P武装和平，中档（195m planning floor） & 111 & 74（英I 57.5、A12 120） &
.902 & .876 \\
法P、英维持W战备 & 111 & 140 & .902 & .731 \\
P低档（165m） & 84 & --- & .769 & .747 \\
P高档（220--260m，F3压力带） & 121--140 & --- & 高于.902 & --- \\
\end{longtable}
}

红队纠正了此前综合的误读：\textbf{195m对应P中档111
A12对应.902/.876；165m对应84
A12对应.769/.747}------两组数字各自匹配。真正未决的是P4
L31所述220--260m的层级：它是F3的\literalquote{}晚P压力带\literalquote{}（长期保132艘H及129k人），P4明言\literalquote{}P高端是财政尚未闭合的技术包络，不是获批计划\literalquote{}，且\textbf{不在F4任何平衡案内}。

范围：W舰体.552--.926、效能.282--.644；P舰体.769--1.012、效能.747--.974。敏感性（P4，英方保持W基线）：新下水−30\%效能.321；人员补入−30\%
.305；整备率−10个百分点
.390；\textbf{失安特卫普（初始−12舰、年下水−3）吨.444、效能.267}------这一行是\literalquote{}单一冲击、英国反应保持W基线\literalquote{}，不是P基线的结果。

\subsubsection{三、安特卫普把③和K3绑在一起}\label{ux4e09ux5b89ux7279ux536bux666eux628aux2462ux548ck3ux7ed1ux5728ux4e00ux8d77}

第五十二章已论证：英国\literalquote{}不败而接受\literalquote{}（K3）的价码，按1814年巴黎和约第XV条的先例，很可能包含\literalquote{}Antwerp
shall for the future be solely a commercial
port\literalquote{}。F3的安特卫普贡献：1807--13年五十二艘下水中十八艘；1810年夏帝国控制的50艘战列舰体中安特卫普12、1812年72中20、1814年74中21；1811年3月5日命令\literalquote{}18船台、舰留三年\literalquote{}；橡木流域62\%的合规格大橡木。P4／F3的\literalquote{}失安特卫普\literalquote{}分支给1830年A12主线50（W）／71（P）。

两条终局带（本书作者的联算，量级近似）：

{\def\LTcaptype{none} % do not increment counter
\begin{longtable}[]{@{}
  >{\raggedright\arraybackslash}p{(\linewidth - 10\tabcolsep) * \real{0.1429}}
  >{\raggedright\arraybackslash}p{(\linewidth - 10\tabcolsep) * \real{0.1429}}
  >{\raggedleft\arraybackslash}p{(\linewidth - 10\tabcolsep) * \real{0.1905}}
  >{\raggedleft\arraybackslash}p{(\linewidth - 10\tabcolsep) * \real{0.1905}}
  >{\raggedleft\arraybackslash}p{(\linewidth - 10\tabcolsep) * \real{0.1905}}
  >{\raggedright\arraybackslash}p{(\linewidth - 10\tabcolsep) * \real{0.1429}}@{}}
\toprule\noalign{}
\begin{minipage}[b]{\linewidth}\raggedright
终局
\end{minipage} & \begin{minipage}[b]{\linewidth}\raggedright
条件
\end{minipage} & \begin{minipage}[b]{\linewidth}\raggedleft
1830法A12
\end{minipage} & \begin{minipage}[b]{\linewidth}\raggedleft
舰体比
\end{minipage} & \begin{minipage}[b]{\linewidth}\raggedleft
效能比
\end{minipage} & \begin{minipage}[b]{\linewidth}\raggedright
K3
\end{minipage} \\
\midrule\noalign{}
\endhead
\bottomrule\noalign{}
\endlastfoot
\textbf{α 接受＋差距不缩} &
安特卫普商业港化；法保布雷斯特／土伦／罗什福尔／洛里昂＋荷兰合作船厂（弗利辛恩、阿姆斯特丹）；英P态深裁（I
45--70） & 71 & ≈.58（71/111×.902，比例近似） &
≈.56（同法折算.876；若英保W则≈.27） & 成立（α和约） \\
\textbf{β 不接受＋追赶} &
法保安特卫普军港并扩建；英视之为升档触发器保持W战备（B3：法国60艘若为当季敌对压力，英至少78--97任务舰）
& 111 & .902 & .731 & 不成立（武装共存） \\
\end{longtable}
}

\textbf{因此对分问③的回答分两句}：若英国接受（α），1830年法英战列舰效能比约.55--.60（舰体比约.58），差距\literalquote{}缩到三分之二以下\literalquote{}；若英国不接受而法国追赶（β），效能比约.73（舰体.90），但没有和约。\textbf{旧综合的\literalquote{}.902/.876\literalquote{}是把α的和约与β的舰队拼在一起的结果，本书撤回。}这一联算是比例近似而非P4模型重跑，改判证据：以71舰输入重跑P4的P基线；1806年英法谈判与1813--14卡斯尔雷训令中安特卫普条款的原文；瓦尔赫伦远征训令中对安特卫普船厂的目标定义。

\subsubsection{四、边界与不可比}\label{ux56dbux8fb9ux754cux4e0eux4e0dux53efux6bd4}

\begin{itemize}
\tightlist
\item
  H／I／A12／当季出海不是同一指标；比值不是海战胜率或全球制海权。
\item
  James的tons
  burthen不是法国排水量；\literalquote{}Cruisers\literalquote{}已含Line；Built是毛建成非净增。
\item
  英国低战备需低威胁条件；法国压力升高使英国升档（B3明列）；英1830
  P中档I≈57.5、A12＝120隐含一年召回62.5艘，缺逐舰校准；F3合格人员因子与P4训练因子可能重复折价。
\item
  木材主约束是通道、船厂地理、人员、财政，不是\literalquote{}法国没森林\literalquote{}（E1／F3）。
\item
  1831--48蒸汽、爆炸弹（Paixhans
  1822提案→1824试验）、螺旋桨、铁甲使帆舰吨位线失效------P4不把1830参数机械年推至1848；强法国可在合适港链更快试用炮弹炮，英国有同等工业与召回能力。1848年的蒸汽舰队未数值拟合。
\item
  财政：海军170--220m／年＝主线档；1830年高档需220--260m------与铁路年5,800--9,300万竞争同一笔钱（P5）；陆300＋海195只在高收入＋盟约转移下接近闭合（第五十四章第五节）。
\end{itemize}

\textbf{主线}：法国有成为强大第二海权的路径，没有靠兼并沿海＋多造船自动取代皇家海军的路径；最大行省化≠最大有效海权。

\begin{center}\rule{0.5\linewidth}{0.5pt}\end{center}

\subsection{第六十四章　西班牙的终局}\label{ux7b2cux516dux5341ux56dbux7ae0-ux897fux73edux7259ux7684ux7ec8ux5c40}

\subsubsection{一、前提}\label{ux4e00ux524dux63d0}

S2本义的西班牙是B4b：1808年3--5月不召巴约讷、承认费尔南多七世、联姻去条件化、撤四要塞以外驻军、保留同盟。它的直接后果是\textbf{没有总起义}（SP2三条件缺一：无王位处置公开化）、没有半岛战争、没有加的斯议会、没有1812年宪法------也没有1814年的费尔南多复辟与1820--23年的自由三年。

\subsubsection{二、终局概率带（1848视点）}\label{ux4e8cux7ec8ux5c40ux6982ux7387ux5e261848ux89c6ux70b9}

{\def\LTcaptype{none} % do not increment counter
\begin{longtable}[]{@{}
  >{\raggedright\arraybackslash}p{(\linewidth - 4\tabcolsep) * \real{0.3000}}
  >{\raggedleft\arraybackslash}p{(\linewidth - 4\tabcolsep) * \real{0.4000}}
  >{\raggedright\arraybackslash}p{(\linewidth - 4\tabcolsep) * \real{0.3000}}@{}}
\toprule\noalign{}
\begin{minipage}[b]{\linewidth}\raggedright
终局
\end{minipage} & \begin{minipage}[b]{\linewidth}\raggedleft
S2／S5（B4b）
\end{minipage} & \begin{minipage}[b]{\linewidth}\raggedright
备注
\end{minipage} \\
\midrule\noalign{}
\endhead
\bottomrule\noalign{}
\endlastfoot
波旁绝对主义--大臣改革制（1815--30型延长） & \textbf{0.40--0.55（主线）}
& 财政破产（vales
18.9亿rv、贴水63\%）与美洲战争厌战仍在；大臣改革（Consolidación的延续、教产变现）在王朝存续下推进 \\
波旁立宪（钦定宪章／复召Cortes） & 0.30--0.45 &
宪政晚到（1815--30窗口）、温和化（传统Cortes／两院宪章）、更内战化（缺反法黏合剂）；1820年革命动力学不需要占领变量 \\
波拿巴替换／波拿巴立宪 & 0.05--0.15 &
只在法国改行废立时上升到0.35--0.55（半岛军事全胜前提） \\
共和国 & ≤0.02 &
硬上界≤0.08：需双王朝失信＋无王朝焦点的战争延长＋美洲共和样板回传，史实1868--73才凑齐；1848年前至多凑两项 \\
\end{longtable}
}

\textbf{红队SP1项的处置}：这些带不是互斥且已校准的世界概率表，是给定\literalquote{}保波旁路线及秩序存续\literalquote{}的条件带。费尔南多七世的人格（1807年求婚信、1808年1月的\literalquote{}deshonrado\literalquote{}、1814年撕毁宪法、1823年后的十年）在B4b世界里未经流亡与囚禁------其对联姻与法国监护的反应本身是一个未估的变量。

\subsubsection{三、三根敏感轴与三个断层}\label{ux4e09ux4e09ux6839ux654fux611fux8f74ux4e0eux4e09ux4e2aux65adux5c42}

SP2的总律：西班牙社会对法国政策的敏感轴只有三根------王朝符号、税与征发、教会；三轴齐触则总起义，触一轴则区域抵抗，不触则无感。B4b不触王朝轴；S2的保护性关税不触税轴（对西班牙而言封锁本来就不可执行）；教会轴取决于罗马开关------保留波旁下西班牙教会可挤压可合作（1798--1808教皇简函已合作卖教产）。1815--48三断层恒在：土地／señorío清算（1813年反封建暴动已点燃）、教会财产与修会（1798先例→1834
matanza结构）、中央化vs
foral。任何政权只改变爆发顺序与镇压者身份；管理的单点＝教会协约（堂区司铎薪俸国家化则大众保王主义失去中介神经）与军饷。

\textbf{卡洛斯战争的类比}：大众保王主义是地区结构变量而非全国常数（\literalquote{}The
poorest areas of Spain... were lost causes for Don
Carlos\literalquote{}）；北方三区可养万人级反自由派军队数年，政府需二十万以上军队才能压平。在B4b世界里，1833年式的继承争端（若费尔南多仍无男嗣）照常可能出现------它的第三方是法国：法西联合宪兵是IT2\literalquote{}外军平叛依赖律\literalquote{}在西班牙的版本。

\subsubsection{四、加泰罗尼亚与埃布罗}\label{ux56dbux52a0ux6cf0ux7f57ux5c3cux4e9aux4e0eux57c3ux5e03ux7f57}

B4c（兼并至埃布罗）在史实中已被试验并失败（1810年2月至1812年2月的事实剥离，1812年敕令不敢刊公报，1813年民政自撤）；纳瓦拉榨取1.52亿reales仍征收崩溃（1811--12实收\textless29\%）。\textbf{A层南界＝比利牛斯}。加泰是唯一可能的\literalquote{}合作化兼并区\literalquote{}（棉业市场利益＋反马德里积怨，低中置信）------但这是二十年尺度的B层安排，不是省。

\begin{center}\rule{0.5\linewidth}{0.5pt}\end{center}

\subsection{第六十五章　罗马王与意大利：分治还是统一}\label{ux7b2cux516dux5341ux4e94ux7ae0-ux7f57ux9a6cux738bux4e0eux610fux5927ux5229ux5206ux6cbbux8fd8ux662fux7edfux4e00}

\subsubsection{一、两条路径的前提不同}\label{ux4e00ux4e24ux6761ux8defux5f84ux7684ux524dux63d0ux4e0dux540c}

IT1的三分支概率带（C分治0.55--0.65在位期／B统一0.35--0.45继承后／A行省化0.10--0.20）以史实意大利王国为对象：含威尼托（1806，普雷斯堡）、马尔凯（1808）、特伦蒂诺（1810）；那不勒斯已被征服（1806）并交约瑟夫→缪拉；1805年art.3的分冠触发条件是\literalquote{}外军撤出那不勒斯／爱奥尼亚／马耳他\literalquote{}。红队i\_da49cd2638：这些前提在路径B下不成立。

\textbf{S2本义}：普雷斯堡与那不勒斯征服是史实段，王国的疆域与藩属结构与史实同；IT1的带直接适用。art.3的触发条件在α和约下部分成立（那不勒斯无外军------缪拉体制；爱奥尼亚在1807年提尔西特归法；马耳他------若α和约按亚眠第X条式安排为中立多国保证则成立，若英国保留则不成立）。\textbf{因此统一分支的法理通道以对英终局为前提}：α和约含马耳他中立化→art.3自动执行→王冠移交合法男嗣并永久分冠（art.4）；β武装共存→art.3不触发→王冠转移需新的元老院令（皇帝仍可以1810年式重排完成，但失去\literalquote{}自动条款\literalquote{}的正统性外壳）。

\textbf{路径B}：威尼托留奥、那不勒斯波旁存续、无马尔凯（教廷未被并）------王国≈1805年的意大利共和国疆域（约360--400万人口）；art.3的\literalquote{}那不勒斯外军\literalquote{}指英军在西西里与那不勒斯的存在------不成立；北中意统一的对象缩小，且缺继承后\literalquote{}托斯卡纳／罗马转授意大利王冠\literalquote{}的礼制解法（罗马未并）。本书给路径B的统一带0.20--0.30（继承后），行省化0.05--0.10，分治0.55--0.70。

\subsubsection{二、罗马王本人}\label{ux4e8cux7f57ux9a6cux738bux672cux4eba}

第六十章已论：奥婚在S2本义下0.45--0.60。若成，罗马王1811年生、1829年成年、拿破仑活到该日≈0.50--0.65；若不成，继承人为另一位公主之子，\literalquote{}罗马王\literalquote{}称号仍可授（1810年法令第7条是称号而非血统），但奥地利的保护动机消失（A卡：新君非其外孙则皇后摄政终止）。

\textbf{拿破仑是否愿意让罗马王统一中南意大利？}同期文书的答案是：他把意大利王冠腾出给亲生子（1810年1月3日致阿尔迪尼），把两冠永分写进宪制（1805年art.4），把罗马加冕写进法令（1810年art.10）------\textbf{北中意王国归皇子是设计，不是恩典}；那不勒斯的并入没有同期文书，仅存回溯层（\literalquote{}次子\literalquote{}摄政方案，Corr.
XXX
p.549）。IT2：那不勒斯直接省制净负、藩属尚可；含那不勒斯的全统一需缪拉线被移除（＋王朝条件），带0.10--0.18内含于B分支。IT1的判语：\literalquote{}分治→统一是顺流，分治→行省化是逆流。\literalquote{}

\subsubsection{三、1815--48的校准器}\label{ux4e09181548ux7684ux6821ux51c6ux5668}

奥治伦巴第--威尼西亚：三轮起义（1821／1831／1848）＋1848年30--35k伦-威籍奥军过半仍忠＋行政机器全程运转＋镇压需七万驻军。征服者直辖\literalquote{}可运行、可镇压、不可赢得\literalquote{}。S5意大利的1848形态：分治版要宪章＋意大利关税同盟，统一王国版要求完全主权\literalquote{}毕业\literalquote{}------均非反法独立战争，前提是王冠已落地＋教会已和解；若两前提缺失则1814年米兰\literalquote{}雨伞日\literalquote{}式的精英四裂重演（旧贵族煽动、普里纳被私刑、欧仁候选破产）。

\textbf{主线}：在位期分治藩属（0.55--0.65）；继承后北中意统一王国（S2本义0.35--0.45；路径B
0.20--0.30）；含那不勒斯0.10--0.18；行省化\literalquote{}第二法国\literalquote{}是惩罚态非设计态（0.10--0.20）。改判证据：Du
Casse十卷中1813--14是否正式许欧仁世袭；Rath
1941；α和约中马耳他条款的实际形态。

\begin{center}\rule{0.5\linewidth}{0.5pt}\end{center}

\subsection{第六十六章　普鲁士与奥地利：谁更强}\label{ux7b2cux516dux5341ux516dux7ae0-ux666eux9c81ux58ebux4e0eux5965ux5730ux5229ux8c01ux66f4ux5f3a}

\subsubsection{一、S2本义：残普，奥强}\label{ux4e00s2ux672cux4e49ux6b8bux666eux5965ux5f3a}

提尔西特疆界＋西里西亚完整、无1815年西扩、42,000限军十年期（1809--1818）后\literalquote{}rentre
dans le droit
commun\literalquote{}、三要塞法军10,000由普方供宿粮、涅谢尔罗迭1811年方案以普鲁士存续为最本质条款。1848年残普约620--810万人口，常备5.6--9.7万、危机总动员12.4--21.8万、各战区野战合计8.0--15.2万（显式假设算术）；不自动取得1834年式关税领导权（依赖1815年西部领土与1831年黑森加入）。奥地利：普雷斯堡后失威尼西亚与蒂罗尔，但无1809再战（S5-A）则无伊利里亚割让、无1811年破产；常年16--23万军；多战区并行能力与持续养兵的财政深度。\textbf{\literalquote{}奥强于普\literalquote{}的实质不在单场会战兵力，而在多战区并行与财政深度}------单一战线上A卡的9--14万与残普的8.0--15.2万重叠。本书给S2本义\literalquote{}奥地利强于普鲁士（1848视点，军事--财政综合）\literalquote{}\textbf{0.70--0.85}；改革：普鲁士的Stein--Hardenberg改革以1806年败战为触发，S2本义中它发生了（1807--11）；但Landwehr传统（1813）不发生，普鲁士军队是\literalquote{}改革了的旧军\literalquote{}。

\subsubsection{二、路径B：完整普鲁士，未估的情形（本书作者的估计）}\label{ux4e8cux8defux5f84bux5b8cux6574ux666eux9c81ux58ebux672aux4f30ux7684ux60c5ux5f62ux672cux4e66ux4f5cux8005ux7684ux4f30ux8ba1}

路径B下普鲁士完整（1806年疆域约990万人口，含西里西亚与波兰份额）、未败、无改革（无Stein--Hardenberg、无1807年农奴解放、无1808年城市条例、无Landwehr）、无华沙公国在侧、无法军驻要塞、无限军。P7给它\literalquote{}北德邦联式安排的主持权\literalquote{}------红队正确指出这无锚点：1806年普鲁士自己提出的\literalquote{}北德联邦\literalquote{}方案因法国否决与耶拿而死，路径B下法国是否允许它，由法国决定，P6／P7未讨论。

估计（本书裁量，中低置信）： -
数量：1806年普军纸面约25万（含预备），实际野战约14--16万（权重知识的通说量级，未在单位卡内核------列为改判证据）；无改革则Krümper制不存在，1848年常备8--12万、危机动员16--24万（旧制，无Landwehr）。数量上与奥地利16--23万相近甚至略高（单一战区）。
-
质量：1806式旧军制（Grab、Clausewitz对耶拿前普军的判词）；无1813年动员传统；军官团贵族垄断未破。
-
财政：无1806年赔款与1808年削减，1806年债53.5m塔勒（Murau）不膨胀到1820年的180m；但也无关税同盟领导权（1818年关税法的动机------分散领土的内部关税------在完整但未西扩的普鲁士中较弱）。
-
政治：完整普鲁士是1805年第三次同盟的潜在成员（史实1805年12月波茨坦条约后普鲁士已动员），路径B下它未被击败也未被拉拢，是奥俄潜在盟友------\textbf{这正是路径B\literalquote{}三个完整东欧大国的不动≠受约束\literalquote{}的核心。}

\textbf{本书给路径B\literalquote{}奥地利强于普鲁士\literalquote{}0.45--0.60}------比S2本义低20--25个百分点，因为完整普鲁士在数量与财政上更强，但在质量上更弱；\literalquote{}普鲁士领导德意志\literalquote{}在两条路径都不成立（奥地利不恢复对全德指挥，普鲁士无关税同盟，法国控制西部）。改判证据：法国在路径B下对普鲁士北德方案的态度文书（1806年8--9月）；1806年普军实际野战兵力的现代研究；Murau的债务序列。

\subsubsection{三、共同的判定}\label{ux4e09ux5171ux540cux7684ux5224ux5b9a}

\begin{itemize}
\tightlist
\item
  取消普鲁士（B3v2）在两条路径都不是主线：S2本义下它1807夏--1808年2月有对俄价格、1811年后几近等于对俄战争；路径B下没有取消它的战争。
\item
  普奥关系的形态：S2本义＝法国以残普为对俄缓冲、以奥为东南伙伴，两者互不领导德意志；路径B＝三个完整大国的均势，法国的西德霸权只靠邦联与南德四条件维持。
\item
  1848：德意志统一议程因无普鲁士领导核而更弱、因反法情绪而更泛德意志化------两力相消（P7）。
\end{itemize}

\begin{center}\rule{0.5\linewidth}{0.5pt}\end{center}

\subsection{第六十七章　美洲：法国能否助西班牙保住美洲}\label{ux7b2cux516dux5341ux4e03ux7ae0-ux7f8eux6d32ux6cd5ux56fdux80fdux5426ux52a9ux897fux73edux7259ux4fddux4f4fux7f8eux6d32}

\subsubsection{一、能，且唯一路径是通过马德里}\label{ux4e00ux80fdux4e14ux552fux4e00ux8defux5f84ux662fux901aux8fc7ux9a6cux5fb7ux91cc}

法国\literalquote{}助西班牙保美洲\literalquote{}的唯一可行版本是B4b（保波旁）＋B6b（制度化中立船／自由港）：美洲合法性焦点连续→无1810年主权断裂→juntas化推迟5--15年；法国实得白银90--330万pesos／年（中枢190万，四因子连乘：汲取额×送达率×法国分成×交付折扣）＋市场准入，而非领土；派兵版本全部不可行（实测投送上界＝两名专使＋一条双桅船，1808年7月）。

\textbf{红队I项的调和}：SA的\literalquote{}400--800万pesos／年\literalquote{}是毛口径的白银流量上界（未乘送达率与法国分成），MX的\literalquote{}90--330万\literalquote{}是法国净得。两者不可相加也不冲突------本书以MX四因子为承重数，SA的毛口径作总流量参照。对照1803年补贴条约年要求960--1,920万pesos：\textbf{墨西哥白银在任何情景都填不满补贴条约的三分之一。}

\subsubsection{二、拉美格局1848}\label{ux4e8cux62c9ux7f8eux683cux5c401848}

\begin{itemize}
\tightlist
\item
  \textbf{新西班牙}：自治王国存续至1820s
  0.30--0.45、1830--45转事实独立≥0.70（形态更可能君主制、未经十年内战摧毁）；给1848事实独立≥90\%。Consolidación式追索（Yermo政变的触发链）在B4b世界里仍在------法国控制越强、追索越急、精英越离心；莫雷洛斯纲领第20条先验排除法国客户国。阿尔马登的水银是法国唯一的\literalquote{}刀\literalquote{}（3--7年把年铸币砍到600--1,200万），但受害者排序是英国＞西班牙＞墨西哥，法国无所得。
\item
  \textbf{南美}：两变量（合法国王×出海通道）；S2地图中位4个主权实体（3--7），\literalquote{}帝国联邦化\literalquote{}0.40--0.55（含新西班牙------与MX冲突，P8裁定全西属美洲保留≤0.10）；跨情景不变量：巴西不回欧洲法制秩序、上秘鲁归属恒为争端、蒙得维的亚--布城港口之争持续、巴拉圭孤立自治。英国自判征服无望（卡斯尔雷1807年5月1日），首选波旁系王子的立宪君主国------\textbf{英法在南美的目标函数意外重叠}（都要一个波旁王朝＋开放贸易），差别在谁的船运货。
\item
  \textbf{古巴}：三重否决均衡；1848归属西0.40--0.55／美0.20--0.35／英0.05--0.12／独立0.03--0.08／法系≤0.05；任何法国主导安排下脱离概率上升、去向美国；杰斐逊禁运是\literalquote{}大陆体系式贸易断绝施于加勒比\literalquote{}的唯一实测（一季内对美贸易腰斩、三分之二糖滞销）。
\item
  \textbf{海地}：C层承认／关税交易（促成签字0.55--0.80）；D层直接领有＝10--15年撷取窗口；命运在马德里决定；法国废奴推迟至1855--75（甜菜替代能力越小机制越强，1813年实产1.1M
  kg＝1788年殖民地糖的1.2\%）。
\item
  \textbf{巴西}：宫廷不回欧洲，入英圈成美洲第二大国，葡语世界永久分流；无1808开港加速器时克里奥尔弧线晚于西属美洲10--20年。
\item
  \textbf{美国}：门罗式原则提前且对象改为法国；扩张由人口决定（24年翻番）；佛罗里达推迟至1825--35、抵太平洋0.60--0.80；S6在美国维度不可观测，仅经\literalquote{}西班牙归属\literalquote{}传导。
\end{itemize}

\subsubsection{三、S6在美洲的代价是确定的}\label{ux4e09s6ux5728ux7f8eux6d32ux7684ux4ee3ux4ef7ux662fux786eux5b9aux7684}

合法性归零、白银与市场同时丧失；P2的行省化净收益须扣除美洲现值；S3在南美的净效果是给英国送钱。\textbf{\literalquote{}控制越强实得越少\literalquote{}在美洲是三次复现的结构定理}（F3海军、F4财政、MX白银）。

\begin{center}\rule{0.5\linewidth}{0.5pt}\end{center}

\subsection{第六十八章　印度：意愿、方法与能力}\label{ux7b2cux516dux5341ux516bux7ae0-ux5370ux5ea6ux610fux613fux65b9ux6cd5ux4e0eux80fdux529b}

\subsubsection{一、意愿有据可查}\label{ux4e00ux610fux613fux6709ux636eux53efux67e5}

四个投送方案的档案锚：1803年德卡恩自请六战列舰＋3,000步＋500骑；1805年1月16日致德克雷信（Corr.
X 8279）三分舰队20,000法军＋3,000西军；1808年5月13日（Corr. XVII
13877）洛里昂20舰4,000兵＋布雷斯特20舰13,000兵；1808年6月10日（Corr.
XVII 14078）\literalquote{}12 ships were to start from Nantes, 13 from Lorient and 31
from Brest, carrying a total of 19,600 soldiers and 10,600
sailors\literalquote{}------西班牙起义后永久放弃。1810年9月17日问德克雷埃及远征四万人\literalquote{}C\textquotesingle est
à mon plan de campagne pour
1812\literalquote{}；加尔达纳使团的对印计划有路线、月份、集结点、参谋勘测分工。\textbf{意愿是真的，且反复出现。}

\subsubsection{二、方法：四闸门与分层带}\label{ux4e8cux65b9ux6cd5ux56dbux95f8ux95e8ux4e0eux5206ux5c42ux5e26}

四闸门（出港×中继×季风窗×岁入基地）1803--15从未同开；十年实送≈1,200人＋五艘巡航舰；英方同期完成两次同量级跨洋远征（1810毛里求斯、1811爪哇）证明屏障是制海、中继与优先级。分层带：骚扰层已达上限自供养；千人级军官--军械层1803--08可行（Allard／Ventura
1822年在拉合尔建Fauj-i-Khas是历史验证）；万人级需S1／S2且四闸门同开；十万级任何情景不可行。陆路：单季无预置≥2万法军穿波斯关闭；分梯队＋四仓预置＋里海航运＋首年1,500--3,000万法郎白银可使2--3万人1.5--2个战役季抵赫拉特（工程边缘可行，中低）------但四政权的政治前提两两互斥（芬肯施泰因第14条排俄×俄要海峡×波斯要格鲁吉亚×奥斯曼要完整）。

\textbf{B4是最大单点改判杆}：1808年6月10日的方案在西班牙不变时可以出港。S2本义（B4b）下这条56舰的方案在1808--09年有出港条件，但仍缺中继（法兰西岛防御仅约两千人、奴隶动员被殖民者否决）、季风窗与岁入基地（马拉塔政体无可签约的集体主体、1803年12月诸约禁欧洲人服役）。\textbf{出港不等于抵达，抵达不等于立足。}

\subsubsection{三、能力：双侧脆弱性与英国的反应函数}\label{ux4e09ux80fdux529bux53ccux4fa7ux8106ux5f31ux6027ux4e0eux82f1ux56fdux7684ux53cdux5e94ux51fdux6570}

IN1：贸易养帝国不成立（1793--1812累计贸易利润£6.29m，年均£0.32m；领地账年缺£1.1--2.7m；印度债务£8m→£32m）；欧洲排除放松而非收紧印度体系约束（偿付能力来自中国与印度洋）；董事会1812年：\literalquote{}It
would be an European war for European
objects\literalquote{}。双侧脆弱性定理：增派国王军→公司军官晋升被挤压（1809马德拉斯哗变）；减欧籍骨干扩土著兵→兵变风险（1806韦洛尔）；任何可信的法国威胁把英国推向第一侧。明托1807年：\literalquote{}20,000
French troops... would not... subdue us---but would create a struggle
which must be attended with great present
calamity\literalquote{}。法国东方压力最确凿的产出是\textbf{英属印度前沿体系提前二十年成形}（1809年三条约：阿姆利则、喀布尔、信德；波斯预备约）。

\subsubsection{四、法属印度的终局}\label{ux56dbux6cd5ux5c5eux5370ux5ea6ux7684ux7ec8ux5c40}

终局谱系：非武装归还（1814／15实局）→武装商站（S2可争）→桥头堡（需1830年前不可得的印度洋海军均势）→军官市场国家化（拉合尔实证）→征服（否定）。\textbf{印度对法国霸权的最优用法是威胁在场消耗英国}（1810--11英为拔巢动用2.4万人两舰队）；α和约的形态是\literalquote{}非军事化商站＋法属印度洋两岛\literalquote{}，英国为此支付永久前沿--舰队溢价；B1b红线：不碰印度（和约存续要求红线份额≤0.10--0.15）。1848：英属印度存在且财政£20--25m（毛岁入1837--38实测£20.86m）、公司纯行政机关、鸦片与茶是不依赖欧洲市场的税源（国库从茶得税1792--1809共£32.29m）；旁遮普吞并进程照史实；波斯若法国不付费则入英轨。\textbf{夺取印度在任何情景超出法国海运与财政。}

\begin{center}\rule{0.5\linewidth}{0.5pt}\end{center}

\section{结论}\label{ux7ed3ux8bba}

\subsection{一、对原初问题的回答}\label{ux4e00ux5bf9ux539fux521dux95eeux9898ux7684ux56deux7b54}

\textbf{回到1803年，拿破仑如何逆转历史、赢得战争、控制整个欧洲并建立长久霸权？}

他不能靠打得更好来做到。本书四十四个行动者的能力包络与反应函数一致指向一个结论：法国在1803年已经拥有的西欧优势可以固化，欧洲大陆的四层控制可以用条约建立，俄国可以被关税与同盟约束，英国可以在不被击败的情况下接受------但这一切的前提不是更多的胜利，而是\textbf{在胜利之后停下来}。具体地说，是三个停下来的选择：把封锁变成关税，把西班牙变成联姻的盟友，把教皇留在罗马；再加两个更贵的：不并荷兰与汉萨，不征俄国。

本书把这条路叫做\textbf{S2本义}：1803--1807年按史实打赢三场战争，取得普雷斯堡、提尔西特与莱茵邦联的真实条约资产；1807年秋至1808年春切换两个开关（大陆体系→保护性财政关税并对俄关税容忍；西班牙→保留波旁并联姻）；1809年与教廷和解；1810年不并荷汉、奥地利联姻；1811年统合公债、不吞奥尔登堡；1812年不征俄，在英国失去大陆理由的1811--14年窗口谈总和约；1813年起降为和平姿态。P6引擎在本书作者的补算下给这条路径全程人力TIGHT、1808年后财政为正、无战区冲突、零外生奇迹；它需要的外生事件与P6原主线相同（英国签约0.40--0.60、俄国十年不战0.35--0.50、奥地利不再入场0.40--0.60），但需要的克制节点少三个。

\textbf{它有多现实？}条件于法国采纳并维持有界目标框架，P6给0.30--0.50，P8给0.06--0.40（强共因0.26--0.50）；从1803年出发的无条件概率，P6给1.8--9.0\%，P8给0.4--7.3\%。这些是裁量带不是频率。最可能的实际终局不是签约的和平，而是\textbf{英国始终不签约的武装共存}------大陆控制成立、俄国不战、K3不成立。这条路的每一个构件在当时都被提出过（塔列朗、特里亚农、费尔南多的求婚信、教廷的关港书、路易的条约、涅谢尔罗迭的方案、1811年4月的裁军函），它是\literalquote{}合理\literalquote{}的；账算得平，它是\literalquote{}可行\literalquote{}的；它要求一个\literalquote{}靠自己爬上来的军人\literalquote{}在胜利之后做他在1802--03与1810--11两个史实和平窗口里都没做的事，所以它不是\literalquote{}现实\literalquote{}的------如果\literalquote{}现实\literalquote{}的意思是\literalquote{}拿破仑会那样做\literalquote{}。

\subsection{二、最重要的问题：行省化}\label{ux4e8cux6700ux91cdux8981ux7684ux95eeux9898ux884cux7701ux5316}

\textbf{字面版（法国省制覆盖大陆）：不存在合理可行的路径。}
置信高。不是\literalquote{}很难\literalquote{}，是账算不平：人力（S6-max缺12.2万、full缺30.2万在场兵）、财政（丙层年缺532.6m）、治安（外圈6--10万驻军吃掉全部富余）、盟军（吞南德＝33k盟军换60--90k驻军）、整合（新省第十年才转正且只在最优口径下）、海军（最大行省化≠最大有效海权）、白银（控制越强实得越少）、货币（兼并西班牙＋双金属＝自造通缩）、教会（A层可扩性的第一证伪器）、思想（天主教正统诉讼＞继承真空＞民族主义）------十本账互相独立到足以互证。Q1给出结构判据：A层的地理边界是英奥俄军事投送的边界，不是文化边界。

\textbf{边界版（A层限于内圈）：存在，且已被史实部分实现。}
置信中高。约39.5--40.5百万人口的直辖核心（旧法国＋比利时、左莱茵、皮埃蒙特、利古里亚、莱芒；帕尔马另列；托斯卡纳、瓦莱条件直辖；贝格须获真实市场准入），1806--10巩固、1825--35制度稳态。它也不免费：陆300＋海195只在高收入＋盟约转移下接近闭合，中心案缺60--115m；成立条件是盟约兑现与海军节奏让步。

\textbf{乙层（北意、德西、伊利里亚、加泰，56.7m人口）：技术可闭合、政治未证。}
低置信支线，并列保留不作主线：P2的联合乐观反例证明它不能以\literalquote{}纯技术不可能\literalquote{}结案，但闭合条件（100\%保留原B军并自养、低驻军、350m陆军、海军让步）在政治上等于B层而非A层。

\textbf{控制版：以S2本义的形态存在（有条约支撑），以路径B的形态只达到primacy。}

行省化的正确问法不是\literalquote{}能并到哪里\literalquote{}，而是\literalquote{}并了之后还能得到什么\literalquote{}。本书的答案：内圈之外，每多并一块，法国失去的（自养盟军、市场、白银、教会合作、英国的接受）多于得到的（第十年才转正的兵、预期而非实收的税）。

\subsection{三、分问}\label{ux4e09ux5206ux95ee}

{\def\LTcaptype{none} % do not increment counter
\begin{longtable}[]{@{}
  >{\raggedright\arraybackslash}p{(\linewidth - 6\tabcolsep) * \real{0.2500}}
  >{\raggedright\arraybackslash}p{(\linewidth - 6\tabcolsep) * \real{0.2500}}
  >{\raggedright\arraybackslash}p{(\linewidth - 6\tabcolsep) * \real{0.2500}}
  >{\raggedright\arraybackslash}p{(\linewidth - 6\tabcolsep) * \real{0.2500}}@{}}
\toprule\noalign{}
\begin{minipage}[b]{\linewidth}\raggedright
分问
\end{minipage} & \begin{minipage}[b]{\linewidth}\raggedright
主线
\end{minipage} & \begin{minipage}[b]{\linewidth}\raggedright
备选
\end{minipage} & \begin{minipage}[b]{\linewidth}\raggedright
置信
\end{minipage} \\
\midrule\noalign{}
\endhead
\bottomrule\noalign{}
\endlastfoot
② 秩序如何运作 &
以巴黎为清算中心的分层准入体系；浅铸币同盟（白银区，命运绑在马德里）；节点基建（莱茵自由航行提前6--16年、安特卫普商港、帝国干线铁路3,500km）；三层农村（改革深度与军事价值成反比）；\literalquote{}带盔甲的七月王朝\literalquote{}（1830后）
& 有限关税同盟（1848视点0.25--0.35）；深货币同盟（0.10--0.20，继承后） &
中高 \\
③ 海军差距 &
α（英接受＋安特卫普商业港化）：1830效能比≈.55--.60、舰体≈.58；β（不接受＋追赶）：.73／.90但无和约
& P高档（220--260m）不在任何财政平衡案内 & 中 \\
④ 西班牙终局 &
波旁绝对主义--大臣改革制0.40--0.55；波旁立宪0.30--0.45；波拿巴0.05--0.15；共和≤0.02
& 若法国改行废立：波拿巴0.35--0.55、共和0.03--0.08、割据0.10--0.20 &
中 \\
⑤ 罗马王与意大利 &
在位期分治藩属0.55--0.65；继承后北中意统一王国（S2本义0.35--0.45，路径B
0.20--0.30）；含那不勒斯0.10--0.18；统一分支法理上以对英终局（马耳他）为前提
& 行省化\literalquote{}第二法国\literalquote{}0.10--0.20（惩罚态） & 中 \\
⑥ 普奥消长 &
S2本义（残普）：奥强0.70--0.85；路径B（完整普鲁士）：奥强0.45--0.60；两路径普鲁士均不领导德意志
& 取消普鲁士1807夏--1808年2月有对俄价格、此后≈对俄战争 & 中／中低 \\
⑦ 美洲 &
能，唯一路径是通过马德里（B4b＋B6b）；法国实得白银90--330万pesos／年；墨西哥1848事实独立≥90\%；南美中位4个主权实体、帝国联邦化0.40--0.55；古巴留西0.40--0.55；巴西入英圈；海地命运在马德里
& 全西属美洲保留≤0.10（P8） & 中 \\
⑧ 印度 &
意愿有据（四个投送方案）；能力任何情景不足（四闸门从未同开）；最优用法是威胁在场消耗英国；夺取印度超出法国海运与财政
& S2本义下1808年6月方案可出港但不可立足 & 中高 \\
\end{longtable}
}

\subsection{四、备选结论群}\label{ux56dbux5907ux9009ux7ed3ux8bbaux7fa4}

\begin{enumerate}
\def\labelenumi{\arabic{enumi}.}
\tightlist
\item
  \textbf{路径B（1803年即切换）}：内部自洽的primacy固化线，离开单位卡估计域；若1802--03窗口的克制先验（0.05--0.15）被新证据上调（拿破仑对马耳他仲裁的书面态度），它成为比S2本义更便宜的主线。
\item
  \textbf{路径A（入侵）}：不是主线但不构成逻辑排除；两个物理奇迹（D+2完整上岸、跨海补给）若被登陆战局的新校准（蒙特耶尔运载表、两潮／六潮装卸量）提高到\textgreater0.25，A重回候选。
\item
  \textbf{β终局（武装共存）作为默认}：P6自认最可能；若英国的安特卫普要价被证明不可让（1806／1813--14文书），则K3在任何路径都不成立，\literalquote{}控制整个欧洲\literalquote{}只能以\literalquote{}英国不败而共存\literalquote{}的形态达成------本项目允许这一形态。
\item
  \textbf{乙层政治协约}：若IT1／G1／SP1的合法转隶文书与1812年在岗实员出土，乙层可能从低置信支线升为主线的一部分。
\item
  \textbf{无胜利帝国的合法性侵蚀}：若1810--13年军官团文书显示和平前景引发系统性密谋，S5存续概率应再下调3--8个百分点以上。
\end{enumerate}

\subsection{五、改判证据（全书汇总，按改判幅度排序）}\label{ux4e94ux6539ux5224ux8bc1ux636eux5168ux4e66ux6c47ux603bux6309ux6539ux5224ux5e45ux5ea6ux6392ux5e8f}

\begin{enumerate}
\def\labelenumi{\arabic{enumi}.}
\tightlist
\item
  英国对中率关税与安特卫普非军事化的当时反应文书（1806年谈判法方条款、1812年撤令条件、1813--14卡斯尔雷训令）------决定K3与③。
\item
  拿破仑对1803年马耳他仲裁、1808年联姻去条件化、1810年荷兰保留、1812年俄方照会的书面显示偏好------决定C1／C5／C7／C9的先验。
\item
  1812年同口径在岗文官与宪兵实员；荷／汉／罗年度征兵面板------决定行政硬阈值与外圈判定。
\item
  F4逐年表（Branda 2007、Gaudin附录、Compte
  général）------决定贡赋依赖度的峰值年。
\item
  全国逐年到营／退出／净存量；役期与退出率的F2／F5口径合账------决定S5人力账。
\item
  汉堡agio月度序列------决定封锁伤金属还是伤信用（白银区论证）。
\item
  1830年P基线以71舰重跑P4------决定α终局的精确数字。
\item
  Du Casse十卷与Rath 1941------决定意大利继承。
\item
  1806年普军实际野战兵力与法国对普鲁士北德方案的态度------决定路径B的普奥。
\item
  1808年前西班牙本土共和纲领（PARES／BNE）------决定共和上界。
\end{enumerate}

\subsection{六、未决问题（诚实清单）}\label{ux516dux672aux51b3ux95eeux9898ux8bdaux5b9eux6e05ux5355}

\begin{itemize}
\tightlist
\item
  三个C类节点（对英不加价、不并荷汉、统合公债）恰在奇迹阈值上，\literalquote{}零奇迹\literalquote{}不等于高实现概率。
\item
  F2／F5役期与退出率口径未合；10.6万／年的维持式补充是否越过社会容忍带待对账。
\item
  引擎\texttt{delta\_prussia\_takeover\_mid=6.0}与注释4.5--7.5万的十倍差异未定案（只影响B3v2分支）。
\item
  奥地利在1805受辱后、无西班牙泥潭时1809再战概率（0.25--0.40）为本书裁量。
\item
  奥婚在S2本义下的概率（0.45--0.60）为本书裁量；不成则罗马王链断。
\item
  α终局的舰队数字是比例近似，非P4重跑。
\item
  乙层的政治条件、丙层的行政阈值、外圈的年度面板均未实证。
\item
  白银两口径（MX净得／SA毛流量）已分列但未逐年调和。
\item
  非英语二手文献（俄／波／波斯／土／阿拉伯语）经译本或转引进入，未直读。
\item
  已派子节点的运行状态未由工具确认停止；本书不依赖其任何产出。
\end{itemize}

\subsection{七、最后一句}\label{ux4e03ux6700ux540eux4e00ux53e5}

把欧洲行省化的路径不存在，把欧洲控制住的路径存在。前者的账算不平，后者的账算得平但要拿破仑不是拿破仑。历史学不能替一个人做选择，但可以把选择的价格算出来------这是本书做的事。

\begin{center}\rule{0.5\linewidth}{0.5pt}\end{center}

\section{附录}\label{ux9644ux5f55}

\subsection{附录一　材料总目}\label{ux9644ux5f55ux4e00-ux6750ux6599ux603bux76ee}

本目只列正文实际援引的著作与文书，按访问状态分四类：\textbf{亲读}（本轮在下载件或页图上读到所引段落）；\textbf{页图核}（据扫描页图核对表格或原文）；\textbf{转引}（经另一著作或书评取得，正文已标）；\textbf{未获}（正文提及其存在或据以指明改判证据，但本轮未取得）。工作区696份被引用文件的机器清单另存\texttt{附录一\_材料总目\_工作区文件.csv}（按被引卡数排序），供追溯；本目不以文件名代替引用。

\subsubsection{一手文书与刊本}\label{ux4e00ux624bux6587ux4e66ux4e0eux520aux672c}

{\def\LTcaptype{none} % do not increment counter
\begin{longtable}[]{@{}
  >{\raggedright\arraybackslash}p{(\linewidth - 6\tabcolsep) * \real{0.2500}}
  >{\raggedright\arraybackslash}p{(\linewidth - 6\tabcolsep) * \real{0.2500}}
  >{\raggedright\arraybackslash}p{(\linewidth - 6\tabcolsep) * \real{0.2500}}
  >{\raggedright\arraybackslash}p{(\linewidth - 6\tabcolsep) * \real{0.2500}}@{}}
\toprule\noalign{}
\begin{minipage}[b]{\linewidth}\raggedright
文书
\end{minipage} & \begin{minipage}[b]{\linewidth}\raggedright
出处／版本
\end{minipage} & \begin{minipage}[b]{\linewidth}\raggedright
状态
\end{minipage} & \begin{minipage}[b]{\linewidth}\raggedright
主要使用章
\end{minipage} \\
\midrule\noalign{}
\endhead
\bottomrule\noalign{}
\endlastfoot
1814年5月30日巴黎和约第XV条（\literalquote{}Antwerp shall for the future be solely a
commercial port\literalquote{}／\literalquote{}Dorénavant le port d\textquotesingle Anvers sera
uniquement un port de commerce\literalquote{}） & 英文刊本与法文刊本网页 & 亲读 &
五十二、六十三 \\
1802年亚眠条约第VI、X、XIII--XIV条 & 英文刊本 & 亲读 & 六、十一 \\
1805年3月17日意大利王国第一宪制法令art.1、3、4；1810年2月17日元老院法令art.7、10
& 法文刊本页；意文重排本 & 亲读 & 二十三、六十五 \\
1805年意大利王国财政总法第13条 & 刊本 & 亲读 & 三十八 \\
1803年9月27日法瑞军事协约；1803年调停法 & 法文刊本 & 亲读 & 二十一 \\
1806年7月12日莱茵邦联条约；1808年5月巴伐利亚宪法§1 & 德文刊本 & 亲读 &
十九 \\
1807年10月27日枫丹白露条约（La Parra转录） & 转引 & 转引 & 二十六 \\
1808年3月5日cuestiones proponibles（十八条） & La Parra转录 & 转引 &
二十六 \\
1808年1月10日拿破仑致卡洛斯四世（\literalquote{}Consiento de buen grado\ldots\literalquote{}） & La
Parra转录 & 转引 & 二十六、五十一 \\
1808年3月30日拜永条约（缪拉） & Davis转述 & 转引 & 二十四 \\
1808年9月8日法普巴黎公约正约第VII、VIII、XIII条与分立条款 &
法文／德文刊本 & 亲读（两版冲突并列） & 十八 \\
1809年10月14日申布伦条约秘密条款II、第V条 & 法文刊本 & 亲读 & 十七 \\
1809年9月4日奥方照会（波希米亚三圈） & 刊本 & 亲读 & 十七 \\
1812年2月24日法普秘密军事协定第II条 & 德文刊本 & 亲读 & 十八 \\
1812年3月法奥同盟第六条 & 刊本 & 亲读 & 二十九 \\
1807年5月芬肯施泰因条约第3--4、12--14条 & 法文刊本 & 亲读 & 三十一 \\
1808年10月12日埃尔福特公约第8、12、16条 & Vandal转录 & 转引 & 二十九 \\
1809年阿姆利则、喀布尔、信德三条约；1809年波斯预备约第3条 &
Aitchison卷VII & 亲读 & 三十一 \\
1813年10月8日里德条约公开第4条与秘密第5条 & 德文刊本 & 亲读 & 十九 \\
1814年1月14日基尔条约第4、6条 & 丹麦文刊本 & 亲读 & 二十二 \\
1807年10月22日英葡秘密公约；1810年英葡商约 & Manchester转录FO档 & 转引 &
二十八 \\
1811年9月14日普鲁士调整敕令§9、12、37、41、53；1816年5月29日声明第4--5条
& GHDI英译 & 亲读 & 五十九 \\
1804年8月15日莱茵河航行octroi公约第2、5条；1808年10月27日Magistrat du
Rhin敕令 & 法文刊本 & 亲读 & 五十九 \\
1811年12月16日第7644号道路分级法令 & Bulletin des lois 418 &
转引（Arnaud章） & 五十九 \\
1839年比荷条约第9条；1863年斯海尔德赎买条约 & London Gazette 1863-08-04
& 亲读 & 五十九 \\
1865年12月23日四国货币公约；FRUS 1867 doc.291 & 刊本 & 亲读 & 五十九 \\
1813年2月5日摄政法 & 刊本（Bulletin des lois第474／490号原刊未亲验） &
亲读（转录） & 五、六十 \\
1814年4月6日元老院宪法第6条 & 刊本 & 亲读 & 五 \\
1812年马莱伪元老院令 & SHD dossier Malet转录 & 转引 & 五 \\
1803年爱尔兰临时政府宣言第2条 & 刊本 & 亲读 & 十二 \\
1805年海地帝制宪法总则第12条、第50--52条 & 刊本 & 亲读 & 三十五 \\
1813年莫雷洛斯《Sentimientos de la Nación》第1、9、16、20条 & 刊本 &
亲读 & 三十六 \\
1812年6月23日枢密令（对美撤销） & Historical Hansard & 亲读 & 六 \\
1833年废奴法第LXIV--LXV条 & 刊本 & 亲读 & 十一 \\
1823年12月门罗咨文；1823年坎宁--波利尼亚克备忘录 & 刊本 & 亲读 &
六、三十四 \\
拿破仑《通信集》：1810年1月3日致阿尔迪尼；1810年7月9日致勒布伦；1810年7月25日致莫利安；1810年9月17日致德克雷；1810年11月3日瓦莱信；1811年11月27日致欧仁；1811年4月裁军函（17579）；X
8279；XVII 13877、14078；XXX p.549（\literalquote{}次子\literalquote{}方案，回溯层） &
刊本／Gutenberg／转录 & 亲读（部分转引） &
一、二十一、二十三、三十、三十三 \\
1811年2月27日海军部报告；1810年7月13日与1811年3月5日、10月2日海军命令 &
Todorov转录 & 转引 & 三 \\
1806年fédératif训令；1808年致路易法币统一信；尚帕尼1810年12月8日报告 &
刊本 & 亲读 & 一 \\
1810年12月25日、1811年1月31日、1812年4月1日亚历山大致恰尔托雷斯基三信；1813年1月13日信
& 恰尔托雷斯基回忆录附录 & 亲读 & 十三、十六 \\
1811年10月涅谢尔罗迭报告；1811年6月5日科兰古召对 &
Vandal卷三转录；科兰古回忆录 & 亲读（Vandal）／回溯性 &
十三、十五、十八 \\
斯佩兰斯基1811年财政报告 & 俄文刊本 & 亲读 & 十四 \\
1808年1月华沙公国预算草案（AE CP 324, 343） & 波兰文转引 & 转引 &
十六 \\
1810年9月16日科隆商会致Ladoucette书（AN F12 549--550） & Schmidt
1905附录E & 转引 & 二十 \\
1810年10月2日Bacher报告 & Heckscher转述 & 转引 & 二十 \\
1812年11月10日斯特拉斯堡商会统计表第10号 & Ellis附录D & 页图核 &
五十九 \\
1811年5月20日预算辩论；1813年11月海军辩论；1815年6月14日万西塔特发言；1823年、1830年、1854年下院发言；1842年、1848年宪章请愿点验
& Historical Hansard & 亲读 & 七、八、十一、四十一 \\
东印度公司1812年董事会函、1813年请愿书、Cartwright 1812年3月4日总账 &
刊本 & 亲读 & 三十二 \\
卡斯尔雷1807年5月1日内阁备忘录；波帕姆1803年备忘录；科克伦--米兰达1806年条款
& Castlereagh Correspondence vol.VII & 亲读 & 三十七 \\
明托1807年函；1809年11月12日陆军部长宪兵报告 & 刊本／转录 & 亲读 &
三十二、五十七 \\
1806年12月30日雅茅斯勋爵叙述 & Hansard & 亲读 & 六 \\
教皇1813年撤回函；1808年春教廷关港声明 & Pacca回忆录1833年法文版 & 亲读
& 二十五 \\
《Mirari Vos》 & 刊本 & 亲读 & 二十五、六十 \\
塔列朗1805年10月17日斯特拉斯堡备忘录 & Dard转录 & 亲读（Dard） & 一 \\
梅特涅1829年追述德累斯顿谈话 & 梅特涅回忆录卷一 & 回溯性 & 一、六十 \\
\end{longtable}
}

\subsubsection{二手著作（含回忆录与统计汇编）}\label{ux4e8cux624bux8457ux4f5cux542bux56deux5fc6ux5f55ux4e0eux7edfux8ba1ux6c47ux7f16}

{\def\LTcaptype{none} % do not increment counter
\begin{longtable}[]{@{}
  >{\raggedright\arraybackslash}p{(\linewidth - 6\tabcolsep) * \real{0.2500}}
  >{\raggedright\arraybackslash}p{(\linewidth - 6\tabcolsep) * \real{0.2500}}
  >{\raggedright\arraybackslash}p{(\linewidth - 6\tabcolsep) * \real{0.2500}}
  >{\raggedright\arraybackslash}p{(\linewidth - 6\tabcolsep) * \real{0.2500}}@{}}
\toprule\noalign{}
\begin{minipage}[b]{\linewidth}\raggedright
作者
\end{minipage} & \begin{minipage}[b]{\linewidth}\raggedright
著作
\end{minipage} & \begin{minipage}[b]{\linewidth}\raggedright
状态
\end{minipage} & \begin{minipage}[b]{\linewidth}\raggedright
主要使用章
\end{minipage} \\
\midrule\noalign{}
\endhead
\bottomrule\noalign{}
\endlastfoot
Aaslestad, K. \& Joor, J. (eds.) & Revisiting Napoleon\textquotesingle s
Continental System (2015)，含Rowe、Aaslestad、Spaulding各章 &
亲读（印p.212等） & 三十九、四十四 \\
Aksan, V. & Ottoman Wars 1700--1870 & 亲读 & 二十九、三十 \\
Albion, R. G. & Forests and Sea Power (1926) & 亲读（pp.356、368） &
八 \\
Alder, K. & Engineering the Revolution & 亲读 & 四十 \\
Antipa, P. & \literalquote{}How Fiscal Policy Affects Prices:
Britain\textquotesingle s First Experience with Paper Money\literalquote{} JEH
2016；Antipa \& Chamley 2017 & 亲读 & 四十七 \\
Axtmann, R. & 1991年奥地利财政文（附录转Beer 1877） & 页图核 & 十七 \\
Barton, H. A. & Scandinavia in the Revolutionary Era & 亲读 & 二十二 \\
Beck, K. & 1890年据Eberstein遗档论邦联宪法 & 亲读（开篇长段） & 二十 \\
Berding, H. & Napoleonische Herrschafts- und Gesellschaftspolitik im
Königreich Westfalen & 转引（手稿摘录） & 十八 \\
Bordo, M. \& White, E. & NBER WP 3517 (1990) & 亲读（表2） & 七 \\
Brading, D. \& Cross, H. & 论殖民地白银与水银 & 转引 & 三十六 \\
Branda, P. & 战争融资附表（缓存页逐格）；Le prix de la gloire (2007) &
页图核（附表）／未获（专著） & 四 \\
Broers, M. & The Politics of Religion in Napoleonic Italy &
亲读（pp.179--182） & 四十一、六十 \\
Brun, J.-F. & Des armes et de la poudre & 亲读 & 二 \\
Buist, M. G. & At Spes Non Fracta: Hope \& Co. 1770--1815 (1974) &
亲读（全文） & 三十四、三十六、三十八 \\
Chaptal, J.-A. & De l\textquotesingle industrie françoise (1819) &
亲读（OCR） & 五十九 \\
Chilcott, T. & 2006年英国本土防御研究 & 亲读（pp.242--244） & 六、九 \\
Consalvi, E. & Mémoires (1866) & 亲读 & 二十五 \\
Crafts, N. & 2004年蒸汽动力文（Kanefsky系列） & 转引 & 四十 \\
Crosby, A. W. & America, Russia, Hemp, and Napoleon (1965) &
亲读（全文） & 十四 \\
Czartoryski, A. & Memoirs (含附录信件) & 亲读 & 十三、十六 \\
Dard, É. & Napoléon et Talleyrand & 亲读 & 一 \\
Davey, J. & 2009年博士论文（波罗的海给养） & 亲读（全文） &
八、十四、二十二 \\
Davis, J. A. & Naples and Napoleon & 亲读 & 二十四 \\
Desbrière, É. & Projets et tentatives de débarquement &
亲读（pp.442--444） & 九、十二（转引路径） \\
Driault, É. & La politique orientale de Napoléon (1904)；Mohamed Aly et
Napoléon (1925) & 亲读 & 三十、三十一 \\
Dutt, R. & Economic History of India (1906) 第XXIII章 & 转引 & 三十二 \\
Einaudi, L. & ifo DICE Report 2018/3 & 亲读 & 五十九 \\
Ellis, G. & Napoleon\textquotesingle s Continental Blockade: The Case of
Alsace (1981) & 亲读（附录D、G、I、J；印p.202、286） & 三十九、五十九 \\
Elphinstone, M. & Account of the Kingdom of Caubul (1815) & 亲读 &
三十一 \\
Elting, J. & Swords Around a Throne & 亲读（第三章） & 二 \\
Esdaile, C. & Popular Resistance in the French Wars & 亲读 &
二十五、二十七 \\
Fahmy, K. & All the Pasha\textquotesingle s Men & 亲读 & 三十 \\
Ferguson, N. & The House of Rothschild & 亲读 & 三十八 \\
Forrest, A. & Conscripts and Deserters & 亲读（pp.91--92） & 五 \\
Fortescue, J. & County Lieutenancies and the Army 1803--1814 & 亲读 &
九 \\
Fraser, R. & Napoleon\textquotesingle s Cursed War & 亲读 & 二十七 \\
Geggus, D. & 1979年黄热病文 & 亲读 & 三十五 \\
Glenthøj, R. \& Ottosen, M. N. & Experiences of War and Nationality in
Denmark and Norway & 亲读（pp.186--190） & 二十二、四十一 \\
Gordon, S. & The Marathas & 亲读 & 三十三 \\
Grab, A. & Napoleon and the Transformation of Europe
(2003)；2013年征兵文 & 亲读 &
二、十六、十八、二十、二十三、四十四、四十五 \\
Gregory, D. & Sicily: The Insecure Base & 亲读 & 二十四 \\
Hamnett, B. & IEP 2000（秘鲁） & 亲读 & 三十七 \\
Heckscher, E. & The Continental System &
亲读（印pp.210、214、227、244--245、293、295--306） &
二十、三十五、三十九、四十 \\
Herr, R. & Rural Change and Royal Finances in Spain & 亲读 & 二十六 \\
Hicks, P. & 论教廷--帝国升级链 & 亲读 & 五 \\
Houdaille, J. & 帝国陆军来源统计 & 亲读（印p.49） & 二 \\
Humboldt, A. von & Essai politique sur le royaume de la Nouvelle-Espagne
& 亲读 & 三十六 \\
INED / Bonneuil & 1806年法国生命表重建 & 亲读 & 五 \\
INEHRM & Yermo研究（印pp.54--59） & 亲读 & 三十六 \\
James, W. & Naval History, Abstract 17--23 & 转录（未逐年影像对校） &
八 \\
Jarrige, F. & 卢德运动研究（pp.1--3） & 亲读 & 四十一 \\
Jeremy, D. & 1977年英国技术封锁文 & 亲读 & 四十 \\
Johnston, R. M. & The Napoleonic Empire in Southern Italy & 亲读 &
二十四 \\
Joor, J. & Soldiers, Citizens, Civilians (2009) & 亲读（全文；印p.202）
& 二十一、四十五 \\
Kaye, J. & Life of Malcolm & 亲读 & 三十一 \\
Keep, J. & Soldiers of the Tsar & 亲读（摘录） & 十五 \\
Keller, W. \& Shiue, C. & 2014年关税同盟文 & 亲读 & 五十九 \\
Kopsidis, M. \& Bromley, D. & 2016年南德领主制文 & 亲读 &
十九、五十九 \\
Kukiel, M. & Wojna 1812 & 未获（转引） & 十六 \\
La Parra, E. & 戈多伊传；费尔南多七世传 & 亲读 & 二十六 \\
Lawrence, M. & 论1823与卡洛斯战争 & 亲读 & 二十七 \\
Lentz, T. & Correspondance générale T11导言 & 转引 & 一 \\
Lieven, D. & Russia Against Napoleon & 亲读（部分）／书评转引 &
十三、十五 \\
Lignereux, A. & 旺代1815；宪兵研究书评；PUR 2013并合省rébellions专题 &
亲读（书评）／未获（PUR专题） & 五、四十五、五十七 \\
Linch, K. & Britain and Wellington\textquotesingle s Army & 亲读 & 九 \\
Lugli, A. et al. & 2007年拿破仑死因文 & 亲读 & 一 \\
Madelin, L. & La Rome de Napoléon (1906) & 亲读（印p.375；脚注引AN AF IV
1509为转引谱系） & 四十四 \\
Manchester, A. K. & British Preëminence in Brazil & 亲读（全文） &
二十八 \\
Marichal, C. & Bankruptcy of Empire & 亲读 & 二十六、三十四、三十六 \\
Marion, M. & Histoire financière de la France IV (1925) &
页图核（pp.321、325、419） & 四、二十一、四十三 \\
Marsot, A. L. & Egypt in the Reign of Muhammad Ali & 亲读 & 三十 \\
Marshall, J. & Digest of All the Accounts (1833) & 亲读（印p.74） &
十 \\
Marzagalli, S. & 论汉堡走私 & 转引（1999全书未获） & 三十九 \\
Melchior-Bonnet, B. & Napoléon et le Pape & 转引（p.87） & 五 \\
Mercader Riba, J. & 论加泰法治 & 转引 & 四十四 \\
Milburn, W. & Oriental Commerce & 亲读 & 三十二 \\
Murau, S. & 2023年普鲁士国债文 & 亲读 & 十八 \\
Murray/Prinsep & 1831--33年锡克研究 & 亲读 & 三十一 \\
Oman, C. & History of the Peninsular War III附录VII、IV附录XVIII &
亲读（downloads/f2\_oman\_vol3／4） & 二、二十八 \\
Pacca, B. & Mémoires (1833年法文版) & 亲读 & 二十五 \\
Pedreira, J. & HAHR 2000 & 亲读 & 二十八 \\
Perkins, B. & Prologue to War (1961) & 亲读 & 三十四 \\
Petiteau, N. & 显贵跨政权研究 & 转引 & 五、六十 \\
Pigeard, A. & \literalquote{}La conscription sous le Premier Empire\literalquote{} (RSN 420, 1998) &
亲读（网页） & 二 \\
Pitkin, T. & Statistical View (1835) & 亲读 & 三十四 \\
Planert, U. & Der Mythos vom Befreiungskrieg (2007) & 转引（两篇书评） &
四十一 \\
Ploeckl, F. & 2010年关税同盟博弈文 & 亲读 & 五十九 \\
Porter, G. R. & Progress of the Nation & 亲读 & 四十二 \\
Riehn, R. & 1812: Napoleon\textquotesingle s Russian Campaign & 亲读 &
十五 \\
Rosselli, J. & Lord William Bentinck and the British Occupation of
Sicily & 亲读 & 二十四 \\
Rouanet \& Piano & 2023年逃役反事实文 & 亲读（正文与附录冲突并列） &
四十六 \\
Rowe, M. & From Reich to State & 页图核（表3印p.179；印pp.87、90） &
二、四十三、四十五、四十六 \\
Sen, S. P. & The French in India & 亲读（全文） & 三十三 \\
Shaw, S. & Between Old and New & 亲读 & 二十九 \\
Sondhaus, L. & 论伦-威籍奥军 & 亲读 & 二十三、六十五 \\
Tokarz, W. & 1917年华沙公国被占负担研究 & 亲读 & 十六 \\
Todorov, N. & 法国海军舰体研究 & 亲读（PDF p.15） & 三、六十三 \\
Tone, J. L. & The Fatal Knot & 亲读（转de la Torre） & 二十七 \\
Torrente, M. & Historia de la revolución hispano-americana (1829) & 亲读
& 三十七 \\
Troshin, N. & 2015年俄贸易条件文 & 亲读 & 十四 \\
Ullmann, H.-P. & 法兰克福资本市场研究 & 亲读（全文） & 三十八 \\
van Creveld, M. & Supplying War & 亲读（印p.64） & 二 \\
Vandal, A. & Napoléon et Alexandre Ier三卷 & 亲读（Gutenberg） &
十三、十八、二十九 \\
Weis, E. & 1983年巴伐利亚世俗化文 & 亲读 & 十九 \\
Wilson, W. J. & History of the Madras Army & 亲读 & 三十二 \\
Woolf, S. & Napoleon\textquotesingle s Integration of Europe &
亲读（印pp.203、232、243、245） & 四十三、四十五 \\
Ziegler, P. & The Sixth Great Power & 亲读 & 三十八 \\
ADB（Flathe等条目）；HdBG；HLB；DBI；Iranica & 传记辞典条目 & 亲读 &
十八--二十、二十三、三十一 \\
Bank of England & A Millennium of Macroeconomic Data
v3.1（A27、A31、A40表） & 亲读 & 七、十、四十七 \\
Bairoch, P. & 1982年人均工业化指数 & 转引（NBER WP6904转录） & 五十九 \\
Maddison Project 2020 & 人均GDP & 亲读 & 五十九 \\
\end{longtable}
}

\subsubsection{未获而应获（改判证据所在）}\label{ux672aux83b7ux800cux5e94ux83b7ux6539ux5224ux8bc1ux636eux6240ux5728}

Branda 2007逐年表；Gaudin回忆录附录；各年Compte général；Mitchell \&
Deane 1962 p.442；Sherwig补贴表；Danson 1894；Crouzet两卷；Glover
1967；Winfield--Roberts原书；Rath 1941；Du Casse十卷；Finley
1994；Fugier两卷；Kukiel《Wojna
1812》；Żółtowski决算；Dundulis；Pivec-Stelè；Bundy；Mercader
Riba原书；Lignereux PUR 2013；Marzagalli
1999；Dufraisse原文；Tarlé；Planert／Hagemann原书；Chadwick／Hales／Latreille／Schwarzfuchs／Broers
2002／Caiani 2021；Lynch／Halperín
Donghi／Adelman／Anna／Street；Dubois／Girard／Blackburn／Drescher／Ferrer／Régent／Popkin；Bowen／Peers／Philips／Webster；Amini／Atkin／Yapp／Ingram／Kelly；al-Jabartī原文／Ghorbal／Douin／Bowring；Landes／Pollard／Henderson／Horn／Chassagne／Daumas；Harnisch
1984／Koselleck 1967；Shchepetil\textquotesingle nikov；Saumarez
Papers原件；Herries回忆录正文；Gille／Lévy-Leboyer／Bouvier；Soetbeer
1855／Staatsarchiv Hamburg agio序列；ICPSR微观复制包；HC 1855
Return；Kilmainham／WO 17逐月表；AN AF IV 1123--24；Almanach impérial
1810／1812逐省名录；1806年英法谈判法方书面条款；1813--14卡斯尔雷对安特卫普训令。

\begin{center}\rule{0.5\linewidth}{0.5pt}\end{center}

\subsection{附录二　口径冲突表}\label{ux9644ux5f55ux4e8c-ux53e3ux5f84ux51b2ux7a81ux8868}

{\def\LTcaptype{none} % do not increment counter
\begin{longtable}[]{@{}
  >{\raggedright\arraybackslash}p{(\linewidth - 6\tabcolsep) * \real{0.2500}}
  >{\raggedright\arraybackslash}p{(\linewidth - 6\tabcolsep) * \real{0.2500}}
  >{\raggedright\arraybackslash}p{(\linewidth - 6\tabcolsep) * \real{0.2500}}
  >{\raggedright\arraybackslash}p{(\linewidth - 6\tabcolsep) * \real{0.2500}}@{}}
\toprule\noalign{}
\begin{minipage}[b]{\linewidth}\raggedright
\#
\end{minipage} & \begin{minipage}[b]{\linewidth}\raggedright
冲突
\end{minipage} & \begin{minipage}[b]{\linewidth}\raggedright
两方
\end{minipage} & \begin{minipage}[b]{\linewidth}\raggedright
本书处置
\end{minipage} \\
\midrule\noalign{}
\endhead
\bottomrule\noalign{}
\endlastfoot
1 & Branda\literalquote{}Savings\literalquote{}352 vs 383 & 正文352／附表A＋B可复算383 &
采383并注 \\
2 & Branda西班牙行\literalquote{}350 350 0\literalquote{}第四格 & 空白或丢失 &
不写\literalquote{}两列皆0\literalquote{}，实质结论不变 \\
3 & Marion 1812并合地税 & 预期毛342.26／净226.39 & 一律标\literalquote{}预期非实收\literalquote{} \\
4 & F4经常收入T & 650／750／800 &
中心750；按路径疆域重基714--739（第五十四章第五节） \\
5 & 陆300＋海195闭合 & 高收入＋盟约案 vs 中心案缺120m &
两案分列，不互冒 \\
6 & P4预算层级 & 195m→111 A12→.902/.876；165→84→.769/.747；220--260高档
& 红队纠正后口径；高档不在平衡案 \\
7 & 失安特卫普 & P4 .444/.267为W基线单一冲击 &
P基线比例近似≈.58/.56，标近似 \\
8 & SA白银400--800万 vs MX 90--330万 & 毛流量 vs 法国净得 &
以MX四因子为承重，SA为参照，不相加 \\
9 & 甜菜糖1813 & 展望350万kg／334厂 vs 实产110万kg／158厂 &
用实产；1.2\%而非3.7\% \\
10 & HT死亡率 & 6\%低年／13\%全期均值含疫年／17--30\%疫年 &
不合成；13\%已含疫年 \\
11 & HT在场5,000--8,000 & 在场非在册；q\_C 0.65--0.80 &
在册6,250--12,308 \\
12 & P7 MAN-01 & 需求／供给标签互换 & 已改正读法 \\
13 & 外籍兵员 & Grab 1812战役比例／Elting
1807战区比例／Houdaille累计来源 & 三口径不互加 \\
14 & Rowe表3 & 征前逃避33.65\%（分母在册）vs
入伍后逃亡3.00\%（分母到营） & 不相乘 \\
15 & 1813年10月到营56.71\% & 单批截止日 & 不当全国全年 \\
16 & 粮秣 & F2 2,000吨／日（12.5万马）vs E1 1,750（10万马） &
差额由马群解释；均为说明案 \\
17 & van Creveld p.64 & \literalquote{}Even if\literalquote{}算例非莫斯科实发；吨制未明 &
1.5kg为公吨近似输入 \\
18 & 巴黎公约对奥援军 & 法文10,000／德文16,000 &
并列；数据表只用42,000与12,000 \\
19 & 普鲁士国债 & Murau 1806 53.5m／1820 180m vs 克拉克转载1806
35m／1810 66m & 并列 \\
20 & 拜伦特购价 & Weis 20m fl. vs HdBG 23m法郎 & 并列 \\
21 & 拜永条约缪拉出兵 & 18,000／3,000／25 vs 16,000／2,500／20 & 并列 \\
22 & 缪拉军1815 & Johnston 35,000＋5,000 vs Castellano 17,000 & 并列 \\
23 & 西班牙教士总量 & 20万 vs 14.8--17万 & 并列 \\
24 & 西班牙人口损失 & Fraser 350--510k战前／215--375k战时 vs
Esdaile\literalquote{}可能超百万\literalquote{} & 并列 \\
25 & 巴约讷洪塔首会 & 65／75 & 并列 \\
26 & 卡杜达尔行刑日；莱内表决；三次\literalquote{}八万令\literalquote{}；公投反对票 &
6-25/6-28；223:31/223:51；5,657/5,740 & 备案 \\
27 & 俄军1812 & Lieven一线242k／Riehn
190k／488k／428k；茹拉夫斯基131.8万满编 & 禁互填；茹拉夫斯基禁入总装 \\
28 & Troshin 0.74 & 价格贸易条件指数 & 非出口量或收入恢复 \\
29 & 戴维表29 & 加总79.88／99.28≠100 & 缺格不当零 \\
30 & 基尔\literalquote{}一百万\literalquote{} & 军团补贴 vs 秘密补偿 & 并列 \\
31 & 丹麦炮艇47艘 & 丹麦炮艇 vs 挪威双桅 & 并列 \\
32 & 瑞士团入俄 & 7,000--10,000／生还700--1,300三源 & 并列 \\
33 & 华沙军费占比 & Grab预算三分之二 vs Kowalczyk九成 &
无决算系列不裁 \\
34 & 印度Cartwright总账 & OCR 38,614,174 vs 还原28,614,174 &
按51,182,127−22,567,953还原 \\
35 & 毛里求斯远征 & 6,848／11,306／16,000 & 并列 \\
36 & 三舰派遣日期 & 两说 & 并列 \\
37 & 强征美国水手 & 门罗6,257／联邦党\textless800 & 并列 \\
38 & 半岛人口新西班牙 & 7,906 vs 70,000--80,000 &
并列（决定统治层厚度） \\
39 & Torrente秘鲁栏 & 十年合计标为年产 & 不作1806年产量 \\
40 & 卡斯尔雷\literalquote{}Old Spain\literalquote{}段页眉 & \literalquote{}180M\literalquote{}属1807或1808 & 待确认 \\
41 & Rouanet--Piano & 正文3,256 vs 附录2,812 & 保留冲突 \\
42 & 比利时被征率 & 7.135\%非6.12\% & 已改 \\
43 & Fichte销量、Tugendbund人数、烧炭党三口径 &
600册／620--700人／24--30k、50--60k、200k＋ & 并列 \\
44 & 法国关税收入三套 & Heckscher／Marion／AHAD 1807＝77、1812＝109.5 &
量级不合并列 \\
45 & P5旧摘录455／465、365／377 & gen4采Ellis下行465／377 & 采gen4 \\
46 & 1806年断点CI印成1805 & M1原文错误 & 不擅改 \\
47 & C1先验 & P6 0.45--0.60 vs F1 0.05--0.15 & 只属路径B；分歧保留 \\
48 & C9先验 & P6 0.35--0.55 vs F1 0.10--0.30 & 取0.20--0.45 \\
49 & M-B3 & P6 0.45--0.60 vs 本书0.40--0.60（1805受辱） & 取本书带 \\
50 & 引擎delta\_prussia\_takeover\_mid & 6.0 vs 注释4.5--7.5万 &
未定案，只影响B3v2 \\
51 & 1813年汉堡法军撤出日期 & 两说 & 并列 \\
52 & 意大利征兵 & 155k vs 200k＋ & 并列 \\
53 & 磨坊税 & 撤销 vs 减半 & 并列 \\
54 & 英国人口序列 & GB／UK边界断点 & 不直接用GBR长序列 \\
55 & 敖德萨1813 & 网站缺值渲为零 & 不当观测零 \\
\end{longtable}
}

\begin{center}\rule{0.5\linewidth}{0.5pt}\end{center}

\subsection{附录三　红队意见处置表}\label{ux9644ux5f55ux4e09-ux7ea2ux961fux610fux89c1ux5904ux7f6eux8868}

十项实质异议的处置见导论第四节。局部修正：

{\def\LTcaptype{none} % do not increment counter
\begin{longtable}[]{@{}
  >{\raggedright\arraybackslash}p{(\linewidth - 4\tabcolsep) * \real{0.3333}}
  >{\raggedright\arraybackslash}p{(\linewidth - 4\tabcolsep) * \real{0.3333}}
  >{\raggedright\arraybackslash}p{(\linewidth - 4\tabcolsep) * \real{0.3333}}@{}}
\toprule\noalign{}
\begin{minipage}[b]{\linewidth}\raggedright
项
\end{minipage} & \begin{minipage}[b]{\linewidth}\raggedright
红队意见
\end{minipage} & \begin{minipage}[b]{\linewidth}\raggedright
处置
\end{minipage} \\
\midrule\noalign{}
\endhead
\bottomrule\noalign{}
\endlastfoot
P4数字误读 & 195m与.902/.876本身匹配 & 已纠正（第六十三章） \\
白银两口径 & SA 400--800万 vs MX 90--330万未调和 &
分列不相加（第六十七章） \\
Branda精度 & \literalquote{}两列皆0\literalquote{}过度精确 & 已改（第四章） \\
MAN-01 & 标签互换 & 已改读 \\
独立束数 & R1／R2同国；US／MX／X1共用白银材料；Ellis／Heckscher转引 &
T1定理改为\literalquote{}各至少两束独立支撑\literalquote{}（第四十八章） \\
\literalquote{}普遍准入\literalquote{} & 非当时法方提案 & 标MODEL（第五十一章） \\
K3手填 & 引擎无签约状态 & 标手填，给可核代理（第五十二章） \\
引文抽核8处 & 全部命中 & --- \\
复算P2／P8／P4、F2／B2 & 相符 &
复算正确不消除世界线问题（本书全盘接受） \\
未复跑P7三脚本、未回源D2／Q1／X2、未运行E & --- &
E由本书作者运行；P7未重跑（本书以骨架年表替代338条装配） \\
\end{longtable}
}

\begin{center}\rule{0.5\linewidth}{0.5pt}\end{center}

\subsection{附录四　数字总表}\label{ux9644ux5f55ux56db-ux6570ux5b57ux603bux8868}

{\def\LTcaptype{none} % do not increment counter
\begin{longtable}[]{@{}
  >{\raggedright\arraybackslash}p{(\linewidth - 6\tabcolsep) * \real{0.2500}}
  >{\raggedright\arraybackslash}p{(\linewidth - 6\tabcolsep) * \real{0.2500}}
  >{\raggedright\arraybackslash}p{(\linewidth - 6\tabcolsep) * \real{0.2500}}
  >{\raggedright\arraybackslash}p{(\linewidth - 6\tabcolsep) * \real{0.2500}}@{}}
\toprule\noalign{}
\begin{minipage}[b]{\linewidth}\raggedright
项目
\end{minipage} & \begin{minipage}[b]{\linewidth}\raggedright
数值
\end{minipage} & \begin{minipage}[b]{\linewidth}\raggedright
性质
\end{minipage} & \begin{minipage}[b]{\linewidth}\raggedright
出处章
\end{minipage} \\
\midrule\noalign{}
\endhead
\bottomrule\noalign{}
\endlastfoot
法军在册 & 400--550k（中心450） & 包络 & 二 \\
盟军在册 & 150--250k（中心200） & 包络 & 二 \\
在场率 & 0.85 & 参数 & 二 \\
姿态需求（千在场） & S2 540／S5 490／S6-lite 566.25／S6-max
922.25／S6-full 1,102.25 & 模型 & 二、四十八 \\
筛查余额（千） & S2 +55／S5 +147.5／S6-lite +28.75／S6-max
−122.25／S6-full −302.25 & 模型 & 二 \\
西班牙占领 & 29--36万（Oman两切片） & 史实 & 二 \\
年抽人上限 & 9.8--12.7万 & 包络 & 二 \\
和平年征容忍 & ≤8--10万 & 裁量阈值 & 五 \\
1812--13征召 & 1,527,000／十五月 & 史实（Pigeard） & 二 \\
1813-10-09一批到营 & 72,265/127,433＝56.71\% & 史实（单批） & 二 \\
Rowe莱茵1806--10 & 征前逃避33.65\%／入伍后逃亡3.00\% & 史实 & 二 \\
新省净供兵 & hazard6 +0.790‰／五年含训练 −0.463‰ & 制度敏感性 &
二、五十四 \\
经常收入T & 750（650／800） & 构造 & 四 \\
路径疆域重基T & S2本义724--739／路径B 714 & 本书重算 & 五十四 \\
战争附加成本 & 372.5m／年＝156.4＋161.7＋43.7＋10.7 & Branda复算 & 四 \\
贡赋现金依赖 & 16.4\%（剔西班牙12.4\%） & 复算 & 四 \\
和平释放 & 216.1m／年 & 复算 & 四 \\
陆海恒等式 & 750−300−25＝425；海195→陆230；50万×700＝350→缺120 & 恒等式
& 四 \\
Marion 1812并合地 & 预期毛342.26／净226.39 & 预期 & 四 \\
荷兰年金 & 78→26m & 史实 & 二十一 \\
汉堡罚金 & 4,800万法郎；银行7,506,956 Mark Banco & 史实 &
二十、三十八 \\
承包商年净动员上限 & 6,000万--1亿 & 实测 & 三十八 \\
白银摩擦 & ≤23.5\%（1806-05）／≤62.5\%（1807下半年） & 实测 & 三十八 \\
Hope合约价 & 3.75法郎／比索 & 实测 & 三十六 \\
走私保费 & 莱茵1800 6\%→1806--08 26\%→1809 30\%→1809末--1811 50\% & 史实
& 三十九 \\
海关人员 & 1812峰值35,000 & 史实 & 三十九 \\
关税 & 严格15--35／执行28--38.5／税基−60至−110；保护60--90／执行14--24 &
模型 & 三十九、四十八 \\
法国出口体系内比重 & 1806 58.0\%／1810 74.6\%／1817 33.8\% &
史实（Ellis） & 五十九 \\
Chaptal 1819 & 毛1,820m；利润10.00\%；工资46.37\% & 史实 & 五十九 \\
1848大陆／英工业总量比 & S0 1.17／S2 1.55--1.67／S3 1.59--1.74／S5
1.37--1.51／S6 1.22--1.27 & 包络 & 五十九 \\
铁路1848 & 2,500--5,000km（中3,500） & 包络 & 五十九 \\
战列舰体 & 1810 50／1812 72／1814 74 & 史实（Todorov） & 三 \\
1807--13下水 & ≥50艘（安特卫普18） & 逐舰表 & 三 \\
1830法A12 & W84／P111；失安特卫普50／71 & 模型 & 三、六十三 \\
1830比值 & W .742/.449；P .902/.876；法P英W .902/.731；165m
.769/.747；α≈.58/.56 & 模型／近似 & 六十三 \\
英A12 1830 & W 115--145／P 100--140 & 包络 & 八 \\
英1811 consol & 4.67\%（+26.6bp） & 史实 & 七 \\
英债息／收入 & 1810 33.4\%／1811 34.6\% & 史实 & 七 \\
60\%探针 & 1817--23（S3零增长模型） & 模型 & 七 \\
英1848工业 & S2 79--100\%／S3 54--79\%／S5 88--109\%史实 & 包络 & 十 \\
爱尔兰 & 法占需33--66k；汲取≤£5--5.5m & 模型 & 十二 \\
俄军1812 & 在册\textasciitilde60万／一线24.2万／Riehn 488k总、428k入战 &
史实 & 十五 \\
俄庄园损失 & S2中心5.3\%／S3中心22.8\% & 模型 & 十四 \\
冷和平半衰期 & 4年（史实）→8--12年（关税容忍） & 裁量 & 十三 \\
华沙 & 1812
74,722（部长会议）／108k（Kukiel全口径）；被占1813--15提取2.5亿zł & 史实
& 十六 \\
奥1809财政 & 收入31.1／支出86.3／军费66.8（百万CM盾） & 页图 & 十七 \\
申布伦限军 & 150,000（非五万） & 史实 & 十七 \\
普限军 & 42,000十年期；三要塞法军10,000 & 史实 & 十八 \\
残普1848 & 620--810万人；常备5.6--9.7万 & 算术 & 十八 \\
完整普鲁士1848（路径B） & 常备8--12万／动员16--24万 & 本书裁量 &
六十六 \\
奥强于普 & S2本义0.70--0.85／路径B 0.45--0.60 & 裁量 & 六十六 \\
南德 & 1812入俄巴33,000／符15,800（生还300） & 史实 & 十九 \\
荷兰骚动 & 580＋起、80＋起义、法军均15,000 & 史实（Joor） & 二十一 \\
意大利王国 & 征兵155,000／军70,000／俄役27,000生还\textasciitilde1,000 &
史实（Grab） & 二十三 \\
意大利分支 & C 0.55--0.65／B 0.35--0.45（S2本义）0.20--0.30（B）／A
0.10--0.20 & 裁量 & 六十五 \\
那不勒斯 & 税57→73m里拉；1812军费80\%；卡拉布里亚5,421匪 & 史实 &
二十四 \\
西班牙vales & 18.93亿rv；贴水63\% & 史实 & 二十六 \\
西班牙终局 & 0.40--0.55／0.30--0.45／0.05--0.15／≤0.02 & 裁量 &
六十四 \\
纳瓦拉 & 榨取1.52亿rv＝五年产出40\%＋ & 史实（Tone） & 二十七 \\
洪塔构成 & 585人：市政23.9\%／教士23.6\%／军官18.3\% & 史实（Fraser） &
二十七 \\
布拉干萨出海 & 8,000--15,000人、八千万克鲁扎多 & 史实 & 二十八 \\
奥斯曼新军 & 22,685＋1,590（1806） & 史实 & 二十九 \\
埃及征兵 & 1830年代130,000 & 史实（Fahmy） & 三十 \\
波斯军1808 & 骑144,000／步60,000／炮兵实际不存在 & 史实（Gardane） &
三十一 \\
英对波斯 & 200,000 tomans≈£150,000／年 & 史实 & 三十一 \\
EIC & 贸易利润1793--1812累计£6.29m；债£8m→£32m & 史实 & 三十二 \\
印度远征方案 & 1808-06-10：56舰、19,600兵、10,600水手 & 史实 & 三十三 \\
美国转口 & 1806 \$60.28m→1808 \$12.99m & 史实 & 三十四 \\
英美战争概率 & S1 0.05--0.12／S2 0.30--0.45／S3 0.55--0.75 & 裁量 &
三十四 \\
墨西哥白银法国实得 & 90--330万pesos／年（中枢190万） & 四因子 &
三十六 \\
补贴条约要求 & 960--1,920万pesos／年 & 史实换算 & 三十六 \\
加勒比在场 & 5,000--8,000；在册6,250--12,308 & 模型 & 三十五 \\
甜菜糖1813 & 实产1,100,000kg／158厂 & 史实 & 四十 \\
蒸汽 & 英1830 160,000hp；法1839 33,000 & 史实 & 四十 \\
纱锭 & 帝国≈133万 vs 英460--500万 & 史实 & 四十 \\
Tugendbund & ≤700；1809名录620 & 史实 & 四十一 \\
宪章请愿 & 1842近350万；1848点验1,975,496 & 史实 & 四十一 \\
斯特拉斯堡河运 & 1809--11年2.4--2.7万公吨 & 史实（Ellis附录D） &
五十九 \\
敖德萨小麦 & 1815 171,708／1816 188,259公吨 & 史实 & 四十二 \\
内圈整合 & 7--12年机器期；认同不交割 & 判定 & 四十三 \\
Q1 & 22/35大起义；5/35持续抵抗（全部R∧O） & 编码 & 四十五 \\
Q2净兵 & L −0.04至+5.32k／M −7.10至+2.19k／H
−26.73至−6.03k每百万人第10年 & 模型 & 四十六 \\
提尔西特债价 & 法+12.38\%／英−2.15\% & 描述 & 四十七 \\
P2联合账 &
甲+50.4k／8.5m；乙含贝格+33.7k／20.9m；乙四地＋外圈−116.7k／174.2m；丙−479.8k／532.6m
& 模型 & 五十四 \\
P2乙乐观反例 & 需求571,650；余23,598；军费余19.23m & 模型 & 五十六 \\
P2闭合案 & 434,286在册；余17,093 & 模型 & 五十四 \\
路径E引擎 & 1806--07财政−；1808后可得551；1812余3m & 本书补算 & 五十 \\
C类节点乘积 & E 9.6e-5--3.0e-3／B 2.9e-6--2.9e-4 & 序数比较 & 五十 \\
条件成功 & P6 0.30--0.50／P8 J 0.063--0.404 & 裁量 & 五十三 \\
无条件成功 & P6 0.018--0.090／P8 0.0038--0.0727 & 裁量 & 五十三 \\
O1848 & P8 0.17--4.44\%（中1.27\%）；D\textbar C 0.444--0.611 & 裁量 &
六十二 \\
王朝存续1848 & 40--55\%；行政国家≥70\% & 裁量 & 六十 \\
拿破仑寿限 & 1820s死40--50\%／1830s
30--35\%／1840s＋15--25\%；活到1829≈0.50--0.65 & 裁量 & 一、五 \\
\end{longtable}
}

\begin{center}\rule{0.5\linewidth}{0.5pt}\end{center}

\end{document}
